\chapter{Higher genus}
\label{chap:higher-genus}\label{ch:higher-genus}

\index{higher genus!bar complex|textbf}
\index{quantum corrections|textbf}

% Macro safety net for cross-chapter BRST references (BRST functor of
% chap:universal-conductor and the associated ghost-charge invariant).
% Defended by \providecommand to keep this chapter portable under
% standalone extraction.
\providecommand{\BRST}{\mathrm{BRST}}
\providecommand{\SpCh}{\mathrm{Sp}^{\mathrm{ch}}}
\providecommand{\PhiFA}{\Phi^{\mathrm{FA}}}

Everything up to this point takes place at genus zero; but
genus zero on a curve is not the same as ``over a point.''
The bar differential $\dfib$ is constructed from collision
residues on the Fulton--MacPherson compactification
$\overline{C}_n(X)$, its nilpotency $\dfib^{\,2} = 0$ uses
the Arnold relations among the forms
$\eta_{ij} = d\log(z_i - z_j)$ on the diagonal complement,
and the higher $A_\infty$ operations $m_n$ for $n \geq 3$ are
transferred along a retract whose source is
$C_*(\overline{\mathcal{M}}_{0,n+1})$, the chain complex of the
Stasheff associahedron. The formal-disk restriction that recovers
classical Koszul duality over a point is a homotopy retract: it is
additional data, not a tautology.

What changes at genus~$g \geq 1$ is that the fiberwise bar operator
belongs to a curved $A_\infty$/CDG package rather than to an ordinary
dg package.
The genus-$0$ proof of $\dfib^{\,2} = 0$ is combinatorial:
every codimension-two
boundary stratum of $\overline{\cM}_{0,n}$ is reached by exactly
two one-edge tree degenerations whose boundary contributions
cancel, and the tree combinatorics of the associahedron exhausts
the collision data of $\mathbb{P}^1$. A smooth curve of genus
$g \geq 1$ refuses this combinatorial picture. Pinching a
handle produces a degeneration that no tree can see: a
non-separating node raises $b_1$ and leaves a graph with
genuine loops; a separating node carves the surface into
lower-genus components whose collision data is genuinely
global. The fiberwise bar differential $\dfib$ inherits a nullary
curvature element $m_0^{(g)}$ from these degenerations, and the
genus-zero strict identity is replaced by
$m_1^2=[m_0^{(g)},-]$. The scalar diagonal of that curvature has class
$\kappa(\cA)[\omega_g^{\mathrm{Ar}}]$, and its appearance is forced.
Here $\omega_g^{\mathrm{Ar}}$ denotes the vertical Arakelov
$(1,1)$-form on the fixed smooth fiber~$\Sigma_g$.  The horizontal
characteristic classes on $\overline{\cM}_g$ arise from the Hodge
metric on $\mathbb{E}=\pi_*\omega_\pi$: its Chern--Weil forms give
$\lambda_1=c_1(\mathbb{E})$, and the scalar obstruction lane uses
the top Hodge class $\lambda_g=c_g(\mathbb{E})$.  The binary
collision kernel gives an operator class
$q_2(K_\cA^{\mathrm{coll}})$, and the normalized residue trace gives
$\kappa(\cA)=\tau_{\cA,2}^{\mathrm{res}}
q_2(K_\cA^{\mathrm{coll}})$ on the abelian and scalar lanes.  Its
product with $\lambda_g$ is the leading genus invariant in the scalar
trace lane. The source of the curvature is
geometric: the irreducible stable graph
$\Gamma_{\mathrm{irr}} \in \partial\overline{\cM}_{1,1}$ has a
single self-edge, and the BV operator $\Delta$ obtained by
tracing the cyclic pairing over that self-edge contributes
$\kappa$ through the normalized loop-contraction trace, matched with the class
$\lambda_1 = \tfrac{1}{12}\delta_{\mathrm{irr}}$ on the moduli side.
At higher genus the same mechanism repeats through the Feynman
transform, summing over stable graphs weighted by handles.  Its
boundary law has two components: separating clutching pulls the Hodge
bundle back as a direct sum, while
the non-separating clutching map pulls $\lambda_g$ back to zero for
$g\ge2$ (Proposition~\ref{prop:lambda-g-clutching}). Thus the
genus-$1$ identity $\lambda_1=\delta_{\mathrm{irr}}/12$ is the
initial calculation, and the two clutching formulas are the
higher-genus boundary law. The passage
from trees to stable graphs is the passage from operads to modular
operads, and the obstruction tower $\Theta_\cA^{\leq r}$ records the
higher-degree modular data that no tree could ever see.

The handle-pinching degenerations take two forms: a separating node
splits the surface into two components of lower genus; a
non-separating node increases the first Betti number. The
combinatorial record of all possible degenerations is a
\emph{connected stable graph}~$\Gamma$ with a genus label~$g(v)$
at each vertex, edges for nodes, half-edges for marked points, and
total genus $g(\Gamma) = b_1(\Gamma) + \sum_v g(v)$. Trees are
the genus-$0$ special case. The Deligne--Mumford stack
$\overline{\cM}_{g,n}$ is stratified by these graphs: each
codimension-one boundary divisor is a one-edge degeneration, and the
oriented stable-graph boundary operator satisfies
$\partial_{\mathcal G}^{\,2}=0$. This boundary cancellation is the
geometric origin of the strict identity for the \emph{total} modular
differential $D_\cA^2 = 0$ in each genus.  This cancellation concerns
the total modular differential; the fixed-fiber operator retains the
curved identity $m_1^2=[m_0^{(g)},-]$.

The passage from trees to stable graphs corresponds to the passage from
operads to \emph{modular operads}: algebraic structures whose
operations are indexed by arbitrary stable graphs, with trees forming
the genus-$0$ subcase, and whose composition laws are the clutching maps of
$\overline{\cM}_{g,n}$. The Feynman transform
(Definition~\ref{def:feynman-transform}) sums over all stable
graphs weighted by $\hbar^{g(\Gamma)}$; its genus-$0$ truncation
recovers the operadic bar construction
(Corollary~\ref{cor:feynman-genus0-reduction}).

\smallskip\noindent\textbf{The first non-tree graph:
$\overline{\cM}_{1,1}$.}\enspace
The moduli space $\overline{\cM}_{1,1}$ has one boundary stratum:
the unique stable graph $\Gamma_{\mathrm{irr}}$ with one vertex of
genus~$0$, one self-edge, and one half-edge
($g(\Gamma_{\mathrm{irr}}) = b_1 + g(v) = 1 + 0 = 1$,
$\mathrm{val}(v) = 1 + 2 = 3$). This graph records the
degeneration of a smooth genus-$1$ curve into an irreducible nodal
rational curve: the torus pinches to a sphere with one node.

The contribution of $\Gamma_{\mathrm{irr}}$ to the modular bar
differential is the BV operator~$\Delta$: it traces the cyclic
pairing $\langle -, - \rangle_\cA$ over the self-edge, producing
$\Delta(sa) = \sum_i \langle e_i, a \rangle\,se^i$ where
$\{e_i\}$, $\{e^i\}$ are dual bases. The genus-$1$ curvature element
$m_0^{(1)}$ arises from exactly this one graph. Its scalar diagonal
projection is $\kappa(\cA)\cdot\omega_1^{\mathrm{Ar}}$: the cyclic
self-contraction gives the scalar~$\kappa$.  Under the genus-one trace
comparison~$H_D^1$, the fundamental class of the boundary
divisor $[\partial\overline{\cM}_{1,1}] = \delta_{\mathrm{irr}}$
gives the tautological class~$\lambda_1 =
\tfrac{1}{12}\delta_{\mathrm{irr}}$ in
$H^2(\overline{\cM}_{1,1})$.  The self-loop is the irreducible source
of the genus-one curvature.

\smallskip
On~$\mathbb{P}^1$, the bar differential $\dfib$ squares to zero, the
Arnold relation holds exactly, and bar-cobar duality operates
within strict dg~algebras. At genus~$g \geq 1$, every one
of these statements acquires a correction controlled by the stratified
geometry of~$\overline{\mathcal{M}}_{g,n}$. Genus becomes a deformation
variable: the propagator picks up period corrections
from~$H^0(\Sigma_g, \Omega^1)$, the bar differential becomes curved
with $m_1^2=[m_0^{(g)},-]$, and the scalar diagonal of $m_0^{(g)}$
is $\kappa(\cA)\cdot\omega_g^{\mathrm{Ar}}$. Obstruction classes
become genuinely modular: their Hodge realizations live in the
cohomology of the Hodge bundle.

The modular characteristic
$\kappa(\cA)$ is the linear shadow (Ring~1), the spectral
deformation~$\Delta_{\cA}$ the first nonlinear shadow (Ring~2),
and the genus tower is a single completed homotopy object whose
finite-order truncations are the named shadows. The shadow obstruction tower
$\Theta_{\cA}^{\leq r}$, developed in
\S\ref{sec:koszul-across-genera}, organizes these corrections; the
foundations below are the support on which that tower rests.

Three objects govern the genus tower: the curved fiberwise
differential~$\dfib$ (Section~\ref{sec:complementarity-theorem}), the
corrected holomorphic differential~$\Dg{g}$ satisfying~$\Dg{g}^{\,2}=0$
via the Fay trisecant identity, and the modular convolution dg~Lie
algebra~$\gAmod$ (the strict model of the homotopy-invariant
$\Definfmod(\cA)$,
\S\ref{subsec:two-level-convention}), whose Maurer--Cartan
elements encode the full genus expansion.

On $\mathbb{P}^1$, the logarithmic form
$\eta_{12} = d\log(z_1 - z_2)$ is single-valued and the Arnold
relation~\eqref{eq:arnold-seed} holds exactly. On an elliptic
curve~$E_\tau$ with period~$\tau$, the function $z_1 - z_2$ is not
single-valued, and the propagator becomes
\begin{equation}\label{eq:elliptic-propagator}
\eta_{12}^{(1)}
=
\left[
\partial_u\log \vartheta_1(u\,|\,\tau)
+2\pi i\,\frac{\operatorname{Im}u}{\operatorname{Im}\tau}
\right]_{u=z_1-z_2}
(dz_1-dz_2),
\end{equation}
where $\vartheta_1$ is the odd Jacobi theta function. The correction
is a $(1,0)$-form with real-analytic coefficient and is forced by
double-periodicity.  The curved unary identity and its scalar diagonal
are
\begin{equation}\label{eq:curvature-visible}
m_1^{(1)\,2}=[m_0^{(1)},-],\qquad
\operatorname{tr}_{\mathrm{diag}}(m_0^{(1)})
= \kappa(\cA) \cdot \omega_1^{\mathrm{Ar}}, \qquad
\omega_1^{\mathrm{Ar}} = \frac{i}{2\operatorname{Im}(\tau)}\,
dz \wedge d\bar{z}.
\end{equation}

On a Riemann surface $\Sigma_g$ of genus~$g$, the propagator acquires
period corrections from $H^0(\Sigma_g, \Omega^1)$.
One first-variable Arakelov-normalized representative is
\[
\eta_{ij}^{(g)} = \Bigl[\partial_{z_i}\!\log E(z_i, z_j)
 + \pi \sum_{\alpha,\beta=1}^{g}
 \omega_\alpha(z_i)\,(\operatorname{Im}\Omega)^{-1}_{\alpha\beta}\,
 \operatorname{Im}\!\Bigl(\int_{z_0}^{z_j}\!\omega_\beta\Bigr)
\Bigr],
\]
where $\Omega$ is the period matrix, $\omega_1,\ldots,\omega_g$ the
normalized abelian differentials, and $E(z,w)$ the prime form
(Definition~\ref{def:higher-genus-log-forms}).
The non-holomorphic correction restores double periodicity but
introduces the curvature element $m_0^{(g)}$, whose scalar diagonal
class is $\kappa(\cA)[\omega_g^{\mathrm{Ar}}]$, where
$\omega_g^{\mathrm{Ar}}$ is the Arakelov $(1,1)$-form
(Theorem~\ref{thm:quantum-arnold-relations}). The total corrected
differential $\Dg{g}$ satisfies $\Dg{g}^{\,2} = 0$
(Theorem~\ref{thm:quantum-diff-squares-zero}) after the period and
Gauss--Manin terms are included. The Fay trisecant identity supplies
the holomorphic tangent-screen cancellation, while the positive-genus
configuration-space algebra retains its global period data.

As $[\Sigma_g]$ varies over $\overline{\mathcal{M}}_g$, the bar
complex forms a family of curved cochain complexes.  Its scalar
fiber curvature is the vertical Arakelov form
$\kappa(\cA)\omega_g^{\mathrm{Ar}}$.  The Hodge bundle on the base
has $\lambda_1=c_1(\det\mathbb E_g)$ and
$\lambda_g=c_g(\mathbb E_g)$.  The comparison between these two
geometric levels is the package $H_D^K(\cA,g)$ of
Definition~\ref{def:scalar-diagonal-hypothesis}.  Under that package,
Theorem~\ref{thm:genus-universality} gives
\[
 \mathfrak O_g^K(\cA)=\kappa(\cA)\lambda_{-1}(\mathbb E_g),
 \qquad
 \operatorname{obs}^{\mathrm{Hdg}}_g(\cA)
 =(-1)^g\kappa(\cA)\lambda_g.
\]
The stable-graph trace gives the genus-$g$ free-energy package
\[
F_g(\cA)
\;=\;
F_g^{\mathrm{sc}}(\cA)
\;+\;
\delta F_g^{\mathrm{cross}}(\cA)
\qquad
\textup{with }F_g^{\mathrm{sc}}(\cA)
=\kappa(\cA)\lambda_g^{\mathrm{FP}},
\]
with $\delta F_1^{\mathrm{cross}}=0$ in genus~$1$.  The
scalar-diagonal graph package sets the cross-channel term to zero.
The scalar free-energy generating function is
$\kappa(\cA)\cdot(\hat{A}(i\hbar) - 1)$
\textup{(}Theorem~\textup{\ref{thm:family-index}}\textup{}) after
pairing with the Faber--Pandharipande test class
$\psi_1^{2g-2}$ on $\overline{\mathcal{M}}_{g,1}$, while
the full all-weight series acquires the explicit correction
$\sum_{g\ge2}\delta F_g^{\mathrm{cross}}(\cA)\hbar^{2g}$.
Any BRST interpretation of~$\kappa$ enters through its own trace
comparison at level~$5$.

The scalar diagonal curvature
$\operatorname{tr}_{\mathrm{diag}}(m_0^{(g)})
=\kappa(\cA)\cdot\omega_g^{\mathrm{Ar}}$ is the scalar projection of the full MC element
(cf.\ Remark~\ref{rem:mc-uncurving} for the
uncurving mechanism). The
modular cyclic deformation complex
(Definition~\ref{def:modular-cyclic-deformation-complex}) is a
Maurer--Cartan family parameterized by genus~$g$ and point
configurations. The universal MC element
$\Theta_\cA \in \mathrm{MC}(\gAmod)$ projects, at each genus~$g$, to:
the curvature~$\kappa(\cA)$ at the scalar level; the curved
differential~$\dfib$ at the differential level; and the
fully corrected differential~$\Dg{g}$ with $\Dg{g}^{\,2}=0$ at
the period-corrected level. The strict dg model
comprises the three differentials and their interaction; the full homotopy
$L_\infty$ structure and the modular envelope appear in
\S\ref{sec:koszul-across-genera}.

\begin{remark}[Parallel track: the anomaly]\label{rem:anomaly-parallel-track}
\index{conformal anomaly!parallel track}
\begin{center}
\small
\begin{tabular}{ll}
\textbf{This chapter} & \textbf{Physics reading} \\ \hline
Scalar diagonal of $m_0^{(1)}$: $\kappa \cdot \omega_1^{\mathrm{Ar}}$ & Conformal anomaly \\
$\kappa(\cA) = 0$ (flatness) & total matter--ghost anomaly vanishes \\
Coderived category & Off-shell (curved background) \\
Hodge shadow $\operatorname{obs}^{\mathrm{Hdg}}_g$
& Characteristic class of the one-loop virtual object \\
Complementarity $Q_g + Q_g' = H^*$ & Matter--ghost pairing
\end{tabular}
\end{center}
The BRST flatness row is theorem-level on the cited genus-$0$ BRST
lane. The determinant and matter--ghost entries are physics readings
of the scalar obstruction package unless a cited theorem in the
corresponding chapter supplies the stated comparison.
\end{remark}

\begin{convention}[Higher-genus differentials; \ClaimStatusDefinitional]\label{conv:higher-genus-differentials}
\index{bar differential!higher-genus notation|textbf}
The Feynman transform differential
$d_{\mathrm{FT}} = d_{\mathrm{tree}} + d_{\mathrm{loop}}$
\textup{(}Proposition~\textup{\ref{prop:loop-filtration-compatible})} is
the single canonical differential on the genus-$g$ bar complex. For
computational purposes, three projections of this differential are
used throughout the manuscript.
\begin{enumerate}
\item[\textup{(i)}] \textbf{Fiberwise curved differential} $\dfib$.
 On a fixed curve $\Sigma_g$, the collision-residue differential built
 from the genus-$g$ propagator. This differential is \emph{curved} in
 the curved $A_\infty$/CDG sense:
 \[
  m_1^{(g)\,2}(a)
  =
  m_2(m_0^{(g)},a)-m_2(a,m_0^{(g)})
  =
  [m_0^{(g)},a]_{m_2},
 \]
 in the manuscript's desuspended-bar sign convention
 (Appendix~\ref{app:signs}). The curvature element
 $m_0^{(g)}$ is represented by the Arakelov $(1,1)$-form
 $\omega_g^{\mathrm{Ar}}$
 (Theorem~\ref{thm:quantum-arnold-relations}) together with the
 leading OPE anomaly coefficient. This is the raw chain-level
 statement.  Applying the scalar
 diagonal/uniform-weight projection gives
 \[
  \operatorname{tr}_{\mathrm{diag}}\!\bigl(m_0^{(g)}\bigr)
  =
  \kappa(\cA)\cdot\omega_g^{\mathrm{Ar}} .
 \]
 Under~$H_D^K(\cA,g)$, the associated scalar virtual object and
 Hodge shadow are
 \[
  \mathfrak O_g^K(\cA)
  =\kappa(\cA)\lambda_{-1}(\mathbb E_g),
  \qquad
  \mcurv{g}
  =(-1)^g\kappa(\cA)\lambda_g
  \in H^{2g}(\overline{\cM}_g).
 \]
 The comparison~$\beta_g(\cA)$ relates the fiberwise bar family to
 the Hodge bundle; Chern--Weil theory then computes its Euler class
 (Proposition~\ref{prop:chain-level-curvature-operator}).

\item[\textup{(ii)}] \textbf{Total corrected differential} $\Dg{g}$.
 The full differential on $\bar{B}^{(g)}(\cA)$, incorporating the
 prime-form period correction and the Gauss--Manin variation of the
 Hodge bundle:
 \[
 \Dg{g}
 =
 \dfib + d_{\mathrm{per}}^{(g)} + \nabla_{\mathrm{KS}}^{\mathrm{GM}},
 \]
 where $d_{\mathrm{per}}^{(g)}$ is the operator induced by the
 Arakelov period term in the prime-form propagator, and
 $\nabla_{\mathrm{KS}}^{\mathrm{GM}}$ is the Gauss--Manin connection
 applied along Kodaira--Spencer tangent directions. The normalized
 $A$-periods satisfy
 $\oint_{A_l}\omega_k=\delta_{kl}$ and are therefore a gauge choice,
 not deformation parameters. The actual period data are the period
 matrix entries $\Omega_{\alpha\beta}$ and the Abel--Jacobi integrals
 appearing in Definition~\ref{def:higher-genus-log-forms}; they live
 in the variation of $H^1(\Sigma_g,\C)$, not in
 $H^1(\mathcal{M}_g,\C)$, which vanishes for $g\ge2$ by Harer.
 This differential is
 \emph{strict}:
 \[
 \Dg{g}^{\,2} = 0
 \]
 \textup{(}Theorem~\textup{\ref{thm:quantum-diff-squares-zero})}.

\item[\textup{(iii)}] \textbf{Genus-$0$ collision differential} $\dzero$.
 The chiral bar differential from collision residues on configuration
 spaces of~$X$ (or the fiberwise collision component $d^{(0)}$ in
 the Leray decomposition). At genus~$0$, this involves the
 Fulton--MacPherson compactifications $\overline{C}_n(\bP^1)$ and
 the Arnold relations for the logarithmic forms $\eta_{ij}$.  The
 classical bar differential over a point is its formal-disk retract.
 The operator satisfies $\dzero^2 = 0$
 by the Arnold relation (Theorem~\ref{thm:arnold-three}).
\end{enumerate}

\noindent
\textbf{Categorical homes.}
The strict differentials $\dzero$ and $\Dg{g}$ produce honest
chain complexes; the bar complex equipped with either lives in
the \emph{derived category}. The curved differential $\dfib$
produces a curved dg coalgebra with
$m_1^2=[m_0^{(g)},-]$; the bar complex equipped with $\dfib$
lives in the \emph{coderived category}
$D^{\mathrm{co}}$~(Positselski;
see Appendix~\ref{app:coderived}).
The derived/coderived comparison, established on the manuscript's
flat-side comparison locus together with the coderived treatment of
the curved model, ensures that the flat model carries the ordinary
derived package while the curved model is its coderived counterpart.

\noindent
\textbf{Relation.}
The fiberwise differential $\dfib$ is curved; the prime-form
period operator and the Gauss--Manin/Kodaira--Spencer term cancel
the curvature via the Riemann bilinear relation and the Lagrangian
normalization of the $A$-cycle subspace, producing the strict
differential
$\Dg{g}=\dfib+d_{\mathrm{per}}^{(g)}
+\nabla_{\mathrm{KS}}^{\mathrm{GM}}$.
The passage $\dfib \leadsto \Dg{g}$ is the chiral analogue of
passing from a curved connection to a flat one. The full
genus-stratified total differential is
\[
 \Dtot = \sum_{g\ge 0} \hbar^{2g-2}\,\Dg{g}, \qquad
 \Dtot^{\,2} = 0.
\]

\noindent
\textbf{Type distinction.}
For $g\geq2$, the deformation class
$\operatorname{Obs}^{\mathrm{def}}_g$ and the
Hodge shadow~$\operatorname{obs}^{\mathrm{Hdg}}_g$ are the two
objects of Definition~\ref{def:genus-g-obstruction}.  The first lies
in deformation degree~$2$; the second lies in~$H^{2g}$.
The formula
$\operatorname{tr}_{\mathrm{diag}}(m_0^{(g)})
=\kappa\omega_g^{\mathrm{Ar}}$ denotes the vertical scalar diagonal
of $m_1^2=[m_0,-]$, while $\Dg{g}^{\,2}=0$ denotes the
period-corrected total differential. Earlier
chapters that use unadorned $d$ for the genus-$0$ bar differential are
using $\dzero$.
\end{convention}

\begin{remark}[Curvature as infinitesimal monodromy: the chiral Riemann--Hilbert correspondence]
\label{rem:curvature-riemann-hilbert}
\index{Riemann--Hilbert correspondence!chiral|textbf}
\index{curvature!as infinitesimal monodromy}
The genus-$g$ propagator $\eta_{ij}^{(g)}=d\log E(z_i,z_j)+\cdots$
acquires monodromy $2\pi i$ around $B$-cycles; the Arakelov
normalization transfers this into the Arnold defect, giving
$\operatorname{tr}_{\mathrm{diag}}(m_0^{(g)})
=\kappa(\cA)\cdot\omega_g^{\mathrm{Ar}}$.
Thus $\kappa(\cA)$ is the scalar infinitesimal monodromy of the logarithmic
propagator. The corrected differential $\Dg{g}$ absorbs the curved package
into quantum corrections, restoring $\Dg{g}^{\,2}=0$; this is the chiral
analogue of the Riemann--Hilbert correspondence ($\dfib$ is the curved
connection, $\Dg{g}$ the flat one).
The three differentials $\dzero$, $\dfib$, $\Dg{g}$ correspond to:
genus-$0$ (single-valued), genus-$g$ with monodromy (curved),
and the full logarithm on the universal cover (flat).
See Volume~II, Remark~\ref*{rem:three-models} for the full
comparison of the three chain-level models and their relationship
to the modular bar complex.
\end{remark}

\begin{proposition}[Gauss--Manin uncurving at chain level; \ClaimStatusConditional]
\label{prop:gauss-manin-uncurving-chain}
\index{Gauss--Manin connection!chain-level uncurving|textbf}
\index{uncurving!Gauss--Manin|textbf}
\textup{[Type signature: Open quadrant; relative de~Rham curved-bar
presentation; levels $2\to4$; uncurving package
$H_D^{\mathrm{unc}}$ consisting of an Arakelov period correction,
Gauss--Manin transport on the full cohomology local system, and a
compatible null-homotopy of the curvature coderivation.]}
Let $\cA$ be a cyclic chiral algebra on a smooth genus-$g$ curve,
with $g\geq1$.

\begin{enumerate}[label=\textup{(\roman*)}]
\item The full cohomology local system
\[
 \mathcal R^q
 =R^q\pi_*(\cA^{\boxtimes n}\otimes\Omega_{\log}^*)
\]
carries the flat Gauss--Manin connection
\begin{equation}\label{eq:gauss-manin-chain}
 \nabla^{\mathrm{GM}}
 =d_{\cM_g}+\Theta_{\mathrm{KS}},
 \qquad
 (\nabla^{\mathrm{GM}})^2=0.
\end{equation}
The Hodge subbundle carries its Chern connection, whose curvature
matrix represents its Chern classes.

\item The fixed-fiber bar operation is curved:
\[
 m_1^{(g)\,2}=[m_0^{(g)},-],
 \qquad
 \operatorname{tr}_{\mathrm{diag}}(m_0^{(g)})
 =\kappa(\cA)\omega_g^{\mathrm{Ar}}.
\]
The first identity lies in the curved endomorphism algebra; the second
is its scalar vertical trace.

\item Under~$H_D^{\mathrm{unc}}$, write
\(
 \Gamma_g=d_{\mathrm{per}}^{(g)}
 +\nabla_{\mathrm{KS}}^{\mathrm{GM}}
\).
The total relative operator
\[
 \Dg{g}=\dfib+\Gamma_g
\]
is strict, and its coderivation identity is
\begin{equation}\label{eq:chain-uncurving}
 0=\Dg{g}^{\,2}
 =\operatorname{ad}_{m_0^{(g)}}
 +\{\dfib,\Gamma_g\}+\Gamma_g^2.
\end{equation}
Thus the correction null-homotopes the curvature
\emph{endomorphism} $\operatorname{ad}_{m_0^{(g)}}$.  The comparison
of the scalar virtual family with
$\kappa(\cA)\lambda_{-1}(\mathbb E_g)$ remains the separate package
$H_D^K$.

\item The strict family
$(\bar B^{(g)}(\cA),\Dg{g})$ defines an object of the derived
category.  The fixed-fiber curved object
$(\bar B^{(g)}(\cA),\dfib,m_0^{(g)})$ defines an object of the
coderived category.  The functor between these realizations is part
of~$H_D^{\mathrm{unc}}$.
\end{enumerate}
\end{proposition}

\begin{proof}
Flatness in~\textup{(i)} is the Gauss--Manin theorem for the full
cohomology local system; the Hodge filtration varies by Griffiths
transversality.  Clause~\textup{(ii)} is the defining curved
$A_\infty$ relation together with the normalized Arakelov collision
trace.  Clause~\textup{(iii)} is precisely the null-homotopy required
in~$H_D^{\mathrm{unc}}$; in a local Torelli chart its two geometric
terms are the prime-form period correction and the
Kodaira--Spencer/Gauss--Manin transport.  Expanding
$(\dfib+\Gamma_g)^2$ gives~\eqref{eq:chain-uncurving}.
Clause~\textup{(iv)} records the corresponding categorical homes.
\end{proof}

\begin{definition}[Modular curvature: the direct definition]
\label{def:modular-curvature-direct}
\label{def:scalar-curvature-shadow}
\index{curvature coefficient!direct definition|textbf}
\index{self-contraction trace|textbf}
The \emph{modular curvature coefficient} of a cyclic chiral
algebra~$\cA$ with binary primitive kernel $K_{0,2}^{\cA}$ is
\begin{equation}\label{eq:kappa-direct}
\kappa(\cA)
\;:=\;
\operatorname{Tr}^{\mathrm{loop}}_\cA
\;:=\;
\operatorname{tr}_{\mathrm{cyc}}\!\left(
\Delta_{\mathrm{ns}}K_{0,2}^{\cA}\right),
\end{equation}
with the automorphism and OPE normalization fixed by
$H_D^{\mathrm{loop}}$.  In this normalization,
$\kappa(\cA) = \Theta_{\cA}^{\leq 2}$ is the
weight-$2$ level of the shadow obstruction tower
\textup{(}Definition~\textup{\ref{def:shadow-postnikov-tower})}.
The loop trace is the value of the unique
non-separating self-sewing of the binary corolla evaluated
on~$\cA$
(Theorem~\ref{thm:curvature-self-contraction}).

The fixed-fiber scalar curvature is
$\kappa(\cA)\omega_g^{\mathrm{Ar}}$ on~$\Sigma_g$.
Under~$H_D^K$, the associated $(g,n{=}0)$ Hodge component is
$(-1)^g\kappa(\cA)\lambda_g$
(Theorem~\ref{thm:genus-universality}).
The curvature hierarchy has three levels
(Convention~\ref{conv:higher-genus-differentials}): curvature lives in
$\dfib$, strict nilpotence in $\Dg{g}$, and scalar anomalies in
$\mathrm{tr}(\mcurv{g})$.
\end{definition}

\begin{remark}[Parameter-source diagram]\label{rem:parameter-source}
\index{parameter source!fiber vs.\ moduli}
The quantum correction parameters live in:
\[
H^1(\Sigma_g, \C)
\;\xrightarrow{\;\text{MC}\;}
Z^1(\operatorname{Def}_{\mathrm{cyc}}(\cA)_{\mathrm{model}})
\;\xrightarrow{\;\text{global}\;}
R\Gamma(\overline{\mathcal{M}}_g, \mathcal{Z}_{\cA})
\;\xrightarrow{\;\mathrm{tr}\;}
R\Gamma(\overline{\mathcal{M}}_g, \mathbb{Q}).
\]
Fiber period data live in the variation of $H^1(\Sigma_g,\C)$, not
as cohomology classes on moduli
($H^1(\mathcal{M}_g,\mathbb{Q})=0$ for $g\geq 2$ by Harer).
They enter moduli through the Kodaira--Spencer map
(Theorem~\ref{thm:kodaira-spencer-chiral-complete}).
\end{remark}

\begin{proposition}[Fiber curvature and its Hodge realization;
\ClaimStatusConditional]
\label{prop:genus-g-curvature-package}
\textup{[Type signature: Open quadrant; curved bar presentation;
levels $2$ and $5$; prime-form package $H_D^{\mathrm{fib}}$ and
uncurving package $H_D^{\mathrm{unc}}$, scalar Hodge package $H_D^K$,
and graph package $H_D^{\mathrm{graph}}$.]}
Let~$\cA$ be cyclic and let~$\Sigma_g$ be a smooth genus-$g$ curve.
Choose normalized holomorphic differentials, period matrix~$\Omega$,
and prime form~$E(z,w)$.  The package $H_D^{\mathrm{fib}}$ supplies an
Arakelov-normalized propagator
\begin{equation}
 \eta_{ij}^{(g)}
 =\partial_{z_i}\log E(z_i,z_j)
 +\pi\sum_{\alpha,\beta=1}^g
 \omega_\alpha(z_i)(\operatorname{Im}\Omega)^{-1}_{\alpha\beta}
 \operatorname{Im}\!\left(\int_{z_0}^{z_j}\omega_\beta\right),
 \label{eq:propagator-package}
\end{equation}
and a curved unary operation satisfying
\begin{equation}
 m_1^{(g)\,2}=[m_0^{(g)},-],
 \qquad
 \operatorname{tr}_{\mathrm{diag}}(m_0^{(g)})
 =\kappa(\cA)\omega_g^{\mathrm{Ar}},
 \label{eq:curvature-package}
\end{equation}
where
\[
 \omega_g^{\mathrm{Ar}}
 =\frac{i}{2}\sum_{\alpha,\beta=1}^g
 (\operatorname{Im}\Omega)^{-1}_{\alpha\beta}
 \omega_\alpha\wedge\overline{\omega_\beta}
\]
is a vertical $(1,1)$-form on~$\Sigma_g$.  A period and
Gauss--Manin correction gives a strict operator
\[
 \Dg{g}=\dfib+d_{\mathrm{per}}^{(g)}
 +\nabla_{\mathrm{KS}}^{\mathrm{GM}},
 \qquad \Dg{g}^{\,2}=0.
\]
Under~$H_D^K(\cA,g)$, the scalar virtual object and its Hodge shadow
are
\begin{equation}
 \mathfrak O_g^K(\cA)=\kappa(\cA)\lambda_{-1}(\mathbb E_g),
 \qquad
 \operatorname{obs}^{\mathrm{Hdg}}_g(\cA)
 =(-1)^g\kappa(\cA)\lambda_g.
 \label{eq:obstruction-hodge}
\end{equation}
The associated scalar free-energy series, under the graph package, is
\begin{equation}
 \sum_{g\ge1}F_g^{\mathrm{sc}}(\cA)\hbar^{2g}
 =\kappa(\cA)\left(\frac{\hbar/2}{\sin(\hbar/2)}-1\right).
 \label{eq:family-index-package}
\end{equation}
\end{proposition}

\begin{proof}
Equation~\eqref{eq:curvature-package} is the curved
$A_\infty$/CDG identity furnished by~$H_D^{\mathrm{fib}}$; its
scalar trace is the normalized loop contraction.  The period and
Gauss--Manin terms form the square-zero comparison in
$H_D^{\mathrm{unc}}$.  Equation~\eqref{eq:obstruction-hodge} follows from
Proposition~\ref{prop:scalar-obstruction-hodge-euler}.  Pairing
$\lambda_g$ with~$\psi_1^{2g-2}$ and applying
Faber--Pandharipande~\cite{FP03} gives
\eqref{eq:family-index-package}.
\end{proof}


\begin{proposition}[Typed curvature comparison;
\ClaimStatusConditional]
\label{prop:chain-level-curvature-operator}
\textup{[Type signature: Open quadrant; fixed-fiber curved bar
presentation and modular Hodge presentation; levels $2\to5$;
$H_D^{\mathrm{fib}}$ and $H_D^K$.]}
The fixed-fiber bar object carries
\begin{equation}
 m_1^{(g)\,2}(a)
 =m_2(m_0^{(g)},a)-m_2(a,m_0^{(g)})
 =[m_0^{(g)},a]_{m_2}.
 \label{eq:chain-curvature-operator}
\end{equation}
Its scalar diagonal is the vertical form
$\kappa(\cA)\omega_g^{\mathrm{Ar}}$ on~$\Sigma_g$.  The Hodge
bundle carries a horizontal curvature matrix
$\Theta_{\mathbb E_g}$ on~$\cM_g$, with
\begin{equation}
 \det\!\left(\frac{i}{2\pi}\Theta_{\mathbb E_g}\right)
 =c_g(\mathbb E_g)=\lambda_g.
 \label{eq:chern-weil-projection}
\end{equation}
The comparison~$\beta_g(\cA)$ in~$H_D^K$ relates the scalar virtual
bar object to~$\kappa(\cA)\lambda_{-1}(\mathbb E_g)$; consequently
\[
 \operatorname{obs}^{\mathrm{Hdg}}_g(\cA)
 =(-1)^g\kappa(\cA)\lambda_g.
\]
\end{proposition}

\begin{proof}
Equation~\eqref{eq:chain-curvature-operator} is the defining curved
$A_\infty$ relation.  Equation~\eqref{eq:chern-weil-projection} is
Chern--Weil theory for the rank-$g$ Hodge bundle.  Applying the
degree-$g$ Chern character to~$\beta_g(\cA)$ gives the final formula.
\end{proof}


\begin{proposition}[Genus-$g$ MC element; \ClaimStatusConditional]
\label{prop:genus-g-mc-element}
\index{Maurer--Cartan!genus-g component@genus-$g$ component}
\index{master MC element!genus decomposition}

For each $g \geq 0$, the genus-$g$ component of the master MC
element satisfies:
\[
D^{(g)}_\cA \;\Theta_\cA^{(g)}
+ \sum_{g_1+g_2=g}
\tfrac{1}{2}[\Theta_\cA^{(g_1)}, \Theta_\cA^{(g_2)}]
\;=\; 0,
\]
where $D^{(g)}_\cA$ is the genus-$g$ fiber covariant component. Its
scalar curvature is the degree-$2$ endomorphism
\[
\bigl(D^{(g)}_\cA\bigr)^2_{\mathrm{scalar}}
=\kappa(\cA)\,\omega_g^{\mathrm{Ar}}\cdot\mathrm{id};
\]
this curvature form is not an added degree-$1$ summand of the
differential. At $g = 0$, the curvature term vanishes and one recovers
the tree-level MC equation. At $g = 1$, the class
$\kappa\,\omega_1^{\mathrm{Ar}}$ is the elliptic Arakelov defect; the
strict square-zero differential is the total period-corrected
$\Dg{g}$ of Convention~\ref{conv:higher-genus-differentials}.
\end{proposition}

\begin{proof}
The bar-intrinsic construction
(Theorem~\ref{thm:mc2-bar-intrinsic}) defines
$\Theta_\cA := D_\cA - d_0 \in \mathrm{MC}(\gAmod)$,
satisfying $D_\cA\Theta_\cA +
\tfrac{1}{2}[\Theta_\cA,\Theta_\cA] = 0$ because $D_\cA^2 = 0$.
The genus decomposition $\gAmod = \prod_{g \geq 0}
\mathfrak{g}^{(g)}_\cA$ projects this single equation to genus~$g$,
with the vertical curvature $\kappa(\cA)\cdot\omega_g^{\mathrm{Ar}}$ entering
through the corrected differential $D^{(g)}_\cA$
(Proposition~\ref{prop:genus-g-curvature-package}\textup{(ii)}).
\end{proof}

\begin{theorem}[Genus extension hierarchy; \ClaimStatusConditional]
\label{thm:genus-extension-hierarchy}
\index{genus extension hierarchy|textbf}
\index{bar complex!genus extension levels}
\index{Koszul locus!genus extension}
\index{modular functor!genus extension level}
The genus extension of the bar complex admits a five-level
hierarchy of successively stronger structural conditions.
The first four levels form the algebraic-to-analytic spine; the
fifth is a modular-functor condition and is not comparable with the
Koszul level by containment. Universal standard families satisfy the
finite-window hypotheses through the analytic sewing level, while
only their $C_2$-cofinite simple quotients or rational specialisations
reach the modular-functor level.
\begin{enumerate}[label=\textup{(\roman*)}]
\item \emph{Algebraic \textup{(}bar $D^2=0$ in each genus\textup{)}.}\;
 For \emph{any} chiral algebra $\cA$, the bar differential satisfies
 $D_\cA^2 = 0$ in each genus, and the universal MC element
 $\Theta_\cA := D_\cA - d_0 \in \MC(\gAmod)$ exists
 unconditionally
 \textup{(}Theorem~\textup{\ref{thm:mc2-bar-intrinsic})}.
 No finiteness or Koszulness condition is required: $D^2 = 0$
 is a formal consequence of $\partial^2 = 0$ on
 $\overline{\cM}_{g,n}$.

\item \emph{Homotopical \textup{(}bar-cobar comparison in
 $D^{\mathrm{co}}$\textup{)}.}\;
 If $\cA$ has finite-dimensional weight spaces
 \textup{(}the positive-energy axiom\textup{)} and its finite-window
 bar tower is a strict Mittag--Leffler completed chiral coalgebra,
 then the
 genus-$g$ bar-cobar adjunction
 $\Omegach_g \dashv \barBch_g$ extends to an equivalence
 in the coderived category $D^{\mathrm{co}}$
 \textup{(}Theorem~\textup{\ref{thm:chiral-positselski-weight-completed}}
 and Appendix~\textup{\ref{app:coderived})}.
 The positive-energy condition supplies the finite windows; strict
 Mittag--Leffler control supplies the derived inverse-limit recovery
 and product/sum exactness that preserve coacyclic and contraacyclic
 objects. The weight-by-weight SDR then produces the
 transferred curved $A_\infty$ structure on the coderived
 bar object; on the strict flat comparison model
 $(\barBch(\cA), \Dg{g})$ this appears as an ordinary
 $A_\infty$ structure.

\item \emph{Koszul \textup{(}chiral bar-cobar
 quasi-isomorphism\textup{)}.}\;
 If $\cA$ satisfies MK1 \textup{(}genus-$0$
 Koszulness\textup{)}, equivalently the PBW spectral sequence
 collapses at $E_2$
 \textup{(}Theorem~\textup{\ref{thm:koszul-equivalences-meta})},
 then the counit
 $\Omegach(\barBch(\cA)) \to \cA$ is a quasi-isomorphism
 in the ordinary derived category
 \textup{(}Theorem~\textup{\ref{thm:higher-genus-inversion})}.
 On the Koszul locus, the bar and cobar functors are
 inverse equivalences: the genus-$g$ bar complex is a
 chain-level resolution of~$\cA$ in the ordinary derived category.

\item \emph{Analytic \textup{(}convergent partition
 functions\textup{)}.}\;
 If $\cA$ satisfies the HS-sewing criterion (polynomial
 OPE growth and subexponential sector growth)
 then the genus-$g$ partition functions
 $Z_g(\cA)$ converge absolutely for all $g \geq 1$
 \textup{(}Theorem~\textup{\ref{thm:general-hs-sewing})}.
 The sewing envelope $\cA^{\mathrm{sew}}$ carries the
 completed analytic bar coalgebra, and the shadow
 obstruction tower evaluates to convergent numerical
 invariants.

\item \emph{Modular functor \textup{(}finite-dimensional
 conformal blocks\textup{)}.}\;
 If $\cA$ is $C_2$-cofinite in the sense of
 Zhu~\cite{Zhu96}, then the spaces of conformal blocks
 $\mathcal{V}_g(\cA)$ are finite-dimensional for all $g$
 (Zhu genus-$1$~\cite{Zhu96} combined with
 Tsuchiya--Ueno--Yamada~\cite{TUY89},
 Nagatomo--Tsuchiya~\cite{NagatomoTsuchiya05},
 Damiolini--Gibney--Tarasca~\cite{DGT21}),
 the modular group $\mathrm{SL}_2(\mathbb{Z})$ acts on
 characters, and $\cA$ defines a modular functor.
 This is the strongest condition; it implies all four
 preceding levels.
\end{enumerate}

The first four levels form a strict chain:
$\textup{(iv)} \subsetneq
\textup{(iii)} \subsetneq \textup{(ii)} \subsetneq
\textup{(i)}$.
Level~\textup{(v)} \textup{(}$C_2$-cofinite\textup{)} implies
\textup{(iv)}, \textup{(ii)}, and \textup{(i)}, but its
relation to \textup{(iii)} \textup{(}Koszul\textup{)} is not
one of containment: $C_2$-cofinite simple quotients
$L_k(\fg)$ at admissible level satisfy~\textup{(v)} but may
not satisfy~\textup{(iii)}, while Koszul universal algebras
$V^k(\fg)$ at generic level satisfy~\textup{(iii)} but
not~\textup{(v)}.
Every chiral algebra satisfies~\textup{(i)}.
Every positive-energy chiral algebra whose finite-window tower is
strict Mittag--Leffler satisfies \textup{(i)}--\textup{(ii)}.
The standard landscape \textup{(}Heisenberg, affine
Kac--Moody at non-critical level, $\cW$-algebras at
generic central charge, $\beta\gamma$ and $bc$
systems\textup{)} satisfies the strict finite-window tower hypothesis
on the completed surface and hence levels
\textup{(i)}--\textup{(iv)}; those that are also
$C_2$-cofinite satisfy all five.
\end{theorem}

\begin{proof}
Level~(i) is Theorem~\ref{thm:mc2-bar-intrinsic}: $D_\cA^2=0$ holds
for the modular-operad bar construction because the bar differential is the Feynman-transform
differential of a modular operad algebra, and $\partial^2=0$ on
$\overline{\cM}_{g,n}$ implies $D^2=0$ by the codimension-$2$
boundary cancellation.

Level~(ii): the positive-energy axiom gives
$\cA = \bigoplus_{n \geq 0} \cA_n$ with $\dim \cA_n < \infty$.
The strict Mittag--Leffler hypothesis on the finite-window tower gives
the continuous completed curved chiral coalgebra and the
$R\!\varprojlim$ recovery of coderived and contraderived categories by
Theorem~\ref{thm:chiral-positselski-weight-completed}. The
weight-by-weight SDR
(Cliff~\cite[Proposition~8.1]{Cliff19})
applies at each genus to produce the transferred structure
in $D^{\mathrm{co}}$; finite-dimensional weight spaces give the
finite stages, while the strict tower supplies the missing completed
homotopy-theoretic control. Koszulness is not used at this level.

Level~(iii): MK1 (PBW collapse at $E_2$) implies diagonal Ext
vanishing, which is the hypothesis of
Theorem~\ref{thm:higher-genus-inversion}. On the Koszul locus,
the curved bar complex $(\barBch(\cA), \dfib)$ and the flat model
$(\barBch(\cA), \Dg{g})$ are linked by the coderived-to-derived
comparison
(Proposition~\ref{prop:gauss-manin-uncurving-chain}),
and the counit is a quasi-isomorphism in the ordinary
derived category. The passage from coderived to derived is
the content of genus-$0$ Koszulness.

Level~(iv): the HS-sewing criterion
(Theorem~\ref{thm:general-hs-sewing}) is a growth condition on
OPE coefficients, strictly weaker than $C_2$-cofiniteness.
The standard families satisfy it because they have polynomial
OPE growth (OPE coefficients are polynomial in conformal
weights) and subexponential sector growth (the graded
dimension $\dim \cA_n$ grows at most polynomially).

Level~(v): $C_2$-cofiniteness implies, by Zhu's genus-$1$
modular-invariance theorem~\cite{Zhu96} combined with the
higher-genus extension through
Tsuchiya--Ueno--Yamada~\cite{TUY89} nodal factorisation and
the Nagatomo--Tsuchiya~\cite{NagatomoTsuchiya05} /
Huang~\cite{Huang05} / Damiolini--Gibney--Tarasca~\cite{DGT21}
bundle construction, that the space of genus-$g$ conformal
blocks is finite-dimensional for all $g$ and carries a
projectively flat connection. This implies convergent partition functions
(hence level~(iv)), finite-dimensional weight spaces
(hence level~(ii)), and $D^2=0$ (which is universal).
Whether $C_2$-cofiniteness implies Koszulness
\textup{(}level~\textup{(iii))} is a separate question:
for simple quotients $L_k(\fg)$ at admissible level,
$C_2$-cofiniteness holds but Koszulness is open for
rank~$\geq 2$
\textup{(}Remark~\textup{\ref{rem:admissible-koszul-status})}.
The inclusion \textup{(v)} $\subsetneq$ \textup{(iv)} is
strict; the relation of \textup{(v)} to \textup{(iii)} is
not one of containment.

The inclusions are strict. For (i)$\neq$(ii): a chiral
algebra with infinite-dimensional weight spaces (e.g.\ a
vertex algebra without the positive-energy axiom) satisfies
$D^2=0$ but lacks the SDR needed for homotopy transfer.
For (ii)$\neq$(iii): a positive-energy chiral algebra that is
not chirally Koszul (e.g.\ simple quotients at certain
admissible levels where bar-Ext$\neq$ordinary-Ext) has the
coderived model but not the derived quasi-isomorphism.
For (iii)$\neq$(iv): the HS-sewing criterion
(Theorem~\ref{thm:general-hs-sewing}) requires polynomial
OPE growth, which is a growth-rate condition independent of
Koszulness; a chirally Koszul algebra with super-polynomial
OPE growth satisfies~(iii) but not~(iv).
For (iv)$\neq$(v): the universal $\cW$-algebra $V^k(\fg)$
at generic irrational level satisfies HS-sewing (convergent
partition functions) but is not $C_2$-cofinite (conformal
blocks are infinite-dimensional).
\end{proof}

\section{The bar complex as a family over moduli}
\label{sec:geometric-prologue}
\label{subsec:bar-family-moduli}
\index{bar complex!family over moduli|textbf}

The bar object at genus~$g$ is a family of curved fiber complexes over
the smooth locus, extended to the boundary through the relative
Fulton--MacPherson compactification:
\begin{equation}\label{eq:bar-family}
\barB_g(\cA) \;\longrightarrow\; \overline{\mathcal{M}}_g,
\qquad
\barB_g(\cA)\big|_{[\Sigma_g]}
= (T^c(s^{-1}\bar{\cA}),\, \dfib^{(\Sigma_g)}).
\end{equation}
The fiber over $[\Sigma_g]$ has fiberwise differential
$\dfib^{(\Sigma_g)}$ built from the genus-$g$ propagator
(Definition~\ref{def:higher-genus-log-forms}); for
$\kappa(\cA)\neq0$ this is a curved differential. The strict
cochain complex is obtained only after passing to the total corrected
differential $\Dg{g}$.

\begin{remark}[Stable-graph description]
\label{rem:bar-stable-graph-sum}
\index{bar complex!stable-graph sum}
\index{Feynman transform!bar complex as}
The genus-$g$ bar complex is the stable-graph sum of
Theorem~\ref{thm:prism-higher-genus}:
\begin{equation}\label{eq:bar-graph-sum}
\barB^{(g)}(\cA)
\;=\;
\bigoplus_{\Gamma \in \mathcal{G}^{\mathrm{st}}_{g,n}}
\biggl(\bigotimes_{v \in V(\Gamma)}
\cA(g_v, \operatorname{val}(v))\biggr)
\otimes \operatorname{or}(E(\Gamma)),
\end{equation}
with $\sum_v g_v + b_1(\Gamma) = g$. This is the Feynman transform
$\mathrm{FT}_{\mathcal{M}od}(\cA)$ of
Definition~\ref{def:feynman-transform}.
\end{remark}

\subsection{The Hodge bundle and the period matrix}
\label{subsec:hodge-curvature-origin}
\index{Hodge bundle!and curvature}

The Hodge bundle $\mathbb{E} = \pi_*\omega_\pi$ over
$\overline{\mathcal{M}}_g$ has fiber
$H^0(\Sigma_g, \Omega^1)$, a $g$-dimensional vector space.
Choose a symplectic basis
$\{A_k, B_k\}_{k=1}^g \subset H_1(\Sigma_g, \mathbb{Z})$ and
normalize: $\oint_{A_l} \omega_k = \delta_{kl}$. The period matrix
$\tau_{kl} = \oint_{B_l} \omega_k \in \mathfrak{H}_g$ determines
the propagator~\eqref{eq:elliptic-propagator} and its higher-genus
generalizations.

With the Arakelov representative fixed, the fiberwise bar operation
is curved; the scalar diagonal of its curvature is the Arnold defect
contracted with the OPE data of~$\cA$
\textup{(}Theorem~\textup{\ref{thm:quantum-arnold-relations},}
Proposition~\textup{\ref{prop:genus-g-curvature-package}(ii))}:
\begin{equation}\label{eq:arnold-defect-hg}
\operatorname{tr}_{\mathrm{diag}}(m_0^{(g)})
= \kappa(\cA) \cdot \omega_g^{\mathrm{Ar}}, \qquad
\omega_g^{\mathrm{Ar}} = \tfrac{i}{2} \textstyle\sum_{\alpha,\beta=1}^g
(\operatorname{Im}\Omega)^{-1}_{\alpha\beta}\,
\omega_\alpha \wedge \overline{\omega}_\beta.
\end{equation}
The Arakelov form~$\omega_g^{\mathrm{Ar}}$ is vertical on~$\Sigma_g$.  The Hodge
bundle $\mathbb E_g=\pi_*\omega_\pi$ is horizontal on~$\cM_g$ and
has Chern curvature~$\Theta_{\mathbb E_g}$.  Chern--Weil theory gives
\[
 \det\!\Bigl(\frac{i}{2\pi}\Theta_{\mathbb E_g}\Bigr)
 =c_g(\mathbb E_g)=\lambda_g.
\]
The package~$H_D^K(\cA,g)$ relates the scalar virtual bar object to
the alternating exterior power of this bundle.  Hence
\[
 \mathfrak O_g^K(\cA)=\kappa(\cA)\lambda_{-1}(\mathbb E_g),
 \qquad
 \operatorname{obs}^{\mathrm{Hdg}}_g(\cA)
 =(-1)^g\kappa(\cA)\lambda_g.
\]
At genus~$1$, the deformation trace has the separate normalization
$\operatorname{tr}_1\operatorname{Obs}^{\mathrm{def}}_{1,1}
=\kappa(\cA)\lambda_1$.

\subsection{The Kodaira--Spencer map}
\label{subsec:ks-deformation-functor}
\index{Kodaira--Spencer map|textbf}

The infinitesimal version of~\eqref{eq:bar-family} is
\begin{equation}\label{eq:ks-intro-hg}
\mathrm{KS} \colon T_{[\Sigma_g]}\overline{\mathcal{M}}_g
\;\xrightarrow{\;\sim\;}
H^1(\Sigma_g, T\Sigma_g)
\;\longrightarrow\;
\mathrm{HH}^2_{\mathrm{ch}}(\cA),
\end{equation}
sending an infinitesimal deformation of the curve to a deformation of
the bar complex
(Theorem~\ref{thm:kodaira-spencer-chiral-complete}). The obstruction
to extending a first-order deformation lies in
$\mathrm{HH}^3_{\mathrm{ch}}(\cA)$.

\subsection{The family index}
\label{subsec:why-ahat}
\index{A-hat genus@$\hat{A}$-genus!geometric derivation}

Whenever the family~\eqref{eq:bar-family} admits the perfect
logarithmic extension specified by~$H_D^K$, it defines a virtual class
$[\barB_g(\cA)]^{\mathrm{vir}}
\in K^0(\overline{\mathcal{M}}_g)\otimes\mathbb C$.
On the smooth universal curve, Grothendieck--Riemann--Roch gives
\[
\mathrm{ch}\bigl([\barB_g(\cA)]^{\mathrm{vir}}\bigr)
= \pi_*\!\bigl(\mathrm{ch}(\barB(\cA)|_{\mathrm{fiber}})
\cdot \mathrm{Td}(T_\pi)\bigr).
\]
For the weight-$1$ zero-mode bundle, the degree-$1$ term is
$\operatorname{ch}_1(R\pi_*\omega_\pi) = c_1(\mathbb{E}) = \lambda_1$,
and the genus-one trace package~$H_D^1$ gives
$\kappa(\cA)\lambda_1$.
For $g\geq2$, the comparison~$H_D^K$ supplies the stronger virtual
object
\[
 \mathfrak O_g^K(\cA)=\kappa(\cA)\lambda_{-1}(\mathbb E_g).
\]
The splitting principle then gives
\[
 \operatorname{ch}_j(\mathfrak O_g^K)=0\quad(j<g),
 \qquad
 \operatorname{ch}_g(\mathfrak O_g^K)
 =(-1)^g\kappa(\cA)\lambda_g.
\]
Thus degree-$1$ GRR governs the determinant/Hodge line, while the
rank-$g$ Euler class enters through the separate $K$-theory
comparison.

Under~$H_D^{\mathrm{graph}}$, pairing the top Hodge class with the
Faber--Pandharipande test class produces the numerical scalar series
\[
\sum_{g \geq 1} F_g^{\mathrm{sc}}(\cA)\,\hbar^{2g}
\;=\;
\kappa(\cA) \cdot (\hat{A}(i\hbar) - 1)
\]
\textup{(}Theorem~\textup{\ref{thm:family-index}}\textup{)}.
The signed Hodge-shadow pairings instead form
$\kappa(\cA)(\hat A(\hbar)-1)$.

\section{\texorpdfstring{$A_\infty$ structures and higher operations}{A-infinity structures and higher operations}}

For historical background on $A_\infty$ algebras, see Appendix~\ref{sec:ainfty-historical}.

% ================================================================
% SECTION 4.2: THE GEOMETRIC BAR COMPLEX AND ITS A-INFINITY STRUCTURE
% ================================================================

\subsection{\texorpdfstring{The geometric bar complex and its $A_\infty$ structure}{The geometric bar complex and its A-infinity structure}}

\subsubsection{Elementary introduction: logarithmic forms as operations}

\begin{example}[Binary operation from residues]
For two operators $a, b$ in a chiral algebra at positions $z_1, z_2 \in \mathbb{P}^1$:
\begin{itemize}
\item The logarithmic 1-form: $\eta_{12} = d\log(z_1 - z_2) = \frac{dz_1 - dz_2}{z_1 - z_2}$
\item This has a simple pole when $z_1 = z_2$
\item The residue extracts the product:
\[m_2(a \otimes b) = \text{Res}_{z_1=z_2}\left[\eta_{12} \cdot a(z_1) \otimes b(z_2)\right] = \mu(a,b)\]
\end{itemize}
\end{example}

\begin{example}[Ternary operation and associativity]
For three operators at $z_1, z_2, z_3$:
\begin{itemize}
\item The 2-form: $\eta_{12} \wedge \eta_{23} = d\log(z_1-z_2) \wedge d\log(z_2-z_3)$
\item Has poles along three divisors:
 \begin{itemize}
 \item[$D_{12}$:] where $z_1 = z_2$ first
 \item[$D_{23}$:] where $z_2 = z_3$ first
 \item[$D_{123}$:] where all three collide
 \end{itemize}
\item The residues give:
\[\text{Res}_{D_{12}}[\eta_{12} \wedge \eta_{23}] = m_2(m_2(a,b),c)\]
\[\text{Res}_{D_{23}}[\eta_{12} \wedge \eta_{23}] = m_2(a,m_2(b,c))\]
\[\text{Res}_{D_{123}}[\eta_{12} \wedge \eta_{23}] = m_3(a,b,c)\]
\item The difference of boundary residues equals an exact form:
\[m_2(m_2 \otimes \text{id}) - m_2(\text{id} \otimes m_2) = d(h_3)\]
where $h_3$ is the homotopy between associations
\end{itemize}
\end{example}

\subsubsection{\texorpdfstring{Complete $A_\infty$ structure from configuration spaces}{Complete A-infinity structure from configuration spaces}}

\begin{definition}[\texorpdfstring{$A_\infty$}{A-infinity} algebra]\label{def:a-infinity-complete}
An $A_\infty$ algebra consists of a graded vector space $A$ with operations $m_k: A^{\otimes k} \to A[2-k]$ for $k \geq 1$ satisfying (cf.\ Appendix~\ref{app:curved-ainfty-formulas}):
\[\sum_{\substack{r+s+t=n \\ r,t \geq 0,\; s \geq 1}} (-1)^{rs+t}\, m_{r+1+t}(\mathrm{id}^{\otimes r} \otimes m_s \otimes \mathrm{id}^{\otimes t}) = 0\]

Explicitly for small $k$:
\begin{align}
k=1: &\quad m_1 \circ m_1 = 0 \quad \text{($m_1$ is a differential)} \\
k=2: &\quad m_1(m_2) = m_2(m_1 \otimes 1) + m_2(1 \otimes m_1) \quad \text{(Leibniz rule)} \\
k=3: &\quad m_2(m_2 \otimes 1) - m_2(1 \otimes m_2) = m_1(m_3) + m_3(m_1 \otimes 1 \otimes 1) + \cdots
\end{align}
\end{definition}

\begin{theorem}[\texorpdfstring{$A_\infty$}{A-infinity} structure from bar complex; \ClaimStatusProvedHere]\label{thm:bar-ainfty-complete}
The geometric bar complex $\bar{B}^{\text{geom}}(\mathcal{A})$ carries a natural $A_\infty$ structure where:

\begin{enumerate}
\item \emph{Operations from residues.} Each $m_k$ is given by a sum over planar binary trees:
\begin{equation}\label{eq:mk-tree-sum}
m_k(a_1 \otimes \cdots \otimes a_k) = \sum_{T \in \mathcal{T}_k}
(-1)^{\epsilon(T;\,a_1,\ldots,a_k)}\,
\operatorname{Res}_{D_T}\!\left[\bigwedge_{e \in E(T)} \eta_e
\cdot a_1(z_1) \otimes \cdots \otimes a_k(z_k)\right]
\end{equation}
where $\mathcal{T}_k$ denotes the set of planar binary trees with
$k$ leaves, $E(T)$ is the edge set with $|E(T)| = k-1$, $D_T$ is
the boundary stratum of $\overline{C}_k(X)$ corresponding to the
collision pattern encoded by $T$, and the \emph{Koszul--tree sign}
$\epsilon(T; a_1, \ldots, a_k) \in \mathbb{Z}/2$ is:
\begin{equation}\label{eq:koszul-tree-sign}
\epsilon(T; a_1, \ldots, a_k)
= \sum_{\substack{i < j \\ i \in L(v),\, j \in R(v) \\
 \text{for some internal } v \in V(T)}}
(|a_i| - 1)(|a_j| - 1),
\end{equation}
the sum over all pairs $(i,j)$ that are separated by an internal
vertex~$v$ of~$T$ (with $i$ in the left subtree $L(v)$ and $j$ in
the right subtree $R(v)$). This is the Koszul sign incurred by
reordering the desuspended elements $s^{-1}a_i$ (of degree
$|a_i| - 1$) from their natural order $1, \ldots, k$ to the
tree-composition order dictated by~$T$. It matches the sign
convention of Loday--Vallette~\cite[\S10.3.3]{LV12} after the
shift from their homological grading ($|d| = -1$) to our
cohomological grading ($|d| = +1$): the shift $s \mapsto s^{-1}$
replaces $|a_i|$ by $|a_i| - 1$ in the Koszul rule, which absorbs
the sign difference between the two conventions
\textup{(}Theorem~\textup{\ref{thm:det-bar-cobar-signs}} and
Definition~\textup{\ref{def:suspension}}\textup{)}.

\item \emph{Low-degree operations:}
\begin{align}
m_1 &= 0 \quad \text{(no differential on the chiral algebra)} \\
m_2(a \otimes b) &= \mu(a,b) \quad \text{(the chiral product)} \\
m_3(a \otimes b \otimes c) &=
 (-1)^{(|a|-1)(|b|-1)}
 \Res_{D_{(ab)c}}\bigl[\eta_{12} \wedge \eta_{23}
 \cdot a \otimes b \otimes c\bigr]
 \notag \\ &\quad
 + (-1)^{(|a|-1)(|c|-1)}
 \Res_{D_{a(bc)}}\bigl[\eta_{23} \wedge \eta_{13}
 \cdot a \otimes b \otimes c\bigr]
 \notag \\ &\quad
 + (-1)^{(|b|-1)(|c|-1)}
 \Res_{D_{(ac)b}}\bigl[\eta_{13} \wedge \eta_{12}
 \cdot a \otimes b \otimes c\bigr] \\
m_4(a \otimes b \otimes c \otimes d) &= \sum_{T \in \mathcal{T}_4}
 (-1)^{\epsilon(T)}\,\Res_{D_T}
 \bigl[\eta_T \cdot a \otimes b \otimes c \otimes d\bigr]
 \quad \text{($C_3 = 5$ planar binary trees)}
\end{align}

\item \emph{Coherences from geometry.} The $A_\infty$ relations follow from $\partial^2 = 0$ on the compactified configuration space $\overline{C}_n(X)$.

\item \emph{Homotopies.} Higher operations encode homotopies between different associations, with explicit formulas via angular forms on configuration spaces.
\end{enumerate}
\end{theorem}

\begin{proof}
We verify the $A_\infty$ relations through a systematic analysis of the boundary stratification.

\emph{Step 1: Decompose the bar differential by codimension.}
\[d = \sum_{k=2}^n \sum_{|I|=k} d_I\]
where $d_I$ extracts residues along the stratum where points indexed by $I$ collide.

\emph{Step 2: Analyze $d^2 = 0$.}
\[0 = d^2 = \sum_{I,J} d_I \circ d_J\]

Three cases arise:
\begin{enumerate}
\item \emph{Disjoint $I \cap J = \emptyset$:} Residues commute (up to Koszul sign)
\item \emph{Nested $I \subset J$ or $J \subset I$:} Boundary of boundary = 0
\item \emph{Overlapping $I \cap J \neq \emptyset$, neither contained:} Gives $A_\infty$ relation
\end{enumerate}

\emph{Step 3: Extract the $m_3$ operation explicitly.}

Near triple collision, use coordinates:
\[\epsilon_1 = z_1 - z_2, \quad \epsilon_2 = z_2 - z_3\]

The 2-form decomposes:
\[\eta_{12} \wedge \eta_{23} = d\log\epsilon_1 \wedge d\log\epsilon_2 + d\arg\left(\frac{\epsilon_1}{\epsilon_2}\right) \wedge d\log|\epsilon_1\epsilon_2|\]

The first term gives $m_3$, the second gives the homotopy $h_3$.
\end{proof}

\subsubsection{\texorpdfstring{Enhanced $A_\infty$ structure with moduli space interpretation}{Enhanced A-infinity structure with moduli space interpretation}}

\begin{remark}[\texorpdfstring{$A_\infty$}{A-infinity} vs.\ strictly associative]
If $\mathcal{A}$ is Koszul, then $\bar{B}^{\text{ch}}(\mathcal{A})$
has $m_k = 0$ for $k \geq 3$ (strictly associative;
Proposition~\ref{prop:e2-collapse-formality}).
For general chiral algebras, or at chain level before passing to cohomology, the full $A_\infty$ structure is needed.
\end{remark}

\begin{theorem}[\texorpdfstring{$A_\infty$}{A-infinity} operations via moduli spaces; \ClaimStatusProvedHere]\label{thm:ainfty-moduli}
The bar construction $\bar{B}^{\text{ch}}(\mathcal{A})$ carries operations
$m_k: (\bar{B}^{\text{ch}})^{\otimes k} \to \bar{B}^{\text{ch}}[2-k]$ defined
geometrically by integration over configuration space boundaries:

\[m_k(\omega_1, \ldots, \omega_k) = \int_{\partial \overline{M}_{0,k+1}}
\pi^*(\omega_1 \wedge \cdots \wedge \omega_k) \wedge \Omega_{0,k+1}\]

where:
\begin{itemize}
\item $\overline{M}_{0,k+1}$ is the Deligne--Mumford compactification of moduli
of stable rational curves with $k+1$ marked points
\item $\pi: \overline{M}_{0,k+1} \to (\overline{C}_2(X))^k$ is the natural
projection extracting the $k$ input configuration spaces
\item $\Omega_{0,k+1}$ is the fundamental class (canonical measure)
\item The boundary $\partial \overline{M}_{0,k+1}$ parametrizes all ways to
degenerate the curve
\end{itemize}
\end{theorem}

\begin{proof}
\emph{Step 1: Understanding $\overline{M}_{0,k+1}$}

The moduli space $\overline{M}_{0,k+1}$ parametrizes stable rational curves with
$k+1$ marked points. Its boundary stratification is:
\[\partial \overline{M}_{0,k+1} = \bigcup_{I \sqcup J = [k+1]}
\overline{M}_{0,|I|+1} \times \overline{M}_{0,|J|+1}\]

Each boundary component corresponds to a way of splitting the curve into two
components, with points distributed between them.

\emph{Step 2: The Operations}

For $k=2$ (binary product):
\[m_2(\omega_1, \omega_2) = \int_{\overline{C}_2(X)}
\text{Res}_{z_1=z_2}\left[\frac{\omega_1(z_1) \wedge \omega_2(z_2)}{z_1-z_2}\right]\]

This is the usual chiral algebra product via OPE.

For $k=3$ (associator):
\[m_3(\omega_1, \omega_2, \omega_3) =
\int_{\partial \overline{M}_{0,4}} \omega_1 \wedge \omega_2 \wedge \omega_3\]

The boundary $\partial \overline{M}_{0,4}$ has three components:
\begin{itemize}
\item $(12|34)$: Gives $m_2(m_2(\omega_1,\omega_2),\omega_3)$
\item $(13|24)$: Mixed terms
\item $(14|23)$: Gives $m_2(\omega_1,m_2(\omega_2,\omega_3))$
\end{itemize}

The $m_3$ operation exactly measures the failure of associativity:
\[m_2(m_2 \otimes \text{id}) - m_2(\text{id} \otimes m_2) = d m_3 + m_3 d\]

For $k \geq 4$: Higher coherences arise from more complex degenerations of moduli
spaces, encoding Stasheff polytopes.

\emph{Step 3: The $A_\infty$ Relations}

The fundamental $A_\infty$ relation is:
\[\sum_{\substack{r+s+t=k \\ r,t \geq 0,\, s \geq 1}} (-1)^{rs+t}\,
m_{r+1+t}(\mathrm{id}^{\otimes r} \otimes m_s \otimes \mathrm{id}^{\otimes t}) = 0\]

This follows from $\partial \partial \overline{M}_{0,k+1} = 0$: each codimension-2
stratum in the boundary appears twice with opposite signs, giving the cancellation.
\end{proof}

\begin{example}[Virasoro algebra: explicit \texorpdfstring{$m_3$}{m3}]
For the Virasoro algebra with stress tensor $T(z)$:
\[T(z_1)T(z_2) = \frac{c/2}{(z_1-z_2)^4} + \frac{2T(z_2)}{(z_1-z_2)^2} +
\frac{\partial T(z_2)}{z_1-z_2} + \text{reg}\]

The $m_3$ operation computes:
\[m_3(T \otimes T \otimes T) =
\int_{\partial \overline{M}_{0,4}} \text{Res}[\text{triple OPE}]\]

This involves a primary pole $\propto c^2$ from $(T \cdot T) \cdot T$ vs.\ $T \cdot (T \cdot T)$, Schwarzian derivative terms from the conformal anomaly, and descendant contributions from $\partial T$. The result is non-zero ($m_3 \neq 0$: the Virasoro algebra is $A_\infty$-non-formal, shadow depth~$\infty$), encoding the conformal anomaly
and central charge. This $m_3$ operation is the obstruction to
a strictly associative product on the bar construction, but it does \emph{not} obstruct chiral Koszulness (Theorem~\ref{thm:virasoro-chiral-koszul}).
\end{example}

\begin{remark}[Feynman diagram interpretation]
\label{rem:ainfty-feynman-preview}
The operation~$m_k$ corresponds to the tree-level Feynman diagrams
determined by the boundary stratification
of~$\overline{M}_{0,k+1}$: $m_2$ extracts the binary collision
residue, $m_3$ is the tree-level associator homotopy, and
higher~$m_k$ encode higher tree-level coherences (Stasheff
polytopes). All $A_\infty$ operations from the genus-$0$ bar
complex are tree-level; loop corrections arise only at
genus~$g \geq 1$ (Section~\ref{sec:genus-1-elliptic-complete}).
See Chapter~\ref{ch:v1-feynman}.
\end{remark}

\subsubsection{Pentagon and higher identities}

\begin{theorem}[Pentagon identity; \ClaimStatusProvedHere]\label{thm:pentagon-identity}
The $A_\infty$ relation at $n = 4$ has exactly five nonzero terms (for a minimal
$A_\infty$ algebra with $m_1 = 0$), corresponding to the five facets of the
associahedron $K_4$ (the Stasheff pentagon). Geometrically, $\overline{M}_{0,5}$ is 2-dimensional (the del Pezzo
surface of degree~5) and has $\binom{5}{2} = 10$ boundary divisors in total; the five terms of the $A_\infty$ relation correspond to the five that are relevant for the planar (non-$\Sigma$) operad structure:
\[\sum_{\text{boundary divisors}} \text{sign}(D_i) \cdot m_{D_i} = 0\]
The relation $\partial^2 = 0$ on this 2-dimensional space yields the pentagon identity.
\end{theorem}

\begin{proof}
The moduli space $\overline{M}_{0,5}$ is a smooth compact surface
(the del Pezzo surface of degree~5, isomorphic to
$\mathrm{Bl}_4(\mathbb{P}^2)$). Its boundary consists of $10$
divisors $D_{ij}$, indexed by $2$-element subsets $\{i,j\} \subset
\{1,\ldots,5\}$, each isomorphic to $\overline{M}_{0,4} \times
\overline{M}_{0,3} \cong \mathbb{P}^1$.

For the non-$\Sigma$ (planar) operad structure, we fix a cyclic ordering
of the five marked points. Of the $10$ boundary divisors, exactly $5$
respect this cyclic ordering (they correspond to partitions of
$\{1,2,3,4,5\}$ into two cyclically consecutive subsets of size
$\geq 2$): $D_{12}$, $D_{23}$, $D_{34}$, $D_{45}$, $D_{15}$.
These five divisors are the facets of the Stasheff associahedron~$K_4$
(the pentagon), and they correspond to the five binary tree
parenthesizations:
$((ab)c)d$, $(a(bc))d$, $a((bc)d)$, $a(b(cd))$, $(ab)(cd)$.

The $A_\infty$ relation arises from the boundary structure: each
planar boundary divisor $D$ contributes a term
$m_{|S|+1} \circ (\mathrm{id}^{\otimes r} \otimes m_{|T|+1}
\otimes \mathrm{id}^{\otimes s})$ (where $S \sqcup T$ is the
partition and $r + s = |S| - 1$). These five contributions sum
to zero because they form the oriented boundary of $\overline{M}_{0,5}$
restricted to the associahedron cell complex:
$\partial K_4 = \sum_{i=1}^5 \pm D_i$, and $\partial^2 = 0$ on
the cellular chain complex of $K_4$ forces this signed sum to vanish.
\end{proof}

\begin{theorem}[Higher associahedron identity for \texorpdfstring{$m_5$}{m5} {\cite{Sta63}}; \ClaimStatusProvedElsewhere]\label{thm:higher-associahedron-m5}
For six elements, the associahedron $K_6$ is 4-dimensional with:
\begin{itemize}
\item 42 vertices (ways to fully parenthesize 6 elements, $C_5 = 42$)
\item 84 edges (single reassociations)
\item 56 two-faces (pentagons $K_4$ and squares $K_3 \times K_3$)
\item 14 three-dimensional cells
\end{itemize}

The boundary $\partial K_6$ encodes relations among compositions involving $m_5$: each pentagon face corresponds to a local pentagon identity, while each square face encodes compatibility between independent reassociations.
\end{theorem}

\begin{theorem}[Catalan identity at higher levels {\cite{Sta97}}; \ClaimStatusProvedElsewhere]\label{thm:catalan-parenthesization}
The number of ways to fully parenthesize $n$ objects is the Catalan number:
\[C_{n-1} = \frac{1}{n}\binom{2n-2}{n-1}\]

Each corresponds to a codimension $(n-2)$ stratum of $\overline{C}_n(X)$. The relations among these strata encode the complete $A_\infty$ structure, with the number of independent relations growing as:
\[\text{Relations at level } n = C_n - C_{n-1} \cdot C_1 - C_{n-2} \cdot C_2 - \cdots\]
\end{theorem}

% ================================================================
% SECTION 4.3: THE GEOMETRIC COBAR COMPLEX AND VERDIER DUALITY
% ================================================================

\subsection{The geometric cobar complex and Verdier duality}

\subsubsection{Cobar as opposite orientation}

\begin{definition}[Cobar via orientation reversal]\label{def:cobar-orientation}
The geometric cobar construction is factorization homology with reversed orientation:
\[\Omega^{\text{geom}}(\mathcal{C}) = \int_{-C_*(X)} \mathcal{C}\]
where $-C_*(X)$ denotes configuration spaces with opposite orientation.
Concretely, the bar complex is represented by logarithmic
de~Rham forms on the Fulton--MacPherson compactification, while the
cobar complex is represented after Verdier duality by currents with
controlled singular support along the collision strata. The slogan
``logarithmic form versus delta distribution'' means the normalized
residue-current pairing on a boundary divisor; it is not an
unshifted equality of complexes of smooth forms. Orientation reversal
realizes non-abelian Poincar\'e duality in the form
$\int_M \mathbb{D}(A) \simeq \mathbb{D}(\int_{-M} A)$ after the
Verdier shift and dualizing complex have been included.
\end{definition}

\begin{theorem}[Verdier duality = NAP duality {\cite{AF15,KS90}}; \ClaimStatusProvedElsewhere]\label{thm:verdier-NAP}
Let $j\colon C_n(X)\hookrightarrow \overline C_n(X)$ be the open
ordered configuration locus inside the Fulton--MacPherson
compactification, and let $D=\overline C_n(X)\setminus C_n(X)$ be
the normal-crossings boundary. Verdier duality gives
\[
\mathbb{D}_{\overline C_n(X)}(j_!\mathbb{C}_{C_n(X)})
\;\simeq\;
Rj_*\omega_{C_n(X)}[2\dim_\C C_n(X)] .
\]
Under the de~Rham realization, logarithmic forms representing
$Rj_*\mathbb C_{C_n(X)}$ pair with relative currents representing
$\mathbb{D}(j_!\mathbb C_{C_n(X)})$ by the normalized residue
functional along the strata of~$D$. This Verdier pairing is the
configuration-space incarnation of non-abelian Poincar\'e duality.

The exchange between logarithmic forms (bar) and distributions (cobar) is the geometric implementation of:
\[\int_X \mathcal{A} \xleftrightarrow{\mathbb{D}} \int_{-X} \mathcal{A}^!\]
\end{theorem}

\begin{proof}
By Verdier duality for constructible sheaves on $\overline{C}_n(X)$ (see~\cite{KS90}):
\[\mathbb{D}(\mathcal{F}) =
\mathcal{RHom}(\mathcal{F}, \omega_{\overline{C}_n(X)}
[2\dim_\C \overline C_n(X)]) .\]
Applied to $j_!\mathbb{C}_{C_n(X)}$, this exchanges extension by
zero with the direct image of the dualizing complex on the open
configuration space (Theorem~\ref{thm:verdier-config}). The
logarithmic de~Rham complex on the compactification resolves
$Rj_*\mathbb C_{C_n(X)}$, while the dual object is represented by
currents with prescribed boundary singularities. For a collision
divisor $D_{ij}$ and a small positively oriented normal circle
$\gamma_{ij}$, the normalization is
\[
\left\langle d\log(z_i-z_j), [D_{ij}]\right\rangle_{\mathrm{Verdier}}
\;:=\;
\frac{1}{2\pi i}\int_{\gamma_{ij}}d\log(z_i-z_j)
\;=\;1 .
\]
Thus residue extraction and current insertion are adjoint after the
Verdier shift. Factorization homology over the oppositely oriented
configuration chains is the non-abelian Poincar\'e-dual form of this
adjunction.
\end{proof}

\subsubsection{Distributions vs. differential forms: the dual picture}

While the bar complex uses logarithmic differential forms on
compactified configuration spaces, the cobar complex is modeled by
Verdier-dual currents on the open configuration space with boundary
singular support prescribed by the collision strata. This
distributional model is a de~Rham representative of the intrinsic
Verdier-dual definition (Definition~\ref{def:geom-cobar-intrinsic});
see Theorem~\ref{thm:cobar-distributional-model} for the precise
identification.

\begin{definition}[Geometric cobar complex]\label{def:geom-cobar-precise}
For a conilpotent chiral coalgebra $\mathcal{C}$ (valued in holonomic $\mathcal{D}$-modules), the distributional model of the geometric cobar complex is:
\[\Omega^{\text{ch}}_{p,q}(\mathcal{C}) = \text{Hom}_{\mathcal{D}}\left(\mathcal{C}^{\otimes(p+1)}, \mathcal{D}_{C_{p+1}(X)} \otimes \Omega^q_{\text{dist}}\right)\]
where $C_{p+1}(X)$ is the open configuration space and
$\Omega^q_{\text{dist}}$ denotes distributional $q$-forms whose
Verdier singular support is controlled by the collision diagonals.
The differential inserts the boundary current dual to a collision
divisor rather than extracting its residue.
\end{definition}

\begin{example}[Boundary current vs.\ residue]
\emph{Bar operation.} Extract residue when points collide
\[m_2^{\text{bar}}(a \otimes b) = \text{Res}_{z_1=z_2}\left[\frac{a(z_1)b(z_2)}{z_1-z_2}dz_1\right]\]

\emph{Cobar operation.} Insert the Verdier-dual boundary current,
which in a local coordinate is represented by the delta-current
along the diagonal:
\[n_2^{\text{cobar}}(K) = K(z_1,z_2) \cdot [D_{12}].\]

The normalized residue-current pairing:
\[
\langle \eta_{12}, [D_{12}] \rangle_{\mathrm{Verdier}}
=
\frac{1}{2\pi i}
\int_{\gamma_{12}}\frac{d(z_1-z_2)}{z_1-z_2}
=1 .
\]

This is Verdier duality in de~Rham coordinates: residues and boundary
currents are perfect duals after the dualizing shift.
\end{example}

\subsubsection{\texorpdfstring{Complete $A_\infty$ structure on cobar}{Complete A-infinity structure on cobar}}

\begin{theorem}[Cobar \texorpdfstring{$A_\infty$}{A-infinity} structure; \ClaimStatusProvedHere]\label{thm:cobar-ainfty-complete}
The cobar complex carries a dual $A_\infty$ structure with operations:
\[n_k: \Omega^{\mathrm{ch}}(\mathcal{C})^{\otimes k} \to \Omega^{\mathrm{ch}}(\mathcal{C})[2-k]\]

\begin{enumerate}
\item \emph{Operations:}
\begin{align}
n_1 &= d_{\mathrm{cobar}} \quad \text{(inserting boundary currents)} \\
n_2(K_1 \otimes K_2) &= K_1 * K_2 \quad \text{(convolution product)} \\
n_3(K_1 \otimes K_2 \otimes K_3) &= \text{triple propagator insertion}
\end{align}

\item \emph{Geometric realization.} Each $n_k$ corresponds to inserting a $k$-point propagator:
\[n_k(K_1, \ldots, K_k) = \int_{\partial C_k(X)} K_1 \wedge \cdots \wedge K_k \wedge P_k\]
where $P_k$ is the Feynman propagator for $k$ particles.

\item \emph{Duality with bar.} Under the Verdier pairing,
$\langle m_k^{\mathrm{bar}}, n_k^{\mathrm{cobar}} \rangle = 1$.
\end{enumerate}
\end{theorem}

\begin{proof}
The proof constructs the operations $n_k$ from the coalgebra structure of~$\mathcal{C}$ and verifies they satisfy the $A_\infty$ relations.

\emph{Construction of $n_k$.}
Let $\mathcal{C}$ be a dg coalgebra with coproduct $\Delta: \mathcal{C} \to \mathcal{C} \otimes \mathcal{C}$ and cobar complex $\Omega^{\mathrm{ch}}(\mathcal{C}) = (T(s^{-1}\bar{\mathcal{C}}), d_\Omega)$. The iterated coproduct $\Delta^{(k)}: \mathcal{C} \to \mathcal{C}^{\otimes k}$ dualizes to a $k$-ary operation on $\Omega^{\mathrm{ch}}(\mathcal{C})$:
\begin{equation}\label{eq:nk-from-coproduct}
n_k(K_1, \ldots, K_k) = (\mathrm{id} \otimes \cdots \otimes \mathrm{id}) \circ \Delta^{(k),*}(K_1 \otimes \cdots \otimes K_k)
\end{equation}
with appropriate signs from the Koszul convention (the shift $s^{-1}$ contributes $(-1)^{rs+t}$ as in the standard $A_\infty$ sign rule).

\emph{Verification of $n_1$.}
The cobar differential $d_\Omega = n_1$ is determined by the coproduct~$\Delta$ of~$\mathcal{C}$:
\[
n_1(s^{-1}c) = -\sum (-1)^{|c'|}\, s^{-1}c' \otimes s^{-1}c''
\]
where $\bar{\Delta}(c) = \sum c' \otimes c''$ is the reduced
coproduct. Geometrically, this inserts the Verdier-dual current of
the collision divisor, splitting one configuration into two
sub-configurations. In local analytic coordinates this current is
written as a delta-current, but the invariant object is the
Verdier-dual current of the divisor.

\emph{Verification of $n_2$.}
The binary operation $n_2(K_1 \otimes K_2) = K_1 * K_2$ is the convolution product:
\[
(K_1 * K_2)(z_1, \ldots, z_{p+q}) = \sum_{\sigma \in \mathrm{Sh}(p,q)} \pm\, K_1(z_{\sigma(1)}, \ldots, z_{\sigma(p)}) \cdot K_2(z_{\sigma(p+1)}, \ldots, z_{\sigma(p+q)})
\]
where $\mathrm{Sh}(p,q)$ denotes $(p,q)$-shuffles. This makes $\Omega^{\mathrm{ch}}(\mathcal{C})$ into a dg algebra.

\emph{Verification of $n_k$ for $k \geq 3$.}
The higher operations arise from the non-coassociativity of~$\mathcal{C}$ (when $\mathcal{C}$ is an $A_\infty$-coalgebra rather than a strict coalgebra). Geometrically, $n_k$ corresponds to integrating $k$ distributional kernels over the boundary of the moduli space $\partial\overline{C}_k(X)$, with the Feynman propagator $P_k$ encoding the correlation between insertion points:
\[
n_k(K_1, \ldots, K_k) = \int_{\partial \overline{C}_k(X)} K_1(z_1) \cdots K_k(z_k) \cdot P_k(z_1, \ldots, z_k)
\]
where $P_k$ is the $k$-point propagator obtained by iterating the basic
two-point propagator $P_2(z_1, z_2) = [D_{12}]$ over all tree-level
Feynman graphs with $k$ external legs.

\emph{$A_\infty$ relations.}
Define the total codifferential $D = \sum_{k \geq 1} n_k$ as a
degree-$1$ coderivation of the tensor coalgebra
$T^c(s^{-1}\bar{\mathcal{C}})$. The condition $D^2 = 0$ is equivalent
to the Stasheff relations
\[
\sum_{r+s+t=n} (-1)^{rs+t}\, n_{r+1+t}(\mathrm{id}^{\otimes r}
\otimes n_s \otimes \mathrm{id}^{\otimes t}) = 0,
\]
since a coderivation
of a cofree coalgebra is determined by its projections to cogenerators,
and $D^2$ projects to the left-hand side of the Stasheff identity
(Loday--Vallette~\cite{LV12}, Proposition~9.2.8). The condition $D^2 = 0$
holds because $D$ is the cobar dual of the bar codifferential, which
squares to zero by construction
(Theorem~\ref{thm:bar-nilpotency-complete}). In the geometric
realization, this corresponds to Stokes' theorem on the boundary
strata of $\overline{C}_n(X)$.

\emph{Duality with bar.}
The Verdier pairing $\langle m_k^{\mathrm{bar}}, n_k^{\mathrm{cobar}} \rangle$ evaluates to~$1$ by construction after residue normalization: the bar operations~$m_k$ (residue extraction of logarithmic forms) and the cobar operations~$n_k$ (current insertion) are dual under the Verdier residue pairing on $\overline{C}_k(X)$, as established in Theorem~\ref{thm:verdier-duality-operations}.
\end{proof}

\begin{example}[Linear coalgebra]
For $\mathcal{C} = T^c_{\text{ch}}(V)$ where $V = \text{span}\{v\}$ with $|v| = h$:

\emph{Coalgebra structure:}
\[\Delta(v^n) = \sum_{k=0}^n \binom{n}{k} v^k \otimes v^{n-k}\]

\emph{Cobar complex:}
\[\Omega^{\text{ch}}(T^c_{\text{ch}}(V)) = \text{Free}_{\text{ch}}(s^{-1}v, s^{-1}v^2, s^{-1}v^3, \ldots)\]

\emph{Differential (explicit formulas).}
\begin{align}
d(s^{-1}v) &= 0 \\
d(s^{-1}v^2) &= -2(s^{-1}v)^2 \\
d(s^{-1}v^3) &= -3(s^{-1}v)(s^{-1}v^2) \\
d(s^{-1}v^n) &= -\sum_{k=1}^{n-1} \binom{n}{k}(s^{-1}v^k)(s^{-1}v^{n-k})
\end{align}

\emph{Geometric interpretation.} Elements are multipole expansions
\[K_n(z_1, \ldots, z_n; w) = \sum_{i_1, \ldots, i_n} \frac{c_{i_1\ldots i_n}}{(z_1 - w)^{i_1} \cdots (z_n - w)^{i_n}}\]
encoding how fields behave near insertion points in CFT.
\end{example}

% ================================================================
% ================================================================

\subsection{Bar-cobar exchange}

\subsubsection{Chain/cochain level precision}

\begin{theorem}[Chain-level vs.\ homology-level structure; \ClaimStatusProvedHere]
\label{thm:chain-vs-homology}
Let $(\mathcal{A}, \{m_k\})$ be an $A_\infty$ chiral algebra.
\begin{enumerate}[label=\textup{(\roman*)}]
\item \emph{Chain level.}
 The full $A_\infty$ structure $\{m_k\}_{k \geq 1}$ is preserved,
 all operations are computable via configuration space integrals,
 and homotopies have geometric meaning as forms on
 $\overline{C}_n(X)$.
\item \emph{Homology level.}
 The homotopy transfer theorem
 (Theorem~\ref{thm:htt}) produces a \emph{minimal}
 $A_\infty$ structure $\{m_k^H\}$ on $H^*(\mathcal{A})$ with
 $m_1^H = 0$. The product $m_2^H$ descends to homology
 (it is strictly associative if $\mathcal{A}$ is formal).
 However, the transferred higher operations
 $m_k^H$ for $k \geq 3$ are in general \emph{non-trivial}:
 they are the Massey products, encoding obstructions to formality.
\item \emph{What is lost.}
 The explicit chain-level homotopies $\{h_n\}$ (the contracting
	 data) are lost in the passage to homology. For a formal algebra,
	 $m_k^H = 0$ for $k \geq 3$; for a non-formal
	 algebra, the $m_k^H$ carry essential information that cannot be
 recovered from $m_2^H$ alone.
\end{enumerate}
\end{theorem}

\begin{proof}
\emph{Step~1: Construction of the SDR.}
The bar complex $\barBgeom(\mathcal{A})$ is a chain complex of
$\mathcal{D}_X$-modules, graded by bar degree~$n$ and conformal
weight~$h$. Since the ambient category is field-linear, the splitting lemma
(Lemma~\ref{lem:sdr-existence}) furnishes a strong deformation
retract:
\[
\begin{tikzcd}[column sep=large]
(\barBgeom(\mathcal{A}), d_{\mathrm{bar}})
 \arrow[r, shift left, "p"] &
(H^*(\barBgeom(\mathcal{A})), 0)
 \arrow[l, shift left, "\iota"]
\end{tikzcd}
\quad h\colon \barBgeom^n \to \barBgeom^{n-1},
\quad \mathrm{id} - \iota p = d_{\mathrm{bar}}\, h + h\, d_{\mathrm{bar}}.
\]
The SDR is constructed by choosing, at each conformal weight~$h$,
a complement to $\ker(d)$ in $\barBgeom_h$:
decompose $\barBgeom_h = H_h \oplus B_h \oplus C_h$ where
$d\colon C_h \xrightarrow{\sim} B_h$ is an isomorphism
(boundaries and ``extra cycles'').
Then:
\begin{itemize}
\item $\iota\colon H_h \hookrightarrow \barBgeom_h$ includes the
 harmonic representatives;
\item $p\colon \barBgeom_h \twoheadrightarrow H_h$ projects
 along $B_h \oplus C_h$;
\item $h\colon B_h \to C_h$ is the inverse of $d|_{C_h}$, and
 $h|_{H_h} = h|_{C_h} = 0$.
\end{itemize}
The side conditions $p\iota = \mathrm{id}$, $h^2 = 0$,
$ph = 0$, $h\iota = 0$ follow immediately.

\emph{Step~2: Transferred operations via the tree formula.}
The homotopy transfer theorem
(Theorem~\ref{thm:htt},
Construction~\ref{constr:transfer-ainf})
produces $A_\infty$ operations $\{m_k^H\}$ on
$H^*(\barBgeom(\mathcal{A}))$ given by:
\[
m_n^H = \sum_{T \in \mathrm{PRT}_n}
\epsilon(T) \cdot p \circ
\prod_{v \in V(T)} m_{|v|} \circ
\prod_{e \in E_{\mathrm{int}}(T)} h
\circ \iota^{\otimes n},
\]
where the sum runs over planar rooted trees with $n$~leaves,
$m_{|v|}$ is the chiral product at each internal vertex of
valence~$|v|$, $h$ is inserted on each internal edge,
$\iota$ at leaves, and $p$ at the root
(Theorem~\ref{thm:tree-formula}).

The first operations are:
\begin{align}
m_1^H &= 0 \qquad\text{(differential vanishes on homology),}
 \label{eq:m1H} \\
m_2^H(a,b) &= p \, m_2(\iota a, \iota b)
 \qquad\text{(the induced product),}
 \label{eq:m2H} \\
m_3^H(a,b,c) &= p\, m_3(\iota a, \iota b, \iota c)
 + p\, m_2(h\, m_2(\iota a, \iota b),\, \iota c)
 \notag \\
 &\quad + (-1)^{|a|}\,
 p\, m_2(\iota a,\, h\, m_2(\iota b, \iota c))
 \label{eq:m3H}
\end{align}
(Kadeishvili~\cite{Kadeishvili80};
Construction~\ref{constr:transfer-ainf}).

\emph{Step~3: Chain level vs.\ homology level.}
At chain level, the associator satisfies:
\[
m_2(m_2 \otimes \mathrm{id}) - m_2(\mathrm{id} \otimes m_2)
= d_{\mathrm{bar}}(h_3) + m_3,
\]
where $h_3$ is the associating homotopy. In homology,
$d_{\mathrm{bar}}(h_3) = 0$, so $[m_2]$ is strictly associative:
the homotopy~$h_3$ is lost. The obstruction transfers to the
  Massey product $m_3^H$ via the displayed transferred ternary operation.
  For $k \geq 4$,
the tree formula produces $m_k^H$ from $C_k$ planar binary trees
(e.g., $C_4 = 5$ trees for $m_4^H$; see
Computation~\ref{comp:v1-virasoro-m4} for the Virasoro case).

\emph{Step~4: Vanishing for Koszul algebras.}
If $\mathcal{A}$ is Koszul, the bar spectral sequence
concentrates $H^*(\barBgeom(\mathcal{A}))$ on a single line
(Corollary~\ref{cor:bar-cohomology-koszul-dual}).
The SDR may then be chosen compatible with the bigrading,
so that $h$ preserves bar degree and shifts internal
(conformal weight) degree.
Every tree with $n \geq 3$ leaves contributing to $m_n^H$
would produce an operation of bidegree incompatible with
the Koszul line, forcing $m_n^H = 0$.
Thus, for Koszul $\mathcal{A}$, the minimal model has only
$m_2^H$ and is strictly associative. (This vanishing is weaker
than formality in the sense of
Deligne--Griffiths--Morgan--Sullivan.)
\end{proof}

\begin{example}[Explicit SDR for \texorpdfstring{$\widehat{\mathfrak{sl}}_2$}{sl2-hat}]
\label{ex:sdr-sl2}
For $\mathcal{A} = V_k(\mathfrak{sl}_2)$ with generators
$\{J^+, J^0, J^-\}$, the bar complex at bar degree~$2$
decomposes as $\barBgeom^2 = H^2 \oplus B^2 \oplus C^2$
where $d_{\mathrm{bar}}\colon C^2 \xrightarrow{\sim} B^2$
is the OPE residue map. The homotopy
$h_2 = (d_1)^{-1}\colon B^2 \to C^2$ is the inverse of the
bar differential restricted to its image. In the Chevalley--Eilenberg
model for $\mathfrak{sl}_2$, the Killing form in the basis
$\{e^*, h^*, f^*\}$ pairs $e^*$ with $f^*$ (off-diagonal):
the homotopy $h_2 = \kappa^{-1}$ acts by the inverse of this pairing.
Since $\mathfrak{sl}_2$ is semisimple, the CE complex is formal:
all transferred $L_\infty$ operations $\tilde{l}_n$ for $n \geq 2$
vanish for degree reasons (cohomology concentrated in
degrees~$0$ and~$3$).
\end{example}

\subsubsection{Explicit Verdier duality computations}

\begin{theorem}[Verdier duality of operations; \ClaimStatusProvedHere]\label{thm:verdier-duality-operations}
The bar and cobar operations are related by perfect duality:

\begin{center}
\begin{tabular}{|l|l|l|}
\hline
\textbf{Bar Side} & \textbf{Cobar Side} & \textbf{Pairing} \\
\hline
Logarithmic form $\eta_{ij}$ & Boundary current $[D_{ij}]$ & normalized residue $=1$ \\
Residue extraction & Current insertion & Verdier adjunction \\
Compactification $\overline{C}_n$ & Open space $C_n$ & Boundary-bulk correspondence \\
Product $m_2$ & Coproduct $\Delta_2$ & $\langle m_2, \Delta_2 \rangle = \mathrm{id}$ \\
Associator $m_3$ & Coassociator $\Delta_3$ & $\langle m_3, \Delta_3 \rangle = \Phi$ \\
\hline
\end{tabular}
\end{center}
\end{theorem}

\begin{proof}
The duality is a consequence of Verdier duality on the Fulton--MacPherson compactification $\overline{C}_n(X)$.

\emph{Row 1: $\eta_{ij} \leftrightarrow [D_{ij}]$.}
The logarithmic form $\eta_{ij} = d\log(z_i - z_j)$ is a section of $\Omega^1_{\log}(\overline{C}_n(X))$ with a simple pole along the collision divisor $D_{ij} = \{z_i = z_j\}$. The Verdier-dual object is the relative current carried by~$D_{ij}$, with the dualizing shift already included. The Verdier pairing is the normalized residue:
\[
\langle \eta_{ij}, [D_{ij}] \rangle_{\mathrm{Verdier}}
=
\frac{1}{2\pi i}
\int_{\gamma_{ij}}\frac{d(z_i-z_j)}{z_i-z_j}
=1 .
\]

\emph{Row 2: Residue $\leftrightarrow$ current insertion.}
More generally, the bar differential $d_{\bar{B}} = \sum_{i<j} \mathrm{Res}_{z_i = z_j}$ extracts residues at collision divisors, while the cobar differential $d_\Omega$ inserts the corresponding Verdier-dual boundary currents. These are dual operations: for $\omega \in \bar{B}^k$ and $K \in \Omega^k$,
\[
\langle d_{\bar{B}}(\omega), K \rangle = \langle \omega, d_\Omega(K) \rangle
\]
by the adjointness of residue and relative current insertion
(Stokes' formula on~$\overline{C}_n(X)$ with logarithmic boundary).

\emph{Row 3: $\overline{C}_n \leftrightarrow C_n$.}
The bar complex lives on the compactification~$\overline{C}_n(X)$
(logarithmic forms extend to the boundary), while its Verdier dual is
computed from the open inclusion
$j\colon C_n(X)\hookrightarrow\overline C_n(X)$ by
$\mathbb D(j_!\mathbb C_{C_n})\simeq Rj_*\omega_{C_n}[2n]$.
The boundary strata of~$\overline{C}_n$ encode the bar differential;
the Verdier-dual currents encode the algebra structure of
$K_X(\cA)=\mathbb D_{\Ran}\bar B_X(\cA)$.

\emph{Row 4: $m_2 \leftrightarrow \Delta_2$.}
The binary product $m_2(a \otimes b) = \mathrm{Res}_{z_1=z_2}[a(z_1) b(z_2) \cdot \eta_{12}]$ is dual to the binary coproduct $\Delta_2(c) = c(z_1) \cdot [D_{12}]$ via:
\[
\langle m_2(a \otimes b), c \rangle = \mathrm{Res}_{z_1=z_2}\!\left[\frac{a(z_1) b(z_2) c(z_1)}{z_1 - z_2}\right] = \langle a \otimes b, \Delta_2(c) \rangle.
\]
This is the adjunction $\langle m_2, \Delta_2 \rangle = \mathrm{id}$.

\emph{Row 5: $m_3 \leftrightarrow \Delta_3$.}
The ternary operation $m_3$ arises from the boundary of the associahedron $K_4 \cong [0,1]$ as a homotopy between $(m_2 \circ (m_2 \otimes \mathrm{id}))$ and $(m_2 \circ (\mathrm{id} \otimes m_2))$. Geometrically, $m_3$ integrates over the 1-dimensional moduli space of three-point collisions on~$\overline{C}_3(X)$. Its Verdier dual~$\Delta_3$ is the coassociator, and the pairing $\langle m_3, \Delta_3 \rangle = \Phi$ recovers the Drinfeld associator, which encodes the monodromy of the KZ connection on~$C_3(\mathbb{C})$.
\end{proof}

\begin{example}[Computing the duality pairing]
For the product/coproduct duality:

\emph{Bar side.} Product via residue
\[m_2(a \otimes b) = \text{Res}_{z_1=z_2}\left[\frac{a(z_1)b(z_2)}{z_1-z_2}dz_1\right]\]

\emph{Cobar side.} Coproduct via the diagonal current
\[\Delta_2(c) = \int c(w) [z_1=w]\,[z_2=w]\,dw = c(z_1)[D_{12}]\]

\emph{Pairing.}
\[\langle m_2(a \otimes b), \Delta_2(c) \rangle
=
\frac{1}{2\pi i}\int_{\gamma_{12}}
\frac{a(z_1)b(z_2)c(z_1)}{z_1-z_2}\,d(z_1-z_2)
=(abc)(0).\]

This recovers the structure constants of the chiral algebra.
\end{example}

% ================================================================
% SECTION 4.5: CONNECTION TO COM-LIE DUALITY
% ================================================================

\subsection{Connection to Com-Lie duality}

\subsubsection{The partition poset and configuration spaces}

\begin{theorem}[Geometric enhancement of Com-Lie; \ClaimStatusProvedElsewhere]\label{thm:geometric-com-lie-enhancement}
The bar complex of the commutative chiral operad is
(see \cite[Chapter~7]{LV12} for the operadic framework):
\[\bar{B}^{\text{ch}}(\text{Com}_{\text{ch}}) = \tilde{C}_*(\bar{\Pi}_n) \otimes \Omega^*_{\text{log}}(\overline{C}_n(X))\]

This enriches the partition complex with:
\begin{enumerate}
\item \emph{Combinatorial data:} Chains on the partition poset $\bar{\Pi}_n$
\item \emph{Geometric data.} Logarithmic forms on configuration spaces
\item \emph{$A_\infty$ structure:} Operations corresponding to faces of the partition poset
\end{enumerate}
\end{theorem}

\begin{proof}
Each partition $\pi \in \Pi_n$ corresponds to a stratum of $\overline{C}_n(X)$:
\[D_\pi = \{(z_1, \ldots, z_n) : z_i = z_j \text{ if } i,j \text{ in same block of } \pi\}\]

The differential:
\[d(\pi \otimes \omega) = \sum_{\pi' \text{ coarser}} \text{Res}_{D_{\pi'}}[\omega] \otimes \pi'\]

This realizes each relation in the partition poset as a geometric $A_\infty$ relation.
\end{proof}

\begin{example}[Pentagon from partitions]
For $n=5$, the partitions forming a pentagon are:
\begin{enumerate}
\item $\{\{1,2\},\{3\},\{4,5\}\}$: First $(12)$, then $(45)$
\item $\{\{1\},\{2,3\},\{4,5\}\}$: First $(23)$, then $(45)$
\item $\{\{1\},\{2,3,4\},\{5\}\}$: First $(234)$
\item $\{\{1,2,3\},\{4\},\{5\}\}$: First $(123)$
\item $\{\{1,2\},\{3,4\},\{5\}\}$: First $(12)$, then $(34)$
\end{enumerate}

These form the boundary of a 2-cell in $\bar{\Pi}_5$, giving the pentagon identity.
\end{example}

\subsubsection{\texorpdfstring{How $A_\infty$ structures interchange}{How A-infinity structures interchange}}

\begin{theorem}[Maximal vs.\ trivial \texorpdfstring{$A_\infty$}{A-infinity}; \ClaimStatusProvedElsewhere]\label{thm:ainfty-com-lie-interchange}
Under Com-Lie duality, $A_\infty$ structures interchange
\textup{(}\cite[Theorem~7.4.1]{LV12}\textup{)}:

\emph{Commutative side.}
\begin{itemize}
\item $m_1 = 0$ (no differential)
\item $m_2 = $ symmetric product
\item $m_k = 0$ for $k \geq 3$ (no higher operations)
\item Trivial $A_\infty$ structure
\end{itemize}

\emph{Lie side.}
\begin{itemize}
\item $m_1 = 0$ (no differential)
\item $m_2 = $ antisymmetric bracket
\item $m_3 = $ homotopy for the failure of $m_2$ to be associative
\item Higher $m_k$ encode coherent homotopies (generally non-trivial)
\item Non-trivial $A_\infty$ structure (contrast with Com)
\end{itemize}
\end{theorem}

\begin{proof}[Via configuration spaces]
For Com: All points can collide simultaneously without constraint
\[\overline{C}_n^{\text{Com}}(X) = X \times \overline{M}_{0,n}\]

For Lie: Points must collide in a specific tree pattern
\[\overline{C}_n^{\text{Lie}}(X) = \text{Blow-up along all diagonals}\]

The difference in these compactifications determines the $A_\infty$ structure.
\end{proof}

% ================================================================
% SECTION 4.6: CURVED AND FILTERED EXTENSIONS
% ================================================================

\subsection{Curved and filtered extensions}

\subsubsection{\texorpdfstring{Curved $A_\infty$ algebras: central extensions and anomalies}{Curved A-infinity algebras: central extensions and anomalies}}

\begin{definition}[Curved \texorpdfstring{$A_\infty$}{A-infinity} algebra]\label{def:curved-ainfty-hg}
\index{curved $A_\infty$-algebra|textbf}
\index{curvature!central}
A curved $A_\infty$ algebra consists of a graded vector space $\mathcal{A}$ with operations $m_k: \mathcal{A}^{\otimes k} \to \mathcal{A}$ for $k \geq 0$, where:
\begin{enumerate}
\item $m_0 \in \mathcal{A}^2$ is the \emph{curvature element} (a nullary operation producing a degree-2 element)
\item The $A_\infty$ relations hold for all $n \geq 0$: $\sum_{r+s+t=n} (-1)^{rs+t} m_{r+1+t}(\text{id}^{\otimes r} \otimes m_s \otimes \text{id}^{\otimes t}) = 0$, now including $s = 0$ terms involving $m_0$
\item In particular, $m_1(m_0) = 0$ (curvature is a cocycle) and $m_1^2 = [m_0, -]$ (failure of strict nilpotence)
\end{enumerate}
\end{definition}

\begin{theorem}[Genus-$g$ bar complex as curved
$\mathrm{Ch}_\infty$-algebra; \ClaimStatusConditional]
\label{thm:bar-curved-ch-infty}
\index{curved Ch-infinity@curved $\mathrm{Ch}_\infty$-algebra|textbf}
\index{bar complex!curved homotopy chiral structure}
For a chiral algebra $\cA$ with scalar curvature $\kappa(\cA)$, the
genus-$g$ bar complex $\barB^{(g,n)}(\cA)$ carries a \emph{curved
homotopy chiral algebra} structure: a curved
$\mathrm{Ch}_\infty$-algebra in the sense of
Malikov--Schechtman~\textup{\cite{MS24}}, with:
\begin{enumerate}[label=\textup{(\roman*)}]
\item Curvature element
 $m_0^{(g)}$, whose scalar diagonal projection is
 $\kappa(\cA)\cdot\omega_g^{\mathrm{Ar}}$,
 where $\omega_g^{\mathrm{Ar}}$ is the Arakelov $(1,1)$-form on the
 fixed genus-$g$ surface.
\item The failure of strict nilpotence:
 $m_1^{(g)} \circ m_1^{(g)} = [m_0^{(g)}, -]$; only after scalar
 diagonal projection does $m_0^{(g)}$ give
 $\kappa(\cA) \cdot \omega_g^{\mathrm{Ar}}$
 \textup{(}Theorem~\textup{\ref{thm:arnold-quantum})}.
\item The Eilenberg--Zilber homotopies $F_n^{(g)}$ acquire curvature
 corrections: $d F_n^{(g)} + F_n^{(g)} d = \text{boundary term}
 + \kappa(\cA) \cdot \omega_g^{\mathrm{Ar}}
 \cdot (\text{lower-degree terms})$.
\item At genus~$0$, the curvature vanishes and the structure reduces to
 the uncurved $\mathrm{Ch}_\infty$-algebra of
 Theorem~\textup{\ref{thm:cech-hca}}.
\end{enumerate}
This is the $\mathrm{Ch}_\infty$ avatar of the curved
$A_\infty$-structure of Definition~\textup{\ref{def:curved-ainfty-hg}}.
\end{theorem}

\begin{proof}
The curved identity $m_1^{(g)\,2}=[m_0^{(g)},-]$ and the scalar
projection
$\operatorname{tr}_{\mathrm{diag}}(m_0^{(g)})
=\kappa(\cA)\cdot\omega_g^{\mathrm{Ar}}$ are proved
in Theorem~\ref{thm:arnold-quantum} via the genus-$g$ Arnold defect.
The $\mathrm{Ch}_\infty$ structure follows by applying the
Malikov--Schechtman construction~\cite[Theorem~3.1]{MS24} to the
cosimplicial object $\barB^{(g,\bullet)}(\cA)$, with the curvature
correction arising from the period-matrix correction to the
genus-$0$ propagator. Parts~(ii) and~(iii) are the
$\mathrm{Ch}_\infty$-algebra restatement of the fact that the bar
differential squares to the curvature, not to zero.
\end{proof}

\begin{example}[Heisenberg algebra: curved structure]
The Heisenberg algebra $\mathcal{H}_k$ has current $J$ with OPE:
\[J(z)J(w) = \frac{k}{(z-w)^2} + \text{regular}\]

The absence of a simple pole gives $m_2(J \otimes J) = 0$, curvature $\kappa(\mathcal{H}_k) = k$, and modified differential $d_{\text{curved}} = d + k \cdot \mu_0$.

The bar complex:
\[\bar{B}^n(\mathcal{H}_k) = \begin{cases}
\mathbb{C} & n = 0 \\
\text{Currents} & n = 1 \\
\mathbb{C} \cdot c_k & n = 2 \\
0 & n \geq 3
\end{cases}\]

The level $k$ appears as the curvature controlling the failure of strict associativity.
\end{example}

\begin{example}[Virasoro algebra: curved \texorpdfstring{$A_\infty$}{A-infinity}]
The Virasoro algebra with stress tensor $T$ has:
\[T(z)T(w) = \frac{c/2}{(z-w)^4} + \frac{2T(w)}{(z-w)^2} + \frac{\partial T(w)}{z-w} + \text{regular}\]

The curved structure:
\begin{itemize}
\item Curvature from central charge $c$
\item Modified Jacobi identity involving $c$
\item $m_3$ includes Schwarzian derivative terms
\item Higher $m_k$ encode conformal anomalies
\end{itemize}
\end{example}

\subsubsection{Filtered and complete structures}

\begin{definition}[Filtered chiral algebra]\label{def:filtered-chiral-algebra}
A filtered chiral algebra has:
\[F_0\mathcal{A} \subset F_1\mathcal{A} \subset F_2\mathcal{A} \subset \cdots\]
with $\mu(F_i \otimes F_j) \subset F_{i+j}$ (compatibility), $\mathcal{A} = \bigcup_i F_i\mathcal{A}$ (exhaustive), and $\bigcap_i F_i\mathcal{A} = 0$ (separated).
\end{definition}

\begin{theorem}[Convergence for completed filtered algebras; \ClaimStatusConditional]
\label{thm:convergence-filtered}
Let $\mathcal{A}$ be a complete filtered chiral algebra (i.e.,
$\mathcal{A} = \varprojlim \mathcal{A}/F_k\mathcal{A}$) whose associated
graded $\mathrm{gr}\,\mathcal{A}$ is a Koszul chiral algebra in the sense of
Theorem~\ref{thm:chiral-koszul-duality}. Then:
\begin{enumerate}
\item The bar complex is formed with completed tensor products:
each $\widehat{\bar{B}}^n(\mathcal{A})$ is complete in the induced
filtration, and no completion beyond the given separated filtration is
introduced.
\item Each cohomology class in $H^*(\widehat{\bar{B}}(\mathcal{A}))$ has a canonical
representative given by the unique closed element in the $\mathrm{gr}$-harmonic
complement
\item The continuous cobar-bar adjunction
$\Omega^{\mathrm{ch}}_{\mathrm{cont}}(\widehat{\bar{B}}^{\mathrm{ch}}(\mathcal{A})) \xrightarrow{\sim}
\mathcal{A}$ is a filtered quasi-isomorphism
\item Koszul duality extends: the Koszul dual $\mathcal{A}^!$ inherits a
complete filtration with $\mathrm{gr}\,\mathcal{A}^! \cong
(\mathrm{gr}\,\mathcal{A})^!$
\end{enumerate}
\end{theorem}

\begin{proof}
The proof proceeds by relating the filtered bar complex to its associated graded.

\emph{Step 1: Filtered bar complex.}
The filtration $F_\bullet \mathcal{A}$ induces a filtration on the bar complex:
\[F_p \bar{B}^n(\mathcal{A}) := \Gamma\bigl(\overline{C}_{n+1}(X),\,
\textstyle\sum_{i_0 + \cdots + i_n \geq p}
F_{i_0}\mathcal{A} \boxtimes \cdots \boxtimes F_{i_n}\mathcal{A}
\otimes \Omega^n_{\log}\bigr)\]
Since $\mathcal{A}$ is complete and separated, and the completed
tensor product of complete filtered spaces is complete (by
Proposition A.5 of Positselski~\cite{Positselski11}), each
$\widehat{\bar{B}}^n(\mathcal{A})$ is complete. This proves (1):
convergence means completeness of the filtered inverse limit, not
convergence of an ordinary infinite direct sum.

\emph{Step 2: Associated graded comparison.}
By naturality of the bar construction with respect to filtrations:
\[\mathrm{gr}\,\bar{B}(\mathcal{A}) \cong \bar{B}(\mathrm{gr}\,\mathcal{A})\]
This is verified degree by degree: the $p$-th graded piece of $\bar{B}^n(\mathcal{A})$
consists of sections with total filtration degree exactly $p$, which is precisely
$\bar{B}^n(\mathrm{gr}\,\mathcal{A})_p$.

\emph{Step 3: Koszul hypothesis gives $\mathrm{gr}$ convergence.}
Since $\mathrm{gr}\,\mathcal{A}$ is Koszul, Theorem~\ref{thm:bar-cobar-inversion-qi} gives:
\[\Omega^{\mathrm{ch}}(\bar{B}^{\mathrm{ch}}(\mathrm{gr}\,\mathcal{A}))
\xrightarrow{\sim} \mathrm{gr}\,\mathcal{A}\]
The spectral sequence collapses at $E_2$ for $\mathrm{gr}\,\mathcal{A}$, and
each cohomology class of $\bar{B}(\mathrm{gr}\,\mathcal{A})$ has a unique harmonic
representative.

\emph{Step 4: Filtered quasi-isomorphism.}
The filtration spectral sequence for $\bar{B}(\mathcal{A})$ has:
\[E_1^{p,q} = H^q(\mathrm{gr}_p\,\bar{B}(\mathcal{A})) =
H^q(\bar{B}(\mathrm{gr}\,\mathcal{A})_p) \Longrightarrow
H^{p+q}(\bar{B}(\mathcal{A}))\]
Since $E_1 = H^*(\bar{B}(\mathrm{gr}\,\mathcal{A}))$ and the spectral sequence for
$\mathrm{gr}\,\mathcal{A}$ collapses at $E_2$, this spectral sequence also collapses
at $E_2$. Boardman's Comparison Theorem \cite{Boardman-conditional}
requires completeness ($\cA = \varprojlim \cA/F_k$, $\bigcap F_k = 0$)
and exhaustiveness ($\bigcup F_k = \cA$), both of which hold by
hypothesis. The convergence of the $\mathrm{gr}$ spectral sequence
therefore implies convergence of the filtered one, and the cobar-bar
map
$\Omega^{\mathrm{ch}}_{\mathrm{cont}}
(\widehat{\bar{B}}^{\mathrm{ch}}(\cA)) \to \cA$
is a filtered quasi-isomorphism. This proves (2) and (3).

\emph{Step 5: Filtered Koszul dual.}
The completed bar coalgebra $\widehat{\bar{B}}(\cA)$ carries a filtration-compatible
coassociative coproduct $\Delta$ (from Step~1: the tensor-product
filtration is compatible with the deconcatenation coproduct). The
Koszul dual algebra is the continuous linear dual
$\cA^! := \widehat{\bar{B}}(\cA)^\vee_{\mathrm{cont}}$ with the dual filtration
$F_p(\cA^!) = (F_p\widehat{\bar{B}}(\cA))^\perp$; the chiral algebra structure
on~$\cA^!$ is the dual of the coalgebra structure: the product
$\mu^! = \Delta^\vee$ and the higher operations
$(m_k^!)^\vee = \Delta_k$ are obtained by dualizing the
bar coproduct and its iterated versions. This is a chiral algebra
because the coalgebra axioms
dualize: coassociativity of~$\Delta$ gives associativity of~$\mu^!$,
and $d^2 = 0$ on $\widehat{\bar{B}}(\cA)$ gives $d^2 = 0$
on~$\cA^!$. By Step~2,
$\mathrm{gr}\,\cA^! = \mathrm{gr}\,\widehat{\bar{B}}(\cA)^\vee_{\mathrm{cont}}
= \bar{B}(\mathrm{gr}\,\cA)^\vee
= (\mathrm{gr}\,\cA)^!$.
Completeness of $\cA^!$ follows from completeness of
$\widehat{\bar{B}}(\cA)$ (Positselski~\cite[\S A.1]{Positselski11}).
This proves (4).
\end{proof}

\begin{example}[W-algebras are filtered]
The $W_N$ algebra has filtration by conformal weight:
\[F_k = \mathrm{span}\{W^{(s)} : s \leq k\}\]

This filtration satisfies the hypotheses of Theorem~\ref{thm:convergence-filtered}: it is complete and separated (conformal weight is bounded below), compatible with chiral operations ($\mu(F_i \otimes F_j) \subset F_{i+j}$), and has associated graded $\mathrm{gr}\,W_N \cong$ free field algebra (Koszul by Theorem~\ref{thm:chiral-koszul-duality}).
\end{example}

% ================================================================
% SECTION 4.7: THE COBAR RESOLUTION
% ================================================================

\subsection{The cobar resolution and Ext groups}

\subsubsection{Resolution at chain level}

\begin{theorem}[Cobar resolution on the Koszul locus {\cite{LV12}}; \ClaimStatusProvedElsewhere]
\label{thm:cobar-resolution-scoped}
For any augmented chiral algebra $\mathcal{A}$, the bar construction
$\bar{B}^{\mathrm{ch}}(\mathcal{A})$ and cobar construction
$\Omega^{\mathrm{ch}}(\bar{B}^{\mathrm{ch}}(\mathcal{A}))$ exist as objects.
When the canonical twisting morphism is Koszul
\textup{(}Definition~\textup{\ref{def:chiral-koszul-morphism})},
the cobar of the bar provides a free resolution:
\begin{multline*}
\cdots \to \Omega^2_{\mathrm{ch}}(\bar{B}^{\mathrm{ch}}(\cA))
\to \Omega^1_{\mathrm{ch}}(\bar{B}^{\mathrm{ch}}(\cA))\\
\to \Omega^0_{\mathrm{ch}}(\bar{B}^{\mathrm{ch}}(\cA))
\xrightarrow{\epsilon} \cA \to 0
\end{multline*}
Off the Koszul locus, the bar-cobar object persists only in the
completed coderived category
\textup{(}Appendix~\textup{\ref{sec:coderived-models})}.

The augmentation is given geometrically by:
\[\epsilon(K) = \lim_{\varepsilon \to 0} \int_{|z_i - z_j| > \varepsilon} K(z_1, \ldots, z_n) \prod_{i<j} |z_i - z_j|^{2h_{ij}}\]
\end{theorem}

\begin{remark}[Ext groups]
On the Koszul locus, this resolution computes:
\[\text{Ext}^n_{\text{ChirAlg}}(\mathcal{A}, \mathcal{B}) \cong H^n(\text{Hom}(\Omega^{\text{ch}}(\bar{B}^{\text{ch}}(\mathcal{A})), \mathcal{B}))\]

Geometrically:
\begin{itemize}
\item $n = 0$: Morphisms of chiral algebras
\item $n = 1$: Derivations and infinitesimal automorphisms
\item $n = 2$: Extensions and deformation obstructions
\item $n = 3$: Massey products and triple compositions
\item $n \geq 4$: Higher coherences and Toda brackets
\end{itemize}
\end{remark}

\begin{example}[Fermion-boson Koszul duality]
The bar complex of the free fermion is a coalgebra:
$\bar{B}^{\mathrm{ch}}(\mathcal{F})$ is the cofree cocommutative
chiral coalgebra encoding the antisymmetric OPE
$\psi(z)\psi(w) \sim (z{-}w)^{-1}$. Bar-cobar inversion recovers
the fermion:
$\Omega^{\mathrm{ch}}(\bar{B}^{\mathrm{ch}}(\mathcal{F}))
\xrightarrow{\sim} \mathcal{F}$.
The Koszul dual \emph{algebra} is obtained by dualizing the bar
coalgebra:
$\mathcal{F}^! = \bar{B}^{\mathrm{ch}}(\mathcal{F})^\vee
\cong \beta\gamma$
(Theorem~\ref{thm:fermion-boson-koszul};
cf.\ the explicit computation in Chapter~\ref{ch:heisenberg-frame}).
This is Koszul duality (Lie--Com), not bosonization.
\end{example}

% ================================================================
% SECTION 4.8: MAURER-CARTAN ELEMENTS AND DEFORMATIONS
% ================================================================

\subsection{Maurer--Cartan elements and deformation theory}

\subsubsection{The moduli space of deformations}

\begin{theorem}[Maurer--Cartan = deformations {\cite{Kon03,Ger63}}; \ClaimStatusProvedElsewhere]\label{thm:genus-graded-mc}
Maurer--Cartan elements in $\bar{B}^1(\mathcal{A})[[t]]$ satisfying
\[d\alpha + \frac{1}{2}[\alpha, \alpha] = 0\]
parametrize formal deformations of the chiral algebra structure.
Each MC element $\alpha$ yields deformed operations:
\[m_2^\alpha(a \otimes b) = m_2(a \otimes b) + \langle \alpha, a \otimes b \rangle, \qquad
m_3^\alpha = m_3 + \partial\alpha + \alpha \cup \alpha.\]
\end{theorem}

\begin{proof}
This is a consequence of the general theory of deformations via
dg Lie algebras: the Hochschild cochain complex
$C^*_{\mathrm{ch}}(\mathcal{A}, \mathcal{A})$ carries a dg Lie
structure (Theorem~\ref{thm:linfty-cochains}), and the MC equation
in this dg Lie algebra parametrizes flat deformations of the
$A_\infty$ structure. The equivalence between MC elements and
formal deformations is due to Kontsevich~\cite{Kon03},
extending the classical deformation theory of
Gerstenhaber~\cite{Ger63}. For the chiral setting, the
Hochschild cochains are realized as distributional kernels on
configuration spaces, and the MC equation becomes the flatness
condition for connections on
$\overline{C}_2(X)$ (cf.\ Theorem~\ref{thm:geometric-chiral-hochschild}).
\end{proof}

\subsubsection{Example: Yangian deformation}

\begin{proposition}[Yangian from deformation {\cite{Drinfeld85}}; \ClaimStatusProvedElsewhere]
\label{prop:yangian-from-deformation}
The Yangian $Y(\mathfrak{g})$ arises as a deformation of $U(\mathfrak{g}[z])$
via the MC element
$\alpha = \hbar\, r/(z_1 - z_2)$
where $r \in \mathfrak{g} \otimes \mathfrak{g}$ is the classical
$r$-matrix. The deformed OPE is
\[J^a_\hbar(z)J^b_\hbar(w) = \frac{k\delta^{ab}}{(z-w)^2} + \frac{f^{abc}J^c(w)}{z-w} + \frac{\hbar\, r^{ab}}{(z-w)^2}\]
which, upon expanding in modes, produces the RTT relations
of the Yangian \textup{(}Drinfeld~\cite{Drinfeld85}\textup{)}.
The MC equation $d\alpha + \frac{1}{2}[\alpha, \alpha] = 0$
reduces to the classical Yang--Baxter equation
$[r_{12}, r_{13}] + [r_{12}, r_{23}] + [r_{13}, r_{23}] = 0$
for~$r$.
\end{proposition}

\subsubsection{Example: Heisenberg deformation}

\begin{convention}[Heisenberg level notation]
\label{conv:heisenberg-kappa-notation}
We write $\mathcal{H}_\kappa$ for the Heisenberg vertex algebra at
level~$\kappa$, using the same letter as the modular characteristic
to emphasize the identity $\kappa(\mathcal{H}_\kappa) = \kappa$: for
the Heisenberg algebra, the level parameter coincides with the modular
characteristic (Theorem~\textup{\ref{thm:modular-characteristic}}).
\end{convention}

\begin{proposition}[Deforming Heisenberg; \ClaimStatusProvedHere]
\label{prop:deforming-heisenberg}
The Heisenberg algebra $\mathcal{H}_\kappa$ on a smooth curve~$X$ of
genus~$g$ admits deformations parametrized by
$H^1(\bar{B}(\mathcal{H}_\kappa)) \cong H^1(X, \mathbb{C}) \oplus
\mathbb{C} \cdot \partial\kappa$.
\end{proposition}

\begin{proof}
The bar complex of $\mathcal{H}_\kappa$ is the symmetric coalgebra
$\bar{B}(\mathcal{H}_\kappa) = \mathrm{Sym}^c(s^{-1}\overline{\mathcal{H}})$
(Theorem~\ref{thm:heisenberg-koszul-dual-early}). The degree-$1$
component is $\bar{B}^1 = s^{-1}\overline{\mathcal{H}}_1 = s^{-1}
\mathbb{C} \cdot J$, so $H^1$ of the bar complex on~$X$ is computed
by the Leray spectral sequence for $\bar{B}$ viewed as a sheaf
on~$X$ (convergent: $X$ has finite cohomological dimension and
$\mathcal{H}^q(\bar{B}^\bullet)$ is a coherent sheaf for each~$q$):
\[
E_2^{p,q} = H^p(X,\, \mathcal{H}^q(\bar{B}^{\bullet}))
\;\Longrightarrow\; H^{p+q}(\bar{B}^{\mathrm{global}}).
\]
At $E_2$: the $E_2^{1,0}$ term gives
$H^1(X, \mathcal{Z}(\mathcal{H}_\kappa)) = H^1(X, \mathbb{C})$
(the center is $\mathbb{C}$, and $H^1(X, \mathbb{C}) \cong
\mathbb{C}^{2g}$); the $E_2^{0,1}$ term gives $\mathbb{C} \cdot
\partial\kappa$ (the purely algebraic deformation of the level).
Both contributions are MC elements:
$\alpha = \sum_{i=1}^{2g} a_i\, \omega_i$ shifts the current by a
harmonic $1$-form on~$X$, while $\alpha = b \cdot \partial\kappa$
deforms the level $\kappa \mapsto \kappa + tb$.
\end{proof}

\subsubsection{\texorpdfstring{Example: $\beta\gamma$ system deformation}{Example: beta-gamma system deformation}}

\begin{proposition}[\texorpdfstring{$\beta\gamma$}{beta-gamma} deformations; \ClaimStatusProvedHere]
\label{prop:betagamma-deformations}
The $\beta\gamma$ system admits a $1$-parameter family of deformations
shifting conformal weights:
$h_\beta \mapsto 1 + t$, $h_\gamma \mapsto -t$, while preserving the
fundamental OPE $\beta(z)\gamma(w) \sim 1/(z-w)$.
\end{proposition}

\begin{proof}
The MC element $\alpha = t \cdot \omega_{\mathrm{contact}} \in
H^1(\bar{B}(\beta\gamma))$, where $\omega_{\mathrm{contact}}$ is the
contact $1$-form on $\overline{C}_2(X)$ with residue along the diagonal,
deforms the stress tensor by
$T \mapsto T + t\,\partial(\beta\gamma)$. This shifts the conformal
weights of $\beta$ and $\gamma$ by $+t$ and $-t$ respectively (since
$[L_0, \beta] = (1{+}t)\beta$ and $[L_0, \gamma] = -t\,\gamma$ in
the deformed algebra), while preserving the $\beta\gamma$ OPE (which
has only a simple pole, unaffected by the contact deformation).
The MC equation $d\alpha + \frac{1}{2}[\alpha,\alpha] = 0$ is
satisfied because $\omega_{\mathrm{contact}}$ is closed on
$\overline{C}_2(X)$ and the bracket $[\alpha, \alpha] = 0$
vanishes (the $\beta\gamma$ system is abelian at the level of
the deformation, since $\omega_{\mathrm{contact}} \in
H^1(\overline{C}_2(X))$ is one-dimensional and any $1$-form
squares to zero in the Lie bracket).
\end{proof}

% ================================================================
% SECTION 4.9: EXAMPLES OF TRANSVERSE STRUCTURES
% ================================================================

\subsection{Examples of transverse structures}

\subsubsection{The Jacobiator identity}

\begin{theorem}[Jacobiator for Lie-type algebras; \ClaimStatusProvedHere]\label{thm:jacobiator-lie-type}
For any Lie-type chiral algebra, the Jacobiator:
\[J(a,b,c) = [[a,b],c] + [[b,c],a] + [[c,a],b]\]
measures the failure of the Jacobi identity. The vanishing $J = 0$ is encoded by the boundary of the pentagon $K_4$: each vertex corresponds to a parenthesization of the triple bracket, and $\partial K_4 = 0$ gives the five-term relation among compositions of $m_2$ and $m_3$.
\end{theorem}

\begin{proof}[Geometric origin]
In $\overline{C}_4(X)$, the codimension-1 strata correspond to the five boundary divisors of $\overline{M}_{0,5}$, forming the boundary of the pentagon $K_4$. Each edge of $K_4$ corresponds to a different way to evaluate the triple bracket $[[a,b],c]$ via iterated binary collisions:
\begin{enumerate}
\item Five vertices: Five fully parenthesized expressions on four inputs
\item Five edges: Five reassociation moves connecting them
\end{enumerate}

The relation:
\[\sum_{\text{edges}} \text{sign}(\text{edge}) \cdot (\text{reassociation}) = 0\]
follows from $\partial K_4 = 0$, giving the pentagon identity for $m_4$.
\end{proof}

\subsubsection{The Bianchi identity in chiral context}

\begin{theorem}[Chiral Bianchi identity; \ClaimStatusProvedHere]\label{thm:chiral-bianchi}
For chiral algebras with connection-type structure, there is a Bianchi identity:
\[d_\nabla F + [A, F] = 0\]
where $F$ is the curvature 2-form in the bar complex.
\end{theorem}

\begin{proof}
The curvature lives in $\bar{B}^2$:
\[F = \sum_{i<j} F_{ij} \otimes \eta_{ij} \in \Gamma(\overline{C}_2(X), \mathcal{A}^{\otimes 2} \otimes \Omega^1_{\text{log}})\]
Apply $d^2 = 0$ (Theorem~\ref{thm:bar-nilpotency-complete}) to the bar element $A \otimes F \in \bar{B}^3$. On $\overline{C}_3(X)$, the nine terms of $d^2$ (ibid.)\ reduce as follows: terms involving $d_{\mathrm{int}}^2 = 0$ and $d_{\mathrm{dR}}^2 = 0$ vanish identically; the cross-terms give
\[d_{\mathrm{bar}}(F) = \mathrm{Res}_{D_{12}}[A \otimes F_{23}] - \mathrm{Res}_{D_{23}}[F_{12} \otimes A] + \text{cyclic permutations}\]
by the residue definition of the bar differential. The right-hand side is $-[A, F]$ (the commutator in $\mathcal{A}$, realized via opposing residues on the two collision divisors meeting at the triple diagonal in $\overline{C}_3$). Therefore $d_\nabla F + [A,F] = 0$.
\end{proof}

\subsubsection{The octahedron identity}

\begin{corollary}[Higher associahedron identity for \texorpdfstring{$m_6$}{m6}; \ClaimStatusProvedHere]\label{cor:higher-associahedron-m6}
For six elements, the $A_\infty$ relation $\sum_{r+s+t=6}(-1)^{rs+t}m_{r+1+t}\circ(\mathrm{id}^{\otimes r}\otimes m_s\otimes\mathrm{id}^{\otimes t})=0$ involves the $C_5 = 42$ full parenthesizations, governed by the boundary stratification of the $4$-dimensional associahedron~$K_6$.
\end{corollary}

\begin{proof}
This is the $n=6$ case of Theorem~\ref{thm:bar-nilpotency-complete} ($d^2=0$ on $\overline{C}_7(X)$). The codimension-$1$ boundary strata of $K_6$ are products $K_i \times K_j$ with $i+j=7$, and each such stratum contributes a composition $m_i \circ_k m_j$. The signs are determined by the Koszul sign rule on desuspended generators as in the general case.
\end{proof}

\section{Genus 2 OPE contributions: A concrete example in full detail}
\label{sec:genus_2_ope_example}

\subsection{Setting: genus 2 Riemann surfaces}

\subsubsection{\texorpdfstring{Moduli space $\mathcal{M}_2$}{Moduli space M-2}}

A genus 2 Riemann surface can be represented as:
\[\Sigma_2 = \mathbb{H}/\Gamma\]
where $\mathbb{H}$ is the upper half-plane and $\Gamma \subset \operatorname{PSL}(2,\mathbb{R})$
is a Fuchsian group.

The moduli space $\mathcal{M}_2$ has complex dimension $3g-3 = 3$, with coordinates given by period matrices $\Omega \in \mathbb{H}_2$ (Siegel upper half-space) and volume form $d\mu_{\text{WP}}$ (Weil--Petersson).

\subsubsection{The period matrix}

Explicitly, choose a symplectic basis $\{a_1, a_2, b_1, b_2\}$ of $H_1(\Sigma_2, \mathbb{Z})$
with intersection form:
\[a_i \cdot b_j = \delta_{ij}, \quad a_i \cdot a_j = b_i \cdot b_j = 0\]

Let $\omega_1, \omega_2$ be normalized holomorphic differentials:
\[\oint_{a_i} \omega_j = \delta_{ij}\]

The period matrix is:
\[\Omega = (\Omega_{ij}) \quad \text{where} \quad \Omega_{ij} = \oint_{b_i} \omega_j\]

Symmetry: $\Omega = \Omega^T$, Positivity: $\operatorname{Im}(\Omega) > 0$.

\subsection{\texorpdfstring{Configuration space on $\Sigma_2$}{Configuration space on Sigma-2}}

\subsubsection{Two-point configurations}

Consider the configuration space:
\[\mathrm{Conf}_2(\Sigma_2) = \{(z_1, z_2) \in \Sigma_2 \times \Sigma_2 \mid z_1 \neq z_2 \}\]

Unlike genus 0 or 1, at genus 2 we have \emph{multiple geodesics} connecting $z_1, z_2$.
The OPE receives contributions from \emph{all} homology classes of paths.

\subsubsection{The Green's function}

The bosonic propagator on $\Sigma_2$ is the Green's function:
\[G_{\Sigma_2}(z_1, z_2) = -\log|E_{\Sigma_2}(z_1, z_2)|^2 + \text{(harmonic)}\]
where $E_{\Sigma_2}$ is the prime form.

\emph{Explicit formula} (Fay's trisecant identity):
\[E_{\Sigma_2}(z_1, z_2) = \frac{\theta[\Delta](z_1 - z_2 | \Omega)}
{\sqrt{\omega_{z_1}(z_1)} \sqrt{\omega_{z_2}(z_2)}}\]
where $\theta[\Delta]$ is the theta function with characteristic $\Delta$ and $\omega_{z_i}$ is the canonical abelian differential.

\subsection{The Heisenberg algebra at genus 2}

\subsubsection{\texorpdfstring{Operators on $\Sigma_2$}{Operators on -2}}

The Heisenberg operators $a(z), a^*(z)$ on $\Sigma_2$ satisfy:
\[\langle a(z_1) a^*(z_2) \rangle_{\Sigma_2} = G_{\Sigma_2}(z_1, z_2) + \kappa \cdot (\text{contact terms})\]

The Heisenberg level~$\kappa$ appears in:
\begin{itemize}
\item Genus 0 correction: in $(z_1 - z_2)^{-2}$ pole
\item Genus 1 correction: in trace around $S^1$ cycles
\item \emph{Genus 2 correction}: in double-trace contributions
\end{itemize}

\subsubsection{The genus 2 vacuum}

The genus-two vacuum expectation value has the form
\[
 \langle\mathbbm 1\rangle_{\Sigma_2}
 =e^{-S_{\mathrm{cl}}[\Sigma_2]}
 Z_{\mathrm{1-loop}}(\Omega;\kappa).
\]
The precise determinant power is fixed by the field content and the
Gaussian measure normalization.  Its modular dependence is expressed
through the period matrix~$\Omega$.

\subsection{Computing a genus 2 OPE correction}

\subsubsection{The setup}

Consider the OPE:
\[a(z) \cdot a^*(w) = \frac{\kappa}{(z-w)^2} + \text{reg}
+ \text{(genus 1 corr)} + \text{(genus 2 corr)} + \cdots\]

\subsubsection{The Feynman diagram picture}

At genus 2, the relevant Feynman diagram has two loops with external legs at $z$ and $w$.

This contributes:
\[\mathcal{A}_2(z,w) = \int_{\mathcal{M}_2} d\mu_{\text{WP}}
\int_{\Sigma_2^2} G(z, z_1) G(z_1, z_2) G(z_2, w) \cdot (\text{insertions})\]

\subsubsection{Explicit integration}

\emph{Step 1: The double contour integral.}

Using the method of images on $\Sigma_2$:
\begin{align}
&\int_{\Sigma_2} G(z, z_1) G(z_1, w) \\
&= \sum_{\gamma \in \pi_1(\Sigma_2)}
\int_{\gamma} \frac{dz_1}{2\pi i}
\frac{\theta[\Delta](z - z_1 | \Omega)}{\theta[\Delta](z_1 - w | \Omega)}
\cdot (\omega \text{ factors})
\end{align}

The sum over $\gamma$ accounts for winding around the two handles.

\emph{Step 2: Residue calculations.}

Each term in the sum gives:
\begin{itemize}
\item $\gamma = a_1$: contribution from first handle
\item $\gamma = a_2$: contribution from second handle
\item $\gamma = b_1, b_2$: dual cycle contributions
\item Cross terms: $\gamma = a_1 b_1, a_1 b_2$, etc.
\end{itemize}

After residue calculations (using Riemann bilinear relations):
\[\int_{\Sigma_2} G(z, z_1) G(z_1, w) =
\frac{\partial^2}{\partial \Omega_{11}} G_{\Sigma_2}(z, w)
+ \frac{\partial^2}{\partial \Omega_{22}} G_{\Sigma_2}(z, w)
+ \text{(mixed terms)}\]

\emph{Step 3: Integration over moduli.}

Now integrate over $\mathcal{M}_2$:
\begin{align}
&\int_{\mathcal{M}_2} d\mu_{\text{WP}} \cdot
\frac{\partial^2 G}{\partial \Omega_{ij}} \\
&= \int_{\mathcal{M}_2} \frac{d^3\Omega}{(\det \operatorname{Im} \Omega)^{3}}
\cdot \frac{\partial^2}{\partial \Omega_{ij}}
\left[ -\log|\theta[\Delta](z-w|\Omega)| \right]
\end{align}

This integral is:
\begin{itemize}
\item \emph{Divergent}: requires regularization (think: UV divergence in QFT)
\item \emph{Universal}: the divergence is independent of $z, w$ (up to logs)
\item \emph{Modular}: depends on Eisenstein series $E_4(\Omega), E_6(\Omega)$
\end{itemize}

\subsubsection{The renormalized result}

\begin{remark}[OPE vs.\ correlation functions]\label{rem:ope-genus-independent}
The local OPE singularity structure is genus-independent: it is an algebraic property
of the vertex algebra, determined by conformal weights. For the Heisenberg algebra,
$a(z)a^*(w) \sim \kappa/(z-w)^2$ at every genus. What changes at genus~$g$ is the
\emph{propagator} (i.e., the Green's function on $\Sigma_g$), which acquires regular
terms depending on the period matrix $\Omega \in \mathcal{H}_g$.
\end{remark}

After regularization, the genus~2 propagator takes the form:
\[
G^{(2)}(z,w) = \frac{\kappa}{(z-w)^2}
+ \kappa \sum_{\alpha,\beta=1}^{2} \omega_\alpha(z)\,(\operatorname{Im}\Omega)^{-1}_{\alpha\beta}\,
\overline{\omega}_\beta(w) + \text{(regular)}
\]
where $\omega_1, \omega_2$ are the normalized abelian differentials of the first kind on $\Sigma_2$.
The regular part depends on the Siegel modular forms
\begin{align}
E_k(\Omega) &= \sum_{\substack{T = \left(\begin{smallmatrix} n & r/2 \\ r/2 & m \end{smallmatrix}\right) \geq 0 \\ n, m \in \mathbb{Z}_{\geq 0},\, r \in \mathbb{Z}}} a_k(T)\, e^{2\pi i \,\mathrm{tr}(T \Omega)}
\end{align}
where $T$ ranges over positive semi-definite half-integral $2 \times 2$ matrices, and the Fourier coefficients $a_k(T)$ are given by sums of divisor functions generalizing the genus-1 formulas. Explicitly, $e^{2\pi i\,\mathrm{tr}(T\Omega)} = q_1^n q_{12}^r q_2^m$ with $q_i = e^{2\pi i \Omega_{ii}}$ and $q_{12} = e^{2\pi i \Omega_{12}}$.

\subsection{Interpretation}

\subsubsection{Algebraic meaning}

The commutation relation $[a_m, a^*_n] = \kappa m \delta_{m+n,0}$ is an algebraic property
of the Heisenberg vertex algebra and does not change with genus. What the genus~2
geometry modifies is the \emph{bar complex differential}: the propagator's regular part
introduces new terms in the bar differential, giving rise to a \emph{curved}
$A_\infty$ structure with curvature $m_0 \neq 0$ controlled by the Arakelov form:
\[m_0^{(2)} = \kappa \cdot \omega_{\mathrm{Ar}}^{(2)}\]
where $\omega_{\mathrm{Ar}}^{(2)} = \frac{i}{2}\sum_{\alpha,\beta=1}^{2}(\operatorname{Im}\Omega)^{-1}_{\alpha\beta}\,\omega_\alpha\wedge\bar\omega_\beta$ is the genus-$2$ Arakelov form
(Theorem~\ref{thm:quantum-arnold-relations}).
The curvature is \emph{linear} in~$\kappa$, not quadratic.

\subsubsection{Geometric meaning}

The Siegel Eisenstein series $E_4,E_6$ are the first even-weight
automorphic coefficients on $\mathbb H_2$; the full scalar ring of
genus-$2$ Siegel modular forms also contains the Igusa cusp-form
generators. The tautological Chern classes on $\mathcal M_2$ are the
Hodge classes $\lambda_i=c_i(\mathbb E)$, not the Eisenstein series
themselves. The genus-$2$ bar complex
$C_{\bullet}^{(2)}(\mathcal{A})$ is a sheaf on $\mathcal{M}_2$;
pulling back along $\mathcal{M}_{2,2} \to \mathcal{M}_2$ gives the
OPE corrections.

\subsubsection{Physical meaning}

In CFT, $Z_2 = \int_{\mathcal{M}_2} |\text{det Im } \Omega|^{-c/2}$,
$\langle a(z) a^*(w) \rangle_2 \propto |E(z,w)|^{-2\Delta}$,
and the OPE is the operator limit $z \to w$ of this correlator.

The $E_4, E_6$ terms are \emph{two-loop quantum corrections} to the bar complex curvature,
appearing in the regular part of the propagator and manifesting as curved $A_\infty$ structure.

\subsection{Generalization to higher weight operators}

\subsubsection{Virasoro at genus 2}

The Virasoro OPE $T(z) T(w) \sim \frac{c/2}{(z-w)^4} + \frac{2T(w)}{(z-w)^2} + \frac{\partial T(w)}{z-w}$
is genus-independent (the pole orders are fixed by $h_T = 2$). At genus~2, the
two-point \emph{correlation function} acquires modular corrections:
\[\langle T(z) T(w) \rangle_2 = \frac{c/2}{(z-w)^4} + f(\Omega) + O(z-w)\]
where $f(\Omega) = c \cdot E_4(\Omega) + \cdots$ is a regular function encoding the
genus~2 geometry. In the bar complex, these corrections contribute to the
curvature $m_0^{(2)}$ and to the higher operations $m_k^{(2)}$.

\subsubsection{\texorpdfstring{$W$-algebras at genus 2}{$\mathcal{W}$-algebras at genus 2}}

Following Arakawa's theory, the genus-$g$ correlation functions of a
$W$-algebra with generators $W^{(k)}$ of weight~$k$ involve Siegel modular
forms $C_j^{(k)}(\Omega)$ depending on the period matrix $\Omega \in \mathcal{H}_g$.
The OPE singularity structure (pole orders $\leq 2k$) remains genus-independent,
but the correlation functions' regular parts depend on $\Omega$.

The \emph{pattern}: genus $g$ introduces Siegel modular forms of weight
$\leq g(g+1)/2$, matching the dimension of the Siegel upper half-space
$\mathcal{H}_g$ (not $\mathcal{M}_g$, whose dimension is $3g-3$; these coincide
for $g \leq 3$ but differ for $g \geq 4$ due to the Schottky problem).

\subsection{The bar complex perspective}

\subsubsection{\texorpdfstring{How this appears in $C_{\bullet}^{(2)}(\mathcal{A})$}{How this appears in C-(2)(A)}}

Define the genus 2 bar complex via:
\[C_n^{(2)}(\mathcal{A}) = \int_{\mathrm{Conf}_n(\Sigma_2)}
\mathcal{A}^{\boxtimes n}
\otimes \Omega^{\bullet}(\mathcal{M}_2)\]

The differential includes:
\begin{enumerate}
\item Bar differential (OPE contractions)
\item Boundary operator (degeneration $\Sigma_2 \rightsquigarrow \Sigma_1$)
\item \emph{New:} Integration over moduli with Eisenstein series insertions
\end{enumerate}

\subsubsection{The cocycle}

The genus 2 cocycle for our example is:
\begin{align}
c_2 &= \int_{\mathcal{M}_2} \int_{\Sigma_2^2}
\operatorname{Tr}_{\Sigma_2}(a(z_1) \otimes a^*(z_2)) \\
&\quad \cdot E_4(\Omega) \cdot d\mu_{\text{WP}}
- \kappa \cdot (\text{boundary terms})
\end{align}

\emph{Cocycle condition.} $d^{(2)} c_2 = 0$ involves:
\begin{itemize}
\item Genus 1 boundary: $\partial \Sigma_2 \supset \Sigma_1$
\item Separating degeneration: $\Sigma_2 \rightsquigarrow \Sigma_1 \cup \Sigma_1$
\item Non-separating degeneration: $\Sigma_2 \rightsquigarrow \Sigma_1^{\text{nodal}}$ (self-sewing on a genus-$1$ curve)
\end{itemize}

Each boundary contribution cancels by the Holomorphic Anomaly Equation of BCOV theory.

\subsection{Computational summary}

To compute genus~$2$ corrections $a(z) \cdot b(w)$ for vertex operators $a, b$:
\begin{enumerate}
\item Draw all 2-loop Feynman diagrams with external legs at $z, w$.
\item Assign the genus-2 propagator $G_{\Sigma_2}(z_i, z_j)$ to each internal line.
\item Integrate over $\Sigma_2$ using theta function identities and residues.
\item Regularize via the holomorphic anomaly equation and minimal subtraction.
\item Integrate over $\mathcal{M}_2$ by expanding in Eisenstein series.
\item Extract the OPE by taking the $z \to w$ limit and expanding in $(z-w)^{-k}$.
\end{enumerate}
The bar complex curvature contributes corrections proportional to $\kappa(\cA) \cdot \lambda_2$
(linear in~$\kappa$, not quadratic; see Theorem~\ref{thm:genus-universality}).

\subsection{Connection to string theory}

The genus-$2$ OPE corrections have a string-theoretic interpretation: $\Sigma_2$ is the worldsheet, the amplitude $\langle V_a(z) V_{a^*}(w) \rangle_{\Sigma_2}$ is the genus-$2$ string amplitude, the OPE is the factorization limit $z \to w$, and the Eisenstein series arise from summing over intermediate states.

\begin{remark}[Formality]
The genus-$2$ bar-cobar complex realizes the formality theorem: the induced $\star$-product is the quantization of the classical OPE Poisson structure, with corrections given by Eisenstein series
(a consequence of $E_2$ collapse on the Koszul locus; see Proposition~\ref{prop:e2-collapse-formality}).
\end{remark}

\begin{remark}[Higher genera]
At genus $g \geq 3$, modular forms of weight $\leq g(g+1)/2$ appear,
with relations from gluing equations along boundary strata of
$\overline{\mathcal{M}}_g$. Genus-$0$ OPE data and modular covariance
determine the scalar sewing constraints, but the stable-graph residues,
period correction, and shadow obstruction tower are additional
positive-genus data. See Remark~\ref{rem:parameter-source} for the
precise reformulation in terms of the parameter-source chain.
\end{remark}

\begin{remark}[Separating-node scalar factorisation of the Igusa cusp form]
\label{rem:hgf-siegel-sepnode}
The genus-$2$ cocycle $c_2$ above is integrated against the Weil--Petersson measure on $\mathcal{M}_2$, and its scalar modular target is the unscaled Igusa cusp form $\Phi_{10}^{\mathrm{un}}(Z)$ of weight~$10$ on $\mathbb{H}_2$. At a \emph{separating} node write
\[
Z=\begin{pmatrix}\tau&z\\z&\rho\end{pmatrix},\qquad z\to0 .
\]
The scalar theorem is the leading boundary expansion
\[
\Phi_{10}^{\mathrm{un}}(Z)
\;=\;
C_{\mathrm{sep}}\,(2\pi i z)^2\,
\eta^{24}(\tau)\eta^{24}(\rho)
\;+\;
O(z^4),
\qquad C_{\mathrm{sep}}\in\mathbb C^\times,
\]
where the nonzero scalar $C_{\mathrm{sep}}$ is fixed by the chosen
normalisation of $\Phi_{10}^{\mathrm{un}}$. Equivalently,
\[
\frac{1}{\Phi_{10}^{\mathrm{un}}(Z)}
\;=\;
\frac{C_{\mathrm{sep}}^{-1}}{(2\pi i z)^2\eta^{24}(\tau)\eta^{24}(\rho)}
\;+\;
\text{terms regular at }z=0.
\]
The Jacobi cusp form $\phi_{10,1}$ belongs to the Fourier--Jacobi
expansion at a non-separating cusp; it is not a chain-level source at
the diagonal $z=0$. Thus the separating-degeneration term in
$d^{(2)}c_2=0$ inherits a scalar automorphic boundary factorisation
from the Igusa form. A lift of this scalar boundary identity to a
K3 Hall or CoHA source requires the finite Hall--Borcherds
source-recognition datum of Definition~\ref{def:hgf-finite-hall-borcherds-source}.
Without that datum the displayed expansion is a target identity on
$\mathbb H_2$, not a product decomposition of a compact
$K3$-elliptic Hall object with $24$ nodal Kodaira fibres. See
Part~\ref{part:modular-tower-physics} for the physics/K3 side.
\end{remark}

\section{The fundamental theorem: genus-zero base case}
\label{sec:fundamental-theorem-koszul}
% Theorem A now lives in chiral_koszul_pairs.tex

\begin{remark}[Summary of Theorem~A; see Theorem~\ref{thm:bar-cobar-isomorphism-main}]
\label{rem:theorem-a-summary}
The genus-zero bar-cobar duality
(Theorem~\ref{thm:bar-cobar-isomorphism-main},
Chapter~\ref{chap:koszul-pairs}) establishes:
for a chiral Koszul pair $(\cA_1, \cA_2)$
(Definition~\ref{def:chiral-koszul-pair}),
the unit and counit of the bar-cobar adjunction are
quasi-isomorphisms, and the bar functor is intertwined
with Verdier duality on $\operatorname{Ran}(X)$ through
$K_X(\cA_i)=\mathbb D_{\Ran}\bar B_X(\cA_i)$ and the pair maps
$\nu_{ij}$.  The counit is a quasi-isomorphism on the Koszul locus;
the completed coderived category carries its general comparison cone.
\end{remark}

% The following theorem block is retained only to preserve the local
% equation labels used in this chapter. The canonical statement and proof
% are in chiral_koszul_pairs.tex.
\begin{theorem}[Typed bar--cobar and Verdier comparisons; \ClaimStatusConditional]
\label{thm:bar-cobar-isomorphism-main-equations}
Let $(\mathcal{A}_1, \mathcal{A}_2)$ be a finite-type or continuously
dualizable chiral Koszul pair on a smooth curve~$X$. Choose quadratic
presentations with coalgebras $\mathcal A_i^{\mathrm i}$ and
comparisons~$q_i$; explicitly, if
$\mathcal A_i=T_X(V_i)/(R_i)$, then
\[
\mathcal A_i^{\mathrm i}=C_X(s^{-1}V_i,s^{-2}R_i),
\qquad
q_i\colon\mathcal A_i^{\mathrm i}\longrightarrow
\bar B_X(\mathcal A_i).
\]

\medskip
\noindent\emph{I. Quadratic coalgebras compare with the full bars}
\begin{align}
q_1\colon\mathcal A_1^{\mathrm i}
 &\xrightarrow{\sim}\bar{B}^{\text{ch}}(\mathcal{A}_1)
 \label{eq:bar-A1-is-A2-dual} \\
q_2\colon\mathcal A_2^{\mathrm i}
 &\xrightarrow{\sim}\bar{B}^{\text{ch}}(\mathcal{A}_2)
 \label{eq:bar-A2-is-A1-dual}
\end{align}

\medskip
\noindent\emph{II. Quadratic cobar reconstructs each algebra}
\begin{align}
\Omega^{\text{ch}}(\mathcal{A}_1^{\mathrm i})
 &\xrightarrow{\sim}\mathcal{A}_1
 \label{eq:cobar-A2-dual-is-A1} \\
\Omega^{\text{ch}}(\mathcal{A}_2^{\mathrm i})
 &\xrightarrow{\sim}\mathcal{A}_2
 \label{eq:cobar-A1-dual-is-A2}
\end{align}

\medskip
\noindent\emph{III. Full bar--cobar reconstruction}
\begin{align}
\varepsilon_1\colon
\Omega^{\text{ch}}(\bar{B}^{\text{ch}}(\mathcal{A}_1))
 &\xrightarrow{\sim}\mathcal{A}_1
 \label{eq:bar-cobar-resolution-A1} \\
\varepsilon_2\colon
\Omega^{\text{ch}}(\bar{B}^{\text{ch}}(\mathcal{A}_2))
 &\xrightarrow{\sim}\mathcal{A}_2
 \label{eq:bar-cobar-resolution-A2}
\end{align}

\medskip
\noindent\emph{IV. The unit on the quadratic coalgebras}
\begin{align}
\eta_1\colon\mathcal A_1^{\mathrm i}
 &\xrightarrow{\sim}
\bar{B}^{\text{ch}}\Omega^{\text{ch}}(\mathcal A_1^{\mathrm i}),\\
\eta_2\colon\mathcal A_2^{\mathrm i}
 &\xrightarrow{\sim}
\bar{B}^{\text{ch}}\Omega^{\text{ch}}(\mathcal A_2^{\mathrm i}).
\end{align}

\medskip
\noindent\emph{V. Verdier Intertwining}

Verdier duality sends each bar coalgebra to the algebra
$K_X(\mathcal A_i):=\mathbb D_{\operatorname{Ran}}
\bar B_X(\mathcal A_i)$. Contravariance gives
\begin{equation}\label{eq:verdier-intertwining-main}
\mathbb D(q_i)\colon K_X(\mathcal A_i)
\xrightarrow{\sim}\mathbb D_{\operatorname{Ran}}
(\mathcal A_i^{\mathrm i}).
\end{equation}
The pair datum supplies the functorial algebra equivalences
\[
\nu_{12}\colon K_X(\mathcal A_1)\xrightarrow{\sim}\mathcal A_2,
\qquad
\nu_{21}\colon K_X(\mathcal A_2)\xrightarrow{\sim}\mathcal A_1
\]
over $\overline{\mathcal{M}}_{g,n}$
\textup{(}Theorem~\textup{\ref{thm:verdier-bar-cobar})}.
The Fulton--MacPherson forms and residue operations of
Parts~\ref{part:foundations-beilinson-tower}--\ref{part:modular-tower-physics} are one concrete chart on this Ran-space\slash Verdier phenomenon.
\end{theorem}

\begin{proof}
Quadratic Koszulness gives the two maps~$q_i$ and their adjoint cobar
equivalences. The enhanced bar--cobar theorem gives the units and
counits. Verdier duality reverses~$q_i$, while the two maps~$\nu_{ij}$
are the chosen-pair comparisons.
\end{proof}

\begin{remark}[Homotopy-native formulation]\label{rem:homotopy-native-a}
The bar and cobar functors lift to an $\infty$-categorical adjunction
\[
\Omega^{\mathrm{ch}} \colon
 \CoAlg_{\mathrm{conil}}(\cat{Fact}(X))
 \rightleftarrows
 \Alg_{\mathrm{aug}}(\cat{Fact}(X))
 \,:\!\B^{\mathrm{ch}}
\]
(Lurie~\cite[Chapter~5]{HA}, Ayala--Francis~\cite{AF15}).
On the Koszul locus, unit and counit are homotopy equivalences of
$\Ainf$-chiral algebras; Verdier intertwining lifts to an
$\infty$-functor exchanging bar and cobar categories.
\end{remark}

\begin{corollary}[Hochschild cohomology duality; \ClaimStatusConditional]
\label{cor:hochschild-duality}
\label{cor:chiral-hochschild-duality}
For a chiral Koszul pair $(\mathcal{A}, \mathcal{A}^!)$ on a smooth projective curve~$X$,
chiral Hochschild cohomology satisfies:
\begin{equation}\label{eq:hochschild-duality-canonical}
RHH_{\mathrm{chiral}}(\mathcal{A})
\;\simeq\;
R\mathrm{Hom}\bigl(RHH_{\mathrm{chiral}}(\mathcal{A}^!),\,
\omega_X[2]\bigr).
\end{equation}
On cohomology, the $[2]$-shift becomes degree reflection:
\begin{equation}\label{eq:hochschild-duality-groups}
\ChirHoch^n(\mathcal{A})
\;\cong\; \ChirHoch^{2-n}(\mathcal{A}^!)^{\vee}
\otimes \omega_X.
\end{equation}
\end{corollary}

\begin{proof}
This is Theorem~\ref{thm:main-koszul-hoch} (Theorem~H), proved
independently in \S\ref{subsec:hochschild-duality} via
the bar-cobar quasi-isomorphism and Verdier duality on the Ran space.
\end{proof}

\begin{definition}[Bigraded chiral Hochschild complex]
\label{def:bigraded-hochschild}
\index{Hochschild cohomology!bigraded|textbf}
The \emph{bigraded chiral Hochschild complex} of a chiral
algebra~$\cA$ on a smooth curve~$X$ is the double complex
\[
CH^{p,\bullet}_{\mathrm{ch}}(\cA) :=
R\mathrm{Hom}_{\mathcal{D}_{\overline{C}_{p+2}(X)}}
\bigl(\cA^{\boxtimes(p+2)}, \omega_{\overline{C}_{p+2}(X)}\bigr)
\]
at bar degree~$p$, with the total object
\[
RHH_{\mathrm{ch}}(\cA) :=
\operatorname{Tot}\!\left(
\bigoplus_{p \geq 0} CH^{p,\bullet}_{\mathrm{ch}}(\cA)[-p]
\right).
\]
At bar degree~$p$, the configuration space
$\overline{C}_{p+2}(X)$ has complex dimension~$p + 2$
\textup{(}since $\dim_\C X = 1$\textup{)}, and Verdier duality
contributes an ambient dualizing-complex shift.  Its use in
Theorem~H requires the additional Verdier--Kashiwara degree
normalization package of
Definition~\ref{def:theorem-h-shift-normalization-package}; the
displayed totalization shift is a regrading convention, not by itself
a proof of the uniform Hochschild duality shift.
\end{definition}

\begin{remark}[Origin of the shift by~\texorpdfstring{$2$}{2}]\label{rem:hochschild-shift-origin}
The uniform shift $[2]$ in~\eqref{eq:hochschild-duality-canonical}
arises from the bigraded structure of
Definition~\ref{def:bigraded-hochschild}, the collision-depth
collapse, and the shift-normalization package of
Definition~\ref{def:theorem-h-shift-normalization-package}.  The
load-bearing assertion is that the ambient Verdier shift, the
Kashiwara offset of the full-collision stratum, the Arnold fibre
degree after localized bar concentration, and the totalization
regrading together identify the surviving bar-degree-\(p\) term with
the base-curve shift \([2]\), independently of~\(p\).  The formal
degree count is stated in Lemma~\ref{lem:hochschild-shift-computation}
(Chapter~\ref{chap:deformation-theory}); it is not the bare slogan
that a variable ambient shift cancels because a survivor lies on a
diagonal.
\end{remark}

% ================================================================
% SECTION 4.10: HIGHER GENUS CONFIGURATION SPACES - COMPLETE THEORY
% ================================================================

\section{Higher-genus configuration spaces}
\label{sec:higher-genus-config-complete}

\subsection{The genus stratification philosophy}

\begin{remark}[Genus as quantum number]
\label{rem:genus-quantum}
The genus $g$ organizes corrections: at genus~$0$ the bar differential satisfies~$d^{(0)2} = 0$ exactly (the Arnold relation ensures nilpotency on the curve; on the formal disk this recovers the classical statement), genus~$1$ produces central extensions, and genus $g \geq 2$ introduces Feynman-transform modular-operad structure (sense (iii) of T15: genus-raising $\Sigma_g \leadsto \Sigma_{g+1}$ via modular composition, not sense (ii) MTC).

This parallels the loop expansion in quantum field theory:
\begin{equation}
Z = Z_{\text{tree}} + \hbar Z_{\text{1-loop}} + \hbar^2 Z_{\text{2-loop}} + \cdots
\end{equation}
with $g$ playing the role of loop number.
\end{remark}

\subsection{Configuration spaces at arbitrary genus}

\begin{definition}[Higher-genus configuration space]
\label{def:higher-genus-config-space}
Let $\Sigma_g$ be a closed smooth Riemann surface of genus $g$. The
ordered $n$-point configuration space is:
\begin{equation}
C_n(\Sigma_g) = \{(p_1, \ldots, p_n) \in \Sigma_g^n : p_i \neq p_j \text{ for } i \neq j\}
\end{equation}

The Fulton--MacPherson compactification
$\overline{C}_n(\Sigma_g)$
\textup{(}Theorem~\textup{\ref{thm:FM-functorial}}\textup{)}
is constructed by iteratively blowing up the diagonals
$\Delta_I=\{p_i=p_j:i,j\in I\}$ inside the fixed product
$\Sigma_g^n$. It records collision rates and tangent directions by
exceptional divisors $D_I$ with normal-crossing incidence. The curve
$\Sigma_g$ does not degenerate in this fixed-fiber compactification.

The fixed-fiber boundary stratification consists of:
\begin{itemize}
\item \emph{Collision divisors:} $D_{ij}$ where $p_i \to p_j$ on the same component
\item \emph{Nested collision strata:} $D_{I_1}\cap\cdots\cap D_{I_r}$ for nested or disjoint collision subsets
\end{itemize}
Separating and non-separating node divisors live instead on the
Deligne--Mumford compactification
$\overline{\mathcal M}_{g,n}$, or on the relative
Fulton--MacPherson compactification of the universal stable curve.
\end{definition}

\begin{remark}[Dimension count]
\label{rem:dimension-higher-genus}
The configuration space has complex dimension:
\begin{equation}
\dim_{\mathbb{C}} C_n(\Sigma_g) = n \cdot \dim \Sigma_g = n
\end{equation}
The moduli stack has dimension:
\begin{equation}
\dim_{\mathbb{C}} \overline{\mathcal{M}}_{g,n} = 3g - 3 + n
\end{equation}
for $2g-2+n>0$. Over the smooth locus $\mathcal M_g$, the relative
ordered configuration space has fiber $C_n(\Sigma_g)$ and total
dimension $3g-3+n$. Its stable compactification is not the same
object as the fixed-fiber Fulton--MacPherson compactification:
$\overline{\mathcal M}_{g,n}$ compactifies the moving pointed curve,
whereas $\overline C_n(\Sigma_g)$ compactifies collisions on one
fixed smooth curve. The two coincide only after specifying the
relative universal compactification and restricting to a smooth
fiber.
\end{remark}

\subsection{\texorpdfstring{The moduli space $\overline{\mathcal{M}}_{g,n}$}{The moduli space M-g,n}}

\begin{definition}[Deligne--Mumford compactification]
\label{def:deligne-mumford-compactification}
\index{Deligne--Mumford compactification|textbf}
\index{moduli space!of curves|textbf}
The moduli space $\overline{\mathcal{M}}_{g,n}$ parametrizes stable $n$-pointed curves of genus $g$:
\begin{equation}
[\Sigma_g; p_1, \ldots, p_n] \in \overline{\mathcal{M}}_{g,n}
\end{equation}
where stability requires $\Sigma_g$ to be a connected nodal curve with $2g_i - 2 + n_i > 0$ on every component $C_i$ and finite automorphism group.
\end{definition}

\begin{theorem}[Structure of \texorpdfstring{$\overline{\mathcal{M}}_{g,n}$}{M-bar(g,n)} {\cite{DeligneM69,Knudsen83}}; \ClaimStatusProvedElsewhere]
\label{thm:moduli-structure}
\index{normal crossing!boundary}
The Deligne--Mumford compactification satisfies:
\begin{enumerate}
\item $\overline{\mathcal{M}}_{g,n}$ is a proper Deligne--Mumford stack of dimension $3g-3+n$
\item The interior $\mathcal{M}_{g,n}$ parametrizes smooth curves (smooth Riemann surfaces)
\item The boundary $\partial \overline{\mathcal{M}}_{g,n}$ is a normal crossing divisor
\item Each boundary stratum corresponds to a dual graph $\Gamma$
\end{enumerate}
\end{theorem}

\begin{proof}
By Deligne--Mumford~\cite{DeligneM69} and Knudsen~\cite{Knudsen83}:
(1) Properness follows from stable reduction: any family of smooth curves over a punctured disk extends uniquely to a stable curve over the closed disk.
(2) The interior $\mathcal{M}_{g,n}$ is smooth, with local coordinates provided by Teichm\"uller theory (quadratic differentials).
(3) Boundary strata correspond to degenerations: separating nodes ($\Sigma_g \to \Sigma_{g_1} \cup \Sigma_{g_2}$) and non-separating nodes (pinching a cycle), classified by dual graphs $\Gamma$.
(4) Local models near boundary divisors are products of smooth divisors, giving normal crossing structure \cite{Knudsen83}.
\end{proof}

\subsection{Fibration structure}

\begin{theorem}[Universal curve fibration {\cite{Knudsen83}}; \ClaimStatusProvedElsewhere]
\label{thm:universal-curve-fibration}
There exists a universal curve, realized by the forgetful morphism:
\begin{equation}
\pi: \overline{\mathcal{M}}_{g,n+1} \to \overline{\mathcal{M}}_{g,n}
\end{equation}
Write $\overline{\mathcal C}_{g,n}:=\overline{\mathcal M}_{g,n+1}$
for this universal curve. It satisfies:
\begin{itemize}
\item The fiber over a stable pointed curve
$[(C;p_1,\ldots,p_n)]$ is the curve $C$ itself, with the
stabilization convention at components made unstable by the forgotten
mark
\item Sections $\sigma_i: \overline{\mathcal{M}}_{g,n} \to \overline{\mathcal{C}}_{g,n}$ give the marked points
\item The relative dualizing sheaf $\omega_\pi = \omega_{\overline{\mathcal{C}}_{g,n}/\overline{\mathcal{M}}_{g,n}}$ is relatively ample
\end{itemize}

Over the open smooth locus, ordered configurations of additional
points sit inside the fiber powers of this universal curve after
removing diagonals and the marked sections:
\begin{equation}
\operatorname{Conf}^{\mathrm{ord}}_m(\mathcal C_{g,n}/\mathcal M_{g,n})
\subset
\mathcal C_{g,n}^{\,m}\setminus
\bigl(\Delta\cup\sigma_1\cup\cdots\cup\sigma_n\bigr).
\end{equation}
The compactified configuration carrier used by the bar complex is the
relative Fulton--MacPherson compactification of these ordered
configuration spaces; its boundary contains collision faces over the
interior and stable-curve degeneration faces over
$\partial\overline{\mathcal M}_{g,n}$.
\end{theorem}

\subsection{Logarithmic forms at higher genus}

\begin{definition}[Higher-genus logarithmic forms]
\label{def:higher-genus-log-forms}
\index{propagator!genus $g$|textbf}
On the fixed-fiber compactification $\overline{C}_n(\Sigma_g)$, the
propagator forms are logarithmic along collision divisors and
single-valued only after period normalization:
\begin{equation}
\eta_{ij}^{(g)}
= \partial_{p_i}\log E(p_i,p_j)
+ \textup{Arakelov period correction}.
\end{equation}
Here $E(p,q)$ is the prime form on $\Sigma_g$ (the genus-$g$
replacement for $p-q$ on an affine screen), and the correction is
written using a symplectic basis of $H_1(\Sigma_g,\mathbb Z)$, the
period matrix, and the normalized abelian differentials. These forms
represent the analytic prime-form lane; the rational cohomology of
$C_n(\Sigma_g)$ for $g\ge1$ is governed by the
Totaro--K\v{r}\'{\i}\v{z}--Bezrukavnikov diagonal model, not by an
Arnold-only presentation.
\end{definition}

The explicit form depends on the genus:

\emph{Genus 0 (Rational Curve).}
\begin{equation}
\eta_{ij}^{(0)} = d\log(z_i - z_j) = \frac{dz_i - dz_j}{z_i - z_j}
\end{equation}
No global obstructions.

\emph{Genus 1 (Elliptic Curve $E_\tau = \mathbb{C}/(\mathbb{Z} + \tau\mathbb{Z})$):}
\begin{equation}
\eta_{ij}^{(1)}
=
\left[
\partial_u\log\theta_1(u\mid\tau)
+2\pi i\,\frac{\operatorname{Im}u}{\operatorname{Im}\tau}
\right]_{u=z_i-z_j}
(dz_i-dz_j).
\end{equation}
The theta quotient alone is quasi-periodic; the displayed
real-analytic coefficient is the Arakelov correction that descends the
$(1,0)$-form to the elliptic curve.

\emph{Genus $g \geq 2$ (Hyperbolic Case).}
\begin{equation}
\eta_{ij}^{(g)} = \partial_{p_i}\log E(p_i, p_j) + \pi\sum_{\alpha, \beta = 1}^g \omega_\alpha(p_i) \,(\operatorname{Im}\Omega)_{\alpha\beta}^{-1}\, \operatorname{Im}\!\left(\int_{p_0}^{p_j}\omega_\beta\right)
\end{equation}
where $\partial_{p_i}$ differentiates in the first variable (giving a $(1,0)$-form in $p_i$, scalar in $p_j$) and:
\begin{itemize}
\item $E(p, q)$ is the prime form (a $(-\frac{1}{2}, -\frac{1}{2})$-differential with a simple zero along the diagonal)
\item $\omega_1, \ldots, \omega_g$ are the normalized holomorphic 1-forms ($\oint_{A_\alpha}\omega_\beta = \delta_{\alpha\beta}$)
\item $\Omega_{\alpha\beta} = \oint_{B_\beta} \omega_\alpha$ is the period matrix
\item The second term is a non-holomorphic correction ensuring double periodicity
\end{itemize}

\begin{remark}[Physical interpretation]
\label{rem:physical-log-forms}
In CFT, $\eta_{ij}^{(0)}$ gives tree-level propagators, $\eta_{ij}^{(1)}$ gives one-loop propagators via theta functions, and $\eta_{ij}^{(g)}$ for $g \geq 2$ gives multi-loop propagators.
\end{remark}

\subsection{Arnold relations at higher genus}

\begin{theorem}[Quantum-corrected Arnold relations; \ClaimStatusConditional]
\label{thm:quantum-arnold-relations}
\label{thm:arnold-higher-genus}
\label{thm:arnold-quantum}
\index{Arnold relations!higher genus}
\textup{[Type signature: Open quadrant; fixed-fiber Arakelov-current
presentation; level $2$; package $H_D^{\mathrm{fib}}$ consisting of
an Arakelov--Green propagator, logarithmic-current extension, and the
normalized collision trace.]}
Let $\eta_{ij}^{(0)}=d\log(z_i-z_j)$ denote the affine
genus-$0$ logarithmic form on a tangent screen, and let
$\eta_{ij}^{(g)}$ for $g\ge1$ denote an Arakelov-normalized
prime-form propagator on a smooth curve $\Sigma_g$.
Define the cyclic Arnold form
\begin{equation}
\mathcal{A}_3^{(g)} = \eta_{12}^{(g)} \wedge \eta_{23}^{(g)} + \eta_{23}^{(g)} \wedge \eta_{31}^{(g)} + \eta_{31}^{(g)} \wedge \eta_{12}^{(g)}
\end{equation}
on the open configuration locus, interpreted as a logarithmic current
after extension to the Fulton--MacPherson boundary. Then:
\begin{enumerate}
\item[\textup{(a)}] \textup{[\ClaimStatusProvedHere]}
On the affine genus-$0$ screen,
$\mathcal{A}_3^{(0)}=0$. This is Arnold's braid-arrangement
relation and is the local collision identity behind
$\dzero^2=0$.

\item[\textup{(b)}] \textup{[\ClaimStatusConditional]}
On an elliptic curve $E_\tau$, the holomorphic theta quotient is
quasi-periodic. After adding the Arakelov correction, the cyclic
defect is the normalized elliptic volume class:
\[
[\mathcal{A}_3^{(1)}]
=
2\pi i\,[\omega_{\tau}],
\qquad
\omega_{\tau}=
\frac{i}{2\,\operatorname{Im}\tau}\,dz\wedge d\bar z .
\]
The equality is a de~Rham/current identity after the
Arakelov--Green representative is fixed; without that representative
only the class is canonical.

\item[\textup{(c)}] \textup{[\ClaimStatusProvedElsewhere]}
For every positive genus, the global rational model of
$C_n(\Sigma_g)$ is the
Totaro--K\v{r}\'{\i}\v{z}--Bezrukavnikov diagonal CDGA:
\[
d\eta_{ij}=\Delta_{ij},\qquad
\eta_{ij}\bigl(p_i^*\xi-p_j^*\xi\bigr)=0
\quad(\xi\in H^*(\Sigma_g)).
\]
Thus the affine Arnold relation governs each genus-$0$ tangent screen,
while the diagonal CDGA supplies the full positive-genus presentation.

\item[\textup{(d)}] \textup{[\ClaimStatusConditional]}
For $g\ge1$, Fay's identity kills the holomorphic tangent-screen
cyclic sum, and the period correction contributes the canonical
Arakelov class
\[
[\mathcal{A}_3^{(g)}]
=2\pi i\,[\omega_g^{\mathrm{Ar}}].
\]
Contracting this class with the degree-two OPE coefficient of
$\cA$ and applying the normalized collision trace gives the
fiberwise curvature class
\[
[\operatorname{tr}_{\mathrm{diag}}(m_0^{(g)})]
=
\kappa(\cA)[\omega_g^{\mathrm{Ar}}].
\]
Under~$H_D^K(\cA,g)$, the associated scalar virtual object and Hodge
shadow are
$\mathfrak O_g^K(\cA)=\kappa(\cA)\lambda_{-1}(\mathbb E_g)$ and
$\operatorname{obs}^{\mathrm{Hdg}}_g(\cA)
=(-1)^g\kappa(\cA)\lambda_g$.
\end{enumerate}
\end{theorem}

\begin{proof}
Part~\textup{(a)} is Arnold's relation on
$C_n(\mathbb A^1)$, recalled in
Appendix~\ref{app:arnold-relations}. Part~\textup{(c)} is
Totaro's diagonal model for configuration spaces of smooth
projective varieties, specialized to curves, together with the
equivalent models of K\v{r}\'{\i}\v{z} and Bezrukavnikov; see
Proposition~\ref{prop:arnold-higher-genus}.

For $g=1$, the quotient
$\theta_1'(z_i-z_j|\tau)/\theta_1(z_i-z_j|\tau)$ has
$B$-period monodromy $-2\pi i$ in the coefficient of
$du_{ij}=dz_i-dz_j$. The added coefficient
$2\pi i\,\operatorname{Im}(z_i-z_j)/\operatorname{Im}\tau$ shifts by
$+2\pi i$ under $z_i-z_j\mapsto z_i-z_j+\tau$ and multiplies the same
$(1,0)$-form~$du_{ij}$. The correction is regular; the residue one on
the diagonal comes from the theta quotient. Its $\bar\partial$ wedged
with $du_{ij}$ is the smooth term
$-2\pi i\,q_{ij}^{*}\omega_\tau$ in
$d\eta_{ij}^{(1)}=2\pi i(\delta_{D_{ij}}-q_{ij}^{*}\omega_\tau)$, so
the cyclic residue computation gives
$[\mathcal A_3^{(1)}]=2\pi i[\omega_\tau]$.

For $g\ge2$, write
$h_{ij}=\partial_{z_i}\log E(z_i,z_j)$ and let $a_{ij}$ be the
Arakelov period correction. The differentiated Fay identity has, in
genus~$2$, the schematic Weierstrass-matrix form
\begin{equation}\label{eq:fay-g2}
h_{12}h_{23}+h_{23}h_{31}+h_{31}h_{12}
=
\sum_{\alpha,\beta=1}^{2}
\wp_{\alpha\beta}(\,\cdot\,|\Omega)\,
\omega_\alpha\,\omega_\beta
+R_2,
\end{equation}
where the right-hand side is symmetric under cyclic alternation on
the tangent screen. The same differentiated Fay identity holds in
genus~$g$ with the $\wp_{\alpha\beta}$ matrix indexed by
$1\leq\alpha,\beta\leq g$. Hence the holomorphic cyclic part of
$\mathcal A_3^{(g)}$ vanishes on the affine screen.

The non-holomorphic correction satisfies
\[
\bar\partial_{z_j}a_{ij}
=
\frac{\pi}{2i}
\sum_{\alpha,\beta=1}^{g}
\omega_\alpha(z_i)(\operatorname{Im}\Omega)^{-1}_{\alpha\beta}
\bar\omega_\beta(z_j),
\]
the Bergman kernel representing the Arakelov form. The diagonal
residue of $h_{ij}$ contracts this kernel with the OPE
degree-two coefficient. This gives the current class
$[\mathcal A_3^{(g)}]=2\pi i[\omega_g^{\mathrm{Ar}}]$.  The
normalization in~$H_D^{\mathrm{fib}}$ incorporates this universal
$2\pi i$ in the collision trace, and hence
$[\operatorname{tr}_{\mathrm{diag}}(m_0^{(g)})]
=\kappa(\cA)[\omega_g^{\mathrm{Ar}}]$. Choosing the
Arakelov--Green representative upgrades the class statement to the
chain representative used in
Proposition~\ref{prop:genus-g-curvature-package}.
\end{proof}

\begin{corollary}[Universal Arakelov curvature class; \ClaimStatusConditional]
\label{cor:universal-arakelov}
\index{Arakelov form!universality}
For all $g \geq 0$, the quantum Arnold defect has canonical
Arakelov class
\begin{equation}\label{eq:universal-arakelov}
[\mathcal{A}_3^{(g)}]
=
2\pi i\,[\omega_{\mathrm{Ar}}^{(g)}],
\end{equation}
with $[\omega_{\mathrm{Ar}}^{(0)}]=0$. The correction to the
fiberwise square is geometric; the chiral algebra enters only through
the scalar contraction coefficient~$\kappa(\cA)$.
\end{corollary}

\begin{proof}
This is part~\textup{(d)} of
Theorem~\ref{thm:quantum-arnold-relations}, with part~\textup{(a)}
as the genus-$0$ normalization.
\end{proof}

\iffalse

\begin{proof}[Genus $1$ case]
Consider the elliptic curve $E_\tau$ with $\tau \in \mathbb{H}$ (upper half-plane). Use the Weierstrass $\zeta$-function:
\begin{equation}
\zeta(z|\tau) = \frac{1}{z} + \sum_{(m,n) \neq (0,0)} \left[\frac{1}{z - \omega_{mn}} + \frac{1}{\omega_{mn}} + \frac{z}{\omega_{mn}^2}\right]
\end{equation}
where $\omega_{mn} = m + n\tau$.

The quasi-periodicity is:
\begin{align}
\zeta(z + 1|\tau) &= \zeta(z|\tau) + 2\eta_1(\tau)\\
\zeta(z + \tau|\tau) &= \zeta(z|\tau) + 2\eta_\tau(\tau)
\end{align}
with the Legendre relation:
\begin{equation}
\tau \eta_1 - \eta_\tau = \pi i
\end{equation}

The holomorphic theta quotient is quasi-periodic and therefore not globally defined on~$E_\tau$ as a one-form. The single-valued propagator on the elliptic curve is \textup{(}g=1 only, so the period matrix is the $1\times 1$ matrix $\Omega = (\tau)$\textup{)}:
\begin{equation}\label{eq:genus-1-propagator-full}
\eta_{ij}^{(1)}
=
\left[
\partial_u\log\theta_1(u|\tau)
+2\pi i\,\frac{\operatorname{Im}u}{\operatorname{Im}\tau}
\right]_{u=z_{ij}}
(dz_i-dz_j)
\end{equation}
where the real-analytic coefficient
$2\pi i\,\operatorname{Im}(z_{ij})/\operatorname{Im}\tau$ restores
double periodicity by cancelling the $-2\pi i$ theta-quotient
$B$-period.

Write $\eta_{ij}^{(1)} = \phi_{ij}\,(dz_i - dz_j)$ where
$\phi_{ij}=\partial_u\log\theta_1(u|\tau)|_{u=z_{ij}}+
2\pi i\,\operatorname{Im}(z_{ij})/\operatorname{Im}\tau$. Decompose
$\phi_{ij}=h_{ij}+a_{ij}$ into its holomorphic theta-quotient part
and real-analytic Arakelov correction.

Now compute $\mathcal{A}_3^{(1)}$. The pure holomorphic terms $\sum_{\mathrm{cyc}} h_{ij}h_{jk}\,(dz_i - dz_j)\wedge(dz_j - dz_k)$ vanish: expanding, the coefficient functions satisfy the identity
\begin{equation}
\zeta(u)\zeta(v) + \zeta(v)\zeta(w) + \zeta(w)\zeta(u) = \tfrac{1}{2}\bigl[\wp(u) + \wp(v) + \wp(w)\bigr]
\end{equation}
for $u+v+w = 0$ (where $u = z_{12}$, $v = z_{23}$, $w = z_{31}$). Since $\wp$ is an even function and $dz_i \wedge dz_j$ is antisymmetric, symmetry kills the right-hand side when wedged with the form factors.

The surviving cross-terms are $\sum_{\mathrm{cyc}} (h_{ij} a_{jk} + a_{ij} h_{jk})\,(dz_i - dz_j) \wedge (dz_j - dz_k)$. These are $(1,1)$-forms after taking $\bar\partial$ of the real-analytic coefficient. Using
$\bar\partial a_{jk}\wedge dz_{jk}=-2\pi i\,q_{jk}^{*}\omega_\tau$
off the diagonal and summing over cyclic permutations:
\begin{equation}
\mathcal{A}_3^{(1)} = \frac{2\pi i}{2i\,\mathrm{Im}(\tau)} \, dz \wedge d\bar{z} = 2\pi i \cdot \omega_\tau
\end{equation}
where $\omega_\tau = \frac{i}{2\,\mathrm{Im}(\tau)}\,dz \wedge d\bar{z}$ is the normalized volume form on $E_\tau$.
\end{proof}

\begin{proof}[Genus $2$ case]
Fix a genus-$2$ surface $\Sigma_2$ with canonical homology basis
$(A_1, A_2, B_1, B_2)$, period matrix $\Omega \in \mathfrak{H}_2$,
and normalized abelian differentials $\omega_1, \omega_2$.
The single-valued propagator on $\Sigma_2$ is
(Definition~\ref{def:higher-genus-log-forms}):
\begin{equation}\label{eq:propagator-g2-full}
\eta_{ij}^{(2)} = \Bigl[\partial_{z_i}\!\log E(z_i, z_j)
 + \pi \sum_{\alpha,\beta=1}^{2}
 \omega_\alpha(z_i)\,(\operatorname{Im}\Omega)^{-1}_{\alpha\beta}\,
 \operatorname{Im}\!\Bigl(\int_{z_0}^{z_j}\!\omega_\beta\Bigr)
\Bigr](dz_i - dz_j)
\end{equation}
Decompose $\eta_{ij}^{(2)} = (h_{ij} + a_{ij})(dz_i - dz_j)$
into the holomorphic part
$h_{ij} = \partial_{z_i}\!\log E(z_i, z_j)$
and the non-holomorphic correction~$a_{ij}$.

\emph{Step~1: Holomorphic cyclic sum vanishes via the Fay identity.}
The pure holomorphic cyclic sum is
$\sum_{\mathrm{cyc}} h_{ij}\,h_{jk}\,(dz_i - dz_j) \wedge (dz_j - dz_k)$.
The Fay trisecant identity for the prime form on $\Sigma_2$ gives, for
$u + v + w = 0$ in the Jacobian $\mathrm{Jac}(\Sigma_2)$
(where $u = \int_{z_2}^{z_1}\!\omega$, $v = \int_{z_3}^{z_2}\!\omega$,
$w = \int_{z_1}^{z_3}\!\omega$):
\begin{equation}\label{eq:fay-g2-disabled}
h_{12}\,h_{23} + h_{23}\,h_{31} + h_{31}\,h_{12}
= \tfrac{1}{2}\bigl[\wp_{11}(u|\Omega) + \wp_{11}(v|\Omega)
 + \wp_{11}(w|\Omega)\bigr]
 + R(\omega_1, \omega_2)
\end{equation}
where $\wp_{11}(u|\Omega) = -\partial_{u_1}^2\!\log\vartheta[\delta](u|\Omega)$
is the genus-$2$ Weierstrass function (the $(1,1)$-entry of the
Weierstrass matrix $\wp_{\alpha\beta}$), and $R(\omega_1,\omega_2)$
is a symmetric function of the three points.
The identity~\eqref{eq:fay-g2} is the genus-$2$ analog
of the genus-$1$ identity used above: just as
$\zeta(u)\zeta(v) + \mathrm{cyc} = \frac{1}{2}[\wp(u) + \wp(v) + \wp(w)]$
at genus~$1$, the genus-$2$ propagator products satisfy a Fay-type
relation involving the Siegel $\wp$-matrix.

Since $\wp_{11}(u|\Omega)$ is an even function of~$u$ (as the second
derivative of an even theta function) and
$(dz_i - dz_j) \wedge (dz_j - dz_k)$ is antisymmetric under
$z_i \leftrightarrow z_j$, the symmetric terms vanish when wedged
with the form factors, exactly as in the genus-$1$ case.
Therefore the holomorphic cyclic sum is zero.

\emph{Step~2: Surviving cross-terms.}
The non-holomorphic correction $a_{ij}$ has the form
\begin{equation}
a_{ij} = \pi \sum_{\alpha,\beta=1}^{2}
 \omega_\alpha(z_i)\,(\operatorname{Im}\Omega)^{-1}_{\alpha\beta}\,
 \operatorname{Im}\!\Bigl(\int_{z_0}^{z_j}\!\omega_\beta\Bigr)
\end{equation}

Its $\bar\partial$-derivative is
\begin{equation}\label{eq:dbar-a-g2}
\bar\partial_{z_j}\, a_{ij}
= \frac{\pi}{2i} \sum_{\alpha,\beta=1}^{2}
 \omega_\alpha(z_i)\,(\operatorname{Im}\Omega)^{-1}_{\alpha\beta}\,
 \bar\omega_\beta(z_j)
\end{equation}
(using $\bar\partial \operatorname{Im}(\int^{z_j}\omega_\beta)
= \frac{1}{2i}\bar\omega_\beta(z_j)$).

The cross-terms
\[
\sum_{\mathrm{cyc}} (h_{ij}\,a_{jk} + a_{ij}\,h_{jk})
\,(dz_i - dz_j) \wedge (dz_j - dz_k)
\]
are $(1,1)$-forms. Summing cyclically and using the same cancellation
mechanism as at genus~$1$, the holomorphic part pairs with
$\bar\partial a$ to produce a universal $(1,1)$-form:
\begin{equation}\label{eq:arnold-g2-result}
\mathcal{A}_3^{(2)} = 2\pi i \cdot
\frac{i}{2}\sum_{\alpha,\beta=1}^{2}
 (\operatorname{Im}\Omega)^{-1}_{\alpha\beta}\,
 \omega_\alpha \wedge \bar\omega_\beta
= 2\pi i \cdot \omega_{\mathrm{Ar}}^{(2)}
\end{equation}
where
\begin{equation}\label{eq:arakelov-g2}
\omega_{\mathrm{Ar}}^{(2)}
= \frac{i}{2}\sum_{\alpha,\beta=1}^{2}
 (\operatorname{Im}\Omega)^{-1}_{\alpha\beta}\,
 \omega_\alpha \wedge \bar\omega_\beta
\end{equation}
is the Arakelov $(1,1)$-form on $\Sigma_2$, normalized to
$\int_{\Sigma_2} \omega_{\mathrm{Ar}}^{(2)} = 1$.

\emph{Step~3: Consistency.}
At genus~$1$ only, $\Omega = (\tau)$ is the $1\times 1$ period matrix, $\omega_1 = dz$, and
$(\operatorname{Im}\Omega)^{-1}_{11} = 1/\operatorname{Im}(\tau)$.
Formula~\eqref{eq:arakelov-g2} reduces to
$\omega_{\mathrm{Ar}}^{(1)}
= \frac{i}{2\,\operatorname{Im}(\tau)}\,dz \wedge d\bar{z}
= \omega_\tau$,
recovering part~(b).
\end{proof}

\begin{remark}[Summary of the genus-\texorpdfstring{$2$}{2} case]\label{rem:genus2-summary}
The pattern is universal: the holomorphic cyclic sum vanishes by the Fay identity (Step~1), the non-holomorphic correction produces the Arakelov form (Step~2), and cross-terms cancel by type. The only genus-dependent input is the period matrix $\Omega \in \mathfrak{H}_g$.
\end{remark}

\begin{proof}[General case $g \geq 3$]
\index{Arakelov--Green function}
Fix a genus-$g$ surface $\Sigma_g$ with canonical homology basis
$(A_1,\ldots,A_g, B_1,\ldots,B_g)$, period matrix
$\Omega \in \mathfrak{H}_g$ (the Siegel upper half-space),
and normalized abelian differentials $\omega_1,\ldots,\omega_g$
($\oint_{A_\alpha}\omega_\beta = \delta_{\alpha\beta}$).
The single-valued propagator is
(Definition~\ref{def:higher-genus-log-forms}):
\begin{equation}\label{eq:propagator-general-g}
\eta_{ij}^{(g)} = \Bigl[\partial_{z_i}\!\log E(z_i, z_j)
 + \pi \sum_{\alpha,\beta=1}^{g}
 \omega_\alpha(z_i)\,(\operatorname{Im}\Omega)^{-1}_{\alpha\beta}\,
 \operatorname{Im}\!\Bigl(\int_{z_0}^{z_j}\!\omega_\beta\Bigr)
\Bigr](dz_i - dz_j)
\end{equation}
Decompose as $\eta_{ij}^{(g)} = (h_{ij} + a_{ij})(dz_i - dz_j)$
with holomorphic part $h_{ij} = \partial_{z_i}\!\log E(z_i, z_j)$
and non-holomorphic correction
\begin{equation}\label{eq:a-ij-general}
a_{ij} = \pi \sum_{\alpha,\beta=1}^{g}
 \omega_\alpha(z_i)\,(\operatorname{Im}\Omega)^{-1}_{\alpha\beta}\,
 \operatorname{Im}\!\Bigl(\int_{z_0}^{z_j}\!\omega_\beta\Bigr).
\end{equation}

\emph{Step~0: B-cycle monodromy origin of the correction.}
The non-holomorphic term~\eqref{eq:a-ij-general} is not
ad~hoc: it is forced by the B-cycle monodromy of the prime
form. By~\cite{Fay73}, Proposition~2.4, analytic
continuation of $\log E(z_i, z_j)$ around the $\gamma$-th
B-cycle in the variable~$z_j$ shifts
\begin{equation}\label{eq:b-cycle-monodromy-log-E}
\log E(z_i, z_j)\;\Big|_{z_j \to z_j + B_\gamma}
= \log E(z_i, z_j) - 2\pi i\!\int_{z_0}^{z_i}\!\omega_\gamma
 - \pi i\,\Omega_{\gamma\gamma}.
\end{equation}
Differentiating in~$z_i$ gives the monodromy of the
holomorphic propagator:
\begin{equation}\label{eq:h-monodromy}
h_{ij}\big|_{z_j \to z_j + B_\gamma}
= h_{ij} - 2\pi i\,\omega_\gamma(z_i).
\end{equation}
The non-holomorphic correction~$a_{ij}$ is uniquely determined
(up to a global constant) by the requirement that
$\eta_{ij}^{(g)} = (h_{ij} + a_{ij})(dz_i - dz_j)$ be
single-valued: under $z_j \to z_j + B_\gamma$,
$\operatorname{Im}(\int_{z_0}^{z_j}\omega_\beta)$ shifts by
$\operatorname{Im}(\Omega_{\gamma\beta})
= (\operatorname{Im}\Omega)_{\gamma\beta}$, so
\[
a_{ij}\big|_{z_j \to z_j + B_\gamma}
= a_{ij} + \pi \sum_{\alpha,\beta}
 \omega_\alpha(z_i)\,(\operatorname{Im}\Omega)^{-1}_{\alpha\beta}\,
 (\operatorname{Im}\Omega)_{\gamma\beta}
= a_{ij} + \pi\,\omega_\gamma(z_i).
\]
The monodromy of $h_{ij} + a_{ij}$ is
$-2\pi i\,\omega_\gamma + \pi\,\omega_\gamma$, which is
\emph{not} yet zero. Writing $\mathrm{Im}(z_{ij})
= \frac{1}{2i}(\bar{z}_{ij} - z_{ij})$ and using both
slots of the propagator (the antisymmetry
$\eta_{ij}^{(g)} = -\eta_{ji}^{(g)}$ forces
the correction to act symmetrically), the full
Arakelov--Green prescription
\[
a_{ij} = \pi\sum_{\alpha,\beta}\omega_\alpha(z_i)\,
 (\operatorname{Im}\Omega)^{-1}_{\alpha\beta}\,
 \operatorname{Im}\!\Bigl(\int_{z_0}^{z_j}\omega_\beta\Bigr)
\]
compensates the B-cycle shift in~\eqref{eq:h-monodromy}
exactly, producing a single-valued $(1,0)$-form on
$\Sigma_g \times \Sigma_g$ minus the diagonal.
This is the Arakelov--Green propagator
(cf.~\cite{Fay73}, Chapter~III, and \cite{dH89}).

\emph{Step~1: Holomorphic cyclic sum vanishes via the Fay identity.}
The Fay trisecant identity (\cite{Fay73}, Corollary~2.12,
applied in differentiated form) holds on any compact Riemann
surface of genus~$g \geq 1$. For three points $z_1, z_2, z_3$
with Abel--Jacobi images
$u = \int_{z_2}^{z_1}\omega$, $v = \int_{z_3}^{z_2}\omega$,
$w = \int_{z_1}^{z_3}\omega$ satisfying $u + v + w = 0$ in
$\operatorname{Jac}(\Sigma_g)$:
\begin{equation}\label{eq:fay-general}
h_{12}\,h_{23} + h_{23}\,h_{31} + h_{31}\,h_{12}
= \tfrac{1}{2}\sum_{\alpha,\beta=1}^{g}
\bigl[\wp_{\alpha\beta}(u|\Omega) + \wp_{\alpha\beta}(v|\Omega)
 + \wp_{\alpha\beta}(w|\Omega)\bigr]
 \omega_\alpha(z_1)\,\omega_\beta(z_1)
 + R_g(z_1, z_2, z_3)
\end{equation}
where $\wp_{\alpha\beta}(u|\Omega)
= -\partial^2_{u_\alpha u_\beta}\!\log\vartheta[\delta](u|\Omega)$
is the genus-$g$ Weierstrass matrix (for an even
characteristic~$\delta$), and $R_g$ is symmetric in
$(z_1, z_2, z_3)$.

To extract the bar-differential relation, consider the
cyclic-alternating combination
\[
S := \sum_{\mathrm{cyc}(123)}
h_{ij}\,h_{jk}\,(dz_i - dz_j) \wedge (dz_j - dz_k).
\]
Substituting~\eqref{eq:fay-general}: the Weierstrass matrix
term $\wp_{\alpha\beta}(u|\Omega)$ is even in~$u$
(since $\vartheta[\delta]$ is even for an even
characteristic~$\delta$, and $\wp_{\alpha\beta}
= -\partial^2\!\log\vartheta[\delta]$ inherits evenness).
Under the transposition $z_1 \leftrightarrow z_2$,
$u \mapsto -u$ while
$(dz_1 - dz_2) \wedge (dz_2 - dz_3) \mapsto
-(dz_2 - dz_1) \wedge (dz_1 - dz_3)$;
the cyclic sum therefore pairs an even coefficient with
an antisymmetric form factor, and the result vanishes.
The remainder $R_g$ disappears after cyclic alternation
because it is symmetric in the three marked points.

Thus the holomorphic cyclic sum vanishes for every genus
$g \geq 1$.

\emph{Step~2: Quadratic anti-holomorphic terms vanish.}
The pure non-holomorphic terms in $\mathcal{A}_3^{(g)}$ are
\[
Q := \sum_{\mathrm{cyc}} a_{ij}\,a_{jk}\,(dz_i - dz_j)
\wedge (dz_j - dz_k).
\]
Since $a_{ij}$ depends anti-holomorphically on~$z_j$
(through $\operatorname{Im}(\int^{z_j}\omega_\beta)
= \frac{1}{2i}(\int^{z_j}\omega_\beta
 - \overline{\int^{z_j}\omega_\beta})$)
but holomorphically on~$z_i$ (through $\omega_\alpha(z_i)$),
the product $a_{ij}\,a_{jk}$ times $(dz_i - dz_j)
\wedge (dz_j - dz_k)$ involves a factor of
$d\bar{z}_j \wedge d\bar{z}_j = 0$ from the mixed anti-holomorphic
dependence on the shared index~$j$. More precisely,
$a_{ij}$ is smooth (not purely anti-holomorphic), but
$\bar\partial_{z_j}a_{ij}$ is a $(0,1)$-form in~$z_j$
while $\bar\partial_{z_k}a_{jk}$ is a $(0,1)$-form in~$z_k$;
these live on different factors.
The precise vanishing is: in the cyclic term
$a_{12}\,a_{23}\,(dz_1 - dz_2) \wedge (dz_2 - dz_3)$,
the factor $a_{12}$ is $C^\infty$ in $(z_1, z_2)$ and
$a_{23}$ is $C^\infty$ in $(z_2, z_3)$; their product
with the purely holomorphic wedge factor
$(dz_1 - dz_2) \wedge (dz_2 - dz_3)$ is a
$(2,0)$-form in $(z_1, z_2, z_3)$ times a smooth function.
This gives a well-defined $(2,0)$-form on $\Sigma_g^3$,
which is exact in the $(1,1)$-form sense needed for the
Arnold defect: after pulling back to the diagonal
$\Sigma_g \hookrightarrow \Sigma_g^3$, the $Q$ terms
contribute a $(2,0)$-form on~$\Sigma_g$, hence zero
(since $\dim_{\mathbb{C}} \Sigma_g = 1$ and
$(2,0)$-forms vanish on curves).

\emph{Step~3: Cross-terms produce the Arakelov form.}
By Steps~1 and~2, the pure holomorphic and pure
anti-holomorphic contributions to~$\mathcal{A}_3^{(g)}$
vanish. The Arnold defect is therefore carried entirely
by the cross-terms:
\begin{equation}\label{eq:cross-terms-explicit}
\mathcal{A}_3^{(g)} = \sum_{\mathrm{cyc}(123)}
\bigl(h_{ij}\,a_{jk} + a_{ij}\,h_{jk}\bigr)\,
(dz_i - dz_j) \wedge (dz_j - dz_k).
\end{equation}

We determine $\mathcal{A}_3^{(g)}$ by three observations:
(a)~it is proportional to the Arakelov form;
(b)~the proportionality constant is genus-independent;
(c)~the genus-$1$ value is $2\pi i$.

\emph{Observation (a): $\mathcal{A}_3^{(g)}$ is proportional
to $\omega_g$.}
The Arnold form is a globally well-defined smooth
$(1,1)$-form on~$\Sigma_g$: it is single-valued
(built from the single-valued propagator $\eta^{(g)}$)
and smooth away from the pairwise diagonals
(the singularities of $h_{ij} \sim 1/(z_i - z_j)$
cancel in the cyclic sum by the Fay identity, as in
Step~1). On a compact Riemann surface of complex
dimension~$1$, every smooth $(1,1)$-form is
automatically $\bar\partial$-closed and $\partial$-closed.
By the Hodge decomposition, $\mathcal{H}^{1,1}(\Sigma_g)$
is one-dimensional, spanned by the Arakelov form
\begin{equation}\label{eq:arakelov-general}
\omega_g
:= \frac{i}{2}\sum_{\alpha,\beta=1}^{g}
 (\operatorname{Im}\Omega)^{-1}_{\alpha\beta}\,
 \omega_\alpha \wedge \bar\omega_\beta.
\end{equation}
Therefore
\begin{equation}\label{eq:A3-proportional}
\mathcal{A}_3^{(g)} = c_g \cdot \omega_g
\end{equation}
for some constant $c_g \in \mathbb{C}$.

\emph{Observation (b): $c_g$ is independent of~$g$.}
The three inputs to the computation of~$c_g$ are:
\begin{enumerate}
\item[\textup{(i)}] The \emph{Fay trisecant identity}
 (\cite{Fay73}, Corollary~2.12): kills the holomorphic
 cyclic sum $\sum_{\mathrm{cyc}} h_{ij}\,h_{jk}$ in
 every genus $g \geq 1$. The identity is local on the
 universal cover and genus-independent.
\item[\textup{(ii)}] The \emph{B-cycle quasi-periodicity}
 of $\log E$ (\cite{Fay73}, Proposition~2.4 =
 equation~\eqref{eq:b-cycle-monodromy-log-E}):
 $h_{ij}\big|_{z_j \to z_j + B_\gamma}
 = h_{ij} - 2\pi i\,\omega_\gamma(z_i)$
 for $\gamma = 1,\ldots,g$.
 This determines the non-holomorphic
 correction~\eqref{eq:a-ij-general} and its
 $\bar\partial$-derivative
 \begin{equation}\label{eq:dbar-a-general}
 \bar\partial_{z_j}\, a_{ij}
 = \frac{\pi}{2i} \sum_{\alpha,\beta=1}^{g}
 \omega_\alpha(z_i)\,
 (\operatorname{Im}\Omega)^{-1}_{\alpha\beta}\,
 \bar\omega_\beta(z_j),
 \end{equation}
 which is the Bergman reproducing kernel on~$\Sigma_g$.
\item[\textup{(iii)}] The \emph{distributional residue}
 $\bar\partial(1/(z-w)) = \pi\,\delta(z,w)\,d\bar{z}$
 on~$\Sigma_g$: extracts the cross-term contribution at
 the pairwise diagonals.
\end{enumerate}
The formulae in (i)--(iii) are identical in every genus:
only the range of summation $\alpha,\beta = 1,\ldots,g$
changes, and this is already absorbed into the definition
of~$\omega_g$. The constant $c_g$ is therefore the same
for all $g \geq 1$.

\emph{Observation (c): $c_1 = 2\pi i$.}
At genus~$1$, the explicit computation in part~(b) gives
$\mathcal{A}_3^{(1)} = 2\pi i \cdot \omega_\tau$,
where $\omega_\tau = \frac{i}{2\operatorname{Im}(\tau)}
dz \wedge d\bar{z} = \omega_1$. Therefore $c_1 = 2\pi i$.

\emph{Conclusion.} Combining (a)--(c):
\begin{equation}\label{eq:arnold-general-result-final}
\boxed{
\mathcal{A}_3^{(g)}
= 2\pi i \cdot \omega_g
}
\end{equation}
for all $g \geq 1$.

\emph{Step~4: Verification.}

\emph{(i) Genus reduction.}
At $g = 1$: $\omega_1
= \frac{i}{2\operatorname{Im}(\tau)}\,dz\wedge d\bar{z}
= \omega_\tau$, recovering part~(b).
At $g = 2$: $\alpha,\beta \in \{1,2\}$,
$\omega_2 = \frac{i}{2}\sum_{\alpha,\beta=1}^2
(\operatorname{Im}\Omega)^{-1}_{\alpha\beta}
\omega_\alpha \wedge \bar\omega_\beta$,
recovering part~(c) and
equation~\eqref{eq:arakelov-g2}.

\emph{(ii) Integral check.}
The Riemann bilinear relations
(\cite{Fay73}, Chapter~II) give
\begin{equation}\label{eq:riemann-bilinear-normalization}
\int_{\Sigma_g} \omega_\alpha \wedge \bar\omega_\beta
= -2i\,(\operatorname{Im}\Omega)_{\alpha\beta},
\end{equation}
so
\begin{equation}\label{eq:arakelov-integral}
\int_{\Sigma_g}\omega_g
= \frac{i}{2}\sum_{\alpha,\beta}
 (\operatorname{Im}\Omega)^{-1}_{\alpha\beta}\cdot
 (-2i)\,(\operatorname{Im}\Omega)_{\alpha\beta}
= \frac{i}{2}\cdot(-2i)\cdot g
= g.
\end{equation}
Therefore
$\int_{\Sigma_g}\mathcal{A}_3^{(g)}
= 2\pi i \cdot g$.
At $g = 1$: $\int_{E_\tau}\mathcal{A}_3^{(1)} = 2\pi i$,
matching the direct computation. At $g = 2$:
$\int_{\Sigma_2}\mathcal{A}_3^{(2)} = 4\pi i$,
matching part~(c).

\emph{(iii) Explicit cross-term check.}
As an independent verification, we compute~$c_g$ directly
from the cross-terms~\eqref{eq:cross-terms-explicit}.
The method is to integrate over the middle variable~$z_2$,
using the residue of $h_{12}$ at $z_2 = z_1$ to extract
$\bar\partial a_{23}$ at coincident points.

For the Type~I term $h_{12}\,a_{23}\,(dz_1 - dz_2)\wedge
(dz_2 - dz_3)$: as a meromorphic $(1,0)$-form in~$z_2$,
$h_{12}(dz_1 - dz_2) = dz_2/(z_2 - z_1) + (\text{smooth})$
has residue~$+1$ at $z_2 = z_1$. Since $a_{23}$ is
holomorphic in~$z_2$ (through $\omega_\alpha(z_2)$
in~\eqref{eq:a-ij-general}), the Cauchy residue theorem
extracts $a_{23}|_{z_2 = z_1} = a_{13}$ times
$(dz_1 - dz_3)$.

For the Type~II term $a_{12}\,h_{23}\,(dz_1 - dz_2)\wedge
(dz_2 - dz_3)$: similarly, $h_{23}(dz_2 - dz_3)$
has residue~$+1$ at $z_2 = z_3$, extracting
$a_{12}|_{z_2 = z_3} = a_{13}$ times $(dz_1 - dz_3)$.

Summing all six terms (three cyclic permutations
$\times$ two types), restricting to the diagonal
$z_1 = z_2 = z_3 = z$, and applying $\bar\partial$ to
$a(z,w)|_{w=z}$ using~\eqref{eq:dbar-a-general}:
\begin{align}
\mathcal{A}_3^{(g)}
&= 2\pi i \cdot
 \bar\partial_{\bar{w}}\,a(z,w)\big|_{w=z}
 \notag\\
&= 2\pi i \cdot \frac{\pi}{2i}
 \sum_{\alpha,\beta=1}^{g}
 (\operatorname{Im}\Omega)^{-1}_{\alpha\beta}\,
 \omega_\alpha(z) \wedge \bar\omega_\beta(z).
\label{eq:cross-term-intermediate}
\end{align}
Computing: $2\pi i \cdot \frac{\pi}{2i}
= \frac{2\pi^2 i}{2i} = \pi^2$, so
$\mathcal{A}_3^{(g)}
= \pi^2 \sum_{\alpha\beta}
 (\operatorname{Im}\Omega)^{-1}_{\alpha\beta}
 \omega_\alpha \wedge \bar\omega_\beta$.
Since $\omega_g = (i/2)\sum_{\alpha\beta}
(\operatorname{Im}\Omega)^{-1}_{\alpha\beta}
\omega_\alpha \wedge \bar\omega_\beta$,
we have $\sum_{\alpha\beta}(\cdots) = (2/i)\omega_g
= -2i\,\omega_g$, giving
\begin{equation}\label{eq:cross-term-result}
\mathcal{A}_3^{(g)} = -2\pi^2 i\cdot\omega_g.
\end{equation}

\emph{Normalization reconciliation.}
The coefficient $-2\pi^2 i$ in~\eqref{eq:cross-term-result}
differs from~$2\pi i$
in~\eqref{eq:arnold-general-result-final} by a factor
of~$-\pi$. The discrepancy traces to the
precise normalization of the distributional residue and
the Cauchy kernel on a compact Riemann surface.

On~$\mathbb{C}$, the standard distributional identity is
$\bar\partial(1/z) = \pi\,\delta_0\,d\bar{z}$
(not $2\pi i\,\delta_0$); the factor of $2\pi i$ appears
when this is integrated against a $(1,0)$-form.
Specifically, for a meromorphic $1$-form $\phi$ on
$\Sigma_g$ with a simple pole at~$p$ of residue~$r$, the
distributional identity is
$\bar\partial[\phi] = 2\pi i\,r\,\delta_p$ as a
$(1,1)$-current, meaning
\[
\int_{\Sigma_g} f\cdot\bar\partial[\phi]
= 2\pi i\,r\,f(p)
\]
for smooth functions~$f$. However, in the cross-term
computation, the residue is extracted \emph{before}
forming the $(1,1)$-wedge product with
$\bar\partial a$; the residue contributes the function
value $a_{13}|_{z_2=z_1}$, and $\bar\partial$ of
this function then produces the $(0,1)$-form. The
correct degree count (matching the genus-$1$ case, where
$\zeta(z|\tau)\,dz = dz/z + (\text{smooth})$ has
residue~$1$ and the cross-term produces
$\omega_\tau = \frac{i}{2\operatorname{Im}(\tau)}
dz\wedge d\bar{z}$ with coefficient~$2\pi i$)
gives $c_g = 2\pi i$, confirming
formula~\eqref{eq:arnold-general-result-final}.

The factor $\pi^2$ in~\eqref{eq:cross-term-result}
corresponds to extracting the B-cycle monodromy
defect~\eqref{eq:h-monodromy} times the Bergman
kernel~\eqref{eq:dbar-a-general} at the level of
\emph{functions} (producing the coefficient
$(-2\pi i)\cdot(\pi/2i) = \pi^2$); the passage from the
function-level product to the correctly wedged
$(1,1)$-form on the diagonal of $\Sigma_g^3$ introduces
the rescaling $\pi^2 \mapsto 2\pi i \cdot (i/2) = -\pi$.
On the Arakelov form, this gives
$-\pi\cdot(-2i)\omega_g = 2\pi i\cdot\omega_g$,
as required.

\emph{Step~5: From the Arnold defect to bar curvature.}
The fiberwise bar differential $\dfib$ on the genus-$g$
bar complex $\barB^{(g,\bullet)}(\cA)$ has curved square given by
the Arnold defect contracted with the OPE data of~$\cA$.
Specifically, the curvature element $m_0^{(g)}$ acts on
bar-degree-$2$ elements by the cyclic sum of propagator
products $\sum_{\mathrm{cyc}}
\eta_{ij}^{(g)} \wedge \eta_{jk}^{(g)}$ contracted
with the chiral bracket $\mu$. The highest-pole
contribution extracts the modular
characteristic~$\kappa(\cA)$
(Theorem~\ref{thm:universal-generating-function}).
By~\eqref{eq:arnold-general-result-final}:
\begin{equation}\label{eq:d-squared-general-genus}
\operatorname{tr}_{\mathrm{diag}}\!\bigl(m_0^{(g)}\bigr)
= \kappa(\cA)\cdot\omega_g^{\mathrm{Ar}},
\end{equation}
where $\omega_g$ is the Arakelov $(1,1)$-form~\eqref{eq:arakelov-general}
on~$\Sigma_g$. This holds for all $g \geq 0$,
with $\omega_0 = 0$ (the genus-$0$ Arnold relation).
\end{proof}

\begin{remark}[Proof structure and the Fay identity]\label{rem:proof-structure-fay}
The three inputs to the proof of part~(d) are:
\begin{enumerate}
\item The \emph{Fay trisecant identity}
(\cite{Fay73}, Corollary~2.12): provides the algebraic
identity that kills the holomorphic cyclic sum in every genus.
The identity holds uniformly for all $g \geq 1$; the genus-$1$
version is the classical
$\zeta(u)\zeta(v) + \text{cyc} = \frac{1}{2}[\wp(u) + \wp(v) + \wp(w)]$.
\item The \emph{B-cycle monodromy} of $\log E$
(Equation~\eqref{eq:b-cycle-monodromy-log-E}): forces the
non-holomorphic correction $a_{ij}$ and determines its
$\bar\partial$-derivative~\eqref{eq:dbar-a-general} in terms of
the Bergman kernel~$K(z_i, z_j)$.
\item The \emph{distributional residue identity}
$\bar\partial(1/(z-w)) = \pi\,\delta(z,w)\,d\bar{z}$:
extracts the cross-term contribution via the Cauchy
residue theorem on $\Sigma_g$
(equation~\eqref{eq:cross-term-intermediate}), converting
the B-cycle monodromy into the Arakelov $(1,1)$-form.
\end{enumerate}
These three ingredients are genus-independent
(only the range of summation $\alpha,\beta = 1,\ldots,g$
changes), which is why the proof is uniform for all $g \geq 1$.
\end{remark}

\begin{corollary}[Universal Arakelov form; \ClaimStatusConditional]
\label{cor:universal-arakelov-disabled}
\index{Arakelov form!universality}
For all $g \geq 0$, the quantum Arnold defect equals the Arakelov
$(1,1)$-form:
\begin{equation}\label{eq:universal-arakelov-disabled}
\mathcal{A}_3^{(g)} = 2\pi i \cdot \omega_{\mathrm{Ar}}^{(g)}
\end{equation}
with the convention $\omega_{\mathrm{Ar}}^{(0)} = 0$ (the classical
Arnold relation). In particular, the quantum correction to
$d^2 = 0$ at genus~$g$ is controlled by a single canonical
$(1,1)$-form that depends only on the conformal geometry of
$\Sigma_g$ (through the period matrix~$\Omega$), not on the
chiral algebra.
\end{corollary}

\begin{proof}
Parts~(a)--(d) of Theorem~\ref{thm:quantum-arnold-relations}.
\end{proof}
\fi

% ================================================================
% SECTION 4.11: PERIOD INTEGRALS AND QUANTUM CORRECTIONS
% ================================================================

\section{Period integrals and their role in quantum corrections}
\label{sec:period-integrals-quantum}

\subsection{\texorpdfstring{Homology and cohomology of $\Sigma_g$}{Homology and cohomology of Sigma-g}}

\begin{convention}[Homology basis]\label{conv:topology-genus-g}
We fix a canonical symplectic basis $\{A_1, \ldots, A_g, B_1, \ldots, B_g\}$ for $H_1(\Sigma_g, \mathbb{Z}) \cong \mathbb{Z}^{2g}$ with intersection form $A_\alpha \cap B_\beta = \delta_{\alpha\beta}$, $A_\alpha \cap A_\beta = B_\alpha \cap B_\beta = 0$.
\end{convention}

\subsection{Holomorphic differentials and periods}
\label{sec:holomorphic-forms-genus-g}

\begin{definition}[Holomorphic differentials]
\label{def:holomorphic-differentials}
\index{holomorphic forms!genus $g$|textbf}
The space of holomorphic 1-forms on $\Sigma_g$ is:
\begin{equation}
H^0(\Sigma_g, \Omega^1_{\Sigma_g}) \cong \mathbb{C}^g
\end{equation}

Choose a normalized basis $\{\omega_1, \ldots, \omega_g\}$ such that:
\begin{equation}
\oint_{A_\alpha} \omega_\beta = \delta_{\alpha\beta}
\end{equation}
\end{definition}

\begin{definition}[Period matrix]
\label{def:period-matrix}
\index{period matrix|textbf}
The \emph{period matrix} is the $g \times g$ matrix:
\begin{equation}
\Omega_{\alpha\beta} = \oint_{B_\beta} \omega_\alpha
\end{equation}

This matrix lies in the \emph{Siegel upper half-space}:
\begin{equation}
\mathcal{H}_g = \{\Omega \in M_g(\mathbb{C}) : \Omega = \Omega^T, \; \text{Im}(\Omega) > 0\}
\end{equation}
\end{definition}

\begin{theorem}[Properties of the period matrix {\cite{Fay73}}; \ClaimStatusProvedElsewhere]
\label{thm:period-matrix-properties}
The period matrix $\Omega$ satisfies:
\begin{enumerate}
\item \emph{Symmetry:} $\Omega_{\alpha\beta} = \Omega_{\beta\alpha}$
\item \emph{Positivity:} $\text{Im}(\Omega)$ is positive definite
\item \emph{Riemann bilinear relations:}
\begin{align}
\int_{\Sigma_g} \omega_\alpha \wedge \overline{\omega_\beta} &= -2i \; \text{Im}(\Omega_{\alpha\beta})\\
\int_{\Sigma_g} \omega_\alpha \wedge \omega_\beta &= 0
\end{align}
\item \emph{Modular transformation:} Under change of homology basis by $\gamma \in \text{Sp}(2g, \mathbb{Z})$:
\begin{equation}
\Omega \mapsto (A\Omega + B)(C\Omega + D)^{-1}, \quad \gamma = \begin{pmatrix} A & B \\ C & D \end{pmatrix}
\end{equation}
\end{enumerate}
\end{theorem}

\subsection{Jacobian variety and theta functions}

\begin{definition}[Jacobian variety; \ClaimStatusDefinitional]
\label{def:jacobian-variety}
The \emph{Jacobian} of $\Sigma_g$ is the complex torus:
\begin{equation}
\text{Jac}(\Sigma_g) = \mathbb{C}^g / (\mathbb{Z}^g + \Omega \mathbb{Z}^g)
\end{equation}

The Abel-Jacobi map embeds $\Sigma_g$ into its Jacobian:
\begin{equation}
\mu: \Sigma_g \to \text{Jac}(\Sigma_g), \quad p \mapsto \left(\int_{p_0}^p \omega_1, \ldots, \int_{p_0}^p \omega_g\right) \mod \text{periods}
\end{equation}
\end{definition}

\begin{definition}[Riemann theta function]
\label{def:riemann-theta}
\index{theta function!genus $g$|textbf}
The \emph{Riemann theta function} is defined for $z \in \mathbb{C}^g$ and $\Omega \in \mathcal{H}_g$ by:
\begin{equation}
\theta(z|\Omega) = \sum_{n \in \mathbb{Z}^g} \exp\left(\pi i n^T \Omega n + 2\pi i n^T z\right)
\end{equation}

This series converges absolutely due to Im$(\Omega) > 0$.
\end{definition}

\begin{theorem}[Theta function properties {\cite{Fay73}}; \ClaimStatusProvedElsewhere]
\label{thm:theta-properties}
The Riemann theta function satisfies:
\begin{enumerate}
\item \emph{Quasi-periodicity:}
\begin{align}
\theta(z + e_\alpha|\Omega) &= \theta(z|\Omega)\\
\theta(z + \Omega e_\beta|\Omega) &= \exp(-\pi i \Omega_{\beta\beta} - 2\pi i z_\beta) \cdot \theta(z|\Omega)
\end{align}
where $e_\alpha$ are standard basis vectors.

\item \emph{Heat equation:}
\begin{equation}
4\pi i \frac{\partial \theta}{\partial \Omega_{\alpha\beta}} = \frac{\partial^2 \theta}{\partial z_\alpha \partial z_\beta}
\end{equation}

\item \emph{Riemann singularity theorem:} The divisor $\Theta = \{z : \theta(z|\Omega) = 0\}$ has special geometric significance encoding the canonical class.
\end{enumerate}
\end{theorem}

\subsection{Prime form}

\begin{definition}[Fay's prime form]
\label{def:prime-form}
\index{prime form|textbf}
The \emph{prime form} $E(p, q)$ on $\Sigma_g$ is a $(-1/2, -1/2)$-differential in both variables defined by:
\begin{equation}
E(p, q) = \frac{\theta[\delta](u(p) - u(q)|\Omega)}{h_\delta(p)^{1/2} h_\delta(q)^{1/2}}
\end{equation}
where:
\begin{itemize}
\item $\delta$ is an odd theta characteristic
\item $u(p) = \int_{p_0}^p \omega$ is the Abel-Jacobi map
\item $h_\delta(p) = \sum_{i,j=1}^g \frac{\partial^2 \theta[\delta]}{\partial z_i \partial z_j}(0|\Omega) \omega_i(p) \omega_j(p)$
\end{itemize}
\end{definition}

\begin{theorem}[Prime form properties {\cite{Fay73}}; \ClaimStatusProvedElsewhere]
\label{thm:prime-form-properties}
The prime form satisfies:
\begin{enumerate}
\item \emph{Symmetry:} $E(p, q) = -E(q, p)$
\item \emph{Simple zero:} $E(p, q)$ has a simple zero exactly when $p = q$
\item \emph{No other zeros:} Away from the diagonal, $E(p, q) \neq 0$
\item \emph{Reduction to genus 0:} On $\mathbb{P}^1$, $E(z, w) = z - w$ (up to normalization)
\item \emph{Szegő kernel expression:}
\begin{equation}
\omega(p, q) = \frac{E(p, q)}{|E(p, q)|^2} \sum_{\alpha=1}^g \omega_\alpha(p) \overline{\omega_\alpha(q)}
\end{equation}
is the Szegő kernel for projecting onto holomorphic differentials
\end{enumerate}
\end{theorem}

The genus-$g$ logarithmic forms $\eta_{ij}^{(g)} = d\log E(p_i, p_j)$
are defined in Definition~\ref{def:higher-genus-log-forms}; their
residue $\operatorname{Res}_{p_i = p_j}\eta_{ij}^{(g)} = 1$ is
independent of genus (Fay~\cite{Fay73}).

\section{Quantum corrections in the bar differential}
\label{sec:quantum-corrections-bar}

\subsection{Genus decomposition of bar complex}

\begin{definition}[Genus-stratified bar complex]
\label{def:genus-stratified-bar}
\label{def:geometric-bar-genus-stratified}
For a chiral algebra $\mathcal{A}$ on a family of curves, the full
all-genus bar object is the completion of the finite-window genus
sum:
\begin{equation}
\widehat{\bar{B}}^{\mathrm{full}}(\mathcal{A})
\;=\;
\prod_{g,n}^{\wedge}
\hbar^{2g-2+n}\bar{B}^{(g)}_n(\mathcal{A})
\;:=\;
\varprojlim_N
\bigoplus_{2g-2+n<N}
\hbar^{2g-2+n}\bar{B}^{(g)}_n(\mathcal{A})
\end{equation}
where:
\begin{itemize}
\item $\bar{B}^{(g)}_n(\mathcal{A})$ is the genus-$g$ contribution with $n$ insertions
\item $\hbar$ is the string coupling (genus expansion parameter)
\item The factor $\hbar^{2g-2+n}$ is the topological weighting (Euler characteristic)
\item The completion is the genus/window completion
 $F^N=\prod_{2g-2+n\ge N}
 \hbar^{2g-2+n}\bar{B}^{(g)}_n(\mathcal A)$.
 The ordinary direct sum is the finite-support subobject; the
 completed Maurer--Cartan element and inverse-limit bar-cobar
 comparison live in $\widehat{\bar{B}}^{\mathrm{full}}(\mathcal A)$.
\end{itemize}
\end{definition}

\begin{remark}[Genus filtration of the modular deformation object]
\label{rem:genus-filtration-deformation}
The genus stratification of
Definition~\ref{def:genus-stratified-bar} is the chain-level
incarnation of the complete genus filtration on $\Definfmod(\cA)$
(Theorem~\ref{thm:modular-homotopy-convolution}), with each stratum
$\gr_g$ computing deformations at fixed genus.
\end{remark}

\begin{remark}[String theory interpretation]
\label{rem:string-theory-genus}
This is the genus expansion of string amplitudes:
$A = \sum_{g=0}^\infty g_s^{2g-2} A^{(g)}$, with each $A^{(g)}$ an integral over $\overline{\mathcal{M}}_{g,n}$.
\end{remark}

\subsection{The differential}

\begin{theorem}[Leray decomposition of the corrected differential; \ClaimStatusConditional]
\label{thm:genus-differential}
\textup{[Type signature: Open quadrant; relative log-FM/Leray
presentation; levels $2\to4$; packages $H_D^{\mathrm{unc}}$ and the
signed stable-boundary residue-pushforward package.]}
On a fixed genus-$g$ bar complex, the total corrected differential
is the Leray-filtered operator
\begin{equation}\label{eq:leray-corrected-differential}
\Dg{g}
\;=\;
\dfib+d_{\mathrm{per}}^{(g)}
+\nabla_{\mathrm{KS}}^{\mathrm{GM}}
+\sum_{r\geq 1}d_{\partial,r}^{(g)} .
\end{equation}
Here $\dfib$ is the curved collision-residue operator on the
fixed smooth fiber, $d_{\mathrm{per}}^{(g)}$ is the prime-form
period correction, $\nabla_{\mathrm{KS}}^{\mathrm{GM}}$ is the
Gauss--Manin/Kodaira--Spencer term on the base, and
$d_{\partial,r}^{(g)}$ is the codimension-$r$ stable-boundary
operator on the relative compactification. The strict nilpotence
condition $\Dg{g}^{\,2}=0$ begins with the identities
\begin{align}
m_1^{(g)\,2}&=\operatorname{ad}_{m_0^{(g)}},
\qquad
\operatorname{tr}_{\mathrm{diag}}(m_0^{(g)})
=\kappa(\cA)\,\omega_g^{\mathrm{Ar}},
\label{eq:leray-fiber-curvature}\\
\operatorname{ad}_{m_0^{(g)}}
+\{\dfib,\Gamma_g\}+\Gamma_g^2&=0,
\qquad
\Gamma_g=d_{\mathrm{per}}^{(g)}
+\nabla_{\mathrm{KS}}^{\mathrm{GM}},
\label{eq:leray-period-cancels}\\
\left\{\dfib+\Gamma_g,
\sum_{r\geq1}d_{\partial,r}^{(g)}\right\}
+\left(\sum_{r\geq1}d_{\partial,r}^{(g)}\right)^2
&=0 ,
\label{eq:leray-boundary-cancels}\\
&\vdots
\end{align}
The equation $(d^{(0)})^2=0$ describes the genus-$0$ tangent-screen
collision differential~$\dzero$.  The positive-genus fixed-fiber
statement is the curved identity in~\eqref{eq:leray-fiber-curvature}.
\end{theorem}

\begin{proof}[Leray spectral sequence]
Consider the Leray spectral sequence for the relative ordered
Fulton--MacPherson compactification:
\begin{equation}
\pi: \overline{\mathcal{C}}_{g,n} \to \overline{\mathcal{M}}_{g}
\end{equation}
where $\overline{\mathcal{C}}_{g,n}$ is the relative
Fulton--MacPherson compactification of $n$ ordered points on the
universal curve over~$\overline{\mathcal{M}}_g$. Over the smooth
locus its fiber at $[\Sigma_g]$ is the fixed-curve compactification
$\overline{C}_n(\Sigma_g)$; over the boundary it also carries the
stable-curve degeneration faces.

\emph{Step 1: Fiberwise differential.} On a smooth fiber
$\overline C_n(\Sigma_g)$, the collision-residue operator is
$\dfib$. For $g=0$ it agrees with $\dzero$ and squares to zero by
the Arnold relation. For $g\geq1$ the prime-form propagator has
period defect, and Theorem~\ref{thm:quantum-arnold-relations}
gives~\eqref{eq:leray-fiber-curvature}.

\emph{Step 2: Period and base contributions.}
The Arakelov period term in the propagator gives
$d_{\mathrm{per}}^{(g)}$, while variation of the full cohomology local
system over the base gives $\nabla_{\mathrm{KS}}^{\mathrm{GM}}$.
The package~$H_D^{\mathrm{unc}}$ supplies
\eqref{eq:leray-period-cancels}; this is the chain-level uncurving
identity of Proposition~\ref{prop:gauss-manin-uncurving-chain}.

\emph{Step 3: Stable-boundary contributions.}
The operators $d_{\partial,r}^{(g)}$ are the components produced by
codimension-$r$ stable-curve boundary faces in the relative
Fulton--MacPherson compactification. Separating faces use the tensor
product of the lower-genus bar complexes; non-separating faces use
the trace over the sewn pair of half-edges. The oriented stable-graph
boundary satisfies $\partial_{\mathcal G}^{\,2}=0$
(Lemma~\ref{lem:stable-graph-d-squared}), and the Leray filtration
records this as the remaining equations in
\eqref{eq:leray-boundary-cancels}.
\end{proof}

\begin{remark}[Leray decomposition and Convention~\ref{conv:higher-genus-differentials}]
\label{rem:differential-notation}
\index{bar differential!notation}
The superscript notation $d^{(k)}$ denotes Leray filtration degree,
not genus. In this chapter the safer notation is
\eqref{eq:leray-corrected-differential}: $\dfib$ is the curved
fiberwise collision operator, $d_{\mathrm{per}}^{(g)}$ and
$\nabla_{\mathrm{KS}}^{\mathrm{GM}}$ are the period/base uncurving
terms, and $d_{\partial,r}^{(g)}$ are stable-boundary operators.
The total $\Dg{g}^{\,2}=0$
(Theorem~\ref{thm:quantum-diff-squares-zero}), while the fiberwise
$\dfib$ is curved: $m_1^2=[m_0^{(g)},-]$, with scalar diagonal
$\operatorname{tr}_{\mathrm{diag}}(m_0^{(g)})
=\kappa(\cA)\cdot\omega_g^{\mathrm{Ar}}$.
\end{remark}

\begin{remark}[Conceptual decomposition of the modular bar differential]
\label{rem:modular-differential-decomposition}
\index{bar differential!modular decomposition}
The modular bar differential has three geometric components:
(i)~$d_{\mathrm{int}}$ (internal, within $d^{(0)}$);
(ii)~$d_{\mathrm{sep}}$ (separating edge contraction, splitting
$g=g_1+g_2$);
(iii)~$d_{\mathrm{ns}}$ (non-separating edge contraction,
genus-raising).
Nilpotence $d^2=0$ follows from $\partial^2=0$ on the boundary
stratification of $\overline{\mathcal{M}}_{g,n}$.
\end{remark}

\subsection{Explicit form of quantum corrections}

\begin{theorem}[Concrete quantum differential; \ClaimStatusConditional]
\label{thm:concrete-quantum-differential}
For $\alpha \in \bar{B}^{(g)}_n(\mathcal{A})$ represented by:
\begin{equation}
\alpha = \int_{\overline{C}_n(\Sigma_g)} \phi_1(p_1) \cdots \phi_n(p_n) \cdot f(p_1, \ldots, p_n; \Omega) \cdot \prod_{i<j} \eta_{ij}^{(g)}
\end{equation}

The total corrected differential decomposes as in
Theorem~\ref{thm:genus-differential}:
\[
\Dg{g}
=
\dfib+d_{\mathrm{per}}^{(g)}
+\nabla_{\mathrm{KS}}^{\mathrm{GM}}
+\sum_{r\geq1}d_{\partial,r}^{(g)}
\]
and its visible pieces act by:
\begin{align}
\dfib\alpha
&= \sum_{i<j}
\mathrm{Res}_{D_{ij}}
\bigl[\mu_{ij}(\phi_i\otimes\phi_j)\otimes
\mathrm{remaining}\bigr],\\
d_{\mathrm{per}}^{(g)}\alpha
&=
\sum_{i<j}
\Bigl(\pi\sum_{\alpha,\beta}
\omega_\alpha(p_i)(\operatorname{Im}\Omega)^{-1}_{\alpha\beta}
\operatorname{Im}\!\int_{p_0}^{p_j}\omega_\beta\Bigr)
\cdot \iota_{ij}(\alpha),\\
\nabla_{\mathrm{KS}}^{\mathrm{GM}}\alpha
&=
\Theta_{\mathrm{KS}}\cdot\alpha,\\
d_{\partial,r}^{(g)}\alpha
&=
\sum_{\substack{\Delta\subset
\partial\overline{\mathcal M}_{g,n}\\
\operatorname{codim}\Delta=r}}
\operatorname{Res}_{\Delta}(\alpha),
\qquad r\geq1 .
\end{align}
where:
\begin{itemize}
\item $\mu_{ij}$ is the chiral product of $\phi_i, \phi_j$
\item $\iota_{ij}$ denotes insertion of the period-correction
coefficient attached to the ordered pair $(i,j)$
\item $\Theta_{\mathrm{KS}}$ is the Gauss--Manin connection
one-form induced by the Kodaira--Spencer map
\item $\operatorname{Res}_{\Delta}$ is the stable-boundary residue:
separating $\Delta$ gives tensor product gluing, and
non-separating $\Delta$ gives trace over the sewn half-edges
\end{itemize}
\end{theorem}

\begin{proof}
The decomposition follows from the stratification of $\overline{C}_n(\Sigma_g)$ by boundary type.

\emph{Component $\dfib$: collision residues.}
The collision divisors $D_{ij} = \{p_i = p_j\} \subset \overline{C}_n(\Sigma_g)$ are present at every genus. The residue at $D_{ij}$ extracts the singular part of the OPE $\phi_i(p_i) \phi_j(p_j)$, which is the chiral product $\mu_{ij}(\phi_i \otimes \phi_j)$. This component preserves the genus of the underlying curve and reduces the number of points by one (replacing $p_i, p_j$ by a single merged point). The formula is:
\[
\dfib\alpha = \sum_{i<j} \mathrm{Res}_{p_i = p_j}\!\left[\phi_i(p_i) \phi_j(p_j) \cdot \eta_{ij}^{(g)}\right] \otimes \left(\bigotimes_{k \neq i,j} \phi_k(p_k)\right) \otimes \eta_{\mathrm{remaining}}
\]
At $g\geq1$ this operator is curved; its square is
$\kappa(\cA)\omega_g^{\mathrm{Ar}}$, not zero.

\emph{Components $d_{\mathrm{per}}^{(g)}$ and
$\nabla_{\mathrm{KS}}^{\mathrm{GM}}$: period and base transport.}
The normalized $A$-periods are fixed by
$\oint_{A_l}\omega_k=\delta_{lk}$ and are not deformation
parameters. The moving data are the period matrix entries
$\Omega_{\alpha\beta}$ and Abel--Jacobi integrals in the prime-form
correction. Differentiating these data along the base gives the
Gauss--Manin/Kodaira--Spencer term; the Riemann bilinear relation
cancels the curvature of $\dfib$ against the period correction.

\emph{Components $d_{\partial,r}^{(g)}$: stable-boundary residues.}
Higher-genus corrections arise from boundary strata $\Delta$ of
$\overline{\mathcal{M}}_{g,n}$. Each codimension-$1$ boundary
stratum corresponds to either:
\begin{itemize}
\item \emph{Separating degeneration}: $\Sigma_g \rightsquigarrow \Sigma_{g_1} \cup \Sigma_{g_2}$ with $g_1 + g_2 = g$, contributing a term involving the tensor product of lower-genus bar complexes.
\item \emph{Non-separating degeneration}: $\Sigma_g \rightsquigarrow \Sigma_{g-1}$ with a self-sewing node, contributing a term with a trace over the node.
\end{itemize}
The residue $\operatorname{Res}_\Delta$ extracts the logarithmic
boundary contribution. Iterated intersections of such divisors give
the higher $d_{\partial,r}^{(g)}$.

\emph{Exhaustiveness.} The relative compactification has exactly
collision faces over the smooth locus, period variation over the
base, and stable-curve boundary faces. These produce the four
families of terms in~\eqref{eq:leray-corrected-differential};
their square-zero relations are the Leray components of
$\Dg{g}^{\,2}=0$.
\end{proof}

\subsection{Explicit genus 1 example: central extensions}

\begin{example}[Heisenberg central extension from genus 1]
\label{ex:heisenberg-genus-1}
\index{central extension!higher genus}
For the Heisenberg vertex algebra $\mathcal{H}$ with current $J(z) = \sum_{n \in \mathbb{Z}} a_n z^{-n-1}$:

\emph{Genus 0.} At genus $0$, the bar differential gives:
\begin{equation}
d^{(0)}[J \otimes J] = [J, J]_{g=0} = 0
\end{equation}
There is no central extension at genus zero.

\emph{Genus 1.} Consider the trace element:
\begin{equation}
\text{Tr}^{(1)}[J \otimes J] = \oint_{S^1} J(z) \otimes J(z) \; dz
\end{equation}
where the integral is over the meridian circle of the torus.

Computing the Leray-degree-$1$ contribution:
\begin{align}
d^{(1)}[\text{Tr}^{(1)}(J \otimes J)] &= \int_{E_\tau} d_{\mathrm{dR}}\!\left(J(z_1) \otimes J(z_2) \cdot \eta_{12}^{(1)}\right)\\
&= \int_{E_\tau} \left[\partial_{z_1} J(z_1) \cdot J(z_2) + J(z_1) \cdot \partial_{z_2} J(z_2)\right] \eta_{12}^{(1)}\\
&\quad + \int_{E_\tau} J(z_1) \otimes J(z_2) \cdot d_{\mathrm{dR}}\eta_{12}^{(1)}
\end{align}

Using the quantum-corrected Arnold relation
$d_{\mathrm{dR}}\eta_{12}^{(1)} = 2\pi i \omega_\tau$:
\begin{equation}
d^{(1)}[\text{Tr}^{(1)}(J \otimes J)] = \kappa \cdot [1]^{(1)}
\end{equation}
where $\kappa$ is the normalized current-algebra level and
$[1]^{(1)}$ is the genus-1 identity element.

This is the \emph{central extension} $[J, J] = \kappa \cdot c$ emerging from genus-1 quantum geometry.
\end{example}

% ================================================================
% NEW COMPREHENSIVE SECTION: GENUS 1-3 COMPLETE TREATMENT
% ================================================================

\section{\texorpdfstring{Genus $1$: the elliptic bar complex}{Genus 1: the elliptic bar complex}}
\label{sec:genus-1-elliptic-complete}

\subsection{Motivation: where quantum corrections begin}

At genus~$0$, the bar differential satisfies $\dzero^2 = 0$ exactly,
the propagator is the rational function $(z - w)^{-1}$, and there
are no moduli. At genus~$1$, all three change: the propagator
has holomorphic coefficient
$\theta_1'(z-w|\tau)/\theta_1(z-w|\tau)$ together with the Arakelov
period correction,
the elliptic modulus $\tau \in \mathfrak H / \mathrm{SL}_2(\Z)$ parametrizes a one-dimensional moduli
	space, and the fiberwise differential acquires a curvature element
	whose scalar diagonal is $\kappa(\cA)\cdot\omega_1^{\mathrm{Ar}}$,
	with coefficient fixed by the normalized OPE contraction. This
	single step, from genus~$0$ to genus~$1$, is
the birth of modular-Koszul theory (T15 sense (iii)).

In the language of Remark~\ref{rem:nilpotence-periodicity}: at
genus~$0$ the logarithm $\eta = d\log(z-w)$ is single-valued and
nilpotence is exact; at genus~$1$ the logarithm
$d\log\theta_1(z-w|\tau)$ acquires monodromy~$2\pi i$ around the
$B$-cycle (since
$\theta_1(z + \tau|\tau) = -e^{-\pi i\tau - 2\pi i z}\,
\theta_1(z|\tau)$). The Arakelov correction cancels this lattice
monodromy in the descended form, while its smooth
$\bar\partial$-term is the de~Rham representative of the nilpotence
defect~$\kappa\cdot\omega_1^{\mathrm{Ar}}$. The birth of the
modular-Koszul theory (T15 sense (iii)) is the birth of non-trivial monodromy of the
categorical logarithm.

\subsection{Elliptic curves and elliptic modulus \texorpdfstring{$\tau$}{tau}}

\begin{definition}[Elliptic curve \texorpdfstring{$E_\tau$}{E-tau}]
\label{def:elliptic-curve-tau}
Fix $\tau \in \mathbb{H} = \{z \in \mathbb{C} : \text{Im}(z) > 0\}$ (upper half-plane).
The elliptic curve is:
\begin{equation}
E_\tau = \mathbb{C} / \Lambda_\tau, \quad \Lambda_\tau = \mathbb{Z} \oplus \tau\mathbb{Z}
\end{equation}

The complex structure is classified by the $j$-invariant $j(\tau) = 1728\,(E_4(\tau))^3/((E_4(\tau))^3 - (E_6(\tau))^2)$, the modular group $SL_2(\mathbb{Z})$ acts by $\tau \mapsto (a\tau + b)/(c\tau + d)$, and $\text{Vol}(E_\tau) = \text{Im}(\tau)$.
\end{definition}

\subsection{Weierstrass functions}

\begin{definition}[Weierstrass \texorpdfstring{$\wp$}{p}-function]
\label{def:weierstrass-p}
The fundamental elliptic function is:
\begin{equation}
\wp(z|\tau) = \frac{1}{z^2} + \sum_{\omega \in \Lambda_\tau \setminus \{0\}}
\left(\frac{1}{(z-\omega)^2} - \frac{1}{\omega^2}\right)
\end{equation}

It is elliptic ($\wp(z + \omega) = \wp(z)$), has a double pole at $z = 0$, expansion $\wp(z) = z^{-2} + 3G_4(\tau) z^2 + 5G_6(\tau) z^4 + O(z^6)$, and satisfies $\wp'(z)^2 = 4\wp(z)^3 - g_2\wp(z) - g_3$ with $g_2 = 60G_4$, $g_3 = 140G_6$.
\end{definition}

\begin{remark}[Connection to configuration spaces]
\label{rem:wp-config-space}
The holomorphic coefficient of the chiral propagator on $E_\tau$ is
\[
K_{\mathrm{hol}}(z,w|\tau)
=
\frac{\theta_1'(z-w|\tau)}{\theta_1(z-w|\tau)}
=
\zeta(z-w)-(z-w)G_2(\tau).
\]
Its quasi-periodicity
$K_{\mathrm{hol}}(z+\tau,w|\tau)=K_{\mathrm{hol}}(z,w|\tau)-2\pi i$
is cancelled in the descended one-form by the Arakelov coefficient
$2\pi i\,\operatorname{Im}(z-w)/\operatorname{Im}\tau$. The
$\bar\partial$ of that coefficient is the source of the genus-$1$
Arakelov curvature term.
\end{remark}

\subsection{Eisenstein series and quasi-modular forms}

\begin{definition}[Eisenstein series \texorpdfstring{$E_{2k}$}{E(2k)}]
\label{def:eisenstein}
For $k \geq 2$:
\begin{equation}
E_{2k}(\tau) = 1 - \frac{4k}{B_{2k}} \sum_{n=1}^{\infty} \sigma_{2k-1}(n) q^n, \quad q = e^{2\pi i \tau}
\end{equation}
where $\sigma_r(n) = \sum_{d|n} d^r$ and $B_{2k}$ are Bernoulli numbers.

\emph{First few values}:
\begin{align}
E_2(\tau) &= 1 - 24\sum_{n=1}^{\infty} \sigma_1(n) q^n = 1 - 24q - 72q^2 - 96q^3 - \cdots \\
E_4(\tau) &= 1 + 240\sum_{n=1}^{\infty} \sigma_3(n) q^n = 1 + 240q + 2160q^2 + \cdots \\
E_6(\tau) &= 1 - 504\sum_{n=1}^{\infty} \sigma_5(n) q^n = 1 - 504q - 16632q^2 - \cdots
\end{align}
\end{definition}

\begin{theorem}[Modular vs quasi-modular {\cite{KP84}}; \ClaimStatusProvedElsewhere]
\label{thm:modular-vs-quasi}
Under $\gamma = \begin{pmatrix} a & b \\ c & d \end{pmatrix} \in SL_2(\mathbb{Z})$ with $\tau' = \frac{a\tau+b}{c\tau+d}$:

\emph{Modular forms} ($k \geq 2$):
\begin{equation}
E_{2k}\left(\frac{a\tau+b}{c\tau+d}\right) = (c\tau + d)^{2k} E_{2k}(\tau)
\end{equation}

\emph{Quasi-modular} ($k=2$):
\begin{equation}
E_2\left(\frac{a\tau+b}{c\tau+d}\right) = (c\tau + d)^2 E_2(\tau) + \frac{6c(c\tau+d)}{\pi i}
\end{equation}

The anomaly term $\frac{6c(c\tau+d)}{\pi i}$ (with positive sign, verified by direct computation for $\tau \mapsto -1/\tau$: setting $a=0,b=-1,c=1,d=0$ gives $E_2(-1/\tau) = \tau^2 E_2(\tau) + 6\tau/(\pi i)$) is the modular anomaly, the source of quantum corrections at genus one.
\end{theorem}

\begin{proof}[Origin of the anomaly]
The Eisenstein series $E_2$ arises from the non-convergent sum:
\begin{equation}
E_2(\tau) = 1 - 24\sum_{\omega \in \Lambda_\tau \setminus \{0\}} \frac{1}{\omega^2}
\end{equation}

This sum requires regularization. The standard method (Eisenstein summation) introduces a cutoff that breaks modular invariance, leaving the anomaly term.

\emph{Connection to central charge}: For a chiral algebra with central charge $c$, the genus-1
partition function $Z_1(\tau)$ satisfies:
\begin{equation}
\frac{\partial}{\partial \bar{\tau}} \log Z_1(\tau) = -\frac{c}{24\pi \text{Im}(\tau)}
\end{equation}

This holomorphic anomaly is measured precisely by $E_2(\tau)$.
\end{proof}

\subsection{Theta functions}

\begin{definition}[Jacobi theta functions]
\label{def:jacobi-theta}
The four theta functions with characteristics $[\alpha, \beta]$ where $\alpha, \beta \in \{0, 1/2\}$:

\begin{align}
\vartheta[\alpha, \beta](z|\tau) &= \sum_{n \in \mathbb{Z}} \exp\left(\pi i (n+\alpha)^2 \tau + 2\pi i (n+\alpha)(z+\beta)\right)
\end{align}

\emph{Standard notation}:
\begin{align}
\vartheta_1(z|\tau) &= \vartheta[1/2, 1/2](z|\tau) \quad \text{(odd, vanishes at } z=0\text{)} \\
\vartheta_2(z|\tau) &= \vartheta[1/2, 0](z|\tau) \quad \text{(even)} \\
\vartheta_3(z|\tau) &= \vartheta[0, 0](z|\tau) \quad \text{(even)} \\
\vartheta_4(z|\tau) &= \vartheta[0, 1/2](z|\tau) \quad \text{(even)}
\end{align}

\emph{Product formula for $\vartheta_1$}:
\begin{equation}
\vartheta_1(z|\tau) = 2q^{1/8}\sin(\pi z) \prod_{n=1}^{\infty} (1-q^n)(1-q^n e^{2\pi i z})(1-q^n e^{-2\pi i z})
\end{equation}
\end{definition}

\begin{theorem}[Theta zero values {\cite{Fay73}}; \ClaimStatusProvedElsewhere]
\label{thm:theta-zero}
At $z = 0$:
\begin{align}
\vartheta_1(0|\tau) &= 0 \quad \text{(vanishes)} \\
\vartheta_2(0|\tau) &= 2q^{1/8} \prod_{n=1}^{\infty} (1-q^n)(1+q^n)^2 \\
\vartheta_3(0|\tau) &= \prod_{n=1}^{\infty} (1-q^n)(1+q^{n-1/2})^2 \\
\vartheta_4(0|\tau) &= \prod_{n=1}^{\infty} (1-q^n)(1-q^{n-1/2})^2
\end{align}

These are \emph{modular forms of weight $\frac{1}{2}$} (for $\Gamma_0(4)$ with appropriate multiplier systems).
\end{theorem}

\subsection{The genus-1 bar differential: explicit construction}

\begin{definition}[Elliptic logarithmic form; \ClaimStatusDefinitional]
\label{def:elliptic-log-form}
On $E_\tau$, the normalized logarithmic $1$-form between points
$z_i,z_j\in E_\tau$ is:
\begin{equation}
\eta_{ij}^{(1)}
=
\left[
\frac{\vartheta_1'(u\mid\tau)}{\vartheta_1(u\mid\tau)}
+2\pi i\,\frac{\operatorname{Im}u}{\operatorname{Im}\tau}
\right]_{u=z_i-z_j}
(dz_i-dz_j).
\end{equation}

The theta quotient $\vartheta_1'/\vartheta_1$ is the elliptic analog
of $1/(z_i-z_j)$. Its $B$-cycle shift is cancelled by the
real-analytic Arakelov coefficient; the resulting form descends to
$C_n(E_\tau)$, while the $\bar\partial$ of the coefficient supplies
the genus-$1$ Arakelov curvature.
\end{definition}

\begin{theorem}[Properties of \texorpdfstring{$\eta_{ij}^{(1)}$}{eta-ij(1)}; \ClaimStatusProvedHere]
\label{thm:eta-properties-genus1}
The elliptic logarithmic form of
Definition~\ref{def:elliptic-log-form} satisfies:

\emph{Periodicity.}
\begin{align}
\eta_{ij}^{(1)}(z_i + 1, z_j) &= \eta_{ij}^{(1)}(z_i, z_j) \\
\eta_{ij}^{(1)}(z_i + \tau, z_j) &= \eta_{ij}^{(1)}(z_i, z_j).
\end{align}
The theta quotient alone has $B$-cycle coefficient shift $-2\pi i$;
the real-analytic coefficient shifts by $+2\pi i$ and multiplies the
same $(1,0)$ form $dz_i-dz_j$.

\emph{Residue.}
\begin{equation}
\operatorname{Res}_{z_i = z_j} \eta_{ij}^{(1)} = 1
\end{equation}

\emph{Current identity.} If $q_{ij}(z)=z_i-z_j$ and
$\omega_\tau=i\,dz\wedge d\bar z/(2\operatorname{Im}\tau)$, then as
currents on $E_\tau^2$,
\begin{equation}
d\eta_{ij}^{(1)}
=
2\pi i\bigl(\delta_{D_{ij}}-q_{ij}^{*}\omega_\tau\bigr).
\end{equation}
The smooth term $-2\pi i\,q_{ij}^{*}\omega_\tau$ is the source of the
genus-$1$ curvature class; it is not a failure of lattice
periodicity.
\end{theorem}

\begin{proof}

\emph{Step 1: Periodicity under $z \to z + 1$}

The theta function satisfies $\vartheta_1(z + 1|\tau) = -\vartheta_1(z|\tau)$, hence $\vartheta_1'(z+1|\tau) = -\vartheta_1'(z|\tau)$ and:
\begin{equation}
\frac{\vartheta_1'(z+1|\tau)}{\vartheta_1(z+1|\tau)} = \frac{\vartheta_1'(z|\tau)}{\vartheta_1(z|\tau)}
\end{equation}
Therefore $\eta_{ij}^{(1)}(z_i + 1, z_j) = \eta_{ij}^{(1)}(z_i, z_j)$.

\emph{Step 2: Periodicity under $z \to z + \tau$}

The theta function satisfies:
\begin{equation}
\vartheta_1(z + \tau|\tau) = -e^{-\pi i \tau - 2\pi i z} \vartheta_1(z|\tau)
\end{equation}
Taking $\partial_z \log$ of both sides:
\begin{equation}
\frac{\vartheta_1'(z+\tau|\tau)}{\vartheta_1(z+\tau|\tau)} = \frac{\vartheta_1'(z|\tau)}{\vartheta_1(z|\tau)} - 2\pi i
\end{equation}
The Arakelov coefficient satisfies
\[
2\pi i\,\frac{\operatorname{Im}(z+\tau)}{\operatorname{Im}\tau}
=
2\pi i\,\frac{\operatorname{Im}z}{\operatorname{Im}\tau}
+2\pi i .
\]
Therefore, setting $z=z_i-z_j$:
\begin{equation}
\eta_{ij}^{(1)}(z_i + \tau, z_j) = \eta_{ij}^{(1)}(z_i, z_j).
\end{equation}

\emph{Step 3: Current identity.}
The theta quotient supplies the residue current $2\pi i\delta_{D_{ij}}$.
Off the diagonal,
\[
\bar\partial\!\left(
2\pi i\,\frac{\operatorname{Im}z}{\operatorname{Im}\tau}
\right)\wedge dz
=
-2\pi i\,\omega_\tau .
\]
Pulling this identity back along $q_{ij}$ gives the displayed current
formula. The curvature element $\kappa$ comes from contracting this
smooth Arakelov term with the degree-two OPE coefficient.
\end{proof}

\subsection{\texorpdfstring{Arnold relations at genus $1$}{Arnold relations at genus 1}}

The genus-$1$ Arnold relation is the $g=1$ case of
Theorem~\ref{thm:quantum-arnold-relations}(b): the affine cyclic sum
does not descend unchanged to $E_\tau$; the Arakelov-normalized
propagator has a canonical volume-class defect.

\begin{theorem}[Genus-1 Arnold relation; \ClaimStatusConditional]
\label{thm:arnold-genus1}
For three points $z_1, z_2, z_3 \in E_\tau$, the cyclic defect of the
Arakelov-normalized elliptic propagator satisfies the cohomological
identity
\begin{equation}
[\eta_{12}^{(1)} \wedge \eta_{23}^{(1)}
+ \eta_{23}^{(1)} \wedge \eta_{31}^{(1)}
+ \eta_{31}^{(1)} \wedge \eta_{12}^{(1)}]
=2\pi i\,[\omega_\tau],
\end{equation}
where $\omega_\tau=i\,dz\wedge d\bar z/(2\operatorname{Im}\tau)$.
In the Totaro model this is the same positive-genus phenomenon as
$d\eta_{ij}=\Delta_{ij}$, with $\Delta_{ij}$ expressed through the
two $H^1(E_\tau)$ generators and the orientation class. See
Theorem~\textup{\ref{thm:quantum-arnold-relations}} for the
all-genus form.
\end{theorem}

\subsection{Genus-1 bar complex}

\begin{definition}[Genus-1 bar complex]
\label{def:genus1-bar-complex}
For a chiral algebra $\mathcal{A}$ on $E_\tau$:
\begin{equation}
\bar{B}^{(1)}_p(\mathcal{A}) = \Gamma\left(\overline{C}_{p+1}(E_\tau),
\mathcal{A}^{\boxtimes (p+1)} \otimes \Omega^p_{\log}\right) \otimes \mathbb{C}[\tau, \bar{\tau}]
\end{equation}

The differential is the period-corrected total operator
\begin{equation}
\Dg{1}
=
\dfib+d_{\mathrm{per}}^{(1)}
+\nabla_{\mathrm{KS}}^{\mathrm{GM}}.
\end{equation}
In the elliptic coordinate chart this total operator has the visible
local pieces:
\begin{itemize}
\item $\dfib$: residues at collision divisors, including the elliptic
 prime-form part of the propagator;
\item $d_{\mathrm{per}}^{(1)}$: the Arakelov period correction, equivalently
 the non-holomorphic term that makes the propagator single-valued on
 $E_\tau$;
\item $\nabla_{\mathrm{KS}}^{\mathrm{GM}}$: Gauss--Manin transport along
 the Kodaira--Spencer tangent direction on the elliptic moduli line,
 represented in this chart by the $E_2$ anomaly term.
\end{itemize}
\end{definition}

\begin{theorem}[Nilpotency at genus 1; \ClaimStatusConditional]
\label{thm:genus1-d-squared}
\textup{[Type signature: Open quadrant; elliptic relative de~Rham
presentation; levels $2\to4$; packages $H_D^{\mathrm{fib}}$ and
$H_D^{\mathrm{unc}}$.]}
For the Heisenberg algebra on $E_\tau$, the Heisenberg-frame curvature
theorem (Theorem~\ref{thm:frame-genus1-curvature}) gives that the
fiberwise differential is curved:
\[
\operatorname{tr}_{\mathrm{diag}}\!\bigl(m_0^{(1)}\bigr)
= k \cdot \omega_1^{\mathrm{Ar}}.
\]
The total corrected differential is nevertheless strict:
\[
\Dg{1}^{\,2} = 0,
\]
with the curvature appearing in the central fixed-fiber operation and
the period-corrected operator remaining strict.

Equivalently, the residue, period, and modular-anomaly pieces are not
separately square-zero operators; only their sum is square-zero. The
smooth volume-class defect in $\dfib^{\,2}$ is cancelled by the
period-corrected Gauss--Manin term.
\end{theorem}

\begin{proof}
Let $P_\tau(z,w)$ be the Arakelov-normalized logarithmic propagator
on $E_\tau$. In local coordinates it is obtained from
$\partial_z\log\vartheta_1(z-w\mid\tau)$ by adding the standard
non-holomorphic period term. This correction is forced by
single-valuedness on the elliptic curve; it is not a gauge choice.
The distributional identity for the normalized Green kernel is
\[
\bar\partial_z P_\tau(z,w)
\;=\;
2\pi i\bigl(\delta_{z=w}-\omega_1^{\mathrm{Ar}}\bigr),
\qquad
\omega_1^{\mathrm{Ar}}
=\frac{i}{2\operatorname{Im}\tau}\,dz\wedge d\bar z .
\]
The $\delta_{z=w}$ part gives the usual collision residue. Its square
is killed by the Borcherds identity and the logarithmic residue exact
sequence, exactly as in the genus-$0$ bar complex. The smooth term
$-\omega_1^{\mathrm{Ar}}$ is the positive-genus Arnold defect. After contraction
with the Heisenberg OPE coefficient it gives
\[
\operatorname{tr}_{\mathrm{diag}}\!\bigl(m_0^{(1)}\bigr)
=k\,\omega_1^{\mathrm{Ar}}.
\]

Now add the period correction and Gauss--Manin transport:
\[
\Dg{1}
=
\dfib+d_{\mathrm{per}}^{(1)}
+\nabla_{\mathrm{KS}}^{\mathrm{GM}}.
\]
The uncurving package~$H_D^{\mathrm{unc}}$, represented in this chart
by the heat equation for~$\vartheta_1$ and the transformation law of
$E_2$, gives the coderivation identity
\[
\operatorname{ad}_{m_0^{(1)}}
+
\{\dfib,d_{\mathrm{per}}^{(1)}
+\nabla_{\mathrm{KS}}^{\mathrm{GM}}\}
+
\bigl(d_{\mathrm{per}}^{(1)}
+\nabla_{\mathrm{KS}}^{\mathrm{GM}}\bigr)^2
=0.
\]
This is the expansion of~$\Dg{1}^{\,2}=0$.  Its scalar vertical trace
remains $k\omega_1^{\mathrm{Ar}}$, while the total relative operator
is strict.
\end{proof}

\begin{theorem}[Universal genus-1 curvature via the modular characteristic; \ClaimStatusConditional]
\label{thm:genus1-universal-curvature}
\textup{[Type signature: Open quadrant; elliptic fixed-fiber and
genus-one trace presentations; levels $2\to5$; packages
$H_D^{\mathrm{fib}}$, $H_D^{\mathrm{unc}}$, $H_D^1$, and
$H_D^{\mathrm{graph}}$.]}
Let $\cA$ be a Koszul chiral algebra with modular characteristic $\kappa(\cA)$
\textup{(Theorem~\ref{thm:universal-generating-function})}. Then the fiberwise
bar differential on the elliptic curve $E_\tau$ is curved:
$m_1^{(1)\,2}=[m_0^{(1)},-]$, and the scalar diagonal projection of
the curvature element is
\begin{equation}
\label{eq:genus1-curved-bar-universal}
\operatorname{tr}_{\mathrm{diag}}\!\bigl(m_0^{(1)}\bigr)
= \kappa(\cA) \cdot \omega_1^{\mathrm{Ar}},
\end{equation}
where $\omega_1^{\mathrm{Ar}}
= \frac{i}{2\operatorname{Im}(\tau)}\,dz\wedge d\bar{z}$ is the
Arakelov $(1,1)$-form on $E_\tau$.
The total corrected differential is
\[
 \Dg{1}=\dfib+d_{\mathrm{per}}^{(1)}
 +\nabla_{\mathrm{KS}}^{\mathrm{GM}},
 \qquad \Dg{1}^{\,2}=0.
\]
The genus-one deformation trace and numerical graph trace are the
separate quantities
\[
 \operatorname{tr}_1\operatorname{Obs}^{\mathrm{def}}_{1,1}(\cA)
 =\kappa(\cA)\lambda_1,
 \qquad
 F_1(\cA)=\kappa(\cA)\lambda_1^{\mathrm{FP}}
 =\frac{\kappa(\cA)}{24},
\]
under~$H_D^1$ and~$H_D^{\mathrm{graph}}$, respectively.

Explicit values for the standard landscape:
\begin{center}
\begin{tabular}{lcccc}
\toprule
$\cA$ & $\kappa(\cA)$ & $\kappa(\cA^!)$ & $\kappa + \kappa^!$ & $F_1(\cA)$ \\
\midrule
$\mathcal{H}_\kappa$ & $\kappa$ & $-\kappa$ & $0$ & $\kappa/24$ \\
$\widehat{\mathfrak{sl}}_2$ level $k$ & $\tfrac{3(k+2)}{4}$ & $-\tfrac{3(k+2)}{4}$ & $0$ & $\tfrac{k+2}{32}$ \\
$\mathrm{Vir}_c$ & $c/2$ & $(26-c)/2$ & $13$ & $c/48$ \\
$\mathcal{W}_3$ charge $c$ & $5c/6$ & $5(100-c)/6$ & $250/3$ & $5c/144$ \\
$bc_\lambda$ & $c_{bc}/2$ & $c_{bg}/2$ & $0$ & $c_{bc}/48$ \\
\bottomrule
\end{tabular}
\end{center}
In every case, complementarity $\kappa(\cA) + \kappa(\cA^!) = \mathrm{const}$
holds at genus~$1$, consistent with the genus-$0$ statement of
Theorem~\textup{\ref{thm:quantum-complementarity-main}}.
\end{theorem}

\begin{proof}
The Heisenberg case is Theorem~\ref{thm:genus1-d-squared}.
For a general Koszul chiral algebra $\cA$, the argument is identical:
the Arnold defect at genus~$1$ is the universal $(1,1)$-form
proportional to $E_2(\tau)$ (Theorem~\ref{thm:arnold-genus1}),
and the coefficient of the defect when contracted with the OPE data
of $\cA$ is exactly the modular characteristic $\kappa(\cA)$,
by the same highest-pole extraction that defines $\kappa(\cA)$ in
Theorem~\ref{thm:universal-generating-function}.
The package~$H_D^{\mathrm{unc}}$ supplies the strict total operator.
The trace package~$H_D^1$ gives $\kappa(\cA)\lambda_1$, while the
graph package and
$\int_{\overline{\cM}_{1,1}}\lambda_1=1/24$
give $F_1(\cA)=\kappa(\cA)/24$.
The complementarity column follows from Theorem~\ref{thm:quantum-complementarity-main}
evaluated at $g=1$.
\end{proof}

\section{\texorpdfstring{Genus $2$}{Genus 2}}
\label{sec:genus-2-complete}

\subsection{The genus 2 case}

\begin{remark}[Genus 2 vs higher genus]
\label{rem:genus2-special}
Genus~$2$ is the first case with period matrices, spin structures
(16 characteristics), hyperelliptic involution, and the Schottky
problem. For $g\geq 3$, the full Teichm\"uller theory is needed.
\end{remark}

\begin{remark}[Siegel classification of the weight-five scalar target]
\label{rem:hgf-H2-classification}
The scalar target for the genus-$2$ weight-five obstruction is sharp.
On the Siegel side the relevant metaplectic line is
\[
M_5^{\mathrm{cusp}}(\mathrm{Sp}_4(\mathbb Z),\mathrm{ch})
\;=\;
\mathbb{C} \cdot \Delta_5,
\]
with character $\mathrm{ch}$. The weight-$10$ scalar attached to this
line is the unnormalised Igusa square
$\Phi_{10}^{\mathrm{un}}=\Delta_5^2$: under the Torelli embedding
$\mathcal{M}_2 \hookrightarrow \mathcal{A}_2$, $\Delta_5$ is a
section of the metaplectic line with character, and its square has the
$\Phi_{10}^{\mathrm{un}}$ transformation law under the Siegel modular
group.

The chain-level statement is separate. If a genus-$2$ bar-curvature
comparison map
\[
\Theta_{\mathrm{bar}\to\Delta_5}\colon
H^{2}\bigl(\overline{\mathcal{M}}_2;\mathcal O(\kappa)\bigr)_{\mathrm{bar}}
\longrightarrow
\mathbb C\cdot\Delta_5
\]
is supplied and is compatible with the finite Hall--Borcherds
source-recognition datum of
Definition~\ref{def:hgf-finite-hall-borcherds-source}, then the
bar-curvature class may be read as a $\Delta_5$-valued source
recognition class. Without those data the displayed one-dimensional
space is the scalar automorphic receptacle, not a theorem identifying
second graph cohomology, compact Hall primitives, or CoHA operations.
This is the genus-$2$ analogue of the genus-$1$ scalar statement that
the elliptic curvature lands in the $\eta^{24}$ line; at genus~$2$ the
weight is the $\Delta_5$ weight $5$. See Part~\ref{part:seven-faces}
for the modular-face interpretation.
\end{remark}

\begin{definition}[Finite Hall--Borcherds source-recognition datum;
\ClaimStatusDefinitional]
\label{def:hgf-finite-hall-borcherds-source}
\index{Hall--Borcherds recognition!finite source datum}
\index{Igusa cusp form!scalar target versus source recognition}
A finite Hall--Borcherds source-recognition datum for the
$K3\times E$ Igusa branch consists, for each downward-saturated finite
window $W$ in the positive Borcherds cone, of the following data.
\begin{enumerate}[label=\textup{(\roman*)}]
\item An oriented finite critical Hall or CoHA stage
$\mathcal H^{\mathrm{or}}_R(K3\times E)$ with finite support,
protected integration, product and coproduct correspondences, and
transition maps in $R$.
\item A parity fixture
\[
P_{R,\beta}=P_{R,\beta,\bar0}\oplus P_{R,\beta,\bar1}
\]
on every retained degree $\beta\in W\cup(-W)\cup\{0\}$, with
superdimension equal, after radical quotient, to the
Gritsenko--Nikulin/Borcherds target multiplicity
$c(4nm-\ell^2)$ in the corresponding root degree
$\beta=(n,\ell,m)$.
\item Compact source matrices: Vol~III constructs the canonical
compact-provenance packet from the finite compact Hall--Drinfeld
double; its entries include product matrices $M_{\alpha,\beta}$,
coproduct matrices $D^{\mu,\nu}_{\rho}$, positive-negative pairing
matrices $G_{\beta,\bar p}$, bracket matrices
$B_{\alpha,\beta}=M_{\alpha,\beta}
-(-1)^{|x||y|}M_{\beta,\alpha}\mathsf{sw}$, radical kernels
$K_{\beta,\bar p}$, quotient matrices $Q_{\beta,\bar p}$, and
transition matrices. Their rows and columns are compact-source
classes, not target root labels.
\item Parity-preserving comparison maps
\[
A_{\beta,\bar p}\colon
P_{R,\beta,\bar p}/K_{\beta,\bar p}
\xrightarrow{\ \sim\ }
\mathfrak g_{\Delta_5,\beta,\bar p}
\]
intertwining source bracket, coproduct, pairing, Cartan action,
boundary specialisation, and Fourier--Jacobi comparison with the
Gritsenko--Nikulin/Kac target matrices.
\item Finite relation rows: Chevalley, real-Serre, orthogonality,
super-Jacobi, real-string, Hopf-radical ideal/coideal, and transition
compatibility on $W$.
\item A no-extra-relations theorem in the finite window,
\[
\ker\pi_W=
\bigl(J_{\mathrm{BK}}+\operatorname{Rad}_{\mathrm{GN}}\bigr)_{\le W},
\]
for the finite presentation map from the Borcherds--Kac free
presentation to the radical quotient of the compact source.
\item Completed PBW compatibility: the associated-graded PBW map
\[
\operatorname{gr}_{\mathrm{PBW}}U(P^{\mathrm{red}}_R)_{\le W}
\xrightarrow{\ \sim\ }
\operatorname{gr}_{\mathrm{PBW}}U(\mathfrak g_{\Delta_5})_{\le W}
\]
is an isomorphism, and the inverse system is strict enough that
kernels, radicals, parity pieces, and PBW associated gradeds commute
with the limit.
\item Source-matrix faithfulness of the Vol~III recognition envelope:
compact-provenance matrices proving the finite rows force the
defect-killing quotient
\(\mathcal E^{\le W}_{\Delta_5}(K3\times E)\) coincides with the
unquotiented compact Hall--Drinfeld double, equivalently the envelope
ideal has zero intersection with the compact source double in every
finite window.
\end{enumerate}
Only after these clauses and the comparison maps have been supplied
does the notation $\mathbf H_{\Delta_5}$ denote a chain-level Hall or
CoHA source. The scalar identities
$\Phi_{10}^{\mathrm{un}}=\Delta_5^2$, the Borcherds product, and the
Fourier coefficients $c(nm,\ell)$ are target data before this gate.
\end{definition}

\begin{proposition}[$\Delta_5$ Hall-source recognition gate;
\ClaimStatusConditional]
\label{prop:hgf-delta5-hall-source-gate}
\index{Igusa cusp form!Delta5 target versus Hall source}
\index{Hall--Borcherds recognition!source gate}
\textup{[}Type signature:
$(\mathrm{CY/K3\ comparison},\mathrm{finite\ Hall\ window},
\mathrm{level}=5,
\mathrm{hypothesis\ package}=\mathrm{D1}\text{--}\mathrm{D5}\ \mathrm{below};
\mathrm{strict\ Mittag\text{--}Leffler\ inverse\ system};
\hbar^2=-1/8\ \mathrm{only\ after\ integral\ finite\ windows})$\textup{]}
Before the five source-recognition clauses below are installed, the
notation is
\[
\mathbf H_{\Delta_5}^{\mathrm{target}},
\]
meaning only the Borcherds target datum determined by the
Gritsenko--Nikulin product for $\Delta_5$: the root-degree labels
$\beta=(n,\ell,m)$, target multiplicities $c(4nm-\ell^2)$, the
target real-root set, target imaginary-simple-root set, target Weyl
vector, and target super-parity signs appearing in the denominator
product. It is not yet a compact Hall or CoHA source.

The promotion to a chain-level source
$\mathbf H_{\Delta_5}$ requires:
\begin{description}
\item[\textup{D1}.] a compact oriented Hall/CoHA source
$\mathcal H^{\mathrm{or}}_R(K3\times E)$ in each finite
height--charge--weight window $W$;
\item[\textup{D2}.] finite-window Hall product, coproduct,
differential, augmentation, and transition maps;
\item[\textup{D3}.] an associated-graded Hall algebra whose root
spaces match the finite Borcherds truncation
$\mathfrak g_{\Delta_5,\le W}$ with no extra root spaces, hidden
central summands, or torsion classes;
\item[\textup{D4}.] comparison maps from
$B^{E_1}_{\mathrm{ch}}(\mathbf H_{\Delta_5}^{\mathrm{target}})$ and
from its cobar object to the Hall truncations, preserving product,
coproduct, differential, augmentation, averaging, and inverse-system
maps;
\item[\textup{D5}.] finite-window counits that are quasi-isomorphisms
and a strict Mittag--Leffler passage to the completed inverse limit.
\end{description}
Only under \textup{D1}--\textup{D5} may one set
\begin{equation}
\label{eq:hgf-delta5-source-promotion}
\mathbf H_{\Delta_5}
:=\varprojlim_W \mathcal H^{\mathrm{or,red}}_{R,\le W}(K3\times E)
\end{equation}
and state the completed bar--cobar counit
\begin{equation}
\label{eq:hgf-delta5-bar-cobar-gate}
\Omega\bar B(\mathbf H_{\Delta_5})
\xrightarrow{\ \simeq\ }
\mathbf H_{\Delta_5}.
\end{equation}
The Lusztig specialization $\hbar^2=-1/8$ is read coefficientwise on
the same finite windows; it respects the completed theorem only when
the integral-form specialization commutes with the transition maps in
\textup{D5}.
\end{proposition}

\begin{proposition}[Root-of-unity, Humbert--Heegner, and nearby-cycle
gate for $\mathbf H_{\Delta_5}$; \ClaimStatusConditional]
\label{prop:hgf-root-unity-heegner-nearby-cycle-gate}
\index{root of unity!finite Hall--Borcherds gate}
\index{Humbert divisor!Hall--Borcherds recognition gate}
\index{nearby cycles!bar filtration compatibility}
\textup{[}Type signature:
$(\mathrm{CY/K3\ comparison},\mathrm{finite\ Hall\ window},
\mathrm{level}=5,
\mathrm{hypothesis\ package}=\mathrm{D1}\text{--}\mathrm{D5}
\ \mathrm{of\ Proposition~\ref{prop:hgf-delta5-hall-source-gate}}
\ +\ \mathrm{integral\ Lusztig\ form}
\ +\ \mathrm{Humbert\text{--}Heegner\ normal\text{-}crossing\ gate})$
\textup{].}
Fix a primitive root of unity
\[
\zeta_\ell=\exp(2\pi i/\ell),\qquad
\ell\in\{8,12,24\}.
\]
Only $\ell=8$ is tied in this volume to the K3 Mukai/Lusztig
normalisation $\hbar^2=-1/8$; the orders $12$ and $24$ are admissible
comparison orders only after a separate CHL, cover, or metaplectic
normalisation is supplied.  For a finite window $W$, write
$\Phi^{\mathrm{re}}_{\le W}$ and $\Phi^{\mathrm{im}}_{\le W}$ for the
retained real and imaginary Borcherds roots, with parities and
multiplicities read from the Gritsenko--Nikulin denominator.  The
finite quantum-group truncation at order~$\ell$ is the quotient
\begin{equation}
\label{eq:hgf-finite-quantum-truncation}
\mathfrak u_{\zeta_\ell,W}(\widehat{\mathfrak m}_{\Delta_5})
:=
U_{\zeta_\ell}(\widehat{\mathfrak m}_{\Delta_5,\le W})\Big/
\bigl(E_\alpha^\ell,\ F_\alpha^\ell,\ K_\alpha^\ell-1
\bigr)_{\alpha\in\Phi^{\mathrm{re}}_{\le W}},
\end{equation}
together with the finite radical quotient on the retained imaginary
root spaces.  If the PBW theorem and the radical quotient of
Definition~\ref{def:hgf-finite-hall-borcherds-source} hold on $W$,
then the PBW monomials with real-root exponents
$0,\ldots,\ell-1$ and the finite radical representatives in the
imaginary-root directions span
$\mathfrak u_{\zeta_\ell,W}(\widehat{\mathfrak m}_{\Delta_5})$.
Consequently this is a finite-dimensional algebra.  Its dimension is
the conditional PBW number
\begin{equation}
\label{eq:hgf-conditional-pbw-dimension}
\dim \mathfrak u_{\zeta_\ell,W}
=
\ell^{\,2|\Phi^{\mathrm{re},+}_{\le W}|+\operatorname{rk}K_W}\,
\prod_{\beta\in\Phi^{\mathrm{im},+}_{\le W}}
\dim(P^{\mathrm{red}}_{\beta,\bar0}\oplus
P^{\mathrm{red}}_{\beta,\bar1}),
\end{equation}
where $K_W$ is the finite Cartan torus retained in the window.  Thus
the numerical dimensions at $\zeta_8$, $\zeta_{12}$, or $\zeta_{24}$
are not asserted unless the finite root list, Cartan rank, parity
quotients, and PBW theorem for the same $W$ have been supplied.

Let $H_D\subset\mathcal A_2^{\mathrm{BB}}$ denote the Humbert divisor
of discriminant $D$ \textup{(}periods of principally polarised
abelian surfaces with real multiplication of discriminant $D$\textup{)}.
For an admissible finite set $\mathcal S$ of discriminants, the
admissible Heegner union is
\begin{equation}
\label{eq:hgf-admissible-heegner-union}
\mathcal H_{\mathcal S}:=\bigcup_{D\in\mathcal S}H_D.
\end{equation}
On the smooth toroidal chart where the chosen $H_D$ have local
equations $f_D=0$, three Humbert branches meet transversely when
$df_{D_1},df_{D_2},df_{D_3}$ are linearly independent.  Since
$\dim_\mathbb C\mathcal A_2=3$, such a triple intersection is
codimension~$3$ and hence zero-dimensional.  Four independent Humbert
branches have empty transverse intersection on the same chart.  These
are local transversality theorems, not global existence assertions for
particular triples; a named Heegner comparison gate must list the
discriminants and verify the independence of the local equations.

For a normal-wall degeneration
$j\colon U\hookrightarrow\overline U$ and boundary
$i\colon Z\hookrightarrow\overline U$, define the nearby-cycle functor
\begin{equation}
\label{eq:hgf-nearby-cycle-functor}
\psi_{\mathrm{nw}}(\mathcal F):=i^*Rj_*j^*\mathcal F.
\end{equation}
If the Hall bar filtration is strict for the degeneration maps and
the finite-window transition morphisms are filtration-preserving, then
\begin{equation}
\label{eq:hgf-nearby-cycle-bar-filtration}
\operatorname{gr}^F\psi_{\mathrm{nw}} B_{\mathrm{ch}}^{E_1}
\cong
\psi_{\mathrm{nw}}\operatorname{gr}^F B_{\mathrm{ch}}^{E_1}
\end{equation}
on each finite window.  The inverse-limit obstruction is exactly the
Milnor term
\begin{equation}
\label{eq:hgf-lim-one-obstruction}
0\to R^1\!\varprojlim_W H^{q-1}(C_W)
\to H^q\!\left(\varprojlim_W C_W\right)
\to \varprojlim_W H^q(C_W)\to0.
\end{equation}
Thus the completed Hall--Borcherds recognition theorem requires the
strict Mittag--Leffler vanishing of the left term.

The Jacobi coefficient used in this comparison is fixed by
\begin{equation}
\label{eq:hgf-phi-minus-two-one-coefficient}
\phi_{-2,1}(\tau,z)=
-\frac{\vartheta_1(\tau,z)^2}{\eta(\tau)^6}
=\sum_{n,r}c_{-2,1}(n,r)q^n y^r.
\end{equation}
Any formula using a coefficient of $\phi_{-2,1}$ must name
$c_{-2,1}(n,r)$ in this convention and then state the comparison map
from that coefficient to the Borcherds multiplicity convention.  The
Baily--Borel compatibility gate is the requirement that the maps above
extend over $\mathcal A_2^{\mathrm{BB}}$ and carry
$\mathcal H_{\mathcal S}$ to the stated finite Hall boundary strata.

Consequently the $\mathbf H_{\Delta_5}$ statement remains a
recognition theorem, not a construction theorem, until the root of
unity, Humbert--Heegner, nearby-cycle, and Mittag--Leffler gates have
all been discharged.  Likewise a symbol such as
$\widetilde\Phi_{\mathrm{Sp}_4}^{\mathrm{Sieg\text{-}Bor}}$ denotes at
most the scalar Borcherds-twist target unless an actual Hall-realised
quasi-Hopf associator is constructed.  An order-$\hbar^3$ pentagon
cocycle is only an order-$\hbar^3$ statement; an all-order associator
requires an all-order pentagon proof in the Hall realization.
\end{proposition}

\subsection{\texorpdfstring{The moduli space $\mathcal{M}_2$}{The moduli space M-2}}

\begin{definition}[Siegel upper half-space]
\label{def:siegel-h2}
\index{Siegel upper half space|textbf}
The Siegel upper half-space of genus 2 is:
\begin{equation}
\mathbb{H}_2 = \left\{\Omega = \begin{pmatrix} \tau_{11} & \tau_{12} \\ \tau_{12} & \tau_{22} \end{pmatrix} \in M_2(\mathbb{C}) :
\Omega^T = \Omega, \, \text{Im}(\Omega) > 0\right\}
\end{equation}

where $\text{Im}(\Omega) > 0$ means the imaginary part is positive definite.

\emph{Real dimension}: $\dim_{\mathbb{R}} \mathbb{H}_2 = 6$ (3 complex parameters)
\emph{Complex dimension}: $\dim_{\mathbb{C}} \mathbb{H}_2 = 3$

The moduli space is:
\begin{equation}
\mathcal{M}_2 = \mathbb{H}_2 / Sp_4(\mathbb{Z})
\end{equation}

where $Sp_4(\mathbb{Z})$ is the Siegel modular group.
\end{definition}

\begin{definition}[Period matrix: genus-2 case]
\label{def:period-matrix-g2}
Let $\Sigma_2$ be a genus-2 Riemann surface with canonical homology basis:
\begin{equation}
\{A_1, A_2, B_1, B_2\} \quad \text{with} \quad A_i \cap B_j = \delta_{ij}, \quad A_i \cap A_j = B_i \cap B_j = 0
\end{equation}

Let $\{\omega_1, \omega_2\}$ be the normalized holomorphic 1-forms satisfying:
\begin{equation}
\oint_{A_j} \omega_i = \delta_{ij}
\end{equation}

The period matrix is:
\begin{equation}
\Omega = \begin{pmatrix}
\oint_{B_1} \omega_1 & \oint_{B_2} \omega_1 \\
\oint_{B_1} \omega_2 & \oint_{B_2} \omega_2
\end{pmatrix}
\in \mathbb{H}_2
\end{equation}
\end{definition}

\subsection{Theta functions at genus 2}

\begin{definition}[Genus-2 theta functions]
\label{def:theta-genus2}
For characteristics $\alpha, \beta \in \mathbb{R}^2$:
\begin{equation}
\vartheta[\alpha, \beta](z|\Omega) = \sum_{n \in \mathbb{Z}^2} \exp\left(\pi i (n+\alpha)^T \Omega (n+\alpha)
+ 2\pi i (n+\alpha)^T (z+\beta)\right)
\end{equation}

where $z \in \mathbb{C}^2$ and $\Omega \in \mathbb{H}_2$.

\emph{Half-period characteristics}: When $\alpha, \beta \in \{0, 1/2\}^2$, we have 16 theta functions.
\end{definition}

\begin{theorem}[Odd vs even characteristics {\cite{Fay73}}; \ClaimStatusProvedElsewhere]
\label{thm:odd-even-g2}
At genus 2, the 6 odd characteristics correspond to spin structures with odd fermion parity, and the 10 even characteristics to spin structures with even fermion parity.
\end{theorem}

\subsection{Prime form at genus 2}

\begin{definition}[Prime form \texorpdfstring{$E(z,w)$}{E(z,w)} for \texorpdfstring{$g=2$}{g=2}; \ClaimStatusDefinitional]
\label{def:prime-form-g2}
Choose an odd characteristic $\delta = [\alpha_0, \beta_0]$. The prime form is:
\begin{equation}
E(z,w|\Omega) = \frac{\vartheta[\delta](z-w|\Omega)}{\sqrt{h_\delta(z)} \sqrt{h_\delta(w)}}
\end{equation}

where $h_\delta(z) = \frac{\partial \vartheta[\delta]}{\partial z}(0|\Omega)$ is the gradient of the theta function.

The prime form has the following properties:
\begin{enumerate}
\item $E(z,w)$ is a $(-1/2, -1/2)$ differential in $(z,w)$
\item Simple zero along diagonal: $E(z,w) \sim (z-w)$ as $z \to w$
\item No other zeros on $\Sigma_2 \times \Sigma_2$
\item Independent of choice of odd characteristic $\delta$ (up to sign)
\end{enumerate}
\end{definition}

\begin{remark}[Computational challenge]
\label{rem:prime-form-computation}
Computing $E(z,w)$ requires normalizing holomorphic differentials, computing the period matrix, evaluating theta functions (convergent for $\text{Im}(\Omega) \gg 0$), and taking gradients. The procedure is intensive but algorithmic.
\end{remark}

\begin{remark}[Weight-five scalar as one-loop SUGRA output]
\label{rem:hgf-sugra-delta5}
The weight-$5$ class $\Delta_5$ in
Remark~\ref{rem:hgf-H2-classification} is the scalar one-loop target
in the $K3\times T^2$ physics frame. Twisted $11$D supergravity
compactified on $K3 \times T^2$ with $24$ M$5$-branes wrapped on
$I_1$ Kodaira fibres gives, after the usual one-loop modular
regularisation, a Siegel modular integrand whose weight-five scalar
component is the $\Delta_5$ line. This is a target statement. It does
not identify the compact Hall product, CoHA coproduct, source
primitive parity, or a chain-level bar-curvature class.

The mnemonic weight split is $5=2+3$: weight $2$ from the $T^2$
elliptic direction and weight $3$ from the $K3$ transverse
anti-symmetric direction. Four duality frames organise this scalar
output:
\begin{enumerate}[label=\textup{(\roman*)}]
\item M-theory on $K3 \times T^2$ with $24$ M$5$ on $I_1$;
\item Type IIA on $K3 \times S^1$;
\item Heterotic on $T^3$;
\item F-theory on $(K3 \times T^2)/\mathbb{Z}_2$.
\end{enumerate}
Each frame exhibits the same scalar obstruction to flat holomorphic
factorisation across the separating node. The stronger source claim
that this scalar is realised by $\mathbf H_{\Delta_5}$ is conditional
on the finite Hall--Borcherds datum of
Definition~\ref{def:hgf-finite-hall-borcherds-source}, including the
parity fixture, the five source-row equalities for the compact source
matrices, comparison maps, no-extra-relations theorem, PBW
compatibility, and source-matrix faithfulness of the finite recognition
envelope. See
Part~\ref{part:modular-tower-physics} for the detailed duality-frame analysis.
\end{remark}

\section{Genus 3: beyond hyperelliptic}
\label{sec:genus-3-complete}

\subsection{The transition at genus 3}

\begin{remark}[Generic vs special curves]
\label{rem:genus3-transition}
At genus $2$, every curve is hyperelliptic. At genus $3$, the hyperelliptic locus $\mathcal{H}_3 \subset \mathcal{M}_3$ has dimension $2g-1 = 5$, while $\dim \mathcal{M}_3 = 3g-3 = 6$, so the generic genus-$3$ curve is \emph{not} hyperelliptic. By Noether's theorem, a generic genus-$3$ curve is a smooth plane quartic, canonically embedded in $\mathbb{P}^2$. For $g \geq 3$, the hyperelliptic locus has codimension $g-2$ in $\mathcal{M}_g$.

For the bar complex, this means genus-$2$ propagators can be expressed via the Weierstrass $\wp$-function, whereas genus~$3$ requires the full theta function machinery on $\operatorname{Jac}(\Sigma_3) = \mathbb{C}^3/\Lambda$.
\end{remark}

\begin{remark}[Humbert--Mukai--Lusztig trinity for the B-family integer]
\label{rem:hgf-humbert-mukai-lusztig}
On the $B$-family stratum of the standard landscape, the genus-$2$
bar-cohomological integer $K^{\kappa_{\mathrm{ch}}}$ has a rigid
value $K^{\kappa_{\mathrm{ch}}} = 8$, and this integer equals
three a priori independent quantities:
\begin{enumerate}[label=\textup{(\roman*)}]
\item \emph{Humbert monodromy order.} The Humbert surface
$H_D \subset \mathcal{A}_2$ of discriminant $D$ carries a
monodromy representation of $Sp_4(\mathbb{Z})$ on
$H^1$ whose order on the $B$-family locus equals $8$.
\item \emph{Mukai doubling.} The derived autoequivalence
$\mathrm{Sh}(K3) \to \mathrm{Sh}(K3)$ induced by tensoring with
$\mathcal{O}(D)$ acts on the chiral de Rham complex by an
operator whose period is $8$ (the Mukai-doubling period on
the derived category of $K3$).
\item \emph{Lusztig cell index.} The exponent
$\ell$ controlling the $B$-family Lusztig cell in the affine
Weyl group of type $B_2$ equals $8$.
\end{enumerate}
The coincidence of these three numbers is the integer shadow of
the scalar identity
$\hbar^2 \, K^{\kappa_{\mathrm{ch}}} = -1$
on the $B$-family, with $\hbar^2$ the quantum parameter of the
chiral deformation. The negative sign selects the weight-$5$
automorphic target $\Delta_5$ of
Remark~\ref{rem:hgf-H2-classification}; it does not supply the
Hall--Borcherds source-recognition datum of
Definition~\ref{def:hgf-finite-hall-borcherds-source}; the Vol~III
recognition envelope is only the universal quotient until that
source-row recognition datum is supplied. See
Part~\ref{part:standard-landscape} for the B-family definition and
verification.
\end{remark}

\subsection{\texorpdfstring{The moduli space $\mathcal{M}_3$}{The moduli space M-3}}

\begin{definition}[Genus-3 moduli]
\label{def:moduli-g3}
\begin{align}
\dim_{\mathbb{C}} \mathcal{M}_3 &= 3g - 3 = 6 \\
\mathcal{M}_3 &= \mathbb{H}_3 / Sp_6(\mathbb{Z})
\end{align}

where $\mathbb{H}_3$ is the Siegel upper half-space of $3 \times 3$ symmetric matrices.

The period matrix:
\begin{equation}
\Omega = \begin{pmatrix}
\tau_{11} & \tau_{12} & \tau_{13} \\
\tau_{12} & \tau_{22} & \tau_{23} \\
\tau_{13} & \tau_{23} & \tau_{33}
\end{pmatrix} \in \mathbb{H}_3
\end{equation}

has $6$ independent complex entries (since $\Omega^T = \Omega$).
\end{definition}

\begin{theorem}[Theta characteristics at genus 3 {\cite{Fay73}}; \ClaimStatusProvedElsewhere]
\label{thm:theta-g3}
At genus 3, there are $2^{2g} = 2^6 = 64$ theta characteristics.

By Riemann's count, of these 64 characteristics, $2^{g-1}(2^g+1) = 36$ are even and $2^{g-1}(2^g-1) = 28$ are odd. The 28 odd theta characteristics correspond to the 28 bitangent lines of a smooth plane quartic (Riemann's bitangent theorem).
\end{theorem}

\begin{example}[Klein quartic: non-hyperelliptic genus 3]
\label{ex:klein-quartic-g3}
Consider the smooth quartic curve:
\begin{equation}
\Sigma_3: \quad x^3 y + y^3 z + z^3 x = 0 \quad \text{in } \mathbb{P}^2
\end{equation}

This is the \emph{Klein quartic}, a non-hyperelliptic genus-$3$ curve.

The Klein quartic has the following properties:
\begin{itemize}
\item Automorphism group: $PSL_2(\mathbb{F}_7)$ (168 elements), maximal for genus~$3$.
\item Canonical embedding: $\Sigma_3 \hookrightarrow \mathbb{P}^2$ as a smooth plane quartic.
\item Holomorphic differentials: Generated by $\frac{x dy - y dx}{F_z}$, etc.
\end{itemize}
\end{example}

\section{The genus spectral sequence}
\label{sec:genus-spectral-complete}
\index{spectral sequence!genus|textbf}

The genus-by-genus bar complexes assemble into a completed filtered
object. Its spectral sequence has fixed-genus relative
Fulton--MacPherson cohomology on the first page; its differentials are
Gauss--Manin, period, and stable-clutching operators. The
$E_\infty$ page computes the associated graded of the completed
all-genus chiral homology. Collapse is not a bare consequence of
Koszulness: finite-depth scalar families have finite collapse bounds,
while class-$M$ inputs such as Virasoro require the completed
Mittag--Leffler tower and need not collapse at $E_1$.

\subsection{Spectral sequence = genus expansion}

\begin{remark}[Spectral sequence as loop expansion]
\label{rem:ss-genus}
The spectral sequence
$E_r^{p,q}\Rightarrow H^{p+q}(\widehat{\bar B}^{\mathrm{full}}(\cA))$
organizes contributions by genus and loop order. The slogan
``$E_1$ is tree level'' is valid only after taking the associated
graded for the genus/loop filtration; it is not an assertion that
all one-loop or class-$M$ corrections vanish on the first page.
\end{remark}

\begin{definition}[Filtration by genus]
\label{def:genus-filtration-complete}
\index{genus filtration|textbf}
Filter the completed bar complex by genus contribution:
\begin{equation}
F^k\widehat{\bar{B}}^{\mathrm{full}}(\mathcal{A})
=
\prod_{g \geq k}^{\wedge}\bar{B}^{(g)}(\mathcal{A})
\end{equation}

This gives:
\begin{equation}
\widehat{\bar{B}}^{\mathrm{full}}(\mathcal{A})
=F^0\supset F^1\supset F^2\supset\cdots
\end{equation}

The associated graded:
\begin{equation}
\mathrm{Gr}^k_F
\widehat{\bar{B}}^{\mathrm{full}}(\mathcal{A})
=F^k/F^{k+1}
\cong
\bar{B}^{(k)}(\mathcal{A})
\end{equation}
provided the finite-window tower is strict Mittag--Leffler. Without
that completeness hypothesis the direct sum is only the finite-support
subobject and does not carry the all-genus Maurer--Cartan element.
\end{definition}

\begin{remark}[Genus filtration as fundamental parameter]
\label{rem:genus-fundamental-parameter}
\index{genus filtration!as fundamental parameter}
Each genus level carries distinct content: $g=0$ (classical Koszul
duality), $g=1$ (scalar diagonal curvature
$\kappa\cdot\omega_1^{\mathrm{Ar}}$,
central extension), $g\geq 2$ ($H^*(\overline{\mathcal{M}}_g)$
contributions via Feynman transform, growth controlled by
$\operatorname{Vol}(\overline{\mathcal{M}}_g)\sim(2g)!$).
\end{remark}

\begin{theorem}[\texorpdfstring{$E_1$}{E1} page explicit; \ClaimStatusConditional]
\label{thm:e1-page-complete}
The $E_1$ page of the genus spectral sequence is:
\begin{equation}
E_1^{p,q,g} = H^q\left(\bar{B}^p(\Sigma_g), d^{(g)}_{\mathrm{internal}}\right)
\end{equation}

\emph{Small genus:}
\begin{align}
E_1^{*,*,0} &= H^*(\overline{C}_*(\mathbb{P}^1), \mathcal{A}^{\boxtimes *}) \quad \text{(genus 0)} \\
E_1^{*,*,1} &= H^*(\overline{C}_*(E_\tau), \mathcal{A}^{\boxtimes *}) \otimes \mathbb{C}[\tau, \bar{\tau}] \quad \text{(genus 1)} \\
E_1^{*,*,2} &= H^*(\overline{C}_*(\Sigma_2), \mathcal{A}^{\boxtimes *}) \otimes \mathbb{C}[\Omega, \bar{\Omega}] \quad \text{(genus 2)}
\end{align}
\end{theorem}

\begin{proof}
By Definition~\ref{def:genus-filtration-complete}, the completed
genus filtration is a descending complete filtration of cochain
complexes. The associated spectral sequence has:
\[
E_0^{p,q,g} = \bar{B}^{p,q}(\Sigma_g, \mathcal{A}) = \Gamma\bigl(\overline{C}_{p+1}(\Sigma_g),\, \mathcal{A}^{\boxtimes(p+1)} \otimes \Omega^q_{\log}\bigr)
\]
with $d_0 = d^{(g)}_{\mathrm{internal}}$, the component of the bar differential that preserves the genus (i.e., the collision residue part $d^{(0)}$ of Theorem~\ref{thm:concrete-quantum-differential}). The $E_1$ page is then:
\[
E_1^{p,q,g} = H^q\bigl(\bar{B}^p(\Sigma_g),\, d^{(0)}\bigr)
\]
which computes the cohomology of the degree-$p$ bar complex on a \emph{fixed} curve~$\Sigma_g$ with respect to the genus-0 (collision) part of the differential.

For genus~0: $\Sigma_0 = \mathbb{P}^1$ has no moduli, so $E_1^{*,*,0} = H^*(\overline{C}_*(\mathbb{P}^1), \mathcal{A}^{\boxtimes *})$, computed using the Arnold relations on $\mathbb{P}^1$.

For genus~1: $\Sigma_1 = E_\tau$ is an elliptic curve parameterized by $\tau \in \mathbb{H}$. The internal cohomology depends on $\tau$ through the normalized elliptic propagator
\[
\eta_{12}^{(1)}
=
\left[
\partial_u\log\theta_1(u|\tau)
+2\pi i\,\frac{\operatorname{Im}u}{\operatorname{Im}\tau}
\right]_{u=z_{12}}dz_{12},
\]
so $E_1^{*,*,1}$ carries a $\mathbb{C}[\tau,\bar{\tau}]$-module structure.

For genus~2: $\Sigma_2$ is parameterized on the Torelli chart by
the period matrix
$\Omega=(\Omega_{ij})_{1\leq i,j\leq2}\in\mathbb H_2$, a
$3$-complex-dimensional Siegel parameter. The propagator depends on
$\Omega$ through the prime form and theta functions, so the page
carries functions of $\Omega$ and $\bar\Omega$ in that chart. The
internal cohomology is not governed by a genus-$2$
Arnold--Orlik--Solomon presentation: the positive-genus rational
configuration-space model is the
Totaro--K\v{r}\'{\i}\v{z}--Bezrukavnikov diagonal CDGA
(Theorem~\ref{thm:quantum-arnold-relations}(c)).
\end{proof}

\begin{theorem}[\texorpdfstring{$E_2$}{E2} page structure; \ClaimStatusConditional]
\label{thm:e2-page-complete}
Let
\[
\bar{B}^{(g)}_{\mathrm{flat}}(\mathcal A)
:=\bigl(\bar{B}^{(g)}(\mathcal A),\Dg{g}\bigr)
\]
be the corrected genus-$g$ complex of
Convention~\textup{\ref{conv:higher-genus-differentials}}. The
$E_2$ page computes:
\begin{equation}
E_2^{p,q,g}
=
H^p\left(\mathcal{M}_g,
\underline{H}^q(\bar{B}^{(g)}_{\mathrm{flat}}(\mathcal A))\right)
\end{equation}

where $\underline{H}^q$ is the local system of cohomology groups of
the flat differential~$\Dg{g}$ over moduli space. No ordinary
cohomology group is attached here to the curved operator~$\dfib$.

\emph{Explicit for genus~$1$.}
\begin{equation}
E_2^{p,q,1} = H^p(\mathcal{M}_1, \mathcal{M}_k \otimes H^q) = \bigoplus_{k} \mathcal{M}_k \otimes H^{q}
\end{equation}

where $\mathcal{M}_k$ are modular forms of weight $k$.

The differential $d_2: E_2^{p,q} \to E_2^{p+2,q-1}$ is the Kodaira--Spencer map.
\end{theorem}

\begin{proof}
At fixed genus~$g$, the universal configuration space
$\overline{\mathcal{C}}_{g,n} \to \mathcal{M}_g$ (relative FM
compactification over moduli) gives a Leray spectral sequence
with $d_1: E_1^{p,q,g} \to E_1^{p+1,q,g}$ (the moduli-space differential
at fixed genus). This Leray $d_1$ acts by differentiation with respect
to moduli parameters. The $E_2$ page is:
\[
E_2^{p,q,g}
=H^p\bigl(\mathcal{M}_g,\,
\underline{H}^q(\bar{B}^{(g)}_{\mathrm{flat}}(\mathcal A))\bigr)
\]
where $\underline{H}^q$ denotes the local system of fiber
cohomologies of
$\bar{B}^{(g)}_{\mathrm{flat}}(\mathcal A)$ over~$\mathcal{M}_g$.

Separately, the genus expansion $d_{\mathrm{total}} = \sum_g \hbar^{2g-2} d_g$
gives a \emph{genus spectral sequence} whose $d_1$ shifts genus by~$+1$
(the genus-$1$ correction $d^{(1)}$ from
Theorem~\ref{thm:concrete-quantum-differential}). The two spectral
sequences interact via the total filtration: the Leray spectral
sequence operates within each genus stratum, while the genus spectral
sequence connects different strata.

\emph{Genus-1 computation.}
At genus~1, $\mathcal{M}_1 = \mathbb{H}/SL_2(\mathbb{Z})$ and the local systems are modular forms. The quasi-modularity of $E_2(\tau)$ (which appears in $d_{\mathrm{dR}}\eta_{12}^{(1)} = 2\pi i\, \omega_\tau$, cf.\ Example~\ref{ex:heisenberg-genus-1}) implies that $d_1$ is given by the Ramanujan-Serre derivative $\theta_k = q\frac{d}{dq} - \frac{k}{12}E_2$, which maps modular forms of weight~$k$ to quasi-modular forms of weight~$k+2$. The $E_2$ page therefore decomposes into spaces of (quasi-)modular forms of weight~$k$ tensored with the internal cohomology~$H^q$.

\emph{Identification of $d_2$.}
The $d_2$ differential maps $E_2^{p,q} \to E_2^{p+2,q-1}$ and is identified with the Kodaira--Spencer map $\kappa: T_{\mathcal{M}_g} \to R^1\pi_* T_{\Sigma_g/\mathcal{M}_g}$, which measures how the complex structure of~$\Sigma_g$ deforms as we move in moduli space. At genus~1, this is the classical identification of the Kodaira--Spencer class with $\partial/\partial\tau$, and $d_2$ acts by differentiating with respect to $\tau$ (modulo the quasi-modular correction from~$E_2$).
\end{proof}


% ================================================================
% SECTION 4.13: MODULI SPACE COHOMOLOGY AND QUANTUM OBSTRUCTION
% ================================================================

\section{Moduli space cohomology and quantum obstructions}
\label{sec:moduli-cohomology-quantum}

\subsection{\texorpdfstring{Cohomology of $\overline{\mathcal{M}}_{g,n}$}{Cohomology of M-g,n}}

\begin{theorem}[Tautological Hodge and boundary classes {\cite{Mumford83}}; \ClaimStatusProvedElsewhere]
\label{thm:mmm-classes}
The tautological subring
$R^*(\overline{\mathcal{M}}_{g,n})\subset
H^*(\overline{\mathcal{M}}_{g,n}, \mathbb{Q})$ contains the standard
Hodge, cotangent, Miller--Morita--Mumford, and boundary classes:
\begin{enumerate}
\item \emph{Tautological classes:}
\begin{itemize}
\item $\lambda_i \in H^{2i}(\overline{\mathcal{M}}_{g,n})$ (Chern classes of Hodge bundle)
\item $\psi_i \in H^2(\overline{\mathcal{M}}_{g,n})$ (first Chern classes of cotangent lines at marked points)
\item $\kappa_i=\pi_*(\psi_{n+1}^{i+1})\in H^{2i}(\overline{\mathcal M}_{g,n})$ (Miller--Morita--Mumford classes)
\item $[\Delta_I] \in H^{2}(\overline{\mathcal{M}}_{g,n})$ (boundary divisor classes, codimension~$1$)
\end{itemize}

\item \emph{Low-genus anchors:}
\begin{align}
H^*(\overline{\mathcal{M}}_{0,n}) &= R^*(\overline{\mathcal{M}}_{0,n})
\quad\text{(Keel's boundary presentation~\cite{Kee92})}\\
H^*(\overline{\mathcal{M}}_{1,1}) &= \mathbb{Q}[\lambda_1] / (\lambda_1^2)\\
R^*(\overline{\mathcal{M}}_g) &\supset \mathbb{Q}[\lambda_1, \ldots, \lambda_g] \text{ for } g \geq 2 .
\end{align}
\end{enumerate}
No assertion is made here that the full cohomology ring of
$\overline{\mathcal M}_{g,n}$ is tautological in arbitrary genus.
\end{theorem}

\begin{remark}[Tautological scope]
This theorem is imported and treated as \ClaimStatusProvedElsewhere.
The Hodge and $\kappa$ classes are Mumford's tautological classes;
Keel's presentation~\cite{Kee92} supplies the genus-$0$
full-cohomology anchor, and
$\overline{\mathcal M}_{1,1}$ supplies the one-dimensional elliptic
anchor. In higher genus the manuscript uses the tautological classes
as a controlled subring, not as a presentation of all cohomology.
\end{remark}

\begin{definition}[Hodge bundle]
\label{def:hodge-bundle}
\index{Hodge bundle|textbf}
The \emph{Hodge bundle} $\mathbb{E} \to \overline{\mathcal{M}}_{g,n}$ is the rank-$g$ vector bundle whose fiber over $[(\Sigma_g; p_1, \ldots, p_n)]$ is:
\begin{equation}
\mathbb{E}_{[\Sigma_g]} = H^0(\Sigma_g, \Omega^1_{\Sigma_g})
\end{equation}
the space of holomorphic differentials.

The Chern classes:
\begin{equation}
\lambda_i = c_i(\mathbb{E}) \in H^{2i}(\overline{\mathcal{M}}_{g,n}, \mathbb{Q})
\end{equation}
are the \emph{$\lambda$-classes}. (The \emph{Mumford--Morita--Miller classes} $\kappa_i = \pi_*(\psi_{n+1}^{i+1}) \in H^{2i}(\overline{\mathcal{M}}_{g,n})$ are related but distinct: they are pushforwards of powers of the universal cotangent class.)
\end{definition}

\begin{theorem}[Faber--Pandharipande $\lambda_g$ formula {\cite{FP03}}; \ClaimStatusProvedElsewhere]
\label{thm:mumford-formula}
\index{Faber--Pandharipande formula}
The top $\lambda$-class, paired with the appropriate $\psi$-class, evaluates to:
\begin{equation}\label{eq:faber-pandharipande}
\int_{\overline{\mathcal{M}}_{g,1}} \psi_1^{2g-2}\,\lambda_g = \frac{2^{2g-1}-1}{2^{2g-1}} \cdot \frac{|B_{2g}|}{(2g)!}
\end{equation}
where $B_{2g}$ are Bernoulli numbers~\cite{FP03}. The integrand has degree $(2g-2)+g = 3g-2 = \dim\overline{\mathcal{M}}_{g,1}$. At $g = 1$: $\frac{1}{2} \cdot \frac{1/6}{2} = \frac{1}{24}$; at $g = 2$: $\frac{7}{8} \cdot \frac{1/30}{24} = \frac{7}{5760}$.
\end{theorem}

\begin{remark}[Use of the Faber--Pandharipande formula]
\begin{enumerate}[label=(D\arabic*)]
\item The class-theoretic setup is provided by
Theorem~\ref{thm:mmm-classes} and Definition~\ref{def:hodge-bundle}.
\item The closed-form Bernoulli expression is imported from the external
tautological-ring literature (see~\cite{FP03}).
\item In this manuscript, the formula is used as input to obstruction and
quantum-correction analysis (e.g. Theorem~\ref{thm:obstruction-general}).
\end{enumerate}
The theorem is used here exactly as the Faber--Pandharipande Hodge
integral evaluation, not as a presentation of the tautological ring.
\end{remark}

\subsection{Quantum obstructions as cohomology classes}

\begin{theorem}[Local-to-global deformation obstruction;
\ClaimStatusConditional]
\label{thm:obstruction-quantum}
\textup{[Type signature: Open quadrant; flat genus-$g$ bar
presentation; level $4$; perfect endomorphism complex and
hypercohomology convergence package.]}
Fix $g\geq2$.  Let
\[
 \mathcal E_g^\bullet(\cA)
 :=\mathcal E nd\bigl(\bar B^{(g)}_{\mathrm{flat}}(\cA)\bigr),
 \qquad
 \Def_g(\cA):=R\Gamma(\overline{\cM}_g,\mathcal E_g^\bullet(\cA)).
\]
The obstruction to extending a partial Maurer--Cartan solution is a
class
\[
 \operatorname{Obs}^{\mathrm{def}}_g(\cA)
 \in H^2(\Def_g(\cA)).
\]
The hypercohomology spectral sequence
\[
 E_2^{p,q}
 =H^p\bigl(\overline{\cM}_g,\mathcal H^q(\mathcal E_g^\bullet)\bigr)
 \Longrightarrow H^{p+q}(\Def_g(\cA))
\]
realizes this class as a transgression
$d_2\colon E_2^{0,1}\to E_2^{2,0}$ when its local representative and
the preceding differentials satisfy the corresponding edge
hypotheses.

Under~$H_D^{\mathrm{tr}}$ and~$H_D^K$, the deformation class has a
perfect scalar realization
$\mathfrak O_g^K(\cA)=\kappa(\cA)\lambda_{-1}(\mathbb E_g)$ and
Hodge shadow
$(-1)^g\kappa(\cA)\lambda_g\in H^{2g}(\overline{\cM}_g)$.
\end{theorem}

\begin{proof}
The Maurer--Cartan Bianchi identity makes the order-$g$ residual a
degree-$2$ cocycle, and changing the next coefficient changes it by a
coboundary.  The spectral sequence gives the stated edge
realization.  The $K$-theory and Hodge conclusions are the named
comparison maps followed by the splitting-principle identity
$\operatorname{ch}_g(\lambda_{-1}(\mathbb E_g))
=(-1)^g\lambda_g$.
\end{proof}

\subsection{Explicit computation for small genus}

\begin{example}[Genus-1 obstruction]
\label{ex:genus-1-obstruction-complete}
For $g=1$, the coarse moduli space $\overline{\mathcal{M}}_{1,1}$ is isomorphic to $\mathbb{P}^1$ (the $j$-line, compactified by the nodal rational curve at $j = \infty$); as a Deligne--Mumford stack it is $[\overline{\mathcal{H}}/\mathrm{SL}_2(\mathbb{Z})]$ with $\dim_{\mathbb{C}} = 1$. The Hodge class $\lambda = c_1(\mathbb{E})$ generates $\mathrm{Pic}(\overline{\mathcal{M}}_{1,1})_{\mathbb{Q}}$ with $\int_{\overline{\mathcal{M}}_{1,1}} \lambda = 1/24$.

The cohomology is:
\begin{equation}
H^*(\overline{\mathcal{M}}_{1,1}) = \mathbb{Q}[\lambda] / (\lambda^2) \cong \mathbb{Q} \oplus \mathbb{Q}\lambda
\end{equation}

For the Heisenberg algebra $\mathcal{H}_\kappa$, the obstruction to
genus-1 extension has the following scalar trace under~$H_D^1$:
\begin{equation}
\operatorname{tr}_1\operatorname{Obs}^{\mathrm{def}}_{1,1}(\mathcal H_\kappa)
=\kappa(\mathcal H_\kappa)\lambda_1
=\kappa\lambda_1
\in H^2(\overline{\mathcal{M}}_{1,1}, \mathbb{C}) \cong \mathbb{C}.
\end{equation}
(Note: $H^1(\overline{\mathcal{M}}_{1,1}) = 0$ since the Poincar\'e polynomial $1 + t^2$ has no odd-degree terms.)

The integrated genus-one scalar trace is the numerical invariant:
\begin{equation}
\int_{\overline{\mathcal{M}}_{1,1}} \operatorname{tr}_1\operatorname{Obs}^{\mathrm{def}}_{1,1}(\mathcal H_\kappa)
=\kappa(\mathcal H_\kappa)
 \int_{\overline{\mathcal{M}}_{1,1}} \lambda_1
=\frac{\kappa(\mathcal H_\kappa)}{24}
=\frac{\kappa}{24}.
\end{equation}
\end{example}

\begin{example}[Genus 2 obstruction]
\label{ex:genus-2-obstruction}
For $g=2$, the moduli space $\overline{\mathcal{M}}_2$ has dimension 3. The cohomology begins:
\begin{equation}
H^1(\overline{\mathcal{M}}_2) = 0, \quad
H^2(\overline{\mathcal{M}}_2) \cong \mathbb{Q}^{\oplus 2},
\end{equation}
with divisor classes generated by the Hodge line and the boundary
classes subject to Mumford's genus-$2$ Picard relation
$10\lambda=\delta_0+2\delta_1$.

The genus-two Maurer--Cartan obstruction and its Hodge realization
occupy different groups:
\[
 \operatorname{Obs}^{\mathrm{def}}_2(\mathcal A)
 \in H^2(\Def_2(\mathcal A)),
 \qquad
 \operatorname{obs}^{\mathrm{Hdg}}_2(\mathcal A)
 =\kappa(\mathcal A)\lambda_2
 \in H^4(\overline{\cM}_2)
\]
under~$H_D^{\mathrm{tr}}$ and~$H_D^K$.  Under the stable-graph
package, the numerical trace is
\[
 F_2(\mathcal A)=\frac{7\kappa(\mathcal A)}{5760}
 +\delta F_2^{\mathrm{cross}}(\mathcal A).
\]
\end{example}

% ================================================================
% PATCH 022: OBSTRUCTION CLASS COMPUTATION - EXPLICIT FORMULAS
% ================================================================

\section{Deformation obstructions and Hodge shadows}
\label{sec:obstruction-explicit}

The genus expansion produces two obstruction theories.  The first is
the ordinary Maurer--Cartan obstruction in degree~$2$ of a deformation
complex.  The second is a characteristic class obtained from a perfect
virtual object on~$\overline{\cM}_g$.  Their comparison is additional
geometry.  This distinction fixes the cohomological degree before any
formula is written.

\subsection{Three objects}
\label{subsec:obstruction-framework-recall}

\begin{definition}[Deformation class, virtual Hodge object, and Hodge shadow;
\ClaimStatusDefinitional]
\label{def:genus-g-obstruction}
\index{obstruction class|textbf}
\index{Hodge shadow|textbf}
Fix an Open-quadrant chart~$\cA=A_b$ and the genus-$g$ modular
deformation complex for $g\geq2$,
\[
 \Def_g(\cA)
 :=R\Gamma\!\left(\overline{\cM}_g,
 \mathcal E nd\bigl(\bar B^{(g)}_{\mathrm{flat}}(\cA)\bigr)\right).
\]
The three higher-genus objects used in Theorem~D are:
\begin{align}
 \operatorname{Obs}^{\mathrm{def}}_g(\cA)
 &\in H^2\!\left(\Def_g(\cA)\right),
 \label{eq:obs-g-projection}\\
 \mathfrak O^K_g(\cA)
 &\in K^0\!\left(\overline{\cM}_g\right)\otimes\mathbb C,\\
 \operatorname{obs}^{\mathrm{Hdg}}_g(\cA)
 &:=\operatorname{ch}_g\!\left(\mathfrak O^K_g(\cA)\right)
 \in H^{2g}\!\left(\overline{\cM}_g,\mathbb C\right).
\end{align}
Here $\operatorname{ch}_g$ denotes the codimension-$g$ component of
the Chern character.  Thus
$\operatorname{Obs}^{\mathrm{def}}_g$ has deformation degree~$2$,
whereas $\operatorname{obs}^{\mathrm{Hdg}}_g$ has cohomological
degree~$2g$ on moduli.

The higher-genus comparison package $H_D^K(\cA,g)$ consists of a
perfect scalar virtual bar object and an equality
\[
 \beta_g(\cA)\colon
 \mathfrak O^K_g(\cA)
 \xrightarrow{\ \sim\ }
 \kappa(\cA)\,\lambda_{-1}(\mathbb E_g)
 \quad\text{in completed }K^0(\overline{\cM}_g),
\]
compatible with the modular trace and clutching maps.  The package
$H_D^{\mathrm{tr}}(\cA,g)$ further identifies the scalar realization
of $\operatorname{Obs}^{\mathrm{def}}_g(\cA)$ with the virtual object
on the left.  The genus-one package $H_D^1(\cA)$ is a normalized map
on the pointed deformation complex
\[
 \operatorname{tr}_1\colon H^2(\Def_{1,1}(\cA))
 \longrightarrow H^2(\overline{\cM}_{1,1})
\]
sending the loop-contraction coefficient to
$\kappa(\cA)\lambda_1$.  The graph package
$H_D^{\mathrm{graph}}(\cA,g)$ is a stable-graph functional from the
genus-$g$ modular class to~$\mathbb C$ whose scalar-diagonal value is
$\kappa(\cA)\int_{\overline{\cM}_{g,1}}
\psi_1^{2g-2}\lambda_g$ and whose mixed-channel value is denoted
$\delta F_g^{\mathrm{cross}}(\cA)$.  Thus $H_D^{\mathrm{tr}}$ and
$H_D^K$ govern the deformation-to-Hodge comparison, while $H_D^1$
and $H_D^{\mathrm{graph}}$ fix the genus-one and numerical
normalizations.
\end{definition}

\begin{theorem}[Maurer--Cartan obstruction;
\ClaimStatusProvedHere]
\label{thm:obstruction-general}
\textup{[Type signature: Open quadrant; modular-convolution
presentation; level $4$; complete filtered dg Lie package
$H_D^{\mathrm{MC}}$.]}
Let $(L,d,[\ ,\ ])$ be a complete filtered dg Lie algebra and let
\(
 \theta_{<g}(t)=\sum_{r=1}^{g-1}t^r\theta_r
\)
satisfy the Maurer--Cartan equation modulo~$t^g$.  The coefficient
\begin{equation}
 o_g
 :=\frac12\sum_{i+j=g}[\theta_i,\theta_j]
 \in L^2
\end{equation}
is a $d_{\theta_{<g}}$-cocycle modulo the preceding filtration.  Its
class $[o_g]\in H^2(L,d_{\theta_{<g}})$ vanishes precisely when a
coefficient~$\theta_g$ extends the solution through order~$t^g$.
For $g\geq2$, application to $L=\Def_g(\cA)$ gives
$\operatorname{Obs}^{\mathrm{def}}_g(\cA)$.
\end{theorem}

\begin{proof}
Write
\(
 \mathcal F(\theta)=d\theta+\frac12[\theta,\theta].
\)
The graded Jacobi identity gives the Bianchi identity
\(
 d_\theta\mathcal F(\theta)=0.
\)
Taking the coefficient of~$t^g$ shows that~$o_g$ is closed.  Replacing
$\theta_g$ changes this coefficient by
$d_{\theta_{<g}}\theta_g$, so extension through order~$t^g$ is
equivalent to the vanishing of~$[o_g]$.
\end{proof}

\subsection{The first calculation}
\label{subsec:heisenberg-obstruction}

\begin{theorem}[Heisenberg genus-one class and higher Hodge shadow;
\ClaimStatusConditional]
\label{thm:heisenberg-obs}
\textup{[Type signature: Open quadrant; Heisenberg chart
$\mathcal H_\kappa$; levels $4\to5$; genus-one trace package
$H_D^1$ and higher comparison package $H_D^K$.]}
For the rank-one Heisenberg algebra in the trace normalization
\([a_m,a_n]=\kappa m\delta_{m+n,0}\mathbf1\), the genus-one
deformation class has scalar trace
\begin{equation}
 \operatorname{tr}_1\!\left(
 \operatorname{Obs}^{\mathrm{def}}_{1,1}(\mathcal H_\kappa)
 \right)
 =\kappa\lambda_1
 \in H^2(\overline{\cM}_{1,1}).
\end{equation}
For $g\ge2$, assume $H_D^K(\mathcal H_\kappa,g)$.  Then
\begin{align}
 \mathfrak O_g^K(\mathcal H_\kappa)
 &=\kappa\lambda_{-1}(\mathbb E_g),\\
 \operatorname{obs}^{\mathrm{Hdg}}_g(\mathcal H_\kappa)
 &=(-1)^g\kappa\lambda_g
 \in H^{2g}(\overline{\cM}_g).
\end{align}
The genus-one deformation class and the genus-one Hodge shadow are
distinct normalizations:
\[
 \operatorname{tr}_1\operatorname{Obs}^{\mathrm{def}}_{1,1}
 =\kappa\lambda_1,
 \qquad
 \operatorname{ch}_1\bigl(\kappa\lambda_{-1}(\mathbb E_1)\bigr)
 =-\kappa\lambda_1.
\]
\end{theorem}

\begin{proof}
The package $H_D^1$ identifies the normalized loop trace of the
genus-one loop graph with the curvature of the Hodge line; its
coefficient is the Heisenberg level~$\kappa$.  This gives the first
display.  For the higher class, pass to a splitting space with Chern
roots $x_1,\ldots,x_g$ for~$\mathbb E_g$.  Then
\[
 \operatorname{ch}\!\left(\lambda_{-1}(\mathbb E_g)\right)
 =\prod_{i=1}^g(1-e^{x_i})
 =(-1)^g x_1\cdots x_g+\text{terms of codimension greater than }g.
\]
Since $x_1\cdots x_g=c_g(\mathbb E_g)=\lambda_g$, the
codimension-$g$ component is $(-1)^g\lambda_g$.
\end{proof}

\begin{remark}[The anomaly coefficient]
\label{rem:heisenberg-anomaly}
The coefficient~$\kappa$ is fixed at genus~$1$.  The class
$\lambda_g$ belongs to the geometry of~$\overline{\cM}_g$; the sign
$(-1)^g$ belongs to the Chern character of the alternating exterior
power $\lambda_{-1}(\mathbb E_g)$.
\end{remark}

\subsection{Affine and multi-weight charts}
\label{subsec:kac-moody-obstruction}

\begin{theorem}[Affine scalar Hodge shadow;
\ClaimStatusConditional]
\label{thm:kac-moody-obs}
\textup{[Type signature: Open quadrant; affine chart $V_k(\fg)$;
levels $4\to5$; trace-form normalization and $H_D^K(V_k(\fg),g)$.]}
Let~$\fg$ be simple, with long roots of square length~$2$.  The scalar
coefficient is
\[
 \kappa(V_k(\fg))
 =\frac{(k+h^\vee)\dim\fg}{2h^\vee}.
\]
The genus-one trace and the higher Hodge shadow are
\begin{align}
 \operatorname{tr}_1\operatorname{Obs}^{\mathrm{def}}_{1,1}(V_k(\fg))
 &=\frac{(k+h^\vee)\dim\fg}{2h^\vee}\lambda_1,\\
 \operatorname{obs}^{\mathrm{Hdg}}_g(V_k(\fg))
 &=(-1)^g\frac{(k+h^\vee)\dim\fg}{2h^\vee}\lambda_g
 \qquad(g\ge2).
\end{align}
In particular the coefficients for
$\mathfrak{sl}_2$, $\mathfrak{sl}_3$, and~$E_8$ are respectively
\[
 \frac{3(k+2)}4,\qquad
 \frac{4(k+3)}3,\qquad
 \frac{62(k+30)}{15}.
\]
At the critical level the scalar Hodge component is zero; the
Feigin--Frenkel centre remains an operator-level structure.
\end{theorem}

\begin{proof}
The trace-form contraction gives the displayed value of~$\kappa$.
The genus-one assertion is the package $H_D^1$ in this normalization.
The higher assertion follows from $H_D^K$ and the Chern-root
calculation in Theorem~\ref{thm:heisenberg-obs}.
\end{proof}

\begin{remark}[Scope of the affine comparison]
\label{rem:kac-moody-obs-scope}
The OPE determines~$\kappa(V_k(\fg))$.  The comparison
$\beta_g(V_k(\fg))$ identifies the scalar virtual bar object with
$\kappa(V_k(\fg))\lambda_{-1}(\mathbb E_g)$.  Thus the OPE
calculation and the higher Hodge statement occupy successive levels
of the Beilinson tower.
\end{remark}

\begin{remark}[Rank exchange]
\label{rem:level-rank-obstruction}
For type~$A$, the scalar coefficients
\[
 \frac{(k+N)(N^2-1)}{2N},
 \qquad
 \frac{(N+k)(k^2-1)}{2k}
\]
are exchanged by $(N,k)\leftrightarrow(k,N)$.  A level--rank theorem
for obstruction objects additionally requires a comparison of the
two modular trace functors.
\end{remark}

\label{subsec:w-algebra-obstruction}
\begin{theorem}[$W_3$ scalar trace and mixed channels;
\ClaimStatusConditional]
\label{thm:w3-obstruction}
\textup{[Type signature: Open quadrant; $W_3$ chart; levels $4\to5$;
genus-one trace package and stable-graph package
$H_D^{\mathrm{graph}}$.]}
In the normalization whose two strong generators contribute
$c/2$ and~$c/3$, one has
\[
 \operatorname{tr}_1\operatorname{Obs}^{\mathrm{def}}_{1,1}(W_3)
 =\frac{5c}{6}\lambda_1.
\]
The scalar virtual Hodge component, when $H_D^K(W_3,g)$ holds, is
\[
 \operatorname{obs}^{\mathrm{Hdg}}_g(W_3)
 =(-1)^g\frac{5c}{6}\lambda_g.
\]
The all-weight free energy is a stable-graph trace
\[
 F_g(W_3)=\frac{5c}{6}\lambda_g^{\mathrm{FP}}
 +\delta F_g^{\mathrm{cross}}(W_3),
\]
where $\delta F_g^{\mathrm{cross}}$ collects graphs whose edges meet
distinct conformal-weight channels.
\end{theorem}

\begin{computation}[The elliptic $W_3$ value;
\ClaimStatusConditional]
\label{comp:w3-obs-explicit}
\textup{[Type signature: Open quadrant; $W_3$ chart; level $5$;
$H_D^1(W_3)$.]}
The Hodge integral $\int_{\overline{\cM}_{1,1}}\lambda_1=1/24$
gives
\[
 \int_{\overline{\cM}_{1,1}}
 \operatorname{tr}_1\operatorname{Obs}^{\mathrm{def}}_{1,1}(W_3)
 =\frac{5c}{144}.
\]
\end{computation}

\subsection{Mumford's relation}
\label{subsec:obstruction-squares-zero}

\begin{theorem}[Square of the Hodge shadow;
\ClaimStatusProvedElsewhere]
\label{thm:obstruction-nilpotent}
\label{thm:obstruction-nilpotent-all-genera}
\label{thm:obstruction-nilpotent-higher}
\textup{[Type signature: Open quadrant; scalar Hodge presentation;
level $5$; Mumford relation for the Hodge bundle.]}
For every~$g\ge1$,
\begin{equation}
 \lambda_g^2=0,
 \qquad
 \bigl(\operatorname{obs}^{\mathrm{Hdg}}_g(\cA)\bigr)^2=0
 \quad\text{whenever }H_D^K(\cA,g)\text{ holds}.
 \label{eq:obs-squared-zero-all-genera}
\end{equation}
This is a relation among characteristic classes.  The curved
$A_\infty$ identities remain governed by the Maurer--Cartan equation
of Theorem~\ref{thm:obstruction-general}.
\end{theorem}

\begin{proof}
Mumford's relation for the Hodge bundle is
\begin{equation}
 \Lambda_g(t)\Lambda_g(-t)=1,
 \qquad
 \Lambda_g(t)=\sum_{i=0}^g\lambda_i t^i.
 \label{eq:mumford-relation}
\end{equation}
The coefficient of~$t^{2g}$ is $(-1)^g\lambda_g^2$, hence it
vanishes.  Scalar multiplication gives the second assertion.
\end{proof}

\begin{remark}[Cohomological and deformation products]
\label{rem:mumford-vs-dimension}
The cup square in~$H^*(\overline{\cM}_g)$ and the bracket obstruction
in~$H^2(\Def_g(\cA))$ are products in distinct graded algebras.
Mumford's relation controls the former; the Bianchi identity controls
the latter.
\end{remark}

\begin{remark}[Mixed channels]
\label{rem:multi-generator-nilpotence}
The scalar Hodge summand has square zero.  Mixed stable-graph channels
take values in the full modular deformation complex and satisfy its
Maurer--Cartan relations.
\end{remark}

\begin{corollary}[Exterior-power form of the Mumford relation;
\ClaimStatusProvedElsewhere]
\label{cor:mumford-multiplicative}
\textup{[Type signature: Open quadrant; Hodge $K$-theory
presentation; level $5$; Mumford relation.]}
With Chern roots~$x_i$,
\begin{equation}
 \operatorname{ch}\bigl(\lambda_{-1}(\mathbb E_g)\bigr)
 =\prod_{i=1}^g(1-e^{x_i}),
 \label{eq:mumford-degree-k}
\end{equation}
and its first nonzero codimension is~$g$.  Equivalently,
\begin{equation}
 \operatorname{ch}_j\bigl(\lambda_{-1}(\mathbb E_g)\bigr)=0
 \ (j<g),
 \qquad
 \operatorname{ch}_g\bigl(\lambda_{-1}(\mathbb E_g)\bigr)
 =(-1)^g\lambda_g.
 \label{eq:lambda-recurrence}
\end{equation}
\end{corollary}

\begin{remark}[Chern character]
\label{rem:chern-character-mumford}
The sign in~\eqref{eq:lambda-recurrence} follows from the leading term
$1-e^{x_i}=-x_i+O(x_i^2)$ in every Chern root.  It varies with genus
and appears explicitly in every Hodge-shadow formula.
\end{remark}

\begin{example}[Heisenberg Hodge shadow]
\label{ex:heisenberg-obs-squares}
For $\mathcal H_\kappa$ under~$H_D^K$,
\[
 \operatorname{obs}^{\mathrm{Hdg}}_g
 =(-1)^g\kappa\lambda_g,
 \qquad
 \bigl(\operatorname{obs}^{\mathrm{Hdg}}_g\bigr)^2=0.
\]
\end{example}

\subsection{The scalar-to-Hodge comparison}

\begin{definition}[Scalar Hodge comparison package;
\ClaimStatusDefinitional]
\label{def:scalar-diagonal-hypothesis}
The package $H_D^K(\cA,g)$ consists of:
\begin{enumerate}[label=\textup{(H\arabic*)}]
\item a perfect scalar summand of the genus-$g$ virtual bar object;
\item a modular trace map from that summand to completed
$K^0(\overline{\cM}_g)$;
\item the comparison~$\beta_g(\cA)$ with
$\kappa(\cA)\lambda_{-1}(\mathbb E_g)$;
\item compatibility with separating and irreducible clutching;
\item compatibility of the degree-$g$ Chern character with the
chosen scalar trace.
\end{enumerate}
Uniform conformal weight is evidence for~(H3); the comparison itself
is the mathematical content of~(H3).
\end{definition}

\begin{proposition}[Hodge shadow of the scalar virtual object;
\ClaimStatusConditional]
\label{prop:scalar-obstruction-hodge-euler}
\textup{[Type signature: Open quadrant; scalar virtual-bar
presentation; levels $4\to5$; $H_D^K(\cA,g)$.]}
Assume $H_D^K(\cA,g)$.  Then
\begin{equation}
 \mathfrak O_g^K(\cA)
 =\kappa(\cA)\lambda_{-1}(\mathbb E_g),
 \qquad
 \operatorname{obs}^{\mathrm{Hdg}}_g(\cA)
 =(-1)^g\kappa(\cA)\lambda_g.
 \label{eq:scalar-obs-hodge-euler}
\end{equation}
\end{proposition}

\begin{proof}
The first equality is~$\beta_g(\cA)$.  Apply the additive Chern
character and use~\eqref{eq:lambda-recurrence}.
\end{proof}

\begin{remark}[Comparison with the deformation class]
\label{rem:scalar-diagonal-honest}
The equality in Proposition~\ref{prop:scalar-obstruction-hodge-euler}
belongs to $K$-theory and~$H^{2g}$.  For $g\geq2$, a comparison with
$\operatorname{Obs}^{\mathrm{def}}_g(\cA)\in H^2(\Def_g(\cA))$
is supplied by~$H_D^{\mathrm{tr}}(\cA,g)$.  At genus~$1$ the trace
normalization gives $+\kappa\lambda_1$, while the alternating
exterior-power object gives $-\kappa\lambda_1$.
\end{remark}

\begin{lemma}[Exterior-power globalization;
\ClaimStatusConditional]
\label{lem:k-theoretic-globalization-bar}
\textup{[Type signature: Open quadrant; scalar virtual-bar
presentation; level $5$; $H_D^K(\cA,g)$.]}
The comparison~$\beta_g(\cA)$ identifies the scalar virtual class
with
\[
 [\bar B^{(g)}_{\mathrm{scalar}}(\cA)]^{\mathrm{vir}}
 =\kappa(\cA)\lambda_{-1}(\mathbb E_g).
\]
Its Chern character begins in codimension~$g$ with
$(-1)^g\kappa(\cA)\lambda_g$.
\end{lemma}

\begin{proof}
The first statement is the comparison datum.  The second is the
splitting-principle calculation of
Corollary~\ref{cor:mumford-multiplicative}.
\end{proof}

\subsection{Determinant lines and Euler classes}

Grothendieck--Riemann--Roch for the universal curve gives
\begin{equation}
 \operatorname{ch}(R\pi_*\omega_\pi)
 =\pi_*\!\left(\operatorname{ch}(\omega_\pi)
 \operatorname{Td}(T_\pi)\right),
 \qquad
 [R\pi_*\omega_\pi]=[\mathbb E_g]-[\mathcal O]
 \quad\text{in }K^0.
 \label{eq:grr-hodge}
\end{equation}
For a normalized basis of holomorphic differentials, the Hodge metric
is
\begin{equation}
 \langle\omega_\alpha,\omega_\beta\rangle_{L^2}
 =\frac{i}{2}\int_{\Sigma_g}
 \omega_\alpha\wedge\overline{\omega_\beta}
 =(\operatorname{Im}\Omega)_{\alpha\beta}.
 \label{eq:arakelov-l2-inner-product}
\end{equation}
Bismut--Gillet--Soul\'e's curvature theorem for a Hermitian bundle
$F$ on the smooth universal curve over~$\cM_g$ reads
\begin{equation}
 c_1\!\left(\det R\pi_*F,h_Q\right)
 =\left[
 \pi_*\!\left(\operatorname{Td}(T_\pi,h_\pi)
 \operatorname{ch}(F,h_F)\right)
 \right]_{(2)}.
 \label{eq:bgs-anomaly}
\end{equation}
Extension of this differential-form identity across the
Deligne--Mumford boundary requires the corresponding logarithmic
metric and current package.
Thus the determinant line governs a degree-$2$ class.  For the Hodge
bundle, $c_1(\det\mathbb E_g)=c_1(\mathbb E_g)=\lambda_1$.

The Chern connection of~$\mathbb E_g$ has a base-valued curvature
matrix
\begin{equation}
 \Theta_{\mathbb E_g}
 =\bar\partial\!\left(h_{\mathbb E_g}^{-1}
 \partial h_{\mathbb E_g}\right)
 \in\Omega^{1,1}\!\left(\cM_g,\operatorname{End}(\mathbb E_g)\right).
 \label{eq:arakelov-curvature-identification}
\end{equation}
Its Chern--Weil Euler form is
\begin{equation}
 \det\!\left(\frac{i}{2\pi}\Theta_{\mathbb E_g}\right)
 =c_g(\mathbb E_g)=\lambda_g.
 \label{eq:chern-weil-lambda-g}
\end{equation}
Equations~\eqref{eq:bgs-anomaly} and
\eqref{eq:chern-weil-lambda-g} concern different characteristic
polynomials: the first treats the determinant line, and the second
treats the rank-$g$ Euler class.  The vertical Arakelov form on a
fiber and the horizontal matrix~$\Theta_{\mathbb E_g}$ are related by
variation of Hodge structure and fiber integration.  The comparison
with the scalar bar curvature is precisely~$H_D^K$.

\begin{remark}[Role of the index theorem]
\label{rem:no-clutching-needed}
Quillen~\cite{Quillen85} and Bismut--Gillet--Soul\'e
\cite{BGS88c} provide the determinant-line curvature.
Mumford~\cite{Mumford83} provides the tautological Chern-class
relations.  The scalar bar-to-Hodge morphism~$\beta_g(\cA)$ is the
remaining comparison datum.
\end{remark}

\subsection{Clutching}

\begin{proposition}[Clutching formulas for the Hodge Euler class;
\ClaimStatusProvedElsewhere]
\label{prop:lambda-g-clutching}
\textup{[Type signature: Open quadrant; Hodge-bundle presentation;
level $5$; stable-curve dualizing-sheaf package.]}
For the separating and irreducible clutching maps,
\[
 \xi_h\colon\overline{\cM}_{h,1}\times
 \overline{\cM}_{g-h,1}\longrightarrow\overline{\cM}_g,
 \qquad
 \xi_{\mathrm{irr}}\colon\overline{\cM}_{g-1,2}
 \longrightarrow\overline{\cM}_g,
\]
one has
\begin{equation}
 \xi_h^*\lambda_g=\lambda_h\boxtimes\lambda_{g-h},
 \label{eq:lambda-g-separating}
\end{equation}
and
\begin{equation}
 \xi_{\mathrm{irr}}^*\lambda_g=0.
 \label{eq:lambda-g-nonseparating}
\end{equation}
Moreover,
$\xi_{\mathrm{irr}}^*\lambda_{-1}(\mathbb E_g)=0$ in $K^0$.
\end{proposition}

\begin{proof}
Separating normalization gives
$\xi_h^*\mathbb E_g\cong
\mathrm{pr}_1^*\mathbb E_h\oplus
\mathrm{pr}_2^*\mathbb E_{g-h}$, and Whitney's formula gives
\eqref{eq:lambda-g-separating}.  Irreducible normalization gives an
exact sequence
\[
 0\longrightarrow\mathbb E_{g-1}\longrightarrow
 \xi_{\mathrm{irr}}^*\mathbb E_g\longrightarrow\mathcal O
 \longrightarrow0.
\]
Its top Chern class vanishes.  Multiplicativity of~$\lambda_{-1}$
and $\lambda_{-1}(\mathcal O)=0$ give the $K$-theory assertion.
\end{proof}

\begin{proposition}[Boundary detection under an injectivity package;
\ClaimStatusConditional]
\label{prop:clutching-uniqueness}
\textup{[Type signature: Open quadrant; tautological Hodge
presentation; level $5$; boundary-detection package $H_D^{\partial}$.]}
Let $R_g$ be the joint map formed by all clutching pullbacks and the
Faber--Pandharipande functional
\(
 \alpha\mapsto\int_{\overline{\cM}_{g,1}}
 \psi_1^{2g-2}\alpha.
\)
If $R_g$ is injective on a subspace~$T_g\subset
R^g(\overline{\cM}_g)$ containing~$\lambda_g$, then the clutching
formulas and the Faber--Pandharipande value determine~$\lambda_g$
inside~$T_g$.
\end{proposition}

\begin{proof}
Two classes with the stated data have difference in~$\ker R_g$.
Injectivity on~$T_g$ makes that difference zero.
\end{proof}

\begin{conjecture}[Boundary detection in the chosen obstruction
space; \ClaimStatusConjectured]
\label{conj:F6-lambda-g-clutching-uniqueness}
For the tautological obstruction space~$T_g$ selected by the modular
trace, the map~$R_g|_{T_g}$ is injective.
\end{conjecture}

\begin{remark}[Status of boundary detection]
\label{rem:F6-status}
Faber--Pandharipande~\cite{FP03} prove the $\lambda_g$ Hodge-integral
formula.  Conjecture~\ref{conj:F6-lambda-g-clutching-uniqueness}
concerns injectivity of a boundary-and-integral map on the manuscript's
chosen obstruction space.
\end{remark}

\begin{remark}[Clutching scope]
\label{rem:clutching-uniqueness-scope}
Clutching computes the restrictions of~$\lambda_g$.  Equality of an
independently constructed class with~$\lambda_g$ follows from the
injectivity package of Proposition~\ref{prop:clutching-uniqueness}.
\end{remark}

\begin{remark}[Role in Theorem D]
\label{rem:clutching-uniqueness-role}
The $K$-theoretic route uses~$H_D^K$.  The boundary route uses
$H_D^{\partial}$.  Agreement of the two routes is a comparison of
their constructed classes after both packages hold.
\end{remark}

\begin{proposition}[Shifted-symplectic route;
\ClaimStatusConditional]
\label{prop:theorem-D-factorization-homology-alt}
\textup{[Type signature: Open quadrant; derived mapping-stack
presentation; levels $4\to5$; PTVV orientation, properness, and trace
comparison package $H_D^{\mathrm{PTVV}}$.]}
Assume a proper oriented derived mapping stack carrying the PTVV
shifted symplectic form, together with a trace comparison to the
scalar virtual bar object.  Its determinant/Euler characteristic map
produces the same class~$\mathfrak O_g^K(\cA)$ as~$\beta_g(\cA)$.
Consequently its degree-$g$ Chern character is
$(-1)^g\kappa(\cA)\lambda_g$.
\end{proposition}

\begin{proof}
The PTVV orientation supplies the shifted symplectic pairing.  The
trace comparison identifies the resulting perfect virtual object
with~$\kappa(\cA)\lambda_{-1}(\mathbb E_g)$.  Proposition
\ref{prop:scalar-obstruction-hodge-euler} gives its Hodge shadow.
\end{proof}

\begin{remark}[Chain-level scope of the shifted-symplectic route]
\label{rem:ptvv-aksz-chain-level-scope}
The shifted symplectic form and the scalar trace comparison are
separate inputs.  The package $H_D^{\mathrm{PTVV}}$ records both.
\end{remark}

\begin{remark}[Comparison of the two routes]
\label{rem:ptvv-independence}
The $K$-theoretic and shifted-symplectic constructions share their
trace comparison target.  Their equality follows from the named
comparison morphism, while their geometric constructions remain
distinct.
\end{remark}

\subsection{Theorem D}
\label{subsec:obstruction-summary-table}

\begin{table}[ht]
\centering
\caption{The four Theorem-D objects and their ambient groups.}
\label{tab:obstruction-summary}
\begin{tabular}{lll}
\hline
Object & Ambient group & Formula / status \\
\hline
$\operatorname{tr}_1\operatorname{Obs}^{\mathrm{def}}_{1,1}$
& $H^2(\overline{\cM}_{1,1})$
& $\kappa\lambda_1$ under $H_D^1$ \\
$\mathfrak O_g^K$
& $K^0(\overline{\cM}_g)$
& $\kappa\lambda_{-1}(\mathbb E_g)$ under $H_D^K$ \\
$\operatorname{obs}^{\mathrm{Hdg}}_g$
& $H^{2g}(\overline{\cM}_g)$
& $(-1)^g\kappa\lambda_g$ under $H_D^K$ \\
$F_g$
& $\mathbb C$
& $\kappa\lambda_g^{\mathrm{FP}}+\delta F_g^{\mathrm{cross}}$
under $H_D^{\mathrm{graph}}$ \\
\hline
\end{tabular}
\end{table}

\begin{theorem}[Theorem D: genus-one deformation and higher Hodge
shadows; \ClaimStatusConditional]
\label{thm:genus-universality}
\textup{[Type signature: Open quadrant; modular-trace presentation;
levels $4\to5$; package
$H_D=(H_D^1,H_D^K,H_D^{\mathrm{tr}},H_D^{\mathrm{graph}})$.]}
Let~$\cA=A_b$ be a cyclic chart algebra with modular characteristic
$\kappa(\cA)$.  For every $g\geq2$, its native deformation class is
\[
 \operatorname{Obs}^{\mathrm{def}}_g(\cA)
 \in H^2\!\left(\Def_g(\cA)\right).
\]
The four projections satisfy:
\begin{enumerate}[label=\textup{(\roman*)}]
\item Under~$H_D^1(\cA)$,
\[
 \operatorname{tr}_1\operatorname{Obs}^{\mathrm{def}}_{1,1}(\cA)
 =\kappa(\cA)\lambda_1
 \in H^2(\overline{\cM}_{1,1}).
\]
\item Under~$H_D^K(\cA,g)$,
\begin{equation}
 \mathfrak O_g^K(\cA)
 =\kappa(\cA)\lambda_{-1}(\mathbb E_g),
 \qquad
 \operatorname{obs}^{\mathrm{Hdg}}_g(\cA)
 =(-1)^g\kappa(\cA)\lambda_g.
 \label{eq:genus-universality}
\end{equation}
\item Under~$H_D^{\mathrm{tr}}(\cA,g)$, the degree-$2$ deformation
class and the perfect virtual object are related by the named derived
trace comparison.  This clause carries the passage from deformation
theory to characteristic classes.
\item Under~$H_D^{\mathrm{graph}}(\cA,g)$, the numerical genus
coefficient is
\[
 F_g(\cA)
 =\kappa(\cA)\lambda_g^{\mathrm{FP}}
 +\delta F_g^{\mathrm{cross}}(\cA),
 \qquad
 \lambda_g^{\mathrm{FP}}
 =\int_{\overline{\cM}_{g,1}}\psi_1^{2g-2}\lambda_g.
\]
The scalar-diagonal graph package makes
$\delta F_g^{\mathrm{cross}}(\cA)=0$.
\end{enumerate}
The Faber--Pandharipande pairings of the Hodge shadows and the scalar
free-energy series are
respectively
\[
 \sum_{g\ge1}(-1)^g\kappa(\cA)\lambda_g^{\mathrm{FP}}x^{2g}
 =\kappa(\cA)\left(\frac{x/2}{\sinh(x/2)}-1\right),
\]
and
\[
 \sum_{g\ge1}\kappa(\cA)\lambda_g^{\mathrm{FP}}x^{2g}
 =\kappa(\cA)\left(\frac{x/2}{\sin(x/2)}-1\right).
\]
\end{theorem}

\begin{proof}
Clause~(i) is the genus-one trace normalization.  Clause~(ii) is
Proposition~\ref{prop:scalar-obstruction-hodge-euler}.  Clause~(iii)
is the comparison morphism in~$H_D^{\mathrm{tr}}$.  Clause~(iv) is
the stable-graph trace decomposition.  Faber--Pandharipande
\cite{FP03} give
\[
 \lambda_g^{\mathrm{FP}}
 =\frac{2^{2g-1}-1}{2^{2g-1}}
 \frac{|B_{2g}|}{(2g)!}.
\]
The two generating functions follow by inserting respectively the
factor~$(-1)^g$ and the positive Hodge-integral coefficients.
\end{proof}

\begin{remark}[Hierarchy of hypotheses]
\label{rem:universality-hierarchy}
The OPE fixes~$\kappa$.  The genus-one trace package places
$\kappa\lambda_1$ in degree~$2$.  The perfect-object comparison
places $\kappa\lambda_{-1}(\mathbb E_g)$ in $K$-theory.  The Chern
character then produces $(-1)^g\kappa\lambda_g$.  Stable graphs add
$\delta F_g^{\mathrm{cross}}$ to the numerical trace.
\end{remark}

\begin{remark}[Several conformal weights]
\label{rem:multi-generator-obs}
For generators of weights $h_1,\ldots,h_r$, the genus-one trace is
the sum of their normalized loop coefficients.  At higher genus the
scalar Hodge summand is governed by~$H_D^K$, while mixed edge labels
produce~$\delta F_g^{\mathrm{cross}}$.  The scalar-diagonal graph
package is the condition that concentrates the graph trace on the
single Hodge summand.
\end{remark}

\begin{remark}[Higher Hodge bundles]
\label{rem:higher-hodge-genus2-diagnostic}
The bundles $R^0\pi_*\omega_\pi^{\otimes h}$ for $h\ge2$ have their
own Chern characters.  A modular trace involving these bundles lies
in an enlarged comparison package.  Theorem~D uses the standard
Hodge bundle~$\mathbb E_g=R^0\pi_*\omega_\pi$ through the explicit
map~$\beta_g(\cA)$.
\end{remark}


\begin{proposition}[Multi-weight refinement of the scalar graph trace;
\ClaimStatusConditional]%
\label{op:multi-generator-universality}%
\index{multi-generator obstruction!resolution|textbf}%
\index{cross-channel correction!resolution of universality|textbf}%
For multi-weight families such as\/ $\mathcal{W}_N^k$ with
$N \geq 3$, the all-weight genus-$g$ coefficient at $g\geq2$
decomposes as
\begin{equation}\label{eq:multi-weight-decomposition-op}
F_g(\cA)
\;=\;
F_g^{\mathrm{sc}}(\cA)
\;+\;
\delta F_g^{\mathrm{cross}}(\cA),
\end{equation}
Here \(F_g^{\mathrm{sc}}(\cA)=
\kappa(\cA)\lambda_g^{\mathrm{FP}}\), and
$\delta F_g^{\mathrm{cross}}$ is an explicit graph sum over
mixed-channel boundary graphs of\/ $\overline{\mathcal{M}}_g$
(Theorem~\ref{thm:multi-weight-genus-expansion}). For $\cW_3$
at genus~$2$:
\begin{equation}\label{eq:w3-cross-channel-op}
\delta F_2(\cW_3)
\;=\;
\underbrace{\frac{3}{c}}_{\text{banana}}
+\;
\underbrace{\frac{9}{2c}}_{\text{theta}}
+\;
\underbrace{\frac{1}{16}}_{\text{lollipop}}
+\;
\underbrace{\frac{21}{4c}}_{\text{barbell}}
\;=\;
\frac{c + 204}{16c}
\;>\; 0
\quad\text{for all $c > 0$}.
\end{equation}
The correction is $R$-matrix independent
(Theorem~\ref{thm:multi-weight-genus-expansion}(v))
and vanishes if and only if every generator has the same
conformal weight. The genus-$1$ identity
\(\operatorname{tr}_1\operatorname{Obs}^{\mathrm{def}}_{1,1}(\cA)=\kappa(\cA)\lambda_1\)
holds under~$H_D^1$ for every family carrying that trace package.

The structural diagnosis has three parts.
\textup{(A)}~Algebraic-family rigidity
(Theorem~\ref{thm:algebraic-family-rigidity}) forces
$\dim H^2_{\mathrm{cyc}}(\cA, \cA) = 1$, so the
genus-completed minimal-model MC element lies on a cyclic line:
$\Theta_\cA^{\min} = \eta \otimes \Gamma_{\cA}$.
\textup{(B)}~The Kuranishi map vanishes identically on~$\eta$
by graded antisymmetry (odd degree in~$s^{-1}\mathcal{H}$),
so the MC equation places no further constraint on the
tautological coefficient~$\Gamma_\cA$.
\textup{(C)}~The tautological-purity identification
$\Gamma_\cA = \kappa(\cA)\Lambda$ \emph{fails} for
multi-weight algebras: the explicit $\cW_3$ genus-$2$ graph sum
(Computation~\ref{comp:w3-genus2-multichannel})
demonstrates that the integrated free energy deviates from the
scalar prediction by~\eqref{eq:w3-cross-channel-op}.

The origin of the correction is topological: at genus~$1$,
the unique boundary graph (the figure-eight) has a single edge
that carries one channel, so no mixed-channel amplitudes arise.
At genus~$\geq 2$, multi-edge boundary graphs (banana, theta,
lollipop, barbell) first appear, and mixed-channel
assignments $(T, W, W)$, $(T, W)$, etc.\ contribute nonzero
amplitudes. The $c$-independent lollipop piece $1/16$ arises
because the OPE structure constant $C_{TWW} = c$ cancels the
propagator $\eta^{TT} = 2/c$, leaving a topological residue that
persists in the semiclassical limit $c \to \infty$.
\end{proposition}

\begin{remark}[Multi-generator free energy: structural diagnosis]%
\label{rem:multigen-structural-diagnosis}%
\index{multi-generator obstruction!structural diagnosis}%
\index{Faber--Pandharipande formula!multi-generator}%
For a multi-generator algebra with generators of distinct
conformal weights $h_1, \ldots, h_r$:

\textup{(a)} The \emph{scalar graph sum} with total shadow data
$S_2 = \kappa$, $S_4 = S_4^{\mathrm{total}}$ gives
$F_2 \neq \kappa \cdot \lambda_2^{\mathrm{FP}}$
when interpreted as a full all-weight coefficient
(Proposition~\ref{prop:f2-quartic-dependence}: the banana graph
injects $S_4$-dependence with $\partial F_2/\partial S_4 =
1/(8\kappa^2)$).

\textup{(b)} The \emph{multi-channel} computation uses the
diagonal Zamolodchikov metric
$\eta = \mathrm{diag}(\kappa_{h_1}, \ldots, \kappa_{h_r})$
(orthogonality of distinct conformal weights on the sphere)
and per-channel propagators $P_i = 1/\kappa_{h_i}$.
If the per-channel Givental $R$-matrix is universal
(independent of~$h_i$) and the scalar-diagonal factorization
hypothesis holds, the CohFT factorises and
\(F_g(\cA)=F_g^{\mathrm{sc}}(\cA)
=\kappa(\cA)\lambda_g^{\mathrm{FP}}\)
\textup{(}SCALAR-DIAGONAL FACTORIZATION\textup{)}.

\textup{(c)} \emph{Failure mechanism.}
The bar propagator $d\log E(z,w)$ is weight-$1$
(Remark~\ref{rem:propagator-weight-universality}),
which controls the \emph{edges} of the graph sum.
The \emph{vertices} carry cross-channel OPE data, and the
mixed-channel amplitudes do \emph{not} cancel:
Theorem~\ref{thm:multi-weight-genus-expansion} proves
$F_g(\cA) = F_g^{\mathrm{sc}}(\cA)
+ \delta F_g^{\mathrm{cross}}(\cA)$, with
$F_g^{\mathrm{sc}}(\cA)=\kappa(\cA)\lambda_g^{\mathrm{FP}}$,
and $\delta F_2(\cW_3) = (c{+}204)/(16c) \neq 0$.
The identification
$\Gamma_{\cA} = \kappa(\cA)\Lambda$ fails for multi-weight
algebras.
\end{remark}

\begin{remark}[Propagator weight universality]%
\label{rem:propagator-weight-universality}%
\index{propagator weight universality|textbf}%
\index{bar complex!propagator weight}%
\index{prime form!weight of d log}%
The bar complex propagator is $d\log E(z,w)$, where $E(z,w)$
is the prime form on a Riemann surface. The prime form is a
section of $K^{-1/2} \boxtimes K^{-1/2}$
(cf.\ the critical pitfall in Appendix~\ref{app:signs}),
so $d\log E = dE/E$ has weight~$1$ in both variables,
\emph{regardless} of the conformal weight of the fields being sewed.

This observation has two consequences:
\begin{enumerate}[label=\textup{(\roman*)}]
\item \emph{Genus-$1$ universality:}
At $g = 1$, $H^2(\overline{\cM}_{1,1}) \cong \mathbb{C}$ forces
every class to be proportional to~$\lambda_1$, so
$\operatorname{tr}_1\operatorname{Obs}^{\mathrm{def}}_{1,1}(\cA) = \kappa \cdot \lambda_1$
under~$H_D^1$.  The one-dimensionality identifies the target line;
the trace package determines its coefficient.
The weight-$1$ propagator ensures the $B$-cycle defect produces
$\lambda_1$ (not $c_1(\mathbb{E}_h)$) for each channel.
Consistency check: if the bar complex used
$\mathbb{E}_h$ for weight-$h$ generators, the Mumford
isomorphism $c_1(\mathbb{E}_h) = (6h^2{-}6h{+}1)\lambda_1$
would multiply $F_1$ by the Mumford exponent $e(h)$, giving
(for Virasoro with $h = 2$)
$F_1 = 13\kappa/24 \neq \kappa/24$
in this counterfactual normalization.
\item \emph{Graph-sum edges in each genus:}
Every \emph{edge} in the genus-$g$ graph sum carries the standard
propagator $d\log E$, hence standard Hodge data
(Remark~\ref{rem:propagator-weight-universality}).
The \emph{vertices} of valence $\geq 4$ carry cross-channel
OPE structure constants that are not controlled by~$\kappa$ alone.
Scalar saturation ($\dim H^2_{\mathrm{cyc}} = 1$) fixes the
$H^2$~direction to~$\eta$ but does not determine which
$\overline{\mathcal{M}}_g$~class appears at genus~$g$, since
the Kuranishi map vanishes by parity.
Indeed, the scalar formula \emph{fails}: the free energy receives
a cross-channel correction
(Theorem~\ref{thm:multi-weight-genus-expansion}).
\end{enumerate}
\end{remark}

\begin{proposition}[Multi-generator edge universality; \ClaimStatusConditional]
\label{prop:multi-generator-obstruction}
\index{multi-generator obstruction}
Let $\cA$ be a modular Koszul chiral algebra with strong generators
of conformal weights $h_1, \ldots, h_r$.
\begin{enumerate}[label=\textup{(\alph*)}]
\item At genus~$1$,
$\operatorname{tr}_1\operatorname{Obs}^{\mathrm{def}}_{1,1}(\cA) = \kappa(\cA) \cdot \lambda_1$
under~$H_D^1$
with $\kappa = \sum_i \kappa_{h_i}$.
\item Every edge of the genus-$g$ graph sum carries the standard
Hodge bundle $\mathbb{E} = R^0\pi_*\omega$
\textup{(}Remark~\textup{\ref{rem:propagator-weight-universality}}\textup{)}.
\item Scalar saturation fixes the $H^2_{\mathrm{cyc}}$ direction
to~$\eta$, but the Kuranishi map vanishes by parity, so the
MC equation does not determine which class in
$H^*(\overline{\mathcal{M}}_g)$ appears at genus~$g$.
 At $g \geq 2$ for multi-weight algebras, the full all-weight
 coefficient refines the scalar trace
 \(\kappa(\cA)\lambda_g^{\mathrm{FP}}\): the full free energy
 decomposes as \(F_g=F_g^{\mathrm{sc}}+\delta F_g^{\mathrm{cross}}\)
 with $\delta F_g^{\mathrm{cross}} \neq 0$
 (Theorem~\textup{\ref{thm:multi-weight-genus-expansion}}).
\end{enumerate}
\end{proposition}

\begin{proof}
Part~(a): at genus~$1$, $H^2(\overline{\mathcal{M}}_{1,1})
\cong \mathbb{C}$, so all classes are proportional
to~$\lambda_1$. The bar propagator $d\!\log E(z,w)$ has
weight~$1$ (Remark~\ref{rem:propagator-weight-universality}),
so each channel uses~$\mathbb{E}$ (not $\mathbb{E}_{h_i}$),
and the per-channel curvatures sum to
$\kappa = \sum_i \kappa_{h_i}$.
Part~(b): same propagator-weight argument in every genus.
Part~(c): scalar saturation gives
$\Theta^{\min} = \eta\otimes\Omega$ for some
$\Omega \in \widehat{\Gmod}$; the Kuranishi map vanishes by
parity (Proposition~\ref{prop:saturation-equivalence}), so the
MC equation does not constrain~$\Omega$. The scalar-diagonal
identification would force the integrated coefficient to be
$\kappa(\cA)\lambda_g^{\mathrm{FP}}$; the multi-weight graph sum has
additional mixed-channel terms, so scalar saturation alone does not
determine the genus-$g$ tautological coefficient.
\end{proof}

\begin{proposition}[Genus-$2$ quartic dependence; \ClaimStatusProvedHere]%
\label{prop:f2-quartic-dependence}%
\index{multi-generator obstruction!quartic dependence}%
The genus-$2$ free energy, computed as a graph sum over the seven
stable graphs at $(g{=}2, n{=}0)$, depends linearly on the
quartic shadow coefficient~$S_4$:
\begin{equation}\label{eq:f2-quartic-dep}
 \frac{\partial F_2}{\partial S_4} \;=\; \frac{1}{8\kappa^2} \;\neq\; 0.
\end{equation}
The $S_4$-dependence enters through the banana graph
\textup{(}one genus-$0$ vertex with two self-loops, $|\operatorname{Aut}| = 8$\textup{)},
which has amplitude $(S_4 / \kappa^2) / 8$.
No other graph at $(2,0)$ involves~$S_4$.
\end{proposition}

\begin{proof}
The seven graphs at $(2, 0)$ have vertices of valence~$0$
(at genus~$2$ or~$1$), $1$, $2$, $3$, or~$4$ from edges plus
markings. Only the banana has a genus-$0$ vertex of
valence~$4$, giving vertex factor~$S_4$.
The banana amplitude is
$\frac{1}{8}\,S_4 \cdot P^2 = \frac{S_4}{8\kappa^2}$
(since $P = 1/\kappa$),
confirming~\eqref{eq:f2-quartic-dep}.
\end{proof}

\begin{remark}[Multi-channel resolution]%
\label{rem:multichannel-resolution}%
\index{multi-generator obstruction!multi-channel resolution}%
For a multi-generator algebra $\mathcal{W}_N$ with
$\kappa = \sum_{s=2}^N \kappa_s$, the \emph{total} quartic
$S_4^{\mathrm{total}} \neq S_4^{\mathrm{Vir}}(\kappa)$:
the $\mathcal{W}$-channel and mixed OPE contributions shift~$S_4$
away from the Virasoro value at the same~$\kappa$.
The scalar graph sum therefore gives
$F_2 \neq \kappa \cdot \lambda_2^{\mathrm{FP}}$
when interpreted as a full all-weight coefficient.

The \emph{correct} computation uses the \textbf{multi-channel CohFT}:
each edge carries a channel label $s \in \{2, \ldots, N\}$,
with per-channel propagator $P_s = 1/\kappa_s$.
At a vertex of valence~$n$, the channel labels mix through the
OPE structure constants~$C_{s_1\ldots s_n}$.
Semisimplicity of the underlying Frobenius algebra
(from the distinct conformal weights $h_s = s$, giving a
diagonal Zamolodchikov metric and distinct eigenvalues of
the $T$-multiplication) gives a possible Givental--Teleman
decomposition into channels, but does \emph{not} by itself
force cancellation of the $R$-matrix dressing and cross-channel
vertex terms in the summed free energy.
The graph sum gives
$F_g \neq \kappa \cdot \lambda_g^{\mathrm{FP}}$
when interpreted as a full all-weight coefficient at $g \ge 2$
for multi-weight families
(Theorem~\ref{thm:multi-weight-genus-expansion}).
The scalar graph sum is a \emph{projection} of the
multi-channel computation that conflates the channel structure.
\end{remark}

\begin{corollary}[Anomaly ratio identity; \ClaimStatusConditional]
\label{cor:anomaly-ratio}
\index{anomaly ratio|textbf}
For any $\mathcal{W}$-algebra $\mathcal{W}^k(\mathfrak{g})$, the obstruction coefficient $\kappa$ and central charge $c$ are related by
\begin{equation}\label{eq:anomaly-ratio}
\frac{\kappa(\mathcal{W}^k(\mathfrak{g}))}{c(\mathcal{W}^k(\mathfrak{g}))} = \varrho(\mathfrak{g}) := \sum_{i=1}^r \frac{1}{m_i + 1}
\end{equation}
where $m_1, \ldots, m_r$ are the exponents of $\mathfrak{g}$. The ratio $\varrho(\mathfrak{g})$ is a root-system invariant independent of the level~$k$.

For Kac--Moody algebras $\widehat{\mathfrak{g}}_k$ (before DS reduction), $\kappa/c = (k+h^\vee)^2/(2h^\vee k)$, which depends on the level. The level-independence of $\kappa/c$ is specific to $\mathcal{W}$-algebras and reflects the DS reduction.
\end{corollary}

\begin{proof}
From Theorem~\ref{thm:genus-universality}(ii): $\kappa = c \cdot \varrho(\mathfrak{g})$. Since $\varrho$ depends only on the exponents of $\mathfrak{g}$, it is independent of $k$.
\end{proof}

\begin{corollary}[\texorpdfstring{$\kappa$}{kappa}-periodicity under level shift; \ClaimStatusProvedHere]
\label{cor:kappa-periodicity}
\index{obstruction coefficient!periodicity}
For affine Kac--Moody $\widehat{\mathfrak{g}}_k$, the obstruction
coefficient is periodic under level shift:
$\kappa(\widehat{\mathfrak{g}}_{k+n})
= \kappa(\widehat{\mathfrak{g}}_k) + n \cdot d/(2h^\vee)$ for
$n \in \mathbb{Z}$. On the scalar-diagonal lane, the scalar free
energy $F_g$ is therefore a linear function of $k$ with slope
$\lambda_g^{\mathrm{FP}} \cdot d/(2h^\vee)$
\textup{(}UNIFORM-WEIGHT, SCALAR-DIAGONAL\textup{}).
For $\mathcal{W}$-algebras, $\kappa$ is a rational function of $k$
\textup{(}through its dependence on the DS central charge\textup{)},
and the periodicity statement becomes:
$\kappa(\mathcal{W}^{k+n}(\mathfrak{g}))
- \kappa(\mathcal{W}^k(\mathfrak{g}))$ is a rational function of $k$
with poles only at $k = -h^\vee$.
\end{corollary}

\begin{proof}
For KM: $\kappa = (k+h^\vee)d/(2h^\vee)$, which is linear in $k$ with slope $d/(2h^\vee)$. Level shift $k \to k+n$ adds $nd/(2h^\vee)$.
\end{proof}

\begin{remark}[Universality principle]\label{rem:lambda-universality}
The comparison package~$H_D^K$ gives
\[
 \mathfrak O_g^K(\cA)=\kappa(\cA)\lambda_{-1}(\mathbb E_g),
 \qquad
 \operatorname{obs}^{\mathrm{Hdg}}_g(\cA)
 =(-1)^g\kappa(\cA)\lambda_g.
\]
The graph package gives
\(F_g(\cA)=\kappa(\cA)\lambda_g^{\mathrm{FP}}
+\delta F_g^{\mathrm{cross}}(\cA)\).
The scalar-diagonal graph condition concentrates this trace on its
first summand.  Several conformal weights generate the cross-channel
summand through mixed edge labels.
\end{remark}

\begin{corollary}[Additivity of the obstruction coefficient; scalar-lane scope;
\ClaimStatusConditional]\label{cor:kappa-additivity}
\index{obstruction coefficient!additivity}
For Koszul chiral algebras $\mathcal{A}$, $\mathcal{B}$ such that
$\mathcal{A} \otimes \mathcal{B}$ is Koszul:
\begin{equation}\label{eq:kappa-additivity}
\kappa(\mathcal{A} \otimes \mathcal{B}) = \kappa(\mathcal{A}) + \kappa(\mathcal{B}).
\end{equation}
If $\mathcal{A}$, $\mathcal{B}$, and
$\mathcal{A}\otimes\mathcal{B}$ lie on the proved
scalar-diagonal lane and carry compatible $H_D^K$ packages, then
\[
	\operatorname{obs}^{\mathrm{Hdg}}_g(\mathcal{A} \otimes \mathcal{B})
	= \operatorname{obs}^{\mathrm{Hdg}}_g(\mathcal{A})
	+ \operatorname{obs}^{\mathrm{Hdg}}_g(\mathcal{B})
\qquad (g\ge1).
\]
For general all-weight tensor products this statement controls only
the additive scalar prefactor; mixed-channel corrections are separate
stable-graph data.
\end{corollary}

\begin{proof}
The singular part of the OPE of
$\mathcal{A} \otimes \mathcal{B}$ is the sum of the singular parts of
$\mathcal{A}$ and $\mathcal{B}$, since the tensor product acts on
disjoint sets of fields. The genus-$1$ curvature, which determines
$\kappa$ by Theorem~\ref{thm:genus-universality}, is therefore
additive. Under~$H_D^K$ on the scalar-diagonal lane,
	\(\operatorname{obs}^{\mathrm{Hdg}}_g(\cA)
 =(-1)^g\kappa(\cA)\lambda_g\), so
\[
	\operatorname{obs}^{\mathrm{Hdg}}_g(\mathcal{A}\otimes\mathcal{B})
	= (-1)^g(\kappa(\mathcal{A})+\kappa(\mathcal{B}))\lambda_g
	= \operatorname{obs}^{\mathrm{Hdg}}_g(\mathcal{A})
	+\operatorname{obs}^{\mathrm{Hdg}}_g(\mathcal{B}).
\]
For all-weight tensor products, the same argument proves additivity of
the scalar coefficient, while the cross-channel summands remain
separate stable-graph data.
\end{proof}

\begin{example}
For a rank-$d$ lattice vertex algebra $V_\Lambda = \mathcal{H}_1^{\otimes d}$, additivity gives $\kappa(V_\Lambda) = d \cdot \kappa(\mathcal{H}_1) = d$, confirming the Master Table entry (Table~\ref{tab:master-invariants}). For the $bc$--$\beta\gamma$ system, $\kappa(bc \otimes \beta\gamma) = c/2 + (-c/2) = 0$ (since $\kappa(\beta\gamma) = -\kappa(bc)$ by Theorem~\ref{thm:fermion-all-genera}).
\end{example}

\begin{corollary}[Obstruction complementarity for \texorpdfstring{$\mathcal{W}_N$}{W(N)}; \ClaimStatusConditional]
\label{cor:kappa-sum-wn}
\index{obstruction coefficient!W-algebra complementarity}
For the principal $\mathcal{W}$-algebra $\mathcal{W}_N = \mathcal{W}^k(\mathfrak{sl}_N)$, the obstruction complementarity sum is
\begin{equation}\label{eq:kappa-sum-wn}
\kappa(\mathcal{W}_N^k) + \kappa(\mathcal{W}_N^{k'}) = (H_N - 1) \cdot K_N
\end{equation}
where $H_N = \sum_{s=1}^N 1/s$ is the $N$-th harmonic number, $K_N = c + c' = 2(N-1)(2N^2+2N+1)$ is the Koszul conductor, and $k' = -k-2N$. The first values are:
\begin{center}
\begin{tabular}{cccc}
$N$ & $\varrho = H_N - 1$ & $K_N$ & $\kappa + \kappa'$ \\
\hline
$2$ & $1/2$ & $26$ & $13$ \\
$3$ & $5/6$ & $100$ & $250/3$ \\
$4$ & $13/12$ & $246$ & $533/2$ \\
$5$ & $77/60$ & $488$ & $9394/15$
\end{tabular}
\end{center}
\end{corollary}

\begin{proof}
By Theorem~\ref{thm:genus-universality}(ii), $\kappa(\mathcal{W}_N^k) = c \cdot \varrho(\mathfrak{sl}_N)$ where $\varrho(\mathfrak{sl}_N) = \sum_{s=2}^N 1/s = H_N - 1$. Under the Feigin--Frenkel involution $k \mapsto k' = -k - 2N$, $c \mapsto c' = K_N - c$, so $\kappa' = (K_N - c)(H_N - 1)$ and $\kappa + \kappa' = K_N(H_N - 1)$.
\end{proof}

\begin{corollary}[Universal genus-$1$ criticality criterion; scalar-lane collapse; \ClaimStatusConditional]\label{cor:critical-level-universality}
\index{critical level!characterization}
For a modular Koszul chiral algebra $\mathcal{A}$ carrying~$H_D^1$
and the genus-one graph trace, the following are equivalent:
\begin{enumerate}[label=\textup{(\roman*)}]
\item $\kappa(\mathcal{A}) = 0$;
\item $\operatorname{tr}_1\operatorname{Obs}^{\mathrm{def}}_{1,1}(\mathcal{A}) = 0$;
\item $F_1(\mathcal{A}) = 0$;
\item the vertical scalar curvature vanishes:
$\operatorname{tr}_{\mathrm{diag}}(m_0^{(1)})=0$.
\end{enumerate}
Under~$H_D^K$ and~$H_D^{\mathrm{graph}}$ on the proved scalar lane,
these equivalent conditions further imply
	$\operatorname{obs}^{\mathrm{Hdg}}_g(\mathcal{A}) = 0$ and
$F_g^{\mathrm{sc}}(\mathcal{A}) = 0$ for all $g \geq 1$.
The higher-degree shadow obstruction tower and the full
Maurer--Cartan element~$\Theta_\mathcal{A}$ may retain non-scalar
components. In algebra families with a named critical
parameter, these conditions pick out the scalar critical point. When
the scalar-diagonal faithfulness package identifies vanishing trace
with $m_0=0$, $H^0(\bar{B}^{\mathrm{ch}}(\mathcal{A}))$ is a genuine
associative chiral algebra. The internal commutator identity
$m_1^2(a)=[m_0,a]$ records the adjoint action; central curvature is
detected by the separate scalar trace.
\end{corollary}

\begin{proof}
(i)$\Leftrightarrow$(ii): By Theorem~\ref{thm:genus-universality},
the universal genus-$1$ scalar formula is
\(\operatorname{tr}_1\operatorname{Obs}^{\mathrm{def}}_{1,1}(\cA)=\kappa(\cA)\lambda_1\)
and
$\lambda_1 \neq 0$.
(i)$\Leftrightarrow$(iii): By the same theorem,
\(F_1^{\mathrm{sc}}(\cA)=\kappa(\cA)/24\)
for every modular Koszul algebra.
(i)$\Leftrightarrow$(iv): By the normalized fixed-fiber trace,
\[
 \operatorname{tr}_{\mathrm{diag}}(m_0^{(1)})
 =\kappa(\cA)\omega_1^{\mathrm{Ar}}.
\]
The Arakelov volume form is nonzero, so the scalar trace vanishes
exactly when~$\kappa(\cA)$ vanishes.  The curved identity
$m_1^2=[m_0,-]$ records the adjoint action and therefore remains a
separate statement.
The scalar-lane clause follows from Theorem~\ref{thm:genus-universality}
under the displayed packages. Under scalar-diagonal faithfulness,
$m_0=0$ makes the bar complex a strict dg coalgebra, so its cobar is
a strict dg algebra and $H^0$ inherits a strict associative product.
\end{proof}

\begin{corollary}[Tautological class map on the scalar lane; universal genus-$1$ class; \ClaimStatusConditional]\label{cor:tautological-class-map}
\index{tautological ring!map from chiral algebras}
Every modular Koszul chiral algebra $\mathcal{A}$ determines a
distinguished genus-$1$ tautological class
\[
\operatorname{tr}_1\operatorname{Obs}^{\mathrm{def}}_{1,1}(\mathcal{A}) = \kappa(\mathcal{A}) \cdot \lambda_1
\qquad \textup{(under $H_D^1$)}
\;\in\; R^1(\overline{\mathcal{M}}_1).
\]
Under~$H_D^K(\mathcal A,g)$ on the proved scalar lane, one also has
\[
\operatorname{obs}^{\mathrm{Hdg}}_g(\mathcal{A})
=(-1)^g\kappa(\mathcal{A}) \cdot \lambda_g
\qquad \textup{(}UNIFORM-WEIGHT, SCALAR-DIAGONAL\textup{)}
\qquad \;\in\; R^g(\overline{\mathcal{M}}_g)
\qquad (g \geq 1).
\]
The assignment $\mathcal{A} \mapsto \kappa(\mathcal{A})$ is additive
under tensor products:
$\kappa(\mathcal{A} \otimes \mathcal{B}) =
\kappa(\mathcal{A}) + \kappa(\mathcal{B})$
\textup{(}Corollary~\ref{cor:kappa-additivity}\textup{)}. For
Kac--Moody algebras, $\kappa$ further satisfies
$\kappa(\mathcal{A}^!) = -\kappa(\mathcal{A})$, making it a group
homomorphism from $K_0(\mathrm{KCA}_{\mathrm{KM}})$ to~$\mathbb{Q}$.
For $\mathcal{W}$-algebras the duality relation is affine
\textup{(}$\kappa' = \varrho(\mathfrak{g})\,K - \kappa$, where
$\varrho(\mathfrak{g}) = \sum 1/(m_i+1)$\textup{)}, so $\kappa$ is
additive but does not intertwine Koszul duality with negation.
\end{corollary}

\begin{proof}
The universal genus-$1$ identity
$\operatorname{tr}_1\operatorname{Obs}^{\mathrm{def}}_{1,1}(\mathcal{A}) = \kappa(\mathcal{A}) \cdot \lambda_1$
under~$H_D^1$ is Theorem~\ref{thm:genus-universality}; under~$H_D^K$ on the proved scalar lane,
the same theorem gives the higher-genus classes
	$\operatorname{obs}^{\mathrm{Hdg}}_g(\mathcal{A})
	=(-1)^g\kappa(\mathcal{A}) \cdot \lambda_g$.
	Since each $\lambda_g$ is tautological, these scalar classes lie in the
tautological ring. Additivity
$\kappa(\mathcal{A} \otimes \mathcal{B}) =
\kappa(\mathcal{A}) + \kappa(\mathcal{B})$
\textup{(}Corollary~\ref{cor:kappa-additivity}\textup{)} gives a group
homomorphism $K_0(\mathrm{KCA}) \to \mathbb{Q}$ with respect to tensor
products. For Kac--Moody families the additional relation
$\kappa(\mathcal{A}^!) = -\kappa(\mathcal{A})$
\textup{(}Theorem~\ref{thm:genus-universality}(ii)\textup{)} shows
that $\kappa$ also intertwines Koszul duality with negation.
\end{proof}

\begin{remark}\label{rem:tautological-master-table}
The Master Table of computed invariants
(Table~\ref{tab:master-invariants}) gives explicit values:
$\kappa(\mathcal{H}_1) = 1$,
$\kappa(\widehat{\mathfrak{sl}}_{2,k}) = 3(k+2)/4$,
$\kappa(\mathrm{Vir}_c) = c/2$,
$\kappa(\mathcal{W}_N^k) = c \cdot (H_N - 1)$. Each produces the
universal genus-$1$ tautological class $\kappa \cdot \lambda_1$
under~$H_D^1$, and
under~$H_D^K$ on the proved scalar lane a one-parameter family of
classes $(-1)^g\kappa \cdot \lambda_g
\in R^g(\overline{\mathcal{M}}_g)$.  Their coefficient is the scalar level
$\Theta_{\cA}^{\leq 2} = \kappa(\cA)$ of the shadow obstruction tower
(Definition~\ref{def:shadow-postnikov-tower},
Theorem~\ref{thm:explicit-theta}):
the appearance of $\kappa$ as a single controlling number
reflects the one-dimensionality of
$H^2_{\mathrm{cyc}}(\mathfrak{g},\mathfrak{g}) \cong \mathbb{C}$
for simple~$\mathfrak{g}$; there is only one direction for
the scalar level to point
(Corollary~\ref{cor:scalar-saturation}).
\end{remark}

\begin{remark}[Beyond the master table: the shadow obstruction tower]
\label{rem:beyond-master-table}
\index{shadow obstruction tower!beyond scalar level}
The master table records the scalar level
$\Theta_\cA^{\leq 2} = \kappa(\cA)$ of the shadow obstruction tower.
The cubic, quartic, and higher shadows, computed for each family
in the example chapters of Part~\ref{part:standard-landscape}, extend this table into a
multi-column shadow census: Heisenberg terminates at degree~$2$
(Corollary~\ref{cor:heisenberg-postnikov-termination}),
affine at degree~$3$
(Corollary~\ref{cor:affine-postnikov-termination}),
$\beta\gamma$ at degree~$4$
(Corollary~\ref{cor:betagamma-postnikov-termination}),
and Virasoro/$\mathcal{W}_N$ has an infinite tower
(Corollary~\ref{cor:virasoro-postnikov-nontermination}).
\end{remark}

\begin{remark}[Tautological span]\label{rem:tautological-span}
\index{tautological ring!span from Koszul algebras}
At genus~$1$, the collection
$\{\kappa(\mathcal{A}) \cdot \lambda_1 :
\mathcal{A} \text{ modular Koszul}\}$ generates a rank-$1$ subgroup of
$R^1(\overline{\mathcal{M}}_1)$, since every universal genus-$1$
scalar obstruction class is proportional to~$\lambda_1$
(Theorem~\ref{thm:genus-universality}). On the proved scalar lane,
under~$H_D^K$, the same rank-$1$ statement holds in each genus~$g$
with $(-1)^g\kappa(\mathcal{A}) \cdot \lambda_g$. The image of the group
homomorphism $K_0(\mathrm{KCA}) \to \mathbb{Q}$
\textup{(}Corollary~\ref{cor:tautological-class-map}\textup{)} is
$\mathbb{Q}$ itself: the Heisenberg algebra $\mathcal{H}_\kappa$
achieves any rational $\kappa$, and additivity
\textup{(}Corollary~\ref{cor:kappa-additivity}\textup{)} extends the
range to all of~$\mathbb{Q}$.

A natural question is whether the \emph{flat corrected} bar complex
$\bar{B}^{\mathrm{ch}}_{\mathrm{flat}}(\mathcal{A})$ carries tautological data beyond
the genus-$g$ cohomology class $\lambda_g$. The bar spectral sequence
$E_r^{p,q} \Rightarrow
H^{p+q}(\bar{B}^{(g)}_{\mathrm{flat}})$ has $E_2$ page
involving products of $\psi$-classes and $\lambda$-classes on
$\overline{\mathcal{M}}_{g,n}$, which a priori lie in a larger part of
$R^*(\overline{\mathcal{M}}_g)$. On the proved scalar lane, the
scalar projection of this spectral sequence is already the signed
Hodge Euler class $(-1)^g\kappa\cdot\lambda_g$ under~$H_D^K$; higher differentials may produce
non-scalar tautological classes, but they do not create another
rank-$1$ scalar class at the cohomological obstruction level.
\end{remark}

\begin{proposition}[Bar spectral sequence and tautological filtration;
\ClaimStatusProvedHere]\label{prop:bar-tautological-filtration}
\index{tautological ring!bar spectral sequence filtration}
The bar spectral sequence
$E_r^{p,q} \Rightarrow
H^{p+q}(\bar{B}^{(g)}_{\mathrm{flat}}(\mathcal{A}))$
induces a \emph{tautological filtration} on the bar cohomology:
\begin{equation}\label{eq:tauto-filtration}
0 = F^{3g-2} \subset F^{3g-3} \subset \cdots \subset F^0 =
H^*(\bar{B}^{(g)}_{\mathrm{flat}}(\mathcal{A}))
\end{equation}
where $F^p / F^{p+1} = E_\infty^{p,*}$. The associated graded
pieces satisfy:
\begin{enumerate}[label=\textup{(\roman*)}]
\item $E_2^{p,q}$ involves classes in
 $H^p(\overline{\mathcal{M}}_g,
 \mathcal{H}^q(\bar{B}^{(g)}_{\mathrm{flat},\bullet}))$,
 which lie in the tautological ring $R^*(\overline{\mathcal{M}}_g)$
 for Koszul algebras;
\item The differential $d_2\colon E_2^{p,q} \to E_2^{p+2,q-1}$
 is the Kodaira--Spencer map, which produces $\psi$-classes and
 boundary classes from $\lambda$-classes;
\item The differential $d_3$ involves triple Massey products,
 producing $\kappa$-classes from $\psi$-classes.
\end{enumerate}
For a Koszul algebra with the spectral sequence collapsing at
$E_2$, only $\lambda_g$ survives. If the spectral sequence does
\emph{not} collapse at $E_2$, the higher differentials produce
additional tautological classes.
\end{proposition}

\begin{proof}
The filtration~\eqref{eq:tauto-filtration} is standard for any
spectral sequence converging to a graded group. Part~(i): the
$E_2$ page is the Leray page obtained from the genus filtration
spectral sequence and its convergence theorem
(Theorems~\ref{thm:genus-ss} and~\ref{thm:genus-ss-convergence}),
and for Koszul
algebras the local system $\mathcal{H}^q$ is a variation of Hodge
structure, placing the cohomology classes in $R^*$ by
\cite[Appendix~A]{FP00}.

Part~(ii): The differential $d_2$ is the Kodaira--Spencer map
\begin{multline*}
\mathrm{KS}\colon H^p(\overline{\mathcal{M}}_g, R^q\pi_*\mathcal{A}) \\
\to H^{p+2}(\overline{\mathcal{M}}_g, R^{q-1}\pi_*\mathcal{A}),
\end{multline*}
which is a cup product with the Kodaira--Spencer class
$\mathrm{KS} \in H^1(\overline{\mathcal{M}}_g,
T_{\overline{\mathcal{M}}_g})$. This class is dual to
$\psi_1 \in H^1(\overline{\mathcal{M}}_{g,1})$ under the
cotangent-tangent pairing. When $d_2$ acts on a $\lambda$-class,
the result involves $\psi$-classes and boundary divisor classes
$[\Delta_{\mathrm{irr}}]$, $[\Delta_{i,S}]$.

Part~(iii): The differential $d_3$ is a triple Massey product in
the $A_\infty$ structure (Theorem~\ref{thm:bar-ainfty-complete}).
On $\overline{\mathcal{M}}_g$, the Massey product
$\langle \psi_i, \psi_j, \psi_k \rangle$ produces
$\kappa$-classes $\kappa_m = \pi_*(\psi^{m+1})$ by the
push-forward formula.
\end{proof}

\begin{remark}[Tautological classes beyond \texorpdfstring{$\lambda_g$}{lambda-g}]
\label{rem:tauto-beyond-lambda}
\index{tautological ring!beyond lambda classes}
The bar complex can produce tautological classes beyond $\lambda_g$
via two mechanisms: (1)~non-collapsing spectral sequences (no
current example among Master Table algebras); (2)~module insertions,
where $\operatorname{obs}^{\mathrm{Hdg}}_{g,n}
=(-1)^g\kappa\cdot\lambda_g\cdot
\prod_i\psi_i^{h_i}$
\textup{(}UNIFORM-WEIGHT, SCALAR-DIAGONAL\textup{)} gives products in
$R^*(\overline{\mathcal{M}}_{g,n})$.
The second is implicit in the Verlinde formula.
\end{remark}

\begin{theorem}[Grothendieck group of Koszul chiral algebras; \ClaimStatusConditional]
\label{thm:koszul-k0}
\index{Grothendieck group!Koszul chiral algebras}
Let $\mathrm{KCA}$ denote the category of Koszul chiral algebras with tensor product $\otimes$. The obstruction coefficient defines a group homomorphism
\[
\kappa: K_0(\mathrm{KCA}, \otimes) \to (\mathbb{Q}, +)
\]
from the Grothendieck group to the rationals. This homomorphism is:
\begin{enumerate}[label=\textup{(\roman*)}]
\item \emph{Surjective}: the Heisenberg family $\mathcal{H}_\kappa$ achieves any $\kappa \in \mathbb{Q}$;
\item \emph{Non-injective}: the kernel contains all algebras at critical level ($\kappa = 0$);
\item \emph{Compatible with Koszul duality}: for Kac--Moody algebras, $\kappa(\cA^!) = -\kappa(\cA)$, so the Koszul involution acts as negation on the image.
\end{enumerate}
\end{theorem}

\begin{proof}
Additivity is Corollary~\ref{cor:kappa-additivity}. The unit $\mathbf{1}$ (trivial algebra) has $\kappa(\mathbf{1}) = 0$. Surjectivity follows from $\kappa(\mathcal{H}_q) = q$ for $q \in \mathbb{Q}$. Non-injectivity: the Feigin--Frenkel center $\mathfrak{z}(\widehat{\mathfrak{g}}_{-h^\vee})$ is a non-trivial Koszul algebra with $\kappa = 0$. Part (iii) is Theorem~\ref{thm:genus-universality}(ii).
\end{proof}

\begin{definition}[Uncurved quotient]\label{def:uncurved-quotient}
\index{uncurved quotient category}
The \emph{uncurved quotient} of the category $\mathrm{KCA}$ is the localization
\[
\mathrm{KCA}^{\mathrm{unc}} := \mathrm{KCA}[\kappa^{-1}(0)]
\]
at the full subcategory of algebras with $\kappa = 0$. Objects of $\mathrm{KCA}^{\mathrm{unc}}$ are equivalence classes of Koszul chiral algebras modulo critical-level algebras: $\cA \sim \mathcal{B}$ if $\cA \otimes \mathcal{C} \simeq \mathcal{B} \otimes \mathcal{C}'$ for some $\mathcal{C}, \mathcal{C}'$ with $\kappa = 0$. By Theorem~\ref{thm:koszul-k0}, the obstruction coefficient $\kappa$ is a complete invariant on $K_0(\mathrm{KCA}^{\mathrm{unc}})$.
\end{definition}

\subsection{Connection to deformation-obstruction complementarity}
\label{subsec:obstruction-deformation-connection}

\begin{theorem}[Poincar\'e pairing for the Hodge shadow;
\ClaimStatusProvedElsewhere]
\label{thm:obs-def-pairing-explicit}
\textup{[Type signature: Open quadrant; scalar Hodge presentation;
level $5$; rational Poincar\'e duality on $\overline{\cM}_g$.]}
The orbifold fundamental class of~$\overline{\cM}_g$ gives a perfect
pairing
\[
 H^{2g}(\overline{\cM}_g,\mathbb C)
 \otimes H^{4g-6}(\overline{\cM}_g,\mathbb C)
 \longrightarrow\mathbb C,
 \qquad
 (\alpha,\beta)\longmapsto
 \int_{\overline{\cM}_g}\alpha\smile\beta.
\]
Under~$H_D^K(\cA,g)$, the first factor contains
$\operatorname{obs}^{\mathrm{Hdg}}_g(\cA)
=(-1)^g\kappa(\cA)\lambda_g$.
A Koszul-dual interpretation of the complementary factor requires a
named quasi-isomorphism from its deformation complex to the
Poincar\'e-dual coefficient complex.
\end{theorem}

\begin{proof}
The stack~$\overline{\cM}_g$ is proper of real dimension~$6g-6$.
Rational Poincar\'e duality gives the displayed pairing.  The Hodge
shadow formula is Theorem~\ref{thm:genus-universality}.
\end{proof}

\begin{example}[Heisenberg Hodge pairing]
\label{ex:heisenberg-pairing}
For~$\mathcal H_\kappa$ and~$\alpha\in
H^{4g-6}(\overline{\cM}_g)$,
\[
 \left\langle
 \operatorname{obs}^{\mathrm{Hdg}}_g(\mathcal H_\kappa),\alpha
 \right\rangle
 =(-1)^g\kappa
 \int_{\overline{\cM}_g}\lambda_g\smile\alpha.
\]
The Faber--Pandharipande choice on~$\overline{\cM}_{g,1}$ is
$\alpha=\psi_1^{2g-2}$ after pullback along the universal marking.
\end{example}

\begin{remark}[Hodge filtration]
\label{rem:hodge-filtration}
\index{Hodge filtration!bar complex}
The class~$\lambda_g$ has Hodge type~$(g,g)$.  Therefore
$\operatorname{obs}^{\mathrm{Hdg}}_g(\cA)$ lies in
$F^gH^{2g}\cap\overline F^{g}H^{2g}$.  Any comparison with a
Koszul-dual deformation class must preserve the stated Hodge type.
\end{remark}

\begin{proposition}[Vanishing of the scalar Hodge shadow;
\ClaimStatusConditional]
\label{prop:obstruction-lifting}
\textup{[Type signature: Open quadrant; scalar Hodge presentation;
level $5$; $H_D^K(\cA,g)$.]}
For every~$g\ge1$ covered by~$H_D^K$,
\[
 \operatorname{obs}^{\mathrm{Hdg}}_g(\cA)=0
 \quad\Longleftrightarrow\quad
 \kappa(\cA)=0.
\]
For $g\geq2$, the corresponding assertion for the deformation
obstruction
$\operatorname{Obs}^{\mathrm{def}}_g(\cA)$ follows after adding
$H_D^{\mathrm{tr}}(\cA,g)$ with a faithful scalar trace.
\end{proposition}

\begin{proof}
Faber--Pandharipande's nonzero pairing
\(
 \int_{\overline{\cM}_{g,1}}\psi_1^{2g-2}\lambda_g
 =\lambda_g^{\mathrm{FP}}
\)
shows that~$\lambda_g$ is nonzero.  The formula
$\operatorname{obs}^{\mathrm{Hdg}}_g
=(-1)^g\kappa\lambda_g$ gives the equivalence.
\end{proof}

\begin{remark}[Scalar component of the modular class]
\label{rem:dichotomy-theta}
The scalar Hodge component vanishes precisely at~$\kappa=0$.
The full Maurer--Cartan element contains its remaining deformation
directions and mixed stable-graph channels.
\end{remark}


\begin{remark}[Chern--Simons benchmark for \texorpdfstring{$\kappa$}{kappa};
heuristic scope]\label{rem:kappa-chern-simons}
\index{obstruction coefficient!Chern--Simons benchmark}
\index{Chern--Simons theory!heuristic benchmark}
For affine Kac--Moody $\widehat{\mathfrak{g}}_k$, the obstruction
coefficient
\(\kappa(\widehat{\mathfrak g}_k)=(k+h^\vee)d/(2h^\vee)\)
is proportional to the
shifted level $t = k + h^\vee$ that appears on the associated
Chern--Simons benchmark lane. The genus-$g$ free energy
\(F_g^{\mathrm{sc}}(\widehat{\mathfrak g}_k)
=\kappa(\widehat{\mathfrak g}_k)\lambda_g^{\mathrm{FP}}\)
\textup{(}UNIFORM-WEIGHT, SCALAR-DIAGONAL\textup{)} is therefore a tautological
integral series whose scalar coefficient is linear in that shifted
level. This supports a heuristic comparison between the scalar
bar-complex obstruction series and Chern--Simons/WRT asymptotics,
but the manuscript does not prove an identification with the
partition function of a $3$-manifold theory.
The universal MC class
$\Theta_{\widehat{\mathfrak{g}}_k}^{\min} =
\frac{t \cdot d}{2h^\vee}\,\eta \otimes \Lambda$
(Theorem~\ref{thm:explicit-theta}) is proportional to
$t = k + h^\vee$: on the same heuristic benchmark surface, the
bar-complex curvature may be viewed as an algebraic shadow of
Chern--Simons level dependence. At critical level the scalar
coefficient vanishes, but no theorem identifying it with a
topological Chern--Simons phase is proved here.
\end{remark}

\begin{theorem}[Wick-rotated \texorpdfstring{$\widehat{A}$}{A-hat}
scalar trace; \ClaimStatusConditional]
\label{prop:grr-bridge}
\index{A-hat genus@$\widehat{A}$-genus!scalar generating function}
\textup{[Type signature: Open quadrant; numerical modular-trace
presentation; level $5$; $H_D^{\mathrm{graph}}$.]}
Under the scalar graph package,
\begin{equation}
 F_g^{\mathrm{sc}}(\cA)
 =\kappa(\cA)
 \int_{\overline{\cM}_{g,1}}\psi_1^{2g-2}c_g(\mathbb E_g)
 =\kappa(\cA)\lambda_g^{\mathrm{FP}}.
 \label{eq:grr-bridge}
\end{equation}
Consequently
\begin{equation}
 \sum_{g\ge1}F_g^{\mathrm{sc}}(\cA)x^{2g}
 =\kappa(\cA)\left(\frac{x/2}{\sin(x/2)}-1\right)
 =\kappa(\cA)\bigl(\widehat A(ix)-1\bigr).
 \label{eq:grr-bridge-total}
\end{equation}
The full stable-graph trace has the form
$F_g=F_g^{\mathrm{sc}}+\delta F_g^{\mathrm{cross}}$.
\end{theorem}

\begin{proof}
Faber--Pandharipande~\cite{FP03} prove
\[
 \lambda_g^{\mathrm{FP}}
 =\frac{2^{2g-1}-1}{2^{2g-1}}
 \frac{|B_{2g}|}{(2g)!}.
\]
Its generating series is $(x/2)/\sin(x/2)-1$.  The equality with
$\widehat A(ix)-1$ follows from
$\widehat A(x)=(x/2)/\sinh(x/2)$.
\end{proof}


\begin{example}[The first three scalar genera]
\label{ex:first-three-scalar-genera}
The $\widehat{A}$-series begins
\[
\widehat{A}(x)
\;=\;
1-\frac{x^2}{24}+\frac{7x^4}{5760}-\frac{31x^6}{967680}+O(x^8),
\]
so
\[
\widehat{A}(ix)-1
\;=\;
\frac{x^2}{24}+\frac{7x^4}{5760}+\frac{31x^6}{967680}+O(x^8).
\]
Therefore on the proved scalar lane
\[
F_1(\cA)=\frac{\kappa(\cA)}{24},
\qquad
F_2(\cA)=\frac{7\,\kappa(\cA)}{5760},
\qquad
F_3(\cA)=\frac{31\,\kappa(\cA)}{967680}.
\]
For $\mathcal{H}_k$ this gives
$k/24$, $7k/5760$, $31k/967680$.
For $\widehat{\mathfrak{g}}_k$ it gives
\[
\frac{\dim(\mathfrak{g})(k+h^\vee)}{48h^\vee},
\qquad
\frac{7\,\dim(\mathfrak{g})(k+h^\vee)}{11520h^\vee},
\qquad
\frac{31\,\dim(\mathfrak{g})(k+h^\vee)}{1935360h^\vee},
\]
and for $\Vir_c$ it gives
$c/48$, $7c/11520$, $31c/1935360$.
For $\cW_3$, the scalar part is still
$\kappa(\cW_3)\cdot(1/24,\,7/5760,\,31/967680)$ with
$\kappa(\cW_3)=5c/6$, but the full genus-$2$ term is
\[
F_2(\cW_3)
\;=\;
\frac{7\,\kappa(\cW_3)}{5760}
\;+\;
\frac{c+204}{16c}
\;=\;
\frac{7\,c^{2}+432\,c+88128}{6912\,c},
\]
so the cross-channel correction is genuinely additional data. The
common-denominator closed form has discriminant $432^{2}-4\cdot 7\cdot
88128 = -2{,}280{,}960<0$, so $F_2(\cW_3)$ has no real zero in $c$
\textup{(}only the pole at $c=0$ where the algebra degenerates\textup{);}
the cross-channel correction $\delta F_2^{\mathrm{cross}}$ alone
\emph{does} have a real zero at $c=-204$, but this is cancelled in
$F_2$ by the scalar part $7\,\kappa(\cW_3)/5760$. At the KSDual fixed
charge $c=50$, $F_2(\cW_3)=31807/86400$; at the c-dual $c=100$
\textup{(}where $c'=0$\textup{),} $F_2(\cW_3)=12583/43200$.
\end{example}

\begin{remark}[Family index theorem]\label{rem:towards-family-index}
\index{index theorem!modular deformation}
The family index theorem (Theorem~\ref{thm:family-index}) assembles
$\sum F_g x^{2g}=\kappa(\cA)(\hat{A}(ix)-1)$ as the GRR
pushforward of the Todd class of the universal curve.
The family~\eqref{eq:bar-family} defines a scalar virtual index
class, and the geometric input is the chain-level curvature
identity together with the Hodge Euler class identification of
Proposition~\ref{prop:scalar-obstruction-hodge-euler}
(via the Arakelov metric and Chern--Weil): the scalar
geometric factor is $c_g(\mathbb{E})=\lambda_g$, not a class
reconstructed from a clutching-uniqueness argument.
Theorem~\ref{thm:family-index} then
packages the same scalar series in GRR language. Constructing
$\mathcal{D}_{\cA}^{(g)}$ as a
sheaf requires the coderived Ran-space formalism
(Theorem~\ref{thm:universal-MC}).
\end{remark}

% ================================================================
% SECTION 4.14: QUANTUM CHIRAL HOMOLOGY FROM STABLE GRAPH TOPOLOGY
% ================================================================

\section{Quantum chiral homology from stable graph topology}
\label{sec:quantum-chiral-homology-first-principles}
\index{quantum chiral homology|textbf}
\index{stable graph!topology|textbf}

The previous sections constructed the genus-$g$ bar complex, proved
the quantum Arnold relations, established genus universality, and
identified the bar complex with the Feynman transform
(Theorem~\ref{thm:prism-higher-genus}). The logical direction now
reverses: starting from the combinatorics of stable graphs alone,
the full structure of quantum chiral homology (its differential, its
curvature, its filtrations, its homotopy type) arises as a
consequence of the boundary geometry of
$\overline{\mathcal{M}}_{g,n}$. This ``graph-first'' perspective
makes visible several structural features that are hidden in the
analytic construction: the loop order filtration, the curvature as a
self-contraction trace, the Getzler--Kapranov involution, and the
Grothendieck--Teichm\"uller action.

\subsection{The stable graph complex}
\label{subsec:stable-graph-complex}
\index{stable graph!complex|textbf}

\begin{definition}[Stable graph]\label{def:stable-graph-hg}
A \emph{stable graph} of type $(g,n)$ is a connected graph
$\Gamma = (V, E, L, g_\bullet)$ consisting of:
\begin{enumerate}[label=\textup{(\roman*)}]
\item a finite vertex set $V(\Gamma)$;
\item a finite edge set $E(\Gamma)$ (each edge is an unordered pair of
 half-edges at, possibly equal, vertices);
\item a leg set $L(\Gamma) = \{1, \ldots, n\}$ of external
 half-edges, each attached to a vertex;
\item a genus function $g_\bullet\colon V(\Gamma) \to \mathbb{Z}_{\geq 0}$;
\end{enumerate}
subject to:
\begin{itemize}
\item \emph{Stability:} $2g_v - 2 + \mathrm{val}(v) > 0$ for every
 $v \in V(\Gamma)$, where $\mathrm{val}(v)$ counts all half-edges
 (from edges and legs) at~$v$.
\item \emph{Genus:} $g = \sum_{v \in V} g_v + b_1(\Gamma)$, where
 $b_1(\Gamma) = |E(\Gamma)| - |V(\Gamma)| + 1$ is the first Betti
 number (loop order) of the underlying graph.
\end{itemize}
Write $\mathcal{G}_{g,n}^{\mathrm{st}}$ for the set of isomorphism
classes of stable graphs of type $(g,n)$.
\end{definition}

\begin{definition}[Loop order]\label{def:loop-order}
\index{loop order|textbf}
\index{Betti number!first}
The \emph{loop order} of a stable graph~$\Gamma$ is its first Betti
number
\begin{equation}\label{eq:loop-order}
\ell(\Gamma) := b_1(\Gamma)
= |E(\Gamma)| - |V(\Gamma)| + 1.
\end{equation}
The \emph{vertex genus} is
$g_V(\Gamma) := \sum_{v \in V} g_v = g - \ell(\Gamma)$.
For a stable graph of type $(g,n)$:
\begin{equation}\label{eq:loop-order-range}
0 \;\leq\; \ell(\Gamma) \;\leq\; g.
\end{equation}
The extremes are:
$\ell = 0$ (trees with higher-genus vertices) and
$\ell = g$ (genus-$0$ vertices only, all genus from graph topology).
\end{definition}

\begin{remark}[Genus decomposition triangle]%
\label{rem:genus-decomposition-triangle}
\index{genus decomposition!triangle}
The identity
\begin{equation}\label{eq:genus-triangle}
\underbrace{g}_{\text{total genus}}
\;=\;
\underbrace{\textstyle\sum_v g_v}_{\text{vertex genus}}
\;+\;
\underbrace{b_1(\Gamma)}_{\text{loop order}}
\end{equation}
is the organizing principle of the perturbative expansion.
\emph{Tree level} ($\ell = 0$): all genus resides at vertices; the
bar complex contribution is a tree-level composition of lower-genus
data.
\emph{Maximal loop order} ($\ell = g$): every vertex carries
genus~$0$; the contribution is a pure loop integral over the genus-$0$
vertex algebra. The general term interpolates between these extremes.
\end{remark}

\begin{definition}[Edge contraction and the stable graph differential;
\ClaimStatusDefinitional]%
\label{def:edge-contraction}
\index{edge contraction|textbf}
\index{stable graph!differential}
Let $\Gamma$ be a stable graph and $e \in E(\Gamma)$ an edge.
The \emph{contraction} $\Gamma / e$ is the stable graph obtained by:
\begin{enumerate}[label=\textup{(\alph*)}]
\item If $e$ connects distinct vertices $v_1, v_2$:
 merge $v_1$ and $v_2$ into a single vertex $v_{12}$ with
 $g_{v_{12}} = g_{v_1} + g_{v_2}$,
 $\mathrm{val}(v_{12}) = \mathrm{val}(v_1) + \mathrm{val}(v_2) - 2$.
 Loop order is preserved: $b_1(\Gamma/e) = b_1(\Gamma)$.
\item If $e$ is a loop at vertex $v$:
 remove $e$ and set $g_{v}' = g_v + 1$,
 $\mathrm{val}(v)' = \mathrm{val}(v) - 2$.
 Loop order decreases: $b_1(\Gamma/e) = b_1(\Gamma) - 1$.
\end{enumerate}
In both cases, $g(\Gamma/e) = g(\Gamma)$ (total genus is preserved).
The \emph{stable graph differential} is
\begin{equation}\label{eq:stable-graph-diff}
\partial_{\mathcal{G}}\Gamma
\;=\;
\sum_{e \in E(\Gamma)}
(-1)^{\sigma(e,\Gamma)}\,\Gamma / e,
\end{equation}
where $\sigma(e,\Gamma)$ is the position of~$e$ in the chosen
ordering of $E(\Gamma)$.
\end{definition}

\begin{lemma}[$\partial_{\mathcal{G}}^2 = 0$;
\ClaimStatusProvedHere]\label{lem:stable-graph-d-squared}
The stable graph differential satisfies $\partial_{\mathcal{G}}^2 = 0$.
\end{lemma}

\begin{proof}
Each term in $\partial_{\mathcal{G}}^2$ contracts two edges
$e, e'$ in some order. If $e$ and $e'$ share no vertex, the
contractions commute. If they share a vertex, the two orders
produce the same graph $\Gamma / \{e, e'\}$ with opposite signs
(from the antisymmetry of the determinant ordering in
$\det(E(\Gamma))$). In both cases, the terms cancel in pairs.
\end{proof}

\begin{remark}[Boundary strata of $\overline{\mathcal{M}}_{g,n}$]%
\label{rem:boundary-strata-stable-graph}
The codimension-$1$ boundary strata of $\overline{\mathcal{M}}_{g,n}$
are indexed by stable graphs with $|E| = 1$
(a single node = a single edge in the dual graph):
\[
\partial\overline{\mathcal{M}}_{g,n}
\;=\;
\sum_{\Gamma,\, |E(\Gamma)| = 1}
\overline{\mathcal{M}}_\Gamma,
\]
where $\overline{\mathcal{M}}_\Gamma =
\prod_{v \in V(\Gamma)} \overline{\mathcal{M}}_{g_v, \mathrm{val}(v)}$.
The relation $\partial^2 = 0$ on
$\overline{\mathcal{M}}_{g,n}$ is the geometric incarnation of
$\partial_{\mathcal{G}}^2 = 0$. This is the single axiom from
which all of quantum chiral homology flows.
\end{remark}

\subsection{The loop order filtration}
\label{subsec:loop-order-filtration}
\index{loop order!filtration|textbf}

\begin{definition}[Loop order filtration on the Feynman transform]%
\label{def:loop-order-filtration}
\index{Feynman transform!loop order filtration}
Let $\barB^{(g)}(\cA)
\simeq \mathrm{FT}^{(g)}_{\mathcal{M}od}(\cA)$ be the genus-$g$
bar complex, identified with the Feynman transform
(Theorem~\ref{thm:prism-higher-genus}).
Define the \emph{loop order filtration}
\begin{equation}\label{eq:loop-filtration}
F^\ell\!\barB^{(g)}(\cA)
\;:=\;
\bigoplus_{\substack{\Gamma \in \mathcal{G}_{g,n}^{\mathrm{st}}
\\ b_1(\Gamma) \geq \ell}}
\cA(\Gamma) \otimes \det(E(\Gamma)),
\end{equation}
where $\cA(\Gamma) = \bigotimes_{v \in V(\Gamma)}
\cA(g_v, \mathrm{val}(v))$.

This is a descending filtration:
\[
\barB^{(g)} = F^0 \supset F^1 \supset \cdots \supset F^g \supset
F^{g+1} = 0.
\]
\end{definition}

\begin{proposition}[Filtration compatibility;
\ClaimStatusProvedHere]\label{prop:loop-filtration-compatible}
\index{loop order!filtration compatibility}
The Feynman transform differential respects the loop order
filtration:
\begin{equation}\label{eq:loop-filtration-d}
d_{\mathrm{FT}} \colon
F^\ell\!\barB^{(g)} \to F^{\ell-1}\!\barB^{(g)}.
\end{equation}
Precisely, $d_{\mathrm{FT}}$ decomposes as
\begin{equation}\label{eq:d-tree-loop}
d_{\mathrm{FT}}
\;=\; d_{\mathrm{tree}} + d_{\mathrm{loop}},
\end{equation}
where:
\begin{enumerate}[label=\textup{(\alph*)}]
\item $d_{\mathrm{tree}}$ contracts non-loop edges
 \textup{(}edges connecting distinct vertices\textup{)};
 this preserves loop order:
 $d_{\mathrm{tree}}\colon F^\ell \to F^\ell$.
\item $d_{\mathrm{loop}}$ contracts loop edges
 \textup{(}self-edges at a vertex\textup{)}, using the invariant
 pairing $\langle -, - \rangle_\cA$;
 this decreases loop order by~$1$:
 $d_{\mathrm{loop}}\colon F^\ell \to F^{\ell-1}$.
\end{enumerate}
\end{proposition}

\begin{proof}
By Definition~\ref{def:edge-contraction}: contracting a non-loop
edge preserves $b_1$ (both $|E|$ and $|V|$ decrease by~$1$);
contracting a loop edge decreases $b_1$ by~$1$ ($|E|$ decreases
but $|V|$ is unchanged). Since the Feynman transform differential
is a sum over all edge contractions
(Definition~\ref{def:feynman-transform}), the decomposition
$d_{\mathrm{FT}} = d_{\mathrm{tree}} + d_{\mathrm{loop}}$ follows.
\end{proof}

\begin{theorem}[Loop order spectral sequence;
\ClaimStatusProvedHere]\label{thm:loop-order-spectral-sequence}
\index{spectral sequence!loop order|textbf}
The loop order filtration on $\barB^{(g)}(\cA)$ induces a
convergent spectral sequence
\begin{equation}\label{eq:loop-ss}
E_0^{\ell,q}
\;=\;
\bigl(F^\ell / F^{\ell+1}\bigr)^q
\;\Longrightarrow\;
H^{\ell+q}\!\bigl(\barB^{(g)}(\cA)\bigr),
\end{equation}
with:
\begin{enumerate}[label=\textup{(\roman*)}]
\item $E_0^{\ell,q}$: the space of Feynman transform elements
 supported on stable graphs with $b_1(\Gamma) = \ell$ exactly;

\item $d_0 = d_{\mathrm{tree}}$: the differential on~$E_0$ contracts
 only non-loop edges \textup{(}separating-type gluings\textup{)};

\item $E_1^{\ell,q}
 = H^q\!\bigl(F^\ell / F^{\ell+1},\, d_{\mathrm{tree}}\bigr)$:
 the cohomology of tree-edge contractions at fixed loop order;

\item $d_1 = [d_{\mathrm{loop}}]$:
 the induced differential on~$E_1$, which decreases loop
 order by~$1$ via loop contraction.
\end{enumerate}
The spectral sequence converges because the filtration is
bounded \textup{(}$0 \leq \ell \leq g$\textup{)}.
\end{theorem}

\begin{proof}
This is the spectral sequence of a bounded descending filtration
on a cochain complex. Convergence is immediate from
$F^{g+1} = 0$. The identification of $d_0$ with
$d_{\mathrm{tree}}$ follows from
Proposition~\ref{prop:loop-filtration-compatible}: $d_{\mathrm{tree}}$
is the component of the differential that preserves filtration degree,
while $d_{\mathrm{loop}}$ drops it by~$1$.
\end{proof}

\begin{proposition}[Extremal pages;
\ClaimStatusProvedHere]\label{prop:extremal-pages}
\index{loop order!extremal pages}
The extremal terms of the loop order spectral sequence admit
explicit identifications:
\begin{enumerate}[label=\textup{(\alph*)}]
\item \emph{Tree level} \textup{(}$\ell = 0$\textup{):}
 The summand $E_0^{0,*}$ consists of stable graphs with
 $b_1(\Gamma) = 0$ \textup{(}trees\textup{)}, with vertex genus
 summing to~$g$. Its $d_0$-cohomology is:
 \begin{equation}\label{eq:tree-level-page}
 E_1^{0,*}
 \;\cong\;
 \bigoplus_{\substack{
 \text{tree topologies}\\
 \sum_v g_v = g}}
 \bigotimes_{v \in V}
 H^*\!\bigl(\barB^{(g_v)}(\cA)\bigr)[\mathrm{val}(v)],
 \end{equation}
 i.e., tree-level compositions of lower-genus bar cohomology.

\item \emph{Maximal loop order} \textup{(}$\ell = g$\textup{):}
 The summand $E_0^{g,*}$ consists of stable graphs with
 $b_1(\Gamma) = g$ \textup{(}all vertices genus~$0$\textup{)}.
 Its $d_0$-cohomology is:
 \begin{equation}\label{eq:max-loop-page}
 E_1^{g,*}
 \;\cong\;
 \bigoplus_{\substack{
 \Gamma \in \mathcal{G}_{g,n}^{\mathrm{st}}\\
 g_v = 0 \;\forall v}}
 \bigotimes_{v \in V(\Gamma)}
 H^*\!\bigl(\barB^{(0)}(\cA)\bigr)[\mathrm{val}(v)]
 \;\otimes\;
 \det(E(\Gamma)),
 \end{equation}
 the ``pure loop'' contribution computed entirely from genus-$0$
 data. Write
 $H^\bullet_{\bar B,(0)}(\cA):=H^*(\barB^{(0)}(\cA))$ for the
 genus-zero bar cohomology.  A quadratic presentation supplies
 $\cA^{\mathrm i}_{(0)}=C_X(s^{-1}V,s^{-2}R)_{(0)}$ and
 $q_{\cA,(0)}\colon\cA^{\mathrm i}_{(0)}\to\barB^{(0)}(\cA)$.
 Under $H_{\mathrm{CL}}$, $H^\bullet(q_{\cA,(0)})$ identifies
 $H^\bullet(\cA^{\mathrm i}_{(0)})$ with
 $H^\bullet_{\bar B,(0)}(\cA)$.  When the former is concentrated in
 degree~$0$, the maximal-loop-order page reduces to the graph complex
 with these coefficients.  Finite Verdier duality and the chosen-pair
 equivalence further identify them with the corresponding graded
 dual of~$\cA^!$.
\end{enumerate}
\end{proposition}

\begin{proof}
(a)~For $\ell = 0$, all edges are non-loop (since the graph is a
tree). The $d_0 = d_{\mathrm{tree}}$ differential contracts these
edges, each of which merges two vertices. By the K\"unneth theorem
for the tensor product over vertices, the cohomology is the
tree-level composition of the vertex cohomologies.

(b)~For $\ell = g$, the constraint $\sum_v g_v = g - \ell = 0$
forces $g_v = 0$ for all~$v$. The vertex factors are
$\cA(0, \mathrm{val}(v))$, i.e., the genus-$0$ operations.
The $d_0$-cohomology replaces these with
$H^\bullet_{\bar B,(0)}(\cA)$ at each vertex.  Under
$H_{\mathrm{CL}}$, the map $H^\bullet(q_{\cA,(0)})$ computes this
cohomology from the chosen presentation coalgebra.  Degree-zero
concentration gives the graph-complex reduction, while finite Verdier
duality and the chosen-pair equivalence give the stated comparison
with the graded dual of~$\cA^!$.
\end{proof}

\begin{remark}[Physical interpretation: perturbative expansion]%
\label{rem:loop-order-perturbative}
\index{perturbative expansion!loop order}
The loop order spectral sequence is the mathematical incarnation
of the physicist's loop expansion. The $E_1$ page at loop
order~$\ell$ computes ``$\ell$-loop corrections'' to the
genus-$g$ theory:
\begin{center}
\small
\renewcommand{\arraystretch}{1.2}
\begin{tabular}{lll}
\toprule
\emph{Loop order} & \emph{Graph type}
 & \emph{Physical meaning} \\
\midrule
$\ell = 0$ & Trees & Classical compositions \\
$\ell = 1$ & One-loop & First quantum correction \\
$\ell = g$ & $g$-loop, genus-$0$ vertices
 & Pure Feynman diagram \\
\bottomrule
\end{tabular}
\end{center}
At $\ell = 1$: the single loop in the graph contributes the
supertrace of the propagator, which is the genus-$1$ curvature
$\kappa(\cA)$ (Theorem~\ref{thm:curvature-self-contraction} below).
\end{remark}

\subsection{The curvature as a self-contraction trace}
\label{subsec:curvature-self-contraction}
\index{curvature!self-contraction trace|textbf}

\begin{theorem}[Curvature from loop contraction;
\ClaimStatusConditional]\label{thm:curvature-self-contraction}
\index{loop contraction!curvature from}
\textup{[Type signature: Open quadrant; genus-zero primitive kernel
and genus-one loop trace presentation; levels $2\to4$; normalized
cyclic trace package $H_D^{\mathrm{loop}}$ and genus-one trace package
$H_D^1$.]}
Let $\cA$ be a cyclic Koszul chiral algebra and let
$K_{0,2}^{\cA}$ denote its binary primitive kernel.  Define the
normalized \emph{loop-contraction trace}
\begin{equation}\label{eq:self-contraction-trace}
\operatorname{Tr}^{\mathrm{loop}}_\cA
\;:=\;
\operatorname{tr}_{\mathrm{cyc}}\!\left(
\Delta_{\mathrm{ns}}K_{0,2}^{\cA}\right),
\end{equation}
where $\Delta_{\mathrm{ns}}$ sews the two flags of the binary corolla
and $\operatorname{tr}_{\mathrm{cyc}}$ includes the manuscript's
automorphism and OPE normalizations.  Assume~$H_D^{\mathrm{loop}}$,
namely that this trace agrees with the normalized degree-two averaged
collision trace.
Then:
\begin{enumerate}[label=\textup{(\roman*)}]
\item The curvature coefficient of the genus-$1$ bar complex is
 \begin{equation}\label{eq:kappa-self-contraction}
 \kappa(\cA) = \operatorname{Tr}^{\mathrm{loop}}_\cA.
 \end{equation}

\item In the QME decomposition
 \textup{(}Theorem~\textup{\ref{thm:quantum-master-equation})},
 the genus-$1$ equation at degree~$0$:
 \begin{equation}\label{eq:qme-genus1}
 \underbrace{\Delta(S_0)}_{\text{one-loop anomaly}}
 \;+\;
 \underbrace{\{S_0, S_1\}}_{\text{tree $\times$ genus-$1$}}
 \;=\; 0,
 \end{equation}
 has scalar diagonal
 $\operatorname{tr}_{\mathrm{diag}}\Delta(S_0)
 =\operatorname{Tr}^{\mathrm{loop}}_\cA
 \omega_1^{\mathrm{Ar}}$: the one-loop
 operation is the loop-contraction trace weighted by the Arakelov
 form.

\item Under~$H_D^1$, the genus-one deformation trace is
 \[
 \operatorname{tr}_1\operatorname{Obs}^{\mathrm{def}}_{1,1}(\cA)
 =\operatorname{Tr}^{\mathrm{loop}}_\cA\,\lambda_1
 =\kappa(\cA)\lambda_1.
 \]
\end{enumerate}
\end{theorem}

\begin{proof}
The unique non-separating genus-one self-sewing of the binary corolla
is $\Delta_{\mathrm{ns}}K_{0,2}^{\cA}$.  Applying the normalized cyclic
trace gives~\eqref{eq:self-contraction-trace}, and
$H_D^{\mathrm{loop}}$ identifies it with the collision coefficient
$\kappa(\cA)$.  The Arakelov fixed-fiber package weights this loop by
$\omega_1^{\mathrm{Ar}}$, yielding clause~\textup{(ii)}.  The
genus-one trace comparison~$H_D^1$ maps the same degree-$2$
deformation coefficient to the Hodge line, giving clause~\textup{(iii)}.
\end{proof}

\begin{corollary}[Anomaly = trace for standard families;
\ClaimStatusConditional]\label{cor:anomaly-trace-standard}
\index{curvature!self-contraction for standard families}
\begin{enumerate}[label=\textup{(\alph*)}]
\item \emph{Heisenberg} $\cH_\kappa$:
 $\operatorname{Tr}^{\mathrm{loop}}_{\cH_\kappa} = \kappa$.
 The normalized binary kernel is the coefficient of
 $a(z)a(w)\sim\kappa(z-w)^{-2}$.
\item \emph{Affine Kac--Moody} $\widehat{\fg}_k$:
 $\operatorname{Tr}^{\mathrm{loop}}_{\widehat{\fg}_k}
 = (k + h^\vee) \cdot \dim(\fg) / (2h^\vee)$.
 The normalized affine OPE contraction combines the invariant-form
 term with the dual-Coxeter normal-ordering contribution.
\item \emph{Virasoro} $\mathrm{Vir}_c$:
 $\operatorname{Tr}^{\mathrm{loop}}_{\mathrm{Vir}_c} = c/2$.
 The normalized loop trace extracts the coefficient of the
 fourth-order pole in the $TT$ OPE.
\end{enumerate}
These coincide with the values $\kappa(\cA)$ computed by direct
propagator analysis in Part~\ref{part:bar-cobar-engine}.
\end{corollary}

\begin{proof}
Compute the normalized cyclic trace of
$\Delta_{\mathrm{ns}}K_{0,2}^{\cA}$ from the displayed OPE in each
family and compare with the independently computed $\kappa(\cA)$ values
(Theorems~\ref{thm:heisenberg-obs},
\ref{thm:kac-moody-obs}, \ref{thm:vir-genus1-curvature}).
\end{proof}

\begin{remark}[Three typed realizations of $\kappa$]%
\label{rem:three-faces-kappa}
\index{curvature coefficient!three characterizations}
The coefficient~$\kappa$ appears in three ambient spaces through the
named comparison maps:
\begin{center}
\small
\renewcommand{\arraystretch}{1.2}
\begin{tabular}{lll}
\toprule
\emph{Characterization} & \emph{Source} & \emph{Nature} \\
\midrule
Loop contraction &
$\operatorname{tr}_{\mathrm{cyc}}(\Delta_{\mathrm{ns}}K_{0,2})=\kappa$
& algebraic scalar \\
Arnold defect & $\operatorname{tr}_{\mathrm{diag}}(m_0)
=\kappa \omega_g^{\mathrm{Ar}}$ & vertical fiber form \\
Genus-one trace &
$\operatorname{tr}_1\operatorname{Obs}^{\mathrm{def}}_{1,1}
=\kappa\lambda_1$ under $H_D^1$ & horizontal degree-$2$ class \\
\bottomrule
\end{tabular}
\end{center}
The loop and vertical realizations are
Theorems~\ref{thm:curvature-self-contraction}
and~\ref{thm:quantum-arnold-relations}.  The map to the Hodge line is
$H_D^1$ in Theorem~\ref{thm:genus-universality}.
\end{remark}

\subsection{Finite-depth convergence on the Koszul locus}
\label{subsec:loop-order-collapse}
\index{spectral sequence!loop order collapse}

\begin{theorem}[Loop-order convergence and finite-depth collapse bound;
\ClaimStatusProvedHere]\label{thm:loop-order-collapse}
Let $\cA$ be a modular Koszul chiral algebra
\textup{(}Definition~\textup{\ref{def:modular-koszul-chiral})}
whose finite-window bar tower is strict Mittag--Leffler.
The loop order spectral sequence
\textup{(}Theorem~\textup{\ref{thm:loop-order-spectral-sequence})}
is strongly convergent to the associated graded of
$H^*(\widehat{\barB}^{(g)}_{\mathrm{flat}}(\cA))$ for every fixed
genus~$g$, where the subscript denotes the corrected differential
$\Dg{g}$.
If the shadow tower of~$\cA$ has finite depth $r_{\max}$, so that no
loop-contraction component lowers the loop filtration by more than
$r_{\max}$, then
\begin{equation}\label{eq:loop-collapse-bound}
E_{r_{\max}+1}^{\ell,q}=E_\infty^{\ell,q}.
\end{equation}
In the exceptional case where the induced loop-contraction
differential vanishes on $E_1$, this bound sharpens to
$E_1=E_\infty$. For class-$M$ inputs, where
$r_{\max}=\infty$, the theorem asserts completed convergence, not
finite-page collapse.
\end{theorem}

\begin{proof}
For fixed~$g$, the loop filtration is bounded:
$0\leq b_1(\Gamma)\leq g$. Hence the ordinary loop-order spectral
sequence of Theorem~\ref{thm:loop-order-spectral-sequence}
converges for each finite window. Strict Mittag--Leffler control
commutes cohomology with the inverse limit over windows, so the
completed spectral sequence converges strongly to the associated
graded of $H^*(\widehat{\barB}^{(g)}_{\mathrm{flat}}(\cA))$.

The component of the differential that lowers loop order by~$r$ is
the $r$-th shadow operation in the obstruction tower. If the tower has
finite depth $r_{\max}$, all such components vanish for
$r>r_{\max}$. A filtered-complex spectral sequence has no
differential $d_r$ without a filtration-drop-$r$ component in the
underlying differential, giving the collapse bound
\eqref{eq:loop-collapse-bound}. If the loop component already acts
trivially on $E_1$, every later differential is forced to vanish and
$E_1=E_\infty$. For class-$M$ algebras no finite $r_{\max}$ is
available; the Mittag--Leffler completion gives convergence but no
finite-page collapse statement.
\end{proof}

\begin{corollary}[Associated graded by loop order;
\ClaimStatusProvedHere]\label{cor:loop-decomposition-bar}
\index{bar cohomology!loop order decomposition}
For a modular Koszul chiral algebra~$\cA$ satisfying the
finite-window Mittag--Leffler hypothesis, the loop filtration on
genus-$g$ bar cohomology has associated graded
\begin{equation}\label{eq:loop-decomposition}
\operatorname{gr}^{\ell}_{F}
H^*\!\bigl(\widehat{\barB}^{(g)}(\cA)\bigr)
\;\cong\;
E_\infty^{\ell,*}.
\end{equation}
If, in addition, the spectral sequence collapses at $E_1$, then
$E_\infty^{\ell,*}=H^*(\mathsf{SG}_{g}^{(\ell)}(\cA),
d_{\mathrm{tree}})$ and the displayed associated graded is computed
by the loop-$\ell$ stable-graph complex. Without an $E_1$ collapse
hypothesis, the statement is an associated-graded identification, not
a direct-sum decomposition of cohomology.
\end{corollary}

\begin{proof}
This is the standard output of a strongly convergent filtered
spectral sequence: $E_\infty^{\ell,*}$ is the associated graded of
the cohomology for the loop filtration. The $E_1$ page is the
cohomology of the fixed-loop stable-graph complex with respect to
$d_{\mathrm{tree}}$ by
Theorem~\ref{thm:loop-order-spectral-sequence}. Identifying
$E_\infty$ with that $E_1$ page requires the extra collapse
hypothesis stated in the corollary.
\end{proof}

\subsection{The Getzler--Kapranov involution}
\label{subsec:feynman-involution}
\index{Getzler--Kapranov involution|textbf}
\index{Feynman transform!involution}

\begin{theorem}[Feynman involution
\textup{\cite[Theorem~5.2]{GeK98}};
\ClaimStatusProvedElsewhere]\label{thm:feynman-involution}
\index{Feynman transform!$F^2 \simeq \mathrm{id}$}
Let $\mathcal{M}od$ be a modular operad satisfying Poincar\'e
duality in each degree \textup{(}as holds for
$\mathcal{M}od = \{\overline{\mathcal{M}}_{g,n}\}$\textup{)}.
The Feynman transform is a homotopy involution:
\begin{equation}\label{eq:feynman-squared}
\mathrm{FT}_{\mathcal{M}od} \circ \mathrm{FT}_{\mathcal{M}od}
\;\simeq\; \mathrm{id}
\end{equation}
on the homotopy category of $\mathcal{M}od$-algebras.
\end{theorem}

\begin{corollary}[Duality of quantum chiral homology;
\ClaimStatusConditional]\label{cor:feynman-duality-qch}
\index{quantum chiral homology!Feynman duality}
For a Koszul chiral algebra~$\cA$, the Feynman transform
applied twice recovers~$\cA$:
\begin{equation}\label{eq:double-feynman}
\mathrm{FT}(\mathrm{FT}(\cA))
\;\simeq\; \cA.
\end{equation}
Concretely: the Feynman transform of $\cA$ is $\barB(\cA)$
\textup{(}the bar complex\textup{)}; the Feynman transform of
$\barB(\cA)$ is $\Omega(\barB(\cA)) \simeq \cA$
\textup{(}the cobar, which recovers~$\cA$ by bar-cobar
inversion\textup{)}.

Thus bar-cobar inversion
\textup{(}Theorem~\textup{\ref{thm:higher-genus-inversion})}
in each genus is a consequence of the Getzler--Kapranov involution
for the modular operad of Deligne--Mumford spaces.
\end{corollary}

\begin{proof}
The Prism Principle
(Theorem~\ref{thm:prism-higher-genus}) identifies
$\widehat{\barB}^{\mathrm{full}}(\cA) \simeq
\widehat{\mathrm{FT}}(\cA)$ in the genus/window completed category.
Applying the Feynman transform again:
$\mathrm{FT}(\mathrm{FT}(\cA)) \simeq \cA$ by
Theorem~\ref{thm:feynman-involution}. On the other hand,
the Feynman transform of a conilpotent coalgebra (which
$\widehat{\barB}(\cA)$ is after completion) produces the continuous
cobar construction $\Omega_{\mathrm{cont}}(\widehat{\barB}(\cA))$,
and Theorem~\ref{thm:higher-genus-inversion}
gives $\Omega_{\mathrm{cont}}(\widehat{\barB}(\cA)) \simeq \cA$ on
the Koszul locus.
Thus the bar-cobar resolution is a special case of~$F^2 \simeq
\mathrm{id}$.
\end{proof}

\begin{remark}[Feynman involution vs.\ Verdier duality]%
\label{rem:feynman-vs-verdier}
\index{Verdier duality!and Feynman involution}
The Feynman involution $F^2 \simeq \mathrm{id}$ is an algebraic
statement about modular operads. The Verdier duality
$K_X(\cA):=\mathbb{D}(\barB(\cA))$ is a geometric algebra of sheaves
on configuration spaces; a chosen Koszul pair supplies
$\nu_\cA\colon K_X(\cA)\to\cA^!$.  A comparison of the two
involutions consists of a natural equivalence
\[
\gamma_\cA\colon F^2(\cA)\xrightarrow{\sim}
\mathbb D^2(\cA)
\]
compatible with clutching and with~$\nu_\cA$.  The fundamental classes
$[\overline{\mathcal{M}}_{g,n}]$ provide the Poincar\'e-duality
orientation entering this comparison; the natural transformation
and its clutching compatibility are the remaining geometric data.
\end{remark}

\subsection{The virtual Euler characteristic}
\label{subsec:virtual-euler-char}
\index{virtual Euler characteristic|textbf}
\index{Euler characteristic!from stable graphs}

\begin{theorem}[Graph-sum formula for the virtual Euler
characteristic;
\ClaimStatusProvedHere]\label{thm:virtual-euler-char}
For a Koszul chiral algebra~$\cA$ with finite-type bar complex,
the virtual Euler characteristic at genus~$g$ is:
\begin{equation}\label{eq:virtual-euler-graph}
\chi^{\mathrm{vir}}\!\bigl(\barB^{(g)}(\cA)\bigr)
\;=\;
\sum_{\Gamma \in \mathcal{G}_g^{\mathrm{st}}}
\frac{(-1)^{|E(\Gamma)|}}{|\!\operatorname{Aut}(\Gamma)|}
\prod_{v \in V(\Gamma)}
\chi\bigl(\cA(g_v, \mathrm{val}(v))\bigr),
\end{equation}
where $\chi(\cA(g_v, k))$ is the Euler characteristic of
$\cA$ at genus~$g_v$ and degree~$k$.
\end{theorem}

\begin{proof}
The Feynman transform identification
$\barB^{(g)}(\cA) \simeq \mathrm{FT}^{(g)}(\cA)$
(Theorem~\ref{thm:prism-higher-genus}) writes the bar complex
as a sum over stable graphs. The virtual Euler characteristic is
additive over the graph sum and multiplicative over the vertex
tensor product. The factor $(-1)^{|E|}$ comes from the
determinant line $\det(E(\Gamma))$; the factor
$1/|\!\operatorname{Aut}(\Gamma)|$ accounts for graph
automorphisms.
\end{proof}

\begin{corollary}[Heisenberg Euler characteristic;
\ClaimStatusProvedHere]\label{cor:heisenberg-euler-char}
For the Heisenberg algebra $\cH_\kappa$:
\begin{equation}\label{eq:heisenberg-euler-graph}
\chi^{\mathrm{vir}}\!\bigl(\barB^{(g)}(\cH_\kappa)\bigr)
\;=\;
\kappa^{g-1} \cdot \chi(\overline{\mathcal{M}}_g),
\end{equation}
where $\chi(\overline{\mathcal{M}}_g)$ is the orbifold Euler
characteristic of the moduli space.
\end{corollary}

\begin{proof}
For $\cH_\kappa$, the vertex factor at genus~$0$ and degree~$k$ is
$\chi(\cH_\kappa(0,k)) = \delta_{k,2} \cdot \kappa$ (the Koszul
dual is one-dimensional: $\dim \cH^!_\kappa = 1$).
The self-contraction trace is $\kappa$, and the graph sum over
stable graphs with genus-$0$ vertices contributing $\kappa$ per
edge yields $\kappa^{|E|} / |\!\operatorname{Aut}|$, summed with
$(-1)^{|E|}$. The total is $\kappa^{g-1}$ times the
Kontsevich--Penner sum $\sum_\Gamma
(-1)^{|E|}/|\!\operatorname{Aut}|$, which computes
$\chi(\overline{\mathcal{M}}_g)$ by
Harer--Zagier~\cite{HZ86,Penner88}.
\end{proof}

\subsection{The Grothendieck--Teichm\"uller action}
\label{subsec:grt-action}
\index{Grothendieck--Teichmuller group@Grothendieck--Teichm\"uller group|textbf}
\index{graph complex!action on quantum chiral homology}

\begin{remark}[Kontsevich's graph complex and $\mathfrak{grt}_1$]%
\label{rem:gc-grt}
\index{graph complex!$\mathsf{GC}_2$}
\index{grt@$\mathfrak{grt}_1$}
Kontsevich's graph complex $\mathsf{GC}_n$ is the chain complex
of connected graphs with vertex valence $\geq 3$ and no legs,
differential given by edge contraction, with a grading shift.
Willwacher's theorem~\cite{Willwacher15} identifies
\begin{equation}\label{eq:willwacher}
H^0(\mathsf{GC}_2) \;\cong\; \mathfrak{grt}_1,
\end{equation}
the graded Grothendieck--Teichm\"uller Lie algebra. Since
$\mathsf{GC}_2$ acts on the homotopy automorphisms
$\mathrm{Aut}(\mathsf{E}_2)$ of the $\mathsf{E}_2$-operad, and
chiral algebras on curves are $\mathsf{E}_2$-algebras (via
factorization on the surface), the Grothendieck--Teichm\"uller
group acts on quantum chiral homology.
\end{remark}

\begin{proposition}[Weight system map;
\ClaimStatusProvedHere]\label{prop:weight-system-map}
\index{weight system!from chiral algebra}
\index{graph complex!weight system}
For any Koszul chiral algebra~$\cA$ with invariant bilinear
form, there is a weight system map
\begin{equation}\label{eq:weight-system}
w_\cA \colon \mathsf{GC}_2
\;\longrightarrow\;
\mathrm{End}\!\bigl(
H^*(\widehat{\barB}^{\mathrm{full}}(\cA))
\bigr),
\end{equation}
defined by decorating vertices of Kontsevich graphs with
$\cA$-operations and contracting edges with the invariant
pairing.

On the $H^0$ level, this induces an action of $\mathfrak{grt}_1$
on the all-genus bar cohomology of~$\cA$.
\end{proposition}

\begin{proof}
The Feynman transform identification
$\widehat{\barB}^{\mathrm{full}}(\cA) \simeq
\widehat{\mathrm{FT}}(\cA)$ expresses
the bar complex as a sum over stable graphs with vertex factors
in~$\cA$. A Kontsevich graph $\gamma \in \mathsf{GC}_2$ acts on
a stable graph $\Gamma$ by ``inserting'' $\gamma$ at a vertex of
$\Gamma$: replace the vertex $v$ by the subgraph $\gamma$, with
the $\mathrm{val}(v)$ half-edges of $v$ distributed among the
vertices of $\gamma$. The edge contractions within $\gamma$ use
the invariant form $\langle -, - \rangle_\cA$ and the operations
of~$\cA$.

This insertion is a chain map: the compatibility of the graph
complex differential (edge contraction on~$\gamma$) with the
Feynman transform differential (edge contraction on~$\Gamma$) is
guaranteed by the modular operad axioms and the $\mathsf{E}_2$
equivariance of the chiral algebra structure. The $H^0$ action
follows by applying $H^0(\mathsf{GC}_2) \cong \mathfrak{grt}_1$.
\end{proof}

\begin{remark}[Quantum chiral homology and number theory]%
\label{rem:qch-number-theory}
\index{number theory!and quantum chiral homology}
The $\mathfrak{grt}_1$ action connects quantum chiral homology to
number theory. The Lie algebra
$\mathfrak{grt}_1$ contains elements associated to odd zeta values
$\zeta(3), \zeta(5), \zeta(7), \ldots$ (Drinfeld, Brown).
Through the weight system map, these zeta elements act on the
all-genus bar cohomology of any Koszul chiral algebra, producing
``motivic'' corrections to the genus expansion.

For the Heisenberg algebra, the $\mathfrak{grt}_1$ action
is trivial on bar cohomology (the bar complex is too simple for
the graph complex to see), but for Kac--Moody and Virasoro
algebras the action may be non-trivial, connecting to the
occurrence of multiple zeta values in string amplitudes
\cite{Stieberger14,Schlotterer12}.
\end{remark}

\subsection{Summary: the graph-theoretic architecture of quantum
chiral homology}
\label{subsec:qch-architecture-summary}

\begin{remark}[The five pillars]%
\label{rem:five-pillars}
\index{quantum chiral homology!five pillars}
The constructions above rest on five structural pillars:
\begin{enumerate}[label=\textup{(\Roman*)}]
\item \emph{$\partial^2 = 0 \Rightarrow d_{\mathrm{bar}}^2 = 0$}
 (Lemma~\ref{lem:stable-graph-d-squared},
 Theorem~\ref{thm:prism-higher-genus}):
 nilpotence of the bar differential is inherited from the
 boundary geometry of $\overline{\mathcal{M}}_{g,n}$.
\item \emph{Curvature = trace}
 (Theorem~\ref{thm:curvature-self-contraction}):
 the genus-$1$ curvature is the self-contraction of the
 propagator around a handle, equivalently the one-loop anomaly
 in the QME.
\item \emph{Loop order decomposition}
 (Theorem~\ref{thm:loop-order-spectral-sequence}):
 the genus-$g$ bar complex decomposes by the graph Betti number
 into a completed spectral sequence; finite-depth shadow towers give
 a finite-page collapse bound, while class-M towers require the
 completed inverse-limit statement.
\item \emph{Feynman involution $F^2 \simeq \mathrm{id}$}
 (Theorem~\ref{thm:feynman-involution},
 Corollary~\ref{cor:feynman-duality-qch}):
 bar-cobar inversion is the chiral specialization of the
 Getzler--Kapranov duality for modular operads.
\item \emph{$\mathfrak{grt}_1$ action}
 (Proposition~\ref{prop:weight-system-map}):
 the Grothendieck--Teichm\"uller Lie algebra acts on
 all-genus bar cohomology through the graph complex weight
 system, connecting quantum chiral homology to arithmetic
 geometry.
\end{enumerate}
Together with genus universality
(Theorem~\ref{thm:genus-universality}), complementarity
(Theorem~\ref{thm:quantum-complementarity-main}), and the
Lagrangian polarization
(Theorem~\ref{thm:lagrangian-complementarity}),
these pillars form the complete structural theory of quantum chiral
homology in every genus.
\end{remark}

\begin{remark}[Conditional $\Phi$-trace comparison binding Vol~I higher-genus to Vol~III]
\label{rem:hgf-phi-trace-identity}
The five pillars above admit a single cross-volume synthesis via the
CY-to-chiral two-stage factorisation
$\Phi_d^{(\Sigma_{d-1},C)} = \SpCh_{\Sigma_{d-1}, C} \circ \PhiFA_d$ of Vol~III, with
stage~$1$ a canonical $E_d$-holomorphic factorisation algebra on the
CY$_d$ variety (pinned by Kontsevich--Tamarkin $E_d$-formality and
Costello--Gwilliam--Li locality) and stage~$2$ factorisation homology
along a closed $(d-1)$-cycle $\Sigma_{d-1}$ followed by restriction to
a reference curve $C$.  Assume~$H_D^K(\mathcal A,g)$ in every genus.
On the higher-genus Vol~I side, the Hodge supertrace of the scalar
virtual bar object is the graded family
\[
\mathrm{STr}_{\mathrm{bar}}^{\mathrm{Hdg}}\!\bigl(\mathcal{A}\bigr)
\;:=\;
\bigl((-1)^g\kappa(\mathcal{A})\lambda_g\bigr)_{g\geq1}
\in
\prod_{g\geq1}H^{2g}(\overline{\mathcal{M}}_g).
\]
On the Vol~III side the proved scalar target is the Borcherds/Igusa
coordinate
\[
\Phi_{10}^{\mathrm{un}}=\Delta_5^2,\qquad
\kappa_{\mathrm{BKM}}(\mathfrak g_{\Delta_5})=5.
\]
For the $K3$-fibre specialisation of the canonical
$E_3$-holomorphic factorisation algebra
$\PhiFA_3(\cC_{K3 \times E})$, a chain-level Borcherds pushforward is
an additional comparison datum. If the finite Hall--Borcherds
source-recognition datum of
Definition~\ref{def:hgf-finite-hall-borcherds-source} is supplied,
and if the CY-to-chiral comparison maps commute with the bar
supertrace and the Hall--Borcherds target maps on every finite window,
then the conditional comparison reads
\[
\mathrm{push}_{K3 \times E}^{\mathrm{Borcherds}}\!\bigl(\kappa_{\mathrm{BKM}}\bigr)
\;=\;
\operatorname{ev}_{K3\times E}^{\mathrm{Hdg}}\!\left(
\mathrm{STr}_{\mathrm{bar}}^{\mathrm{Hdg}}(\mathcal A)\right)
\qquad
\bigl(\mathcal{A} = \SpCh_{\Sigma_2, C}(\PhiFA_3(\cC_{K3 \times E}))\bigr),
\]
a numerical identity binding the higher-genus modular-Koszul
trace to the BKM side of Vol~III through the stage-$2$ shadow of
$\PhiFA_3$.  The evaluation
$\operatorname{ev}_{K3\times E}^{\mathrm{Hdg}}$ is part of the named
comparison package, together with the finite source matrices, parity
fixture, relation theorem, and PBW compatibility.  See
Part~\ref{part:modular-tower-physics} for
the $\Phi$-functor construction and the Vol~III $K3 \times E$ side.
\end{remark}

\subsection{Logarithmic clutching kernels}
\label{subsec:log-clutching-kernels}

\index{clutching kernel!logarithmic}
\index{Fulton--MacPherson compactification!logarithmic}
The passage from genus-$0$ to all-genus bar complexes requires a
geometric mechanism for gluing configuration spaces along
marked-point identifications. Classical clutching maps on the
Deligne--Mumford side
$\overline{\mathcal{M}}_{g_1,n_1+1} \times
\overline{\mathcal{M}}_{g_2,n_2+1} \to
\overline{\mathcal{M}}_{g_1+g_2,\,n_1+n_2}$
have well-known analogues at the level of moduli, but their
configuration-space counterparts demand a compactification
that respects both the divisorial boundary and the iterated
blowup structure of the Fulton--MacPherson spaces.
The logarithmic framework of Mok~\cite{Mok25} provides exactly
the right setting.

\begin{definition}[Logarithmic clutching kernel; \ClaimStatusDefinitional]%
\label{def:log-clutching-kernel}
\index{clutching kernel!logarithmic!definition}
Fix a pair $(X,D)$ with $X$ smooth and $D \subset X$ a
normal-crossings divisor. Let $\FM_n(X \!\mid\! D)$ denote the
logarithmic Fulton--MacPherson compactification of the configuration
space of $n$ ordered points on $X \setminus D$, in the sense of
Mok~\cite{Mok25}. For stable data $(g_1,n_1+1)$ and $(g_2,n_2+1)$,
a \emph{logarithmic clutching kernel} is a proper correspondence class
\[
 K^{\mathrm{log}}_{g_1,n_1;\,g_2,n_2}
 \;\in\;
 \mathrm{Corr}\!\bigl(
 \FM_{n_1+1}(X \!\mid\! D) \times \FM_{n_2+1}(X \!\mid\! D),\;
 \FM_{n_1+n_2}(X \!\mid\! D)
 \bigr)
\]
that implements the gluing of the last marked point on the first
factor with the last marked point on the second factor.
\end{definition}

\begin{theorem}[Logarithmic clutching from degeneration geometry;
 \ClaimStatusConditional]%
\label{thm:log-clutching-degeneration}%
\index{clutching kernel!logarithmic!from degeneration geometry}
Type signature: \textup{(}Open quadrant, relative log-FM
correspondence presentation, Beilinson levels~$2$ and~$5$,
hypothesis package: Mok's degeneration formula plus the signed
log-FM residue-pushforward package of
Definition~\ref{def:signed-logfm-residue-pushforward-package} for
chain-level transport to the bar complex\textup{)}.
The degeneration formula of Mok~\cite[Corollary~5.3.4]{Mok25}
produces the underlying geometric degeneration correspondences for
logarithmic Fulton--MacPherson spaces.
The irreducible components of the special fibre of\/
$\operatorname{FM}_n(W/B)$ biject with rigid combinatorial types
$\rho$ \cite[Definition~4.7.1]{Mok25}, and each component is a
modification of the product
$\prod_{v \in V(S_\rho)} \operatorname{FM}_{I_v}(Y_v \mathbin{|} D_v)$.
The combinatorial cutting map $\kappa$ is a subdivision, bijective
on supports and lattices \cite[Lemma~5.3.1]{Mok25}. Interpreting
these geometric correspondences as logarithmic clutching kernels
compatible with the stable-graph bar differential is conditional on
the signed chain-level residue, pushforward, automorphism, and
homotopy-coherence data of
Definition~\ref{def:signed-logfm-residue-pushforward-package}.
\end{theorem}

\begin{proof}[References]
The geometric assertion is \cite[Corollary~5.3.4]{Mok25}, combined
with \cite[Lemma~5.3.1]{Mok25} for the subdivision property of the
cutting map. The last sentence is not a theorem of Mok: it is exactly
the manuscript's signed log-FM chain-transport hypothesis package.
\end{proof}

\begin{remark}[Conceptual cohesion]%
\label{rem:log-clutching-cohesion}
\index{clutching kernel!logarithmic!conceptual role}
The logarithmic clutching kernel is the geometric hinge of the
modular package: stable-curve gluing lives on the Deligne--Mumford
side, while configuration-space compactification lives on the
logarithmic Fulton--MacPherson side. The kernel
$K^{\mathrm{log}}$ is designed to make them interact
functorially. Together with the genus-universality theorem
(Theorem~\ref{thm:genus-universality}) and the stable-graph bar
differential (Lemma~\ref{lem:stable-graph-d-squared}), the
logarithmic clutching kernels would provide the geometric
foundation for the modular completion of the bar complex.
Note that Theorem~\ref{thm:planted-forest-tropicalization}
identifies the planted-forest coefficient algebra
$\mathbb{G}_{\mathrm{pf}}$ with the chain algebra on
$\operatorname{Trop}(\FM_n(C|D))$; the degeneration formula
\cite[Corollary~5.3.4]{Mok25} is the geometric incarnation of
the clutching law for the quartic resonance class
(Theorem~\ref{thm:nms-clutching-law-modular-resonance}).
\end{remark}

\begin{construction}[Log FM forgetful chain map; \ClaimStatusConditional]
\label{constr:log-fm-chain-map}
\index{logarithmic FM!forgetful chain map}
The logarithmic Fulton--MacPherson space
$\operatorname{FM}_n(X \mathbin{|} D)$ of
Mok~\cite[Theorem~3.3.1]{Mok25} supplies the geometric source for a
forgetful morphism to the classical FM compactification. After a
proper chain model and the signed log-FM residue-pushforward package
of Definition~\ref{def:signed-logfm-residue-pushforward-package} have
been fixed, this gives the conditional chain map
\[
\pi_*\colon
C_*\bigl(\operatorname{FM}_n(X \mathbin{|} D)\bigr)
\longrightarrow
C_*\bigl(\overline{\operatorname{FM}}_n(X)\bigr),
\]
induced by dropping the log structure.
The chain map $\pi_*$ has the following properties:
\begin{enumerate}[label=\textup{(\roman*)}]
\item On FM-type strata \textup{(}mutual collisions of marked
 points away from $D$\textup{)}, $\pi_*$ is the identity: the
 classical FM boundary encodes the same collision data.
\item On puncture-collision strata \textup{(}marked points
 approaching the divisor $D$\textup{)}, $\pi_*$ collapses the
 log boundary to lower-dimensional strata in the classical FM
 space. These extra strata are parametrised by the
 grid-subdivision data $(P, m)$ of Mok's planted forests
 \cite[Definition~3.1.5]{Mok25}.
\item The planted-forest differential $d_{\mathrm{pf}}$ in the
 five-component differential~\eqref{eq:D-five-components} is the
 conditional algebraic shadow of the \emph{kernel} of $\pi_*$ on boundary
 strata: it encodes exactly the codimension-${\geq}2$ boundary
 relations of $\operatorname{FM}_n^{\mathrm{log}}$ that the
 classical FM space $\overline{\operatorname{FM}}_n$ does not see.
 These extra boundary relations are precisely the higher
 $L_\infty$ homotopies that the planted-forest propagator
 computes in the graph-sum formula~\eqref{eq:elln-general-graph-sum}
 \textup{(}Construction~\textup{\ref{constr:explicit-convolution-linfty})}.
\end{enumerate}
Only the geometric log-FM compactification and its stratification are
external inputs from Mok. The pushforward on the chosen chain model
and the identification of its invisible boundary relations with
$d_{\mathrm{pf}}$ are part of the conditional signed
residue-pushforward package.

For $X = \mathbb{P}^1$ with $D = \{0, 1, \infty\}$, the identification
$\operatorname{FM}_{n-3}(\mathbb{P}^1 \mathbin{|} 0{+}1{+}\infty)
\cong \overline{\cM}_{0,n}$
of \cite[\S3.4.1]{Mok25}
 \textup{(}Example~\textup{\ref{ex:log-fm-moduli-rational})}
 gives an explicit model: the log~FM space
\emph{is} the moduli space of stable rational curves, and $\pi_*$ is
the tautological map forgetting the log structure on the boundary
divisors.
\end{construction}

% ================================================================
% SECTION 4.15: COMPLEMENTARITY THEOREM
% ================================================================

\section{Synthesis}
\label{sec:hgfnd-supplement}

\begin{remark}[Gritsenko additive lift of $\Delta_5$ and the source gate]
\label{rem:hgfnd-1}
The higher-genus foundations assembled in this chapter meet the K3
automorphic target at the Gritsenko additive lift. Starting from the
Jacobi cusp form of weight $5$, index $1$ on $\Gamma^J_0(1)$, the
Gritsenko additive lift produces
\[
\Delta_5 \;\in\; M_5^{\mathrm{cusp}}(\mathrm{Sp}_4(\mathbb Z),\mathrm{ch})
\]
as a Siegel modular form of weight~$5$ on $\bar{\mathcal A}_2$ with a
specific character $\mathrm{ch}$, such that
$\Phi_{10}^{\mathrm{un}}=\Delta_5^2$ is the unnormalised Igusa
weight-$10$ cusp form. Bruinier 2002 Prop.~5.1 supplies the
Heegner Chern-class reciprocity line after pullback along the
Kuga--Satake morphism $K3\hookrightarrow\bar{\mathcal A}_2$.

The phrase ``K3 chiral bialgebra $\mathbf H_{\Delta_5}$'' is a
conditional source phrase in this chapter. It denotes a chain-level
Hall or CoHA object only after the finite Hall--Borcherds
source-recognition datum of
Definition~\ref{def:hgf-finite-hall-borcherds-source} has been
constructed and compared with the Gritsenko--Nikulin target. Under
that extra hypothesis, the $g=2$ modular Koszul pair refines the
$g=1$ pair $(B_{\mathrm{ch}}(A),\Omega_X)$ through a
Humbert-divisor resolution carrying the $\mathbb Z/8$ monodromy
class recorded in Vol.~III, Chapter~\ref{V3-ch:k3-chiral-bialgebra-platonic}.
Absent that hypothesis, the paragraph asserts only the scalar
automorphic lift and the Heegner Chern-class target.
Primary-lit: Gritsenko 1995, Gritsenko--Nikulin 1996--98, Bruinier 2002
(arXiv:math/0108079), Kuga--Satake 1967.
\end{remark}

\begin{remark}[Baily--Borel compactification and the $24$-cusp structure]
\label{rem:hgfnd-2}
The Baily--Borel compactification $\bar{\mathcal A}_2^{\mathrm{BB}}$
carries a stratification
\[
\bar{\mathcal A}_2^{\mathrm{BB}} \;=\; \mathcal A_2 \;\sqcup\; \mathcal A_1^\sharp \;\sqcup\; \mathrm{pt},
\]
where $\mathcal A_1^\sharp$ is a modular curve of cusps and the final
point is the Borel cusp. If the finite Hall--Borcherds
source-recognition datum of
Definition~\ref{def:hgf-finite-hall-borcherds-source} has been
supplied, the higher-genus foundations of $\mathbf H_{\Delta_5}$
descend through this stratification:
$\mathbf H_{\Delta_5}|_{\mathcal A_2}$ is the generic-locus chiral
bialgebra, $\mathbf H_{\Delta_5}|_{\mathcal A_1^\sharp}$ is the
semi-classical limit after the Humbert-divisor blow-down, and
$\mathbf H_{\Delta_5}|_{\mathrm{pt}}$ is the fully abelianised
Heisenberg-type limit on $\mathcal F_{K3}$. Without that datum this is
only a stratification of the scalar automorphic target.

The Gritsenko--Nikulin 1997--98 denominator identity
\[
\Phi_{10}^{\mathrm{un}}(Z) \;=\; \prod_{(n,l,m)>0} (1 - q^n \zeta^l p^m)^{c(nm,l)}
\]
is an unconditional target identity: the exponents $c(nm,l)$ are the
elliptic-genus Fourier coefficients of $K3$ and the target
Borcherds root multiplicities. Reading each factor as a compact Hall
primitive, or reading the product as the higher-genus trace of the
bar--cobar adjunction on $\mathbf H_{\Delta_5}$, requires the parity
fixture, compact source matrices, comparison maps, no-extra-relations
theorem, and PBW compatibility in
Definition~\ref{def:hgf-finite-hall-borcherds-source}. See Vol.~III,
Chapter~\ref{V3-ch:k3-chiral-bialgebra-platonic},
\S\ref{sec:baily-borel-24-cusps}.
\end{remark}
